% concatenated from V2025.tex
%\documentclass[11pt,a4paper,oneside]{amsart}
%\documentclass[CJK,11pt]{amsart}
\documentclass[11pt,a4paper,oneside]{amsart}
% \usepackage{graphicx}
\usepackage{a4,a4wide}

\usepackage{geometry}
\geometry{top=2.5cm, bottom=2.8cm, left= 2.5cm, right= 2.5cm}

\usepackage{amsmath,amsfonts,amssymb,amsthm}
\usepackage{graphics}
\usepackage{dsfont}
\usepackage{epsfig}
\usepackage{esint}
\usepackage[colorlinks,linkcolor=blue]{hyperref}
%\usepackage{showkeys}
\setlength{\parindent}{0cm}
%\usepackage[inline]{showlabels}
\usepackage{comment}

%Diagram
\usepackage{tikz}
\usetikzlibrary{matrix}
\usetikzlibrary{quotes,angles,positioning}

\numberwithin{equation}{section}



\newcommand{\Bild}[2]{\begin{center}\begin{tabular}{l}
\setlength{\epsfxsize}{#1}
\epsfbox{#2}
\end{tabular}\end{center}}
\newcommand{\lessim}{\stackrel{<}{\sim}}
\newcommand{\gesim}{\stackrel{>}{\sim}}
\newcommand{\ignore}[1]{}
\newcommand{\av}{-\hspace{-2.4ex}\int}
\newcommand{\del}{\delta\hspace{-0.2ex}}
\newtheorem{theorem}{Theorem}[section]
\newtheorem{definition}[theorem]{Definition}
\newtheorem{proposition}[theorem]{Proposition}
\newtheorem{remark}[theorem]{Remark}
\newtheorem{lemma}[theorem]{Lemma}
\newtheorem{corollary}[theorem]{Corollary}
\newtheorem{assumption}[theorem]{Assumption}
\newtheorem{example}[theorem]{Example}
\newtheorem{conjecture}[theorem]{Conjecture}
%%% Some macros

% Mathbb letters
\newcommand{\C}{\mathbb{C}}
\newcommand{\E}{\mathbb{E}}
\newcommand{\N}{\mathbb{N}}
\renewcommand{\P}{{\mathbb P}}
\newcommand{\Q}{{\mathbb Q}}
\newcommand{\R}{\mathbb{R}}
\newcommand{\T}{\mathbb{T}}
\newcommand{\Z}{\mathbb{Z}}
\newcommand{\lt}{\left}
\newcommand{\rt}{\right}
\newcommand{\Boule}{B}
\newcommand{\langl}{\langle}
\newcommand{\rangl}{\rangle}
\newcommand{\pour}{\text{ for }}
\newcommand{\et}{\text{ and }}
\newcommand{\si}{\text{ if }}
\newcommand{\HH}{H}
\newcommand{\Torus}{\T_L}
\newcommand{\pourtout}{\text{for all }}
\newcommand{\loc}{\text{loc}}
\newcommand{\cc}{\text{C}}
\newcommand{\comp}{c}
\newcommand{\abar}{\overline{a}}
\newcommand{\dd}{\mathrm{d}}
\newcommand{\f}{\text{F}}
\newcommand{\Ring}{R}
\newcommand{\ubar}{\overline{u}}
\newcommand{\dans}{\text{ in }}
\newcommand{\Dom}{\text{D}}
\newcommand{\bb}{\text{B}}
\newcommand{\LL}{\mathrm{L}}
\newcommand{\rhs}{r.~h.~s.\,}
\newcommand{\lhs}{l.~h.~s.\,}
\newcommand{\iid}{i.~i.~d.\,}
\newcommand{\QQ}{\mathrm{Q}}
\newcommand{\WW}{\mathrm{W}}
\newcommand{\qq}{\text{Q}}
\newcommand{\pp}{\mathrm{P}}
\newcommand{\fm}[1]{\textcolor{blue}{[F: #1]}}
\newcommand{\nc}[1]{\textcolor{red}{[N: #1]}}
\def\leb{\mathsf{Leb}}
\newcommand{\m}{\mathrm{m}}

% Macros Francesco
\DeclareMathOperator*{\Lip}{Lip}
\DeclareMathOperator*{\dist}{\mathsf{d}}

 %  colorful comments

\usepackage{color}
\definecolor{darkred}{rgb}{0.9,0.1,0.1}
\definecolor{darkblue}{rgb}{0,0,0.7}
\definecolor{darkgreen}{rgb}{0,0.5,0}
\def\comment#1{\marginpar{\raggedright\scriptsize{\textcolor{darkred}{#1}}}}
\def\hwcomment#1{\marginpar{\raggedright\scriptsize{\textcolor{darkblue}{#1}}}}
\newcommand{\hw}[1]{{\color{darkblue}{#1}}}

\newcommand{\diam}{\hspace{-0.4ex}\diamond\hspace{-0.4ex}}

\begin{document}
\title[Annealed quantitative estimates for the 2D-discrete random matching problem]{Annealed quantitative estimates for the quadratic 2D-discrete random matching problem}
\author{  Nicolas Clozeau \address[Nicolas Clozeau]{IST Austria, Austria} \email{nicolas.clozeau@ist.ac.at} \hspace*{0.5cm} Francesco Mattesini \address[Francesco Mattesini]{Universit\"at M\"unster \& MPI Leipzig,  Germany} \email{francesco.mattesini@uni-muenster.de}  } 
\thanks{NC has received funding from the European Research Council (ERC) under the Eu\-ropean Union’s Horizon 2020 research and innovation programme (grant agreement No 948819) \includegraphics[height=\fontcharht\font`\B]{erc_logo.jpg}\,\includegraphics[height=\fontcharht\font`\B]{flag_logo.jpg}}
\thanks{FM is supported by the Deutsche Forschungsgemeinschaft (DFG, German Research Foundation) through the SPP 2265 {\it Random Geometric Systems}. FM has been funded by the Deutsche Forschungsgemeinschaft (DFG, German Research Foundation) under Germany's Excellence Strategy EXC 2044 -390685587, Mathematics M\"unster: Dynamics--Geometry--Structure. FM has been funded by the Max Planck Institute for Mathematics in the Sciences. }
%
%\maketitle
%
\begin{abstract}
%
We study a %bipartite 
random matching problem on closed compact $2$-dimensional Riemannian manifolds (with respect to the squared Riemannian distance), with samples of random points whose common law is absolutely continuous with respect to the volume measure with strictly positive and bounded density. We show that given two sequences of numbers $n$ and $m=m(n)$ of points, asymptotically equivalent as $n$ goes to infinity, the optimal transport plan between the two empirical measures $\mu^n$ and $\nu^{m}$ is quantitatively well-approximated by $\big(\mathrm{Id},\exp(\nabla h^{n})\big)_\#\mu^n$ where $h^{n}$ solves a linear elliptic PDE obtained by a regularized first-order linearization of the Monge-Ampère equation. This is obtained in the case of samples of correlated random points for which a stretched exponential decay of the $\alpha$-mixing coefficient holds and for a class of discrete-time sub-geometrically ergodic Markov chains having a unique absolutely continuous invariant measure with respect to the volume measure.

%of i.~i.~d.~random points and in the case of samples of correlated random points for which $\rho$ is continuous and a stretched exponential decay of the strong mixing coefficient holds.  %Despite that the method employed leads to a sub-optimal rate of convergence, we provide a numerical simulation, based on a formal second-order linearization of the Monge-Ampère equation, leading to the conjecture of an optimal convergence rate of \nc{put what we obtain, surprise !}.  
%
\end{abstract}
%
\maketitle
%
\begin{center}
\textbf{Keywords: }Optimal transport $\cdot$ Matching problem $\cdot$ Quantitative estimates
\end{center}
%

%
\tableofcontents
%
\section{Introduction and statement of the main results}\label{introsec}
%
\subsection{The random matching problem and its asymptotic}\label{TMP}
%
%\nc{Start with basic stuff on the random matching problem, definition applications, ect. Then, recall the known results }




The random matching problem is a popular optimization problem at the interface between analysis and probability with applications in many different fields such as statistical physics \cite{caracciolo2014scaling, MezPar}, computer science \cite{CSMatch} and economics \cite{EcMatch, GalShap}. Within the mathematical literature, it has been subject of intense studies due to its interactions with many areas, including for instance graph theory \cite{lovasz2009matching} and geometric probability \cite{steele1997probability}. In this paper we focus on one of its simple versions. Let $\{X_k\}_{1\leq k\leq n}$ and $\{Y_k\}_{1\leq k\leq m}$ (with possibly $m>n$) be two families of random points on a compact Riemannian manifold $\mathcal{M}$ (endowed with the Riemannian distance $\dd$). We are interested in the quadratic matching problem
\begin{equation}\label{eq:transintr}
\min_{\pi\in\Pi_{nm}} \sum_{i=1}^n \sum_{j=1}^m \pi_{ij}\,\mathrm{d}^2(X_i ,Y_j)%=W_2^2 (\mu^n, \nu^m),
\end{equation}
where
$$\Pi_{nm}:=\bigg\{\pi\in [0,1]^{n\times m}\ \Big\vert\ \sum_{i=1}^n \pi_{ij} = \tfrac{1}{m}\quad\text{and}\quad\sum_{j=1}^{m} \pi_{ij} = \tfrac{1}{n}\bigg\}.$$
Classically, \eqref{eq:transintr} can be phrased in terms of a transport problem. Indeed, letting 
\begin{equation}\label{eq:empmeas}
%
\mu^n:=\frac1n\sum_{i=1}^n \delta_{X_i}\quad\text{and}\quad\nu^m:=\frac1m \sum_{j=1}^m \delta_{Y_j},
%
\end{equation}
be the empirical measures associated with the two point clouds, the linear programming problem \eqref{eq:transintr} amounts to determine the quadratic Wasserstein distance $\ W_2^2(\mu^n,\nu^m)$. 

\medskip
In the special case $n=m$ the Birkhoff-von Neumann  Theorem provides a correspondence between \eqref{eq:transintr} and the usual bipartite matching
\begin{equation}\label{eq:minweightgraph}
%
\min_{\sigma\in \mathcal{S}_{n}} \sum_{i=1}^n \mathrm{d}^2(X_i,Y_{\sigma(i)}),
%
\end{equation}
%
where $\mathcal{S}_{n}$ denotes the set of injective maps $\sigma: \{1, \dots, n\} \rightarrow \{1, \dots, n\}$. 
%
Indeed, since $\Pi_{nn}$ is a convex polytope, minimizers in \eqref{eq:transintr} have to be searched among extremal points. By the Birkhoff-von Neumann Theorem \cite[Lemma 2.1.3]{bapat}, the latter are nothing but permutation matrices (up to a factor $\tfrac{1}{n}$). 

%In this paper, we focus on one of its simple version when trying to match two families of random points $\{X_k\}_{1\leq k\leq n}$ and $\{Y_k\}_{1\leq k\leq m}$ (with possibly $n>m$) on a compact Riemannian manifold $\mathcal{M}$ (endowed with the Riemannian distance $\dd$), while assuming that the cost of matching two points $X_i$ and $Y_j$ depends only on $\mathrm{d}^2(X_i,Y_j)$. Its mathematical formulation reads 
%\begin{equation}\label{eq:transintr}
%\min_{\pi\in\Pi_{nm}} \sum_{i=1}^n \sum_{j=1}^m \pi_{ij}\,\mathrm{d}^2(X_i ,Y_j)%=W_2^2 (\mu^n, \nu^m),
%\end{equation}
%where
%$$\Pi_{nm}:=\bigg\{\pi\in [0,1]^{n\times m}\ \Big\vert\ \sum_{i=1}^n \pi_{ij} = \tfrac{1}{m}\quad\text{and}\quad\sum_{j=1}^{m} \pi_{ij} = \tfrac{1}{n}\bigg\}.$$
%
%In the special case $n=m$, the problem \eqref{eq:transintr} has a one to one correspondence with the bipartite random matching problem
%
%\begin{equation}\label{eq:minweightgraph}
%
%\min_{\sigma\in \mathcal{S}_{n}} \sum_{i=1}^n \mathrm{d}^2(X_i,Y_{\sigma(i)}),
%
%\end{equation}
%
%where $\mathcal{S}_{n}$ denotes the set of injective maps $\sigma: \{1, \dots, n\} \rightarrow \{1, \dots, n\}$. Indeed, since $\Pi_{nn}$ is a convex polytope, minimizers in \eqref{eq:transintr} have to be searched among extremal points. By Birkhoff-von Neumann' theorem \cite[Lemma 2.1.3]{bapat}, the latter are nothing but permutation matrices (up to a factor $\tfrac{1}{n}$). 

%
\medskip
%

%By its definition, \eqref{eq:transintr} has an exact formulation in terms of the Wasserstein cost of order 2 between the two empirical measures
%
%\begin{equation}\label{eq:empmeas}
%
%\mu^n:=\frac1n\sum_{i=1}^n \delta_{X_i}\quad\text{and}\quad\nu^m:=\frac1m \sum_{j=1}^m \delta_{Y_j},
%
%\end{equation}
%
%that is solving \eqref{eq:transintr} amounts to determine the quadratic Wasserstein distance $W^2_2(\mu^n,\nu^m)$. That is the strategy we follow in this paper.

%
\medskip
%

A first natural question is to understand the asymptotics of \eqref{eq:transintr} as $n,m\uparrow\infty$. For the same number of samples $n=m$ and independently and identically distributed (i.i.d.) on the unit square $[0,1]^d$, the scaling of the cost \eqref{eq:transintr} has been well understood in the mathematical and statistical physics literature. A simple heuristic argument, see for instance \cite{MezPar}, suggests that given a point $X_i$, we can find a point $Y_j$ within a volume of order  $O(n^{-1})$ with high probability. For this reason, the typical inter-point distance is of order  $O(n^{-\frac1d})$ suggesting that the scaling of \eqref{eq:transintr} is of order  $O(n^{-\frac2d})$. Although attractive, this heuristic turns out to be unfortunately false in low dimension showing a critical behavior when $d=2$. This critical case is the one on which we focus on in this paper. Ajtai, Koml\'os and Tusn\'ady in \cite{AKT84} were the first to show that, for i.i.d. uniform samples, a logarithmic correction is needed, deriving\footnote{We use the notation $A \lesssim B$ if there exists a global constant $C>0$, which may only depend on $d$, such that $A\le CB$. We write $A\sim B$ if both $A \lesssim B$ and $B \lesssim A$ hold.}
%
\begin{equation}\label{eq:unifmatchingrates}
%
\mathbb{E}\big[W_2^2 (\mu^n, \nu^n)\big]\sim \frac{\log(n)}{n},
%
\end{equation}
%
extended later by Talagrand in \cite{talagrand1992ajtai} for clouds of i.i.d. points which are distributed accordingly to more general common law. A recent breakthrough was obtained within the physics community by Caracciolo, Lucibello, Parisi and Sicuro in \cite{caracciolo2014scaling}, and further developed by Caracciolo and Sicuro in \cite{caracciolo2015scaling} and by Sicuro in \cite{sicuro2017euclidean}, where the asymptotics of the cost are formally derived thanks to a novel PDE approach and optimal transport theory rather than combinatorics. A couple of years later, in general $2$-dimensional compact Riemannian manifolds without boundary, the first-order asymptotic has been rigorously justified by Ambrosio, Stra and Trevisan in \cite{ambrosio2019pde} for i.i.d. uniform samples and recently extended by Ambrosio, Goldman and Trevisan in \cite{ambrosio2021quadratic} for samples distributed accordingly to more general laws which are absolutely continuous (with H\"older continuous density) w.r.t. the volume measure $\dd \mathrm{m}$, leading to 
%
\begin{equation}\label{eq:AsymptoticCost}
%
\lim_{n\rightarrow\infty} \frac{n}{\log(n)} \mathbb{E}\big[W_2^2 (\mu^n,\nu^n)\big] = \frac{\vert\mathcal{M}\vert}{2\pi},
%
\end{equation}
where $|\mathcal{M}|$ denotes the Lebesgue measure of $\mathcal{M}$. The case $n\neq m$ with $n,m\uparrow\infty$ with similar rates is also covered, see \cite[Theorem 1.2]{ambrosio2021quadratic}. %The main novelty of these works is on the originality of the proof techniques, based on a PDE approach and optimal transport theory rather than combinatorics, which allow for an explicit description of the cost and its minimizers. 












\medskip
The novel approach introduced in \cite{caracciolo2014scaling}, later revised in \cite{benedetto2020euclidean}, consists in a linearization of the Monge-Amp\`ere equation that allows for an explicit description of the cost thanks to the linearized proxies (see Section \ref{sec:lin} for more details). The aim of this work is to quantitatively justify the linearization ansatz in terms of convergence of the approximating minimizers of \eqref{eq:transintr} towards the optimal ones. In particular, we are interested in the case where the points are identically distributed with a common law $\rho\,\dd\m$ (where we recall that $\dd\m$ denotes the volume measure) and $\rho$ satisfies for some $\lambda,\Lambda>0$
%The novelty of the approach introduced in \cite{caracciolo2014scaling}, later revised in \cite{benedetto2020euclidean} for non-constant densities, consists on using that minimizers of \eqref{eq:transintr}, thanks to its formulation in terms of the quadratic Wasserstein distance, solve the Monge-Ampère equation. Exploiting the equation and its linearisation, the asymptotics of the cost and approximation of minimizers are then formally derived (we reproduce the formal first-order expansion in Section \ref{MainResultSection}). The aim of this work is to quantitatively justify the linearization ansantz in terms of convergence of the approximating minimizers of \eqref{eq:transintr} towards the optimal ones. In particular, we are interested in the case where the points are identically distributed with a common law $\rho\,\dd\m$ (where we recall that $\dd\m$ denotes the volume measure) and $\rho$ satisfies for some $\lambda,\Lambda>0$
%
\begin{equation}\label{Ellipticity}
%
\lambda\leq\rho\leq\Lambda.
%
\end{equation}
To the best of our knowledge, there are only few results on the asymptotic behavior of the transport map and they are so far limited to the case of i.i.d. uniform samples in the study of the semi-discrete matching problem (that is couplings between $\mu^n$ and $\dd\m$), see the work of Ambrosio, Glaudo and Trevisan in \cite{ambrosio2019optimal}. In connection with this work, quantitative estimates on the optimal map for the matching between the Lebesgue measure and Poisson clouds have been obtained by Goldman, Huesmann and Otto in \cite{GHO1} and Goldman and Huesmann in \cite{GH}. 

%
\medskip
%

Our extension in this paper is fourfold: First, we look at more general distribution of points and we consider the case of general densities $\rho$ satisfying \eqref{Ellipticity}. Second, we do not assume independence and we consider samples which may possess correlations. Third, we do not restrict the analysis to the semi-discrete matching problem and we also investigate the ansatz for the full matching problem \eqref{eq:transintr}. %regularizing the right-hand side of \eqref{eq:ApproxPDE} via the heat-kernel at a time $\sim\frac{1}{n}$ and for weakly correlated points (see Assumptions \ref{Assumptions}), the second item of \eqref{eq:approxcost} holds true with a quantitative convergence rate. 
Finally, we investigate the case where the points are not identically distributed and we extend our result to a class of sub-geometrically ergodic Markov chains.

%
\medskip
%

We finally mention that the effectiveness of the linearization ansatz introduced in \cite{caracciolo2014scaling} is not only limited to the case of i.i.d. distributed points on bounded domains, but it can be employed in many different settings. See for instance \cite{CagliotiPieroni} for an interesting application to the matching on unbounded domains, \cite{Ledoux, Ledoux2, Ledoux3} for an application to Gaussian matching, \cite{Jonas} for an application in random matrix theory, \cite{Wang2, Wang3, Wang, HMTfBm} for an application to a continuous instance of the matching problem, i.e. when the empirical measure is replaced by the occupation measure of a stochastic process. It is worth to further mention that these techniques can also be employed when considering the matching problem with $p$-costs in higher dimension, see \cite{GoldTrev} and when considering a larger class of optimization problems, see \cite{GoldTrevOpt}.
%A possible approach to understand the minimizers in \eqref{eq:transintr} is to justify the linearization ansatz introduced in \cite{caracciolo2014scaling}, later revised in \cite{benedetto2020euclidean} for non-constant densities, that we briefly reproduce here. 

\subsection{Linearization ansatz}\label{sec:lin}
We now briefly reproduce the linearization ansatz introduced in \cite{caracciolo2014scaling}. For simplicity, we consider the case $\mathcal{M} = \mathbb{T}^2$, $n=m$ and i.i.d. samples with common distribution $\rho \, \dd \mathrm{m}$. Let $T^n$ be an optimal transport map (whose existence is ensured by Brenier's Theorem \cite{Brenier}) between $\mu^n$ and $\nu^n$. Based on the transport relation $T^n_{\#}\mu^n=\nu^n$ and a change of variables, $T^n$ solves (formally) the Monge-Amp\`ere equation 
%
\begin{equation}\label{eq:MongAmp}
%
(\nu^n\circ T^n) \det(\nabla T^n) = \mu^n.
%
\end{equation}
%
Since the cost is quadratic, by \cite[Theorem 1.25]{Santambrogio}, there exists a function $h^n$ such that $T^n= \mathrm{Id}+\nabla h^n$. Applying the law of large numbers, we further have the weak convergences $\mu^n,\nu^n {\rightharpoonup}\rho\, \dd \mathrm{m}$ as $n\uparrow \infty$ so that we expect $T^n\approx \mathrm{Id}$ as $n\uparrow\infty$. Thus, this suggests that the correction $\nabla h^n$ is small as $n\uparrow\infty$, allowing to perform (formally) the Taylor expansions
%
\begin{equation}\label{eq:TaylExp}
\nu^n \circ T^n \approx \nu^n + \nabla \nu^n\cdot\nabla h^n\quad\text{and}\quad \det(\nabla T^n) \approx 1 + \Delta h^n.
\end{equation}
%
Plugging \eqref{eq:TaylExp} into \eqref{eq:MongAmp}, neglecting the higher order terms and replacing $\mu^n$ by $\rho$ yields  
%
\begin{equation}\label{eq:ApproxPDE}
\nabla \cdot \rho \nabla h^n = \mu^n-\nu^n.
\end{equation}
%
This formal linearization suggests the following two conjectures 
%
\begin{equation}\label{eq:approxcost}
%
\lim_{n\uparrow\infty}\bigg\vert\mathbb{E}\big[W_2^2 (\mu^n, \nu^n)\big]-\mathbb{E}\bigg[\int_{\mathbb{T}^2} |\nabla h^n|^2\rho\,\dd\m\bigg]\bigg\vert=0\quad\text{and}\quad\lim_{n\uparrow\infty}\int_{\mathbb{T}^2}\vert T^n-(\text{Id}+\nabla h^n)\vert^2=0\text{ a.s.}
%
\end{equation}
%
Unfortunately, \eqref{eq:approxcost} cannot hold as it is, since the solution of \eqref{eq:ApproxPDE} does not belong to $\mathrm{H}^1$ due to the roughness of the source term. To overcome this, following the strategy in \cite{ambrosio2019pde}, a regularization using the heat-semigroup at time $\sim \frac{1}{n}$ (up to logarithmic corrections) is made. Doing so, the first item of \eqref{eq:approxcost} turns out to be true, leading to the result \eqref{eq:AsymptoticCost} (see for instance \cite{ambrosio2019finer} for a convergence rate). 

%{\color{blue} 

%\medskip
%\noindent
%@@@@@@@@@@@@@\\@@@@@@@@@@@@@\\@@@@@@@@@@@@@\\@@@@@@@@@@@@@\\@@@@@@@@@@@@@\\@@@@@@@@@@@@@\\@@@@@@@@@@@@@\\@@@@@@@@@@@@@\\@@@@@@@@@@@@@\\@@@@@@@@@@@@@\\@@@@@@@@@@@@@\\@@@@@@@@@@@@@\\@@@@@@@@@@@@@\\@@@@@@@@@@@@@\\@@@@@@@@@@@@@\\@@@@@@@@@@@@@\\@@@@@@@@@@@@@\\@@@@@@@@@@@@@\\@@@@@@@@@@@@@\\@@@@@@@@@@@@@\\@@@@@@@@@@@@@\\@@@@@@@@@@@@@\\@@@@@@@@@@@@@\\@@@@@@@@@@@@@\\@@@@@@@@@@@@@\\@@@@@@@@@@@@@\\@@@@@@@@@@@@@\\@@@@@@@@@@@@@\\@@@@@@@@@@@@@\\@@@@@@@@@@@@@\\@@@@@@@@@@@@@\\@@@@@@@@@@@@@\\@@@@@@@@@@@@@\\@@@@@@@@@@@@@}

%\medskip
%
%The bipartite random matching problem is a random optimization problem asking to find, among all possible matching in a bipartite weighted graph, the one that minimizes the sum of the costs of the chosen edges. In this paper, we focus on one of its simple version when trying to match two families of random points $\{X_k\}_{1\leq k\leq n}$ and $\{Y_k\}_{1\leq k\leq m}$ (with possibly $m>n$) on a compact Riemannian manifold $\mathcal{M}$ (endowed with the Riemannian distance $\dd$), while assuming that the cost of matching two points $X_i$ and $Y_j$ depends only on $\mathrm{d}^2(X_i,Y_j)$. Its mathematical formulation reads 
%The random matching problem is a random optimization problem at the intersection of analysis, probability and combinatorics which has been subject of intense research in the last 30 years. In this paper we focus on its Euclidean version. Let $\mathcal{M}$ be a closed compact $d$-dimensional Riemannian manifold. Given two finite families of random points $(X_i)_{i=1}^n, (Y_j)_{j=1}^m \subset \mathcal{M}$ with common law $\mu$ we are interested in minimizing the matching cost
%
%\begin{equation}\label{eq:minweightgraph}
%
%\min_{\sigma\in \mathcal{S}_{n,m}} \sum_{i=1}^n \mathrm{d}^2(X_i,Y_{\sigma(i)}),
%
%\end{equation}
%
%where $\mathcal{S}_{n,m}$ denotes the set of injective maps $\sigma: \{1, \dots, n\} \rightarrow \{1, \dots, m\}$. 
%
%The optimization problem \eqref{eq:minweightgraph} arises in many questions of different areas of mathematics, including for instance combinatorics and graph theory \cite{lovasz2009matching} or geometric probability \cite{steele1997probability}.

%
%\medskip
%

%Classically, the minimization problem \eqref{eq:minweightgraph} is often approximated by the Monge-Kantorovich optimal transport problem consisting of replacing $\sigma$ by a coupling (or transport plan), that is a matrix $\{\pi_{ij}\}_{ij}\in [0,1]^{n\times m}$ belonging to the set
%
%$$\Pi_{nm}:=\bigg\{\pi\in [0,1]^{n\times m}\ \Big\vert\ \sum_{i=1}^n \pi_{ij} = \tfrac{1}{m}\quad\text{and}\quad\sum_{j=1}^{m} \pi_{ij} = \tfrac{1}{n}\bigg\}.$$
%
%By its definition, the approximating problem has an exact formulation in terms of the Wasserstein cost of order $2$ between the two empirical measures 
%
%\begin{equation}\label{eq:empmeas}
%
%\mu^n:=\frac1n\sum_{i=1}^n \delta_{X_i}\quad\text{and}\quad\nu^m:=\frac1m \sum_{j=1}^m \delta_{Y_j},
%
%\end{equation}
%
%which reads
%
%\begin{equation}\label{eq:transintr}
%\min_{\pi\in\Pi_{nm}} \sum_{i=1}^n \sum_{j=1}^m \pi_{ij}\,\mathrm{d}^2(X_i ,Y_j)=W_2^2 (\mu^n, \nu^m).
%\end{equation}
%
%In the special case $n=m$, \eqref{eq:minweightgraph} and \eqref{eq:transintr} are equivalent. Indeed, since $\Pi_{nn}$ is a convex polytope, minimizers in \eqref{eq:transintr} have to be searched among extremal points. From Birkhoff-von Neumann's Theorem \cite[Lemma 2.1.3]{bapat}, the latter are nothing but permutation matrices (up to a factor $\tfrac{1}{n}$). From now on, we focus on the formulation of the problem given in \eqref{eq:transintr}. 

%
%\medskip
%

%A first natural question is to understand the asymptotic of \eqref{eq:transintr} as $n,m\uparrow\infty$. For the same number of samples $n=m$ and independently and identically distributed ($i.~i.~d.$) on the unit square $[0,1]^d$, the scaling of the cost \eqref{eq:transintr} has been largely studied in the mathematical and statistical physics literature. A simple heuristic argument, see for instance \cite{MezPar}, suggests that given a point $X_i$ from the first point cloud, we can find a point $Y_j$ within a volume of order  $O(n^{-1})$ with high probability. For this reason, the typical inter-point distance is of order  $O(n^{-\frac1d})$ suggesting that the scaling of \eqref{eq:transintr} is of order  $O(n^{-\frac2d})$. Although attractive, this heuristic turns out to be unfortunately false in low dimension showing a critical behavior when $d=2$, case on which our work is focused. In this case, Ajtai, Koml\'os and Tusn\'ady in \cite{AKT84} were the first to show that, for uniform samples, a logarithmic correction is needed, deriving\footnote{We use the notation $A \lesssim B$ if there exists a global constant $C>0$, which may only depend on $d$, such that $A\le CB$. We write $A\sim B$ if both $A \lesssim B$ and $B \lesssim A$ hold.}
%
%\begin{equation}\label{eq:unifmatchingrates}
%
%\mathbb{E}\big[W_2^2 (\mu^n, \nu^n)\big]\sim \frac{\log(n)}{n},
%
%\end{equation}
%
%extended later by Talagrand in \cite{talagrand1992ajtai} for general laws. A recent breakthrough was obtained by the physical community By Caracciolo, Lucibello, Parisi and Sicuro in \cite{caracciolo2014scaling}, and further developed by Caracciolo and Sicuro in \cite{caracciolo2015scaling} and by Sicuro in \cite{sicuro2017euclidean}, where the asymptotics of the cost are formally derived. A couple of years later, in general $2$-dimensional compact Riemannian manifolds without boundary, the first-order asymptotic has been rigorously justified by Ambrosio, Stra and Trevisan in \cite{ambrosio2019pde} for uniform samples and recently extended by Ambrosio, Goldman and Trevisan in \cite{ambrosio2021quadratic} for more general laws which are absolutely continuous (with H\"older continuous density) w. r. t the volume measure $\dd \mathrm{m}$, leading to 
%
%\begin{equation}\label{eq:AsymptoticCost}
%
%\lim_{n\rightarrow\infty} \frac{n}{\log(n)} \mathbb{E}\big[W_2^2 (\mu^n,\nu^n)\big] = \frac{\vert\mathcal{M}\vert}{2\pi},
%
%\end{equation}
%where $|\mathcal{M}|$ denotes the Lebesgue measure of $\mathcal{M}$. The case $n\neq m$ with $n,m\uparrow\infty$ with similar rates is also covered, see \cite[Theorem 1.2]{ambrosio2021quadratic}. The main novelty of these works is on the originality of the proof techniques, based on a PDE approach and optimal transport theory rather than combinatorics. 

%
%\medskip
%

%A second natural question is to understand the minimizers in \eqref{eq:transintr} as $n,m\uparrow\infty$. We are in particular interested in the case where the points are identically distributed with a common law $\rho\,\dd\m$ where $\dd\m$ denotes the volume measure and $\rho$ satisfies for some $\lambda,\Lambda>0$
%
%\begin{equation}\label{Ellipticity}
%
%\lambda\leq\rho\leq\Lambda.
%
%\end{equation}
%
%A possible approach to understand the minimizers in \eqref{eq:transintr} is to justify the linearization ansatz introduced in \cite{caracciolo2014scaling}, later revised in \cite{benedetto2020euclidean} for non-constant densities, that we briefly reproduce here. For simplicity, we consider the case $\mathcal{M} = \mathbb{T}^2$, $n=m$ and $i.~i.~d.$ samples. Let $T^n$ be an optimal transport map (whose existence is ensured by Brenier's Theorem \cite{Brenier}) between $\mu^n$ and $\nu^n$. Based on the transport relation $T^n_{\#}\mu^n=\nu^n$ and a change of variables, $T^n$ solves (formally) the Monge-Amp\`ere equation 
%
%\begin{equation}\label{eq:MongAmp}
%
%(\mu^n\circ T^n) \det(\nabla T^n) = \nu^n.
%
%\end{equation}
%
%Since the cost is quadratic, by \cite[Theorem 1.25]{Santambrogio} there exists a function $h^n$ such that $T^n= \mathrm{Id}+\nabla h^n$. Applying the law of large numbers, we further have the weak convergences $\mu^n,\nu^n {\rightharpoonup}\rho$ as $n\uparrow \infty$ so that we expect $T^n\approx \mathrm{Id}$ as $n\uparrow\infty$. Thus, this suggests that the correction $\nabla h^n$ is small as $n\uparrow\infty$, allowing to perform (formally) the Taylor expansions
%
%\begin{equation}\label{eq:TaylExp}
%\mu^n \circ T^n \approx \mu^n + \nabla \mu^n\cdot\nabla h^n\quad\text{and}\quad \det(\nabla T^n) \approx 1 + \Delta h^n.
%\end{equation}
%
%Plugging \eqref{eq:TaylExp} into \eqref{eq:MongAmp}, neglecting the higher order terms and replacing $\mu^n$ by $\rho$ yields  
%
%\begin{equation}\label{eq:ApproxPDE}
%-\nabla \cdot \rho \nabla h^n = \mu^n-\nu^n.
%\end{equation}
%
%This formal linearization suggests the following two conjectures 
%
%\begin{equation}\label{eq:approxcost}
%
%\lim_{n\uparrow\infty}\bigg\vert\mathbb{E}\big[W_2^2 (\mu^n, \nu^n)\big]-\mathbb{E}\bigg[\int_{\mathbb{T}^2} |\nabla h^n|^2\bigg]\bigg\vert=0\quad\text{and}\quad\lim_{n\uparrow\infty}\int_{\mathbb{T}^2}\vert T^n-(\text{Id}+\nabla h^n)\vert^2=0\text{ a.s.}
%
%\end{equation}
%
%Unfortunately, \eqref{eq:approxcost} cannot hold as it is since the solution of \eqref{eq:ApproxPDE} does not belong to $\mathrm{H}^1$ due to the roughness of the source term. To overcome this, following the strategy in \cite{ambrosio2019pde}, a regularization using the heat-semigroup at time $\sim \frac{1}{n}$ is made. Doing so, the first item of \eqref{eq:approxcost} turns out to be true, leading to the result \eqref{eq:AsymptoticCost} (see for instance \cite{ambrosio2019finer} for a convergence rate). 
%Without any further assumption on $\rho$ and $\mu^n$ this conjecture turns out to be false. Indeed, in the case of $\rho$ being a constant density \eqref{eq:ApproxPDE} becomes a Poisson equation and it is not hard to show that the right hand side of \eqref{eq:approxcost} diverges. This issue can be fixed via a regularization argument. More precisely, let $\mu^{n,t} := \pp_t \mu^n, \rho_t := \pp_t \rho$ be the evolution at a certain small time $t\ll 1$ of the empirical measure through the heat semigroup. If we argue with $\rho\,\dd m, \mu^n$ replaced by $\rho_t\,\dd m, \mu^{n,t}$ then the ansatz \eqref{eq:approxcost} is sound and Dirichlet energy of the associated solution $h^{n,t}$ well approximates $W_2^2 (\rho\,\dd m, \mu^n)$. 

%Unfortunately, this heuristic turns out to fail in low dimension. Indeed the the asymptotics of \eqref{eq:transintr} are given by
%\begin{equation}\label{eq:unifmatchingrates}
%\mathbb{E}[W_2^2 (\mu^n, \nu^n)]\sim
%n^{-\frac2d}\cdot\begin{cases}
%n & \text{if $d=1$,} \\
%{\ln n} & \text{if $d=2$}, \\
%1 & \text{if $d\ge 3$.}
%\end{cases}
%\end{equation}
%As shown in \eqref{eq:unifmatchingrates}, fluctuations from the heuristic are effective when $d<3$ with a logarithmic correction when $d=2$. In the following of this work we will focus on the ``critical'' case $d=2$. The logarithmic correction appearing in the two dimensional case has been firstly understood in \cite{AKT84}, but it was in recent years that the convergence of the rescaled cost has been shown for $d=2$. Based on an ansatz from statistical physics which arose in \cite{caracciolo2014scaling}, in \cite{ambrosio2019pde} it has been formally shown that
%The main novel of this approach is that is based on partial differential equation (PDE) techniques which allow for explicit description of the cost and the approximating transportation map, see \cite{ambrosio2019finer, ambrosio2019optimal, ambrosio2019pde} for a justification on microscopic scales and \cite{GHO} for a justification on large scales. 


%\medskip
%When $(X_i)_{i=1}^n, (Y_j)_{j=1}^n$ are i.~i.~d. points on $\mathcal{M} \subset \mathbb{R}^2$ with common law $\rho \, \dd x$, a similar approach can be applied to study the convergence of the rescaled cost \eqref{eq:transintr}. Readapting the ansatz of \cite{caracciolo2014scaling} to the case of non-constant densities, in \cite[Conjecture 2]{benedetto2020euclidean}, assuming that $\rho \in C^{0,\alpha}(\mathcal{M})$ is bounded away from $0$, it has been conjectured that when $\mathcal{M}$ is a bounded connected domain
%\[
%\lim_{n\rightarrow\infty} \frac n{\ln n} \mathbb{E} [W_2^2 (\mu^n, \nu^n)] = \frac{|\mathcal{M}|}{2\pi}.
%\]
%Combining PDE techniques together with a subadditivity argument in the spirit of \cite{Barthe2013, Der}, the latter has been rigorously proven in \cite{ambrosio2021quadratic}. As observed in \cite{benedetto2020euclidean} the condition $\rho >0$ intuitively amounts to connectedness of the support of $\rho$ and if violated might change the asymptotics of the transportation cost. For instance, if $\rho$ vanishes somewhere in its support and it is constant in two squares of positive distance then the cost is at least of order $O(1/\sqrt n) \gg O(\ln n / n)$, see \cite[Remark 2]{benedetto2020euclidean} for further details.

%
%\medskip
%

%The aim of this work is to justify the linearization ansatz, that is the second item of \eqref{eq:approxcost}. To the best of our knowledge, there are only few results on the asymptotic behavior of the transport map and they are so far limited to the case of uniformly distributed $i.~i.~d.$ samples in the study of the semi-discrete matching problem (that is couplings between $\mu^n$ and $\dd\m$), see the work of Ambrosio, Glaudo and Trevisan in \cite{ambrosio2019optimal}. In connection with this work, quantitative estimates on the optimal map for the matching between the Lebesgue measure and Poisson clouds have been obtained by Goldman, Huesmann and Otto in \cite{GHO1} and Goldman and Huesmann in \cite{GH}. 

%
%\medskip
%

%Our extension in this paper is threefold: First, we look at more general distribution of points and we consider the case of general densities $\rho$ satisfying \eqref{Ellipticity}. Second, we do not assume independence and we consider samples which may possess correlations. Third, we do not restrict the analysis to the semi-discrete matching problem and we also investigate the ansatz for the full bipartite matching problem \eqref{eq:transintr}. We prove that, after regularizing the right-hand side of \eqref{eq:ApproxPDE} via the heat-kernel at a time $\sim\frac{1}{n}$ and for weakly correlated points (see Assumptions \ref{Assumptions}), the second item of \eqref{eq:approxcost} holds true with a quantitative convergence rate.

%
 %\medskip
%

%The paper is organized as follows: The remaining part of Section \ref{introsec} is devoted to state our main results and possible extensions. Section \ref{StrategyOfProofSection} is devoted to describe the strategy of the proof and list the auxiliary results needed in the proof of the main theorems. The results are then proven in Section \ref{ProofSection}.


%

%The aim of this work is to provide quantitative bounds on the convergence of the approximating map towards the optimal one in the case of non uniformly distributed points. Our primary motivation is to justify the ansatz of \cite{benedetto2020euclidean} on the level of the transport map. To the best of our knowledge (except for the $1$-dimensional case) there are few results on the behavior of the transport map. For example, in \cite{ambrosio2019optimal} quantitative upper bounds for the convergence of the approximating map have been settled in the uniform semi-discrete matching (Lebesgue to empirical measure). In connection with the behavior of the optimal transport map in the Lebesgue to Poisson problem on large scales, quantitative bounds have been obtained in \cite{GHO1, GHO}.  

%
%\medskip
%
%\nc{I would not give the theorem now but more explain in words and where it comes from through the linearization.}
%\textbf{Main result. }Let us assume that the common distribution $\rho$ satisfies the following ellipticity condition, i.~e. there exists some $\lambda,\Lambda >0$ such that
%
%\begin{equation}\label{Ellipticity}
%
%\lambda\leq\rho\leq\Lambda.
%
%\end{equation}
%We will also assume that the density $\rho$ belongs to the fractional Sobolev space $\mathrm{H}^\varepsilon$, for any $\varepsilon>0$,
%
%Note that by the Sobolev embedding Theorem $C^{0,\alpha}(\mathcal{M})\subset \mathrm{H}^\varepsilon$. Moreover, let us introduce the space of Sobolev functions with $\mathrm{L}^2$ gradient and null mean
%\begin{equation}\label{eq:H1}
%\dot{\mathrm{H}}^1:=\Big\{f: \mathcal{M}\rightarrow \mathbb{R}\ \Big\vert\ \int_{\mathcal{M}}\vert \nabla f\vert^2<\infty\quad\text{and}\quad \int_{\mathcal{M}} f=0\Big\}.
%\end{equation}

%\medskip
%In the setting of the $2$-dimensional semidiscrete matching our main result can be stated as follows.
%
%\begin{theorem}\label{thm:sloppyintro}
%
%Let $h^{n,t}\in \dot{\mathrm{H}}^1$ be the weak solution\footnote{where we impose additional Neumann boundary conditions in the case $\mathcal{M}=[0,1]^2$} of
%
%\[
%-\nabla\cdot \rho\nabla h^{n,t} =\mu^{n,t}-\nu^{m,t}.
%\]
%Then there exists a constant $C$ depending on $\rho$ and $\mathcal{M}$ and $\gamma>0$ for which given $t=\frac{\log^{\gamma}(n)}{n}$
%
%\begin{equation}\label{DefApproxPlan}
%
%\gamma^{n,t}:=\big(\mathrm{Id},\exp(\nabla h^{n,t})\big)_\#\mu^{n,t},
%
%\end{equation}
%
%we have: 

%
%\begin{equation}\label{eq:approxtranspomapintro}
%
%\int_\mathcal{M} \mathrm{d}^2\big(T^n,\mathrm{exp}(\nabla f^{n,t})\big)\dd \mathrm{m}\leq \mathcal{C}_n\frac{\log n}{n}\sqrt{\frac{\log\log(n)}{\log(n)}}\quad\text{with }\sup_{n\geq 1}\mathbb{E}[\exp(\tfrac{1}{C}\mathcal{C}_n))]\leq 2.
%
%\end{equation}
%
%\end{theorem}

%We point out that \eqref{eq:approxtranspomapintro} matches the same upper bound obtained in \cite[Theorem 1.3]{ambrosio2019optimal}. \fm{I don't have access to this article do you? I cited the Theorem based on the number I see on the ArXiv version, we need to check also if the number of the quantitative stability result has not changed} Moreover, under suitable $\alpha$ and $\beta$-mixing condition (see \eqref{AlphaMixing} and \eqref{BetaMixing} for a definition) Theorem \ref{thm:sloppyintro} can be extended to the case of \emph{weakly correlated points} (see Theorem \ref{th1} for precise statement). In particular, the case of independently distributed points can be recovered from Theorem \ref{th1} when the $\alpha$-mixing coefficient is zero, which amounts to letting $\eta$ tend to infinity in \eqref{DecayAlphaMixing}. \fm{Is there a better way to say so?} The case of correlated points has been subject of various investigation. In \cite[Section 7]{Fournier2015} general upper bounds for the transportation cost are obtained under mild $\rho$-mixing assumption. More recently, following the same approach of \cite{ambrosio2019pde}, analogous upper bounds for the matching of weakly correlated points have been shown in \cite{borda2021empirical} under the assumption that the sum of the $\beta$-mixing has at most linear growth. More interestingly, Theorem \ref{th1} can be extended to the case of unbalanced matching, i.~e. when the size of the two point clouds is not the same (see Theorem \ref{MainResultMatchingClouds} for a precise statement), providing a novel estimate in terms of convergence of approximating plans and justifying the ansatz of \cite{caracciolo2014scaling} also at the level of couplings of the empirical measures. Finally in Section \ref{Examples} and Appendix \ref{ClassMarkovAppendix} we extend our result to a suitable class of Markov processes, which provides a sequence of non identically distributed random points. 




%
%\medskip
%

%The last equation, at least in distributional sense, makes perfect sense and suggests that the optimal transport map is well approximated by $\text{Id}+\nabla h^n$, where $h^n$ solves \eqref{eq:ApproxPDE}. Thus, it yields a natural estimate of the transportation cost in terms of a negative Sobolev norm, i.\,e. 
%\begin{equation}\label{eq:approxcost}
%W_2^2 (\rho \, \dd m, \mu^n) \approx \int_{\mathbb{T}^2} |\nabla h^n|^2.
%\end{equation}

%\medskip
%Without any further assumption on $\rho$ and $\mu^n$ this conjecture turns out to be false. Indeed, in the case of $\rho$ being a constant density \eqref{eq:ApproxPDE} becomes a Poisson equation and it is not hard to show that the right hand side of \eqref{eq:approxcost} diverges. This issue can be fixed via a regularization argument. More precisely, let $\mu^{n,t} := \pp_t \mu^n, \rho_t := \pp_t \rho$ be the evolution at a certain small time $t\ll 1$ of the empirical measure through the heat semigroup. If we argue with $\rho\,\dd m, \mu^n$ replaced by $\rho_t\,\dd m, \mu^{n,t}$ then the ansatz \eqref{eq:approxcost} is sound and Dirichlet energy of the associated solution $h^{n,t}$ well approximates $W_2^2 (\rho\,\dd m, \mu^n)$. 
%
%\medskip
%
%When $\rho$ is a non constant density there are two sources of \emph{``roughness''} in \eqref{eq:ApproxPDE}. As in the uniform case, one is given by the data and an additional one is given by the operator $- \nabla \cdot \rho \nabla$. The latter provides no regularity for the solution of \eqref{eq:ApproxPDE}, even when the data terms are assumed to be smooth. Indeed, the usual Schauder's theory does not apply for $\rho \in \mathrm{H}^\varepsilon$.  As for the data, this issue can be overcome by mean of a regularization of the operator. Substituting $-\nabla \cdot \rho\nabla$ with $- \nabla \cdot \rho_\delta\nabla$, and solving the PDE
%\[
%-\nabla \cdot \rho_\delta \nabla h^{n,t}_\delta = \rho_t - \mu^{n,t}
%\] 
%yields a smooth solution which is close to the one of \eqref{eq:ApproxPDE} and well approximates the optimal transport map. However, this additional approximation produces another regularization error that, as discussed in further details in Section \ref{StrategyOfProofSection}, is of higher order.  
%
%\medskip
%


%\item \emph{Fluctuations estimate. }It has been showed in \cite{ambrosio2021quadratic} that under independence and the additional regularity assumption $\rho \in C^{0,\alpha}(\mathcal{M})$ the rescaled cost converges to a constant, that is 
%\[
%\sqrt{\frac n{\ln n}} \mathbb{E} [W_2(\mu^n, \leb)] \rightarrow \frac\pi2|\mathcal{M}|.
%\]
%An ultimate goal in the understanding of the convergence of the matching problem would be studying also the convergence of the fluctuations, i.\,e. a result of the type
%\[
%\sqrt{\frac n{\ln n}} |W_2(\mu^n, \leb) - \mathbb{E} [W_2(\mu^n, \leb)]|^2 \rightarrow \xi,
%\]
%where $\xi$ is a Gaussian free field. This will allow to a result of the type (for the random variable)
%\[
%\sqrt{\frac n{\ln n}} W_2(\mu^n, \leb)  \rightarrow \frac\pi2|\mathcal{M}| + \xi. 
%\]

%
%\medskip
%





%Our main result state a quantitative estimate on the error, in $\LL^{2}(\Omega)$, between the transport map $T^n$ and the map $\exp(\nabla f^{n,t})$ \nc{state the equation in the introduction} in $2$-dimensional compact Riemmanian manifold without boundary (or the cube $[0,1]^2$). To state our main result, we need to define two mixing quantities : for a sequence $(X_k)_{k\in\mathbb{N}^*}$ of random points distributed in $\mathcal{M}$, we consider
%
%\begin{itemize}
%
%\item[(i)] the $\alpha$-mixing coefficient defined as, for any $\ell\geq 1$
%

%
%which measures the correlation between points.
%
%\item[(ii)]the $\beta$-mixing coefficient defined as
%
%\begin{equation}\label{BetaMixing}
%
%\beta_{ij}:=\sup_{\vert F\vert\leq 1}\bigg\vert\int_{\mathcal{M}\times\mathcal{M}}\dd (\mathbb{P}_{(X_i,X_j)}-\mathbb{P}_{X_i}\otimes\mathbb{P}_{X_j})F\bigg\vert\quad \text{for any\quad $i,j\geq 1$},
%
%\end{equation}
%
%where the supremum runs over all measurable functions $F : \mathcal{M}\times\mathcal{M}\rightarrow [-1,1]$ and $\mathbb{P}_X$ denotes the distribution of a random variable $X$.

%
%\end{itemize}
%
%\medskip
%

%\nc{Put in the introduction while explaining the main result the definition of the empirical measures.In general, define the objects before.}
%
%Let $(\mathcal{M},\mathrm{d})$ be either a $2$-dimensional compact Riemannian manifold without boundary or the cube $[0,1]^2$, whose volume measure $\mathrm{m}$ is a probability. Given $\rho\in \mathrm{L}^{\infty}$ such that for 




%
%let $(X_k,Y_k)_{k\in\mathbb{N}^*}$ be a sequence of random points in $\mathcal{M}$ with common distribution $\rho\,\dd \m$. We define the empirical measures, for any $n\in \mathbb{N}$ and $\phi(n)\geq n$
%
%\begin{equation}\label{DefEmpiricalMeasures}
%
%\mu^{n}=\frac{1}{n}\sum_{k=1}^n\delta_{X_k}\quad \text{and}\quad \nu^n:=\frac{1}{\phi(n)}\sum_{k=1}^{\phi(n)}\delta_{Y_k}, 
%
%\end{equation}
%
%We define
%

%

%
\medskip
\subsection{Formulation of the main results}
%\subsection{Formulation of the main results}\label{MainResultSection}
%
For the remainder of the paper $\mathcal{M}$ denotes a $2$-dimensional connected and compact Riemannian manifold without boundary (or the square $[0,1]^2$) endowed with the Riemannian distance $\dd$. For $t>0$ we denote by $p_t$ the fundamental solution of the heat operator $\partial_t-\Delta$ on $\mathcal{M}$, where $\Delta$ denotes the Beltrami-Laplace operator. We define the heat semigroup $(\text{P}_t)_{t>0}$ via its action on probability measures $\mu \in \mathcal{P}(\mathcal{M})$ and square integrable functions $f\in\LL^2(\mathcal{M})$
\[
\text{P}_t \mu := \int_\mathcal{M} p_t(\cdot,y)\dd \mu(y) \quad\text{and}\quad \text{P}_t f:= \int_{\mathcal{M}} p_t(\cdot,y) f(y)\dd \m(y).
\]
%
%where $\mathcal{P}(\mathcal{M})$, resp. $\LL^2$, denotes the set of probability measures on $\mathcal{M}$, resp. square integrable function w.r.t $\dd\m$. 

\medskip
We first introduce the class of correlated point clouds that we consider for studying the matching problem \eqref{eq:transintr}. This class concerns point clouds $\{X_i\}_i$ for which the correlations between points decay at an exponential rate, where the correlations are measured in terms of the $\alpha$-mixing coefficient given by, for any $\ell\geq 1$
%
\begin{equation}\label{AlphaMixing}
%
\alpha_\ell:=\sup_{k\geq 1}\sup\bigg\{\frac{\text{cov}(f,g)}{\|f\|_{\LL^{\infty}}\|g\|_{\LL^{\infty}}},\quad f\in\LL^{\infty}(\sigma\{X_j,j\leq k\})\quad \text{and}\quad g\in\LL^{\infty}(\sigma\{X_j,j\geq \ell+k\})\bigg\},
%
\end{equation}
%
and the $\beta$-mixing coefficient given by\footnote{We denote by $\mathbb{P}_X$ the law of a random variable $X$.}
%
\begin{equation}\label{BetaMixing}
%
\beta_{ij}:=\sup_{\vert F\vert\leq 1}\bigg\vert\int_{\mathcal{M}\times\mathcal{M}}F\,\dd (\mathbb{P}_{(X_i,X_j)}-\mathbb{P}_{X_i}\otimes\mathbb{P}_{X_j})\bigg\vert\quad \text{for any $i,j\geq 1$}.
%
\end{equation}
%
\begin{assumption}[Correlated point clouds]\label{Assumptions}
%
We consider point clouds $\{X_i\}_{i}\subset \mathcal{M}$ which are identically distributed according to $\rho\,\dd\m$ where $\rho$ satisfies \eqref{Ellipticity}. We further assume decay of the correlations in the form of 
%
\begin{equation}\label{DecatBetaMixing}
%
\sup_{n\geq 1}\frac{1}{n}\sum_{1\leq i<j\leq n}\beta_{ij}<\infty,
%
\end{equation}
%
and there exist $a,b>0$ and $\eta\in (0,\infty]$ such that 
%
\begin{equation}\label{DecayAlphaMixing}
%
\alpha_\ell\leq a\exp\big(-b\ell^{\eta}\big)\quad\text{for any $\ell\geq 1$.}
%
\end{equation}
%
\end{assumption}
%
Assumption \ref{Assumptions} is made to ensure good concentration properties of the point clouds. On the one hand, under \eqref{DecatBetaMixing}, the cost $W^2_2(\mu^n,\nu^n)$ behaves as in the i.i.d. case \eqref{eq:unifmatchingrates} (cf. \cite[Theorem 2]{borda2021empirical} and Appendix \ref{sec:matchcost}). On the other hand, the sub-exponential decay \eqref{DecayAlphaMixing} of the $\alpha$-mixing coefficient ensures sub-exponential concentration properties (cf. \cite[Theorem 1]{merlevede2011bernstein} and Proposition \ref{BersteinCorrelated}), which is necessary to run our argument. We refer the reader to Section \ref{StrategyOfProofSection} for further technical details.

%
\medskip
%

Our first main result concerns the approximation of transport plans coupling $\{\mu^n\}_n$ and $\{\nu^m\}_m$ defined in \eqref{eq:empmeas}. We justify the formal linearization of the Monge-Amp\`ere equation achieved in Section \ref{sec:lin} in an annealed quantitative way (i.e. in expectation): We show that, a suitable regularization of, the plan $\big(\text{Id},\exp(\nabla h^n)\big)_\#\mu^n$, with $h^n$ defined in \eqref{eq:ApproxPDE}, provides a good approximation when measuring the error with respect to the $W_2$-Wasserstein distance in the product space $\mathcal{M}\times\mathcal{M}$ endowed with the metric 
%
\begin{equation}\label{DistanceTransportPlan}
%
\delta^2 \big((x,y),(z,w)\big):=\dd^2(x,z)+\dd^2(y,w).
%
\end{equation}
%
The density $\rho$ will need further regularity in form of fractional Sobolev spaces defined as, for some $\varepsilon>0$
%
\begin{equation}\label{DefHepsilon}
%
\mathrm{H}^{\varepsilon}:=\Big\{f: \mathcal{M}\rightarrow\mathbb{R}\ \Big\vert \|f\|^2_{\mathrm{H}^\varepsilon}:= \sum_{k\geq 1}\lambda^{2\varepsilon}_k\vert \hat{f}(k)\vert^2<\infty\Big\},
%
\end{equation}
%
where $\{\lambda_k,\phi_k\}_{k}$ denote the eigenvalues and eigenvectors of $-\Delta$ on $\mathcal{M}$ and $\hat{f}(k)=\int_{\mathcal{M}}f\,\phi_k$ denotes the Fourier modes of $f$.
%
Finally, we denote by $\dot{\mathrm{H}}^1$ the $\LL^2$-based Sobolev space 
%
\begin{equation}\label{eq:H1}
%
\dot{\mathrm{H}}^1:=\Big\{f: \mathcal{M}\rightarrow \mathbb{R}\ \Big\vert\ \int_{\mathcal{M}}\vert \nabla f\vert^2<\infty\quad\text{with } \int_{\mathcal{M}} f\dd\m=0\Big\}.
%
\end{equation}
%
\begin{theorem}[Approximation of the transport plan]\label{MainResultMatchingClouds}
%
Let $\rho\in \mathrm{H}^{\varepsilon}$ for some $\varepsilon>0$ satisfying \eqref{Ellipticity} and $\{\mu^n\}_n$ and $\{\nu^{m}\}_m$ be defined in \eqref{eq:empmeas} (for $m=m(n)$ with some given increasing map $m:\mathbb{N}\rightarrow\mathbb{N}$) with point clouds satisfying Assumption \ref{Assumptions} and such that there exists $q\in [1,\infty)$ for which $\frac{m(n)}{n}\underset{n\uparrow\infty}{\rightarrow} q$. We consider\footnote{where we impose additional Neumann boundary conditions in the case $\mathcal{M}=[0,1]^2$} $h^{n,t}\in \dot{\mathrm{H}}^1$ the weak solution of
%
\begin{equation}\label{eq:f^nmt}
\nabla\cdot \rho\nabla h^{n,t} =\mu^{n,t}-\nu^{m,t},
\end{equation}
%
for any $t\in(0,1)$ with $\mu^{n,t}:=\pp_t\mu^n$ and $\nu^{m,t}:=\pp_t\nu^m$.

%
\medskip
%

There exist an exponent $\kappa>0$, a deterministic constant $C$ and a random variable $\mathcal{C}_n$ both depending on $\lambda, \Lambda$ and $\mathcal{M}$ for which given $t=\frac{\log^{\kappa}(n)}{n}$ and
%There exists an exponent $q>1$ such that for any $\gamma_1>0$, there exist $\gamma_2>0$, a constant $C$ depending on $\lambda,\Lambda$ and $\mathcal{M}$ \nc{make precise the dependence} for which given $t=\frac{\log^{\gamma_2}(n)}{n}$, $\delta=\frac{1}{\log^{\gamma_1}(n)}$ and
%
\begin{equation}\label{DefApproxPlan}
%
\gamma^{n,t}:=\big(\mathrm{Id},\exp(\nabla h^{n,t})\big)_\#\mu^{n,t},
%
\end{equation}
%
it holds
%
\begin{equation}\label{ApproximationTransportPlan}
\inf_{\pi}W^2_2(\pi,\gamma^{n,t})\leq \mathcal{C}_n\frac{\log(n)}{n}\sqrt{\frac{\log\log(n)}{\log(n)}}\quad\text{with \,$\sup_{n\geq 1}\mathbb{E}[\tfrac{1}{C}\mathcal{C}_n]\leq 1$,}
\end{equation}
%
where the $\inf$ runs over all optimal transport plans $\pi$ between $\mu^n$ and $\nu^m$.

%
\medskip
%

Furthermore, if \eqref{DecayAlphaMixing} holds with $\eta>2$, the assumption \eqref{DecatBetaMixing} can be dropped and it holds
%
\begin{equation}\label{ApproximationTransportPlanBis}
%
\inf_{\pi}W^2_2(\pi,\gamma^{n,t})\leq \mathcal{C}_n\frac{\log(n)}{n}\sqrt{\frac{\log\log(n)}{\log^{1-\frac{2}{\eta}}(n)}}\quad\text{with }\sup_{n\geq 1}\mathbb{E}[\exp(\tfrac{1}{C}\mathcal{C}^{\frac{1}{2}}_n))]\leq 2,
%
\end{equation}
where the $\inf$ runs over all optimal transport plans $\pi$ between $\mu^n$ and $\nu^m$.
%
%$$a_n:=\log^{\frac{1}{\eta}}(n)\|\rho_\delta-\rho\|^2_{\LL^{q}}+\|\rho_{t+\frac{1}{n}}-\rho_{\frac{1}{n}}\|^2_{\LL^{q}}+\log^{\gamma_1-1}(n)\|\rho_{t+\frac{1}{n}}-\rho_{\frac{1}{n}}\|_{\LL^1}+\sqrt{\tfrac{\log\log(n)}{\log^{1-\frac{1}{\eta}}(n)}}.$$
%
\end{theorem}
%
Our second main result concerns the particular case of the semi-discrete matching problem, i.e. optimal coupling between the common law $\rho\,\dd \m$ and $\{\mu^n\}_n$. We know from McCann's theorem \cite{McCann2001} that there exists a unique optimal transport map $T^n$, that is the optimal transport plan $\pi^n$ can be written as 
%
$$\pi^n=\big(\text{Id},T^n\big)_\#\rho\,\dd m.$$
%
We show that $T^n$ can be approximated in an annealed quantitative way in $\LL^2$ by (a suitable regularized version of) the solution of \eqref{eq:ApproxPDE} with $\nu^m$ replaced by $\rho$.
%
%\begin{corollary}[$W_2$-approximation of the transport plan]\label{ApproximationTransportPlan:Coro}
%
%Under the notations of Theorem \ref{MainResultMatchingClouds}, we define 
%
%$$\gamma^{n,t}:=\big(\mathrm{Id},\exp(\nabla h^{n,t})\big)_\#\mu^{n,t}.$$
%
%It holds for the distance \eqref{DistanceTransportPlan}:
%
%\begin{itemize}
%
%\item If $(X_k,Y_k)_{k\in\mathbb{N}^*}$ are i. i. d, 
%
%\begin{equation}\label{ApproximationTransportPlan2}
%
%\inf_{\pi}W^2_2(\pi,\gamma^{n,t})\leq \mathcal{C}_n\frac{\log(n)}{n}a_n,
%
%\end{equation}
%
%where the $\inf$ is taken over all optimal transport plan from $\mu^n$ to $\nu^n$.

%
%\medskip
%

%\item  If $\rho\in\mathrm{C}^{0}(\mathcal{M})$ and $(X_k,Y_k)_{k\in\mathbb{N}^*}$ are weakly correlated, in the sense of \eqref{Eq:AssumptionsCorrelatedCase}, \eqref{ApproximationTransportPlan2} holds with $q=\infty$ and 
%
%$$\sup_{n\geq 1}\mathbb{E}[\tfrac{1}{C}\mathcal{C}_n]\leq 1.$$
%
%
%\end{itemize}



%
%\end{corollary}
%
\begin{theorem}\label{th1}
%
Let $\rho\in \mathrm{H}^{\varepsilon}$ for some $\varepsilon>0$ satisfying \eqref{Ellipticity} and $\{\mu^{n}\}_n$ be defined in \eqref{eq:empmeas} with a point cloud satisfying Assumption \ref{Assumptions}. We consider $f^{n,t}\in {\rm \dot{H}^1}$ the weak solution of
%
\begin{equation}\label{eq:f^ntBis}
\nabla\cdot \rho\nabla f^{n,t} =\mu^{n,t}-\rho_t,
\end{equation}
%
where, for all $t\in (0,1)$, we recall that $\mu^{n,t}=\pp_t\mu^n$ and $\rho_t=\pp_t\rho$. Finally, we denote by $T^n$ the optimal transport map from $\rho\,\dd \m$ to $\mu^n$. 
%
\medskip
%

There exist an exponent $\kappa>0$, a deterministic constant $C$ and a random variable $\mathcal{C}_n$ both depending on $\lambda,\Lambda$ and $\mathcal{M}$ for which given $t=\frac{\log^{\kappa}(n)}{n}$, it holds
%
\begin{equation}\label{ApproximationTransportMap}
\int_\mathcal{M} \mathrm{d}^2\big(T^n,\mathrm{exp}(\nabla f^{n,t})\big)\dd \mathrm{m}\leq \mathcal{C}_n\frac{\log(n)}{n}\sqrt{\frac{\log\log(n)}{\log(n)}}\quad\text{with \,$\sup_{n\geq 1}\mathbb{E}[\tfrac{1}{C}\mathcal{C}_n]\leq 1$.}
\end{equation}
%
%\begin{equation}\label{Eq78}
%\begin{aligned}
%a_n:=(\log^{\frac{1}{\eta}}(n)+1)(\|\rho_\delta-\rho\|^2_{\LL^{q}}+\|\rho_{t+\frac{1}{n}}-\rho_{\frac{1}{n}}\|^2_{\LL^{q}})+\sqrt{\log^{\gamma_2-1}(n)\|\rho_{t+\frac{1}{n}}-\rho_{\frac{1}{n}}\|_{\LL^1}}+\sqrt{\tfrac{\log\log(n)}{\log(n)}}.
%\end{aligned}
%\end{equation}
%

%
\medskip
%

Furthermore, if \eqref{DecayAlphaMixing} holds with $\eta>2$, the assumption \eqref{DecatBetaMixing} can be dropped and it holds
%
\begin{equation}\label{ApproximationTransportMapBis}
%
\int_\mathcal{M} \mathrm{d}^2\big(T^n,\mathrm{exp}(\nabla f^{n,t})\big)\dd \mathrm{m}\leq \mathcal{C}_n\frac{\log(n)}{n}\sqrt{\frac{\log\log(n)}{\log^{1-\frac{2}{\eta}}(n)}}\quad\text{with }\sup_{n\geq 1}\mathbb{E}[\exp(\tfrac{1}{C}\mathcal{C}_n^\frac12))]\leq 2.
%
\end{equation}
%
%$$a_n:=\log^{\frac{1}{\eta}}(n)\|\rho_\delta-\rho\|^2_{\LL^{q}}+\|\rho_{t+\frac{1}{n}}-\rho_{\frac{1}{n}}\|^2_{\LL^{q}}+\log^{\gamma_1-1}(n)\|\rho_{t+\frac{1}{n}}-\rho_{\frac{1}{n}}\|_{\LL^1}+\sqrt{\tfrac{\log\log(n)}{\log^{1-\frac{1}{\eta}}(n)}}.$$
%
\end{theorem}
%
We finally mention that in the case where the eigenfunctions $\{\phi_k\}_{k}$ admit a uniform bound, the conclusions \eqref{ApproximationTransportPlanBis} and \eqref{ApproximationTransportMapBis} can be improved. We comment on the proof at the end of Sub-section \ref{ProofTransport}.
%
\begin{remark}\label{FlatGeometry}
%
Let $\{\mu^n\}_n$ and $\{\nu^m\}_m$ be as in Theorem \ref{MainResultMatchingClouds}. We assume that the family of eigenfunctions $\{\phi_k\}_k$ satisfies the uniform bound
%
\begin{equation}\label{FlatenessAss}
%
\sup_{k\geq 1}\|\phi_k\|_{\LL^\infty}<\infty.
%
\end{equation}
%
Then, \eqref{ApproximationTransportPlanBis} and \eqref{ApproximationTransportMapBis} hold true for $\eta>1$ with a convergence rate $\frac{\log(n)}{n}\sqrt{\frac{\log\log(n)}{\log^{1-\frac{1}{\eta}}(n)}}$ and the same stochastic integrability. Note that \eqref{FlatenessAss} typically holds when the geometry of $\mathcal{M}$ is flat, see \cite{toth2002riemannian}.
%
\end{remark}




%
\medskip
%

Theorem \ref{MainResultMatchingClouds} and Theorem \ref{th1} are not restricted to the case of identically distributed point clouds and we present in the next section and in Appendix \ref{ClassMarkovAppendix} a possible extension, using the same techniques, to a class of sub-geometrically ergodic discrete-time Markov chains.

%\nc{Example for the correlated case, $\rho$ is assumed to be continuous. }For the correlated case, that is for instance the case for absolutely continuous densities 
%
%$$\vert\rho (x)-\rho(y)\vert\leq \omega(\mathrm{d}(x,y)),$$
%
%where we find 
%
%$$\|\rho_t-\rho\|_{\LL^1(\mathcal{M})}\lesssim \int_{0}^{\infty}\dd r\, e^{-r^2}\omega(r\sqrt{t}),$$
%
%which gives \eqref{Eq81} as soon as
%
%$$\omega(r)\leq \frac{1}{\vert\log^{\gamma_1}(r)\vert}.$$
%
%For an example with non-continuous $\rho$, take $\alpha\in (1,2)$ and $\rho\in \WW^{1,\alpha}(\mathcal{M})$. We then have \eqref{Eq81} using the argument with Lusin-Lipschitz. On deduce easily that
%
%$$\|\rho_t-\rho\|_{\LL^{1}(\mathcal{M})}\lesssim t^{\frac{1}{2}}.$$
%
%\nc{What, quite minimal assumption on $\rho$ avoiding continuity, can we do to assume a quantitative estimate on this ?}
%
%A quantitative result would require to add more regularity on the density $\rho$. As an example, let assume that $\rho\in\mathrm{W}^{1,\alpha}(\mathcal{M})$ for some $\alpha>q$ and $q>2$. In that case, we claim that for any regularized parameter $\tau\in (0,1)$, we have
%
%\begin{equation}\label{Eq81}
%\|\rho_\tau-\rho\|_{\LL^{q}(\mathcal{M})}\lesssim \tau^{\frac{1}{\alpha'}-\frac{1}{2}}.
%\end{equation}
%
%The claim \eqref{Eq81} can be seen using the Lusin-Lipschitz property of Sobolev functions (see \cite{liu1977luzin}) : there exists $g\in\LL^{\alpha}(\mathcal{M})$ such that $\|g\|_{\LL^{\alpha}(\mathcal{M})}\lesssim \|\rho\|_{\WW^{1,\alpha}(\mathcal{M})}$ such that 
%
%$$\vert\rho(x)-\rho(y)\vert\leq (g(x)+g(y))\vert x-y\vert\quad \text{for almost-all $x,y\in \mathcal{M}$}.$$
%
%so that 
%
%\begin{align*}
%\|\rho_{\tau}-\rho\|^2_{\LL^{q}(\mathcal{M})}&= \bigg(\int_{\mathcal{M}}\dd x\Big(\int_{\mathcal{M}}\dd y\,p_\tau(x,y)(\rho(y)-\rho(x))\Big)^{q}\bigg)^{\frac{2}{q}}\\
%&\leq \bigg(\int_{\mathcal{M}}\dd x\Big(\int_{\mathcal{M}}\dd y\,p_\tau(x,y)(g(x)+g(y))\vert x-y\vert\Big)^{q}\bigg)^{\frac{2}{q}}.
%\end{align*}
%
%Then, we have directly using the Green's function estimate \eqref{Eq28}
%
%\begin{align*}
%\bigg(\int_{\mathcal{M}}\dd x\Big(\int_{\mathcal{M}}\dd y\,p_\tau(x,y)g(x)\vert x-y\vert\Big)^{q}\bigg)^{\frac{1}{q}}= \bigg(\int_{\mathcal{M}}\dd x\, g^q(x)\Big(\int_{\mathcal{M}}\dd y\,p_\tau(x,y)\vert x-y\vert\Big)^{q}\bigg)^{\frac{1}{q}}\lesssim \tau,
%\end{align*}
%
%whereas using Hölder's inequality and \eqref{Eq28}
%
%\begin{align*}
%\bigg(\int_{\mathcal{M}}\dd x\Big(\int_{\mathcal{M}}\dd y\,p_\tau(x,y)g(y)\vert x-y\vert\Big)^{q}\bigg)^{\frac{1}{q}}\lesssim \bigg(\int_{\mathcal{M}}\dd x\Big(\int_{\mathcal{M}}\dd y\, p^{\alpha'}_\tau(x,y)\vert x-y\vert^{\alpha'}\Big)^{\frac{q}{\alpha'}}\bigg)^{\frac{1}{q}}\lesssim \tau^{\frac{1}{\alpha'}-\frac{1}{2}}.
%\end{align*}
%
%Estimate \eqref{Eq81} gives
%
%$$a_n\lesssim \delta^{\frac{2}{\alpha'}-1}+\log^{\gamma_1-1}(n)t^{\frac{1}{\alpha'}-\frac{1}{2}}.$$
%
%
\subsection{Extension to a class of sub-geometrically ergodic Markov chains}\label{Examples}
%
%We display here some examples of random processes which fits with our assumption. We also discuss about some extensions to Markov chains. %\nc{I let you do this part}. 
%\nc{I put here an example of Markov chain}
%We provide below a class of Markov chains which satisfies an analogue of Theorem \ref{MainResultMatchingClouds}. 
We first recall some basic facts on discrete-time Markov chains on $\mathcal{M}$. Such a Markov process is described by its initial distribution $\mu_0 \in \mathcal{P}(\mathcal{M})$ and its transition kernel $K$, that is a measurable map from $\mathcal{M}$ to the space of probability measures $\mathcal{P}(\mathcal{M})$. We recall that $K$ acts on $\mathcal{P}(\mathcal{M})$ in form of 
%
\begin{equation}\label{eq:actionkernel}
%
(K\mu)(A) = \int_{\mathcal{M}} K(\cdot,A)\,\dd\mu\quad\text{for every Borel set $A\subset \mathcal{M}$ and $\mu\in \mathcal{P}(\mathcal{M})$,}
%
\end{equation}
%
and likewise on bounded measurable function $\psi$ in form of
%
\begin{equation}\label{eq:actionkernelBis}
%
(K\psi)(x) = \int_{\mathcal{M}} \psi(y) \,K(x, \dd y)\quad\text{for any $x\in\mathcal{M}$.}
%
\end{equation}
%
%Finally we say that a probability measure $\pi$ is an invariant measure for $K$ if it is invariant under its action, that is $K \mu = \mu$. 
%The main difference from the setting in Assumptions \ref{Assumptions} resides in the fact that the elements of a Markov chain are not identically distributed. Indeed, for a state $X_n$ of a Markov chain $(X_n)_{n \ge 0}$ it holds $X_n \sim K^n \lambda$, with $X_1 \sim \lambda$, where we recall $\lambda$ being the initial distribution. Note that by the definition of the action of the kernel $K$ \eqref{eq:actionkernel} it holds for every signed measure $\mu$ and any Borel measurable function $\phi$
Given an initial distribution $\mu_0\in \mathcal{P}(\mathcal{M})$, we recall that the law of a Markov chain $\{X_n\}_{n\geq 0}$ can be expressed with the help of the transition kernel, namely 
%
\begin{equation}\label{LawOfTheChain}
%
X_n\sim K^n\mu_0\quad\text{for any $n\geq 0$,}
%
\end{equation}
%
where $K^n$ stands for the $n^{\mathrm{th}}$-iteration of the kernel $K$. 

%
\medskip
%

We now introduce the class of discrete-time Markov chains we consider.
%
\begin{assumption}\label{ClassMarkov}
%
Let $\mathcal{M}$ be a $2$-dimensional compact Riemannian manifold. Let $\mu_0\in \mathcal{P}(\mathcal{M})$ and $\{X_n\}_{n\geq 1}\subset \mathcal{M}$ be a Markov chain with initial law $\mu_0$. We assume that the chain satisfies the two following conditions:
%
\begin{itemize}
%
\item[(i)]We assume that there exists a measurable function $k : \mathcal{M}\times\mathcal{M}\rightarrow [0,\infty)$ and $\lambda,\Lambda>0$ such that for any Borel set $A\subset\mathcal{M}$
%
\begin{equation}\label{AbsoluteContinuity}
%
K(x,A)=\int_{A}k(x,\cdot)\dd \m\quad\text{with $\lambda\leq k\leq \Lambda$}.
%
\end{equation}
%
\item[(ii)]We assume that there exist an unique invariant measure $\mu_\infty$, i.e. $K\mu_\infty=\mu_\infty$, constants $a,b>0$ and $\eta\in (0,1]$ as well as a map $V : \mathcal{M}\rightarrow [1,\infty)$ such that $\int_\mathcal{M}V\dd\mu_\infty,\int_{\mathcal{M}}V\dd\mu_0<\infty$, for which for any $\ell\geq 1$ and any 
$\phi\in\LL^\infty(\mathcal{M})$
%
\begin{equation}\label{ConvergenceRate}
%
\|K^\ell\phi-\mu_\infty(\phi)\|_V\leq a\exp(-b\ell^{\eta})\|\phi\|_\infty\quad\text{for any $x\in\mathcal{M}$},
%
\end{equation}
%
where 
%
$$\|\phi\|_V:=\sup \frac{\vert \phi\vert}{1+V}\quad\text{and}\quad\mu_\infty(\phi)=\int_{\mathcal{M}}\phi\,\dd\mu_\infty.$$
%
\end{itemize}
%
\end{assumption}
%
We now comment on the consequences of the above assumptions. First, the condition \eqref{AbsoluteContinuity} ensures that the invariant measure $\mu_\infty$ is absolutely continuous with respect to the volume measure with bounded density, namely
%
\begin{equation}\label{AbsolutelyContinuousInvariantMeasure}
%
\mu_\infty=\rho\,\dd \m\quad\text{with $\lambda\leq \rho\leq \Lambda$.}
%
\end{equation}
%
We briefly give the argument for \eqref{AbsolutelyContinuousInvariantMeasure}. Using the second item of \eqref{AbsoluteContinuity}, we have that the operator 
%
$$\mathcal{I}: f\in\LL^1\mapsto \int_{\mathcal{M}}k(\cdot,y)f(y)\dd\m(y)\quad\text{is compact.}$$
%
Furthermore, we note that the closed convex set $\mathcal{C}:=\{f\in \LL^1\,\vert\, \lambda\leq f\leq \Lambda\text{ and }\int_{\mathcal{M}}f=1\}$ is invariant under the action of $\mathcal{I}$. Therefore Schauder's fixed point theorem implies that $\mathcal{I}$ admits a fixed point in $\mathcal{C}$. Given such a fixed point $\rho$, it is clear that $\mu_\infty$ defined in \eqref{AbsolutelyContinuousInvariantMeasure} is an invariant measure according to \eqref{AbsoluteContinuity}. 

%
\medskip
%

Second, for irreducible and aperiodic Markov chains, the condition \eqref{ConvergenceRate} is satisfied when the kernel satisfies the following geometric drift condition: There exist a function $V: \mathcal{M}\rightarrow [1,\infty)$, a petite set $\mathcal{P}$ and a constant $C>0$ such that $\sup_{\mathcal{P}} V<\infty$ and
%
$$KV+\phi\circ V\leq V+C\mathds{1}_\mathcal{P},$$
%
with for large $x\geq 1$, $\phi(x)=c\frac{x}{\log^\alpha(x)}$ for some $\alpha\geq 0$ and $c>0$ (\eqref{ConvergenceRate} is then satisfied with $\eta=\frac{1}{1+\alpha}$), see \cite[Theorem 2.8 \& Section 2.3]{douc2004practical} and the references therein. Moreover, the assumption \eqref{ConvergenceRate} implies the sub-exponential decay of the $\beta$-mixing coefficient \eqref{BetaMixing}, namely there exists a constant $C$ depending on $\lambda$, $\Lambda$ and $\m(\mathcal{M})$ such that
%
\begin{equation}\label{eq:markchainbetamixc}
\beta_{ij}\leq C\exp(-b\vert i-j\vert^\eta)\quad\text{for any $i,j$,}
\end{equation}
% 
which ensures that \eqref{DecatBetaMixing} and \eqref{DecayAlphaMixing} hold and ensures good concentration property of the Markov chain, see Proposition \ref{BersteinCorrelated}. The estimate \eqref{eq:markchainbetamixc} can be seen as a direct consequence of the combination of the estimate on the $\beta$-mixing coefficient in \cite[Proposition 3]{liebscher2005towards} and the assumptions \eqref{AbsoluteContinuity} and \eqref{ConvergenceRate}.
%\begin{theorem}\label{thm:MarkovChain}
%Let $\{X_n\}_{n\ge 1}$ be a Markov chain with initial distribution $\lambda$ belonging to the class $\mathcal{C}$. Let $\pi$ be the invariant measure with distribution $\rho$ as in \eqref{AbsolutelyContinuousInvariantMeasure} and let $\mu^n$ be the empirical measure defined in \eqref{eq:empmeas} associated to the Markov chain. Let $f^{n,t}\in {\rm \dot{H}^1}$ be the weak solution of
%
%\begin{equation}\label{eq:f^ntBis}
%-\nabla\cdot \rho\nabla f^{n,t} =\mu^{n,t}-\rho_t,
%\end{equation}
%
%where we recall that $\mu^{n,t}=\pp_t\mu^n$ and $\rho_t=\pp_t\rho$, and denote by $T^n$ the optimal transport map from $\pi$ to $\mu^n$. Then it holds
%\begin{equation}\label{ApproximationTransportMap}
%\int_\mathcal{M} \mathrm{d}^2\big(T^n,\mathrm{exp}(\nabla f^{n,t})\big)\dd \mathrm{m}\leq \mathcal{C}_n\frac{\log n}{n}\sqrt{\frac{\log\log(n)}{\log(n)}}\quad\text{with \,$\sup_{n\geq 1}\mathbb{E}[\tfrac{1}{C}\mathcal{C}_n]\leq 2$,}
%\end{equation}
%\end{theorem}

%
\medskip
%

Finally, the condition \eqref{ConvergenceRate} quantifies the weak convergence of the law of the Markov chain to its stationary distribution, namely there exists a constant $C>0$ such that for any $f \in \LL^\infty(\mathcal{M})$
%
\begin{equation}\label{eq:lemMark}
|\mathbb{E}[f(X_n)] - \mu_\infty(f)| \leq C\exp(-b n^\eta) \| f\|_{\LL^\infty}.
\end{equation}
%
We shortly give the argument. We first notice that a direct inductive argument together with the semigroup property $K^{n_1+n_2} = K^{n_1}K^{n_2}$ for every $n_1,n_2>0$ and Fubini's theorem gives
%
\begin{equation}\label{eq:swapkern}
%
\int_{\mathcal{M}} f\,\dd K^n\mu_0=  \int_{\mathcal{M}} K^n f\, \dd\mu_0\quad\text{for any $n\geq 1$.}
%
\end{equation}
%
The combination of \eqref{LawOfTheChain}, \eqref{eq:swapkern} and \eqref{ConvergenceRate} gives
%
\begin{align*}
%
|\mathbb{E}[f(X_n)] - \mu_\infty(f)| & \stackrel{\eqref{LawOfTheChain},\eqref{eq:swapkern}}{=} \Big| \int_{\mathcal{M}} (K^n f - \mu_\infty(f))\,\dd\mu_0 \Big| \\ & \stackrel{\eqref{ConvergenceRate}}{\le}  \Big(\int_\mathcal{M} 1+V\dd\mu_0\Big)a\exp(-b n^\eta) \|f\|_\infty.
%
\end{align*}
%
A classical example of a Markov chain satisfying Assumption \ref{ClassMarkov} is given by iterated function systems with additive noise. For simplicity, let $\mathcal{M}=\mathbb{T}^2$. Let $\{\theta_n\}_{n\geq 1}$ be i.i.d. random variables with common law $h\,\dd\m$ for some $h$ satisfying $\lambda\leq h\leq \Lambda$. Let $F:\mathbb{T}^2\rightarrow\mathbb{T}^2$ be a contraction, i.e. there exists a constant $L<1$ such that
%
\begin{equation}\label{eq:lipcont}
\vert F(x)-F(y)\vert \le L \vert x-y\vert\quad\text{for any $x,y\in\mathbb{T}^d$.}
\end{equation}
%
We define the iterated function system $\{X_n\}_{n\geq 1}$ according to the induction
%
$$X_{n+1}=F(X_n)+\theta_n\quad\text{for any $n\geq 1$}.$$
%
The kernel is given by 
%
$$K(x,A)=\int_{\mathbb{T}^2}\mathds{1}_A(F(x)+\theta)h(\theta)\dd\m(\theta),$$
%
so that $K$ satisfies \eqref{AbsoluteContinuity} with 
%
$$k(x,\cdot)=h(\cdot-F(x)).$$
%
Moreover, the condition \eqref{eq:lipcont} ensures the validity of \eqref{ConvergenceRate}, see for instance \cite[Theorem 3.2]{alsm03}. Thus the Markov process $\{X_n\}_{n\ge1}$ satisfies Assumption \ref{ClassMarkov}.

\medskip
We show in Appendix \ref{ClassMarkovAppendix} that the conclusions of Theorem \ref{th1} and Theorem \ref{MainResultMatchingClouds} hold true for Markov chains satisfying Assumption \ref{ClassMarkov}.
%Assume that the Markov chain admits a unique invariant measure $\pi$. In addition, we assume that the transition kernel is given by 
%
%$$P(x,A)=\int_A \dd \m(y)\,\mu(x,y)\quad \text{for any $x\in\mathcal{M}$,}$$
%
%for $\mu$ satisfying \eqref{Ellipticity}. By definition
%
%$$\pi(A)=\int\dd\pi(x)\,P(x,A)\quad\text{for any measurable set $A$,}$$
%
%so that if we look for $\pi$ being absolutely continuous with respect to $\dd\m$, i.e in the form of $\pi=f\dd\m$, a sufficient condition is to solve 
%
%\begin{equation}\label{IntegralEquation}
%
%f-\int_{\mathcal{M}}\dd \m(x)\mu(x,\cdot)f(x)=0\quad\text{with $\int_\mathcal{M} f=1$.}
%
%\end{equation}
%
%Everything is in the study of the integral operator $A: f\mapsto \int_{\mathcal{M}}\dd\m(x)\mu(x,\cdot)f(x)$. To fit we our assumptions, we also require that $\mu$ satisfies the ellipticity condition \eqref{Ellipticity} which transfers into $f$ via \eqref{IntegralEquation}.

%
%\medskip
%

%\nc{A small (necessary) remark on the law of $X_k$. }Using the transition probability and defining $\mu_k$ the law of $X_k$, we have 
%
%$$\mu_{k+1}=\int_{\mathcal{M}} \dd\mu_k(x)P(x,\cdot).$$
%
%Denoting by $\lambda$ the initial distribution, we have by induction that $\mu_k:=f_k\dd \m$ with 
%
%$$f_{k+1}=\int_{\mathcal{M}}\dd \m(x)\, f_k(x)\mu(x,\cdot)\quad\text{with }f_0=\int_{\mathcal{M}}\dd\lambda (x)\mu(x,\cdot),$$
%
%and immediately $f_k$ satisfies \eqref{Ellipticity}.

%
%\medskip
%

%We can solve \eqref{IntegralEquation} in 
%
%$$S:=\Big\{f\in\LL^1(\mathcal{M})\vert \lambda\leq f\leq \Lambda \quad\text{and}\quad \int_{\mathcal{M}}\dd\m\, f=1\Big\},$$
%
%provided that the operator $A$ is compact in $\LL^1$. That is a simple application of Schauder fixed point theorem.
%
\subsection{Open problems}
%
%{\color{red}Here add paragraphs on the possible interesting extensions. First, the question of correlations, mentioning the case of equilibrium of interacting particle systems given by the Gibbs measure : $\{X_i\}_{i\leq N}$ in $\mathbb{T}^2$ has the joint distribution
%%
%$$\mathbb{P}^N=\frac{1}{Z_N}\exp\Big(-\frac{1}{N}\sum_{i\neq j} V(X_i-X_j)\Big).$$
%%
%the mean-field, $\mu^n\rightharpoonup \rho$ a.s where $\rho$ solves
%%
%$$\Delta\rho+\nabla\cdot (\rho(\nabla V\star\rho))=0.$$
%%
%To extend the result, one needs that $\lambda\leq \rho\leq \Lambda$ and quantitative results on the convergence to the mean-field limit and the bound on $\beta_{ij}$. Could be nice if we have an example when this is fulfilled with a conjecture on the result we should obtain without proof.
%
%%
%\medskip
%%
%
%We should also comment on extension of the result to $\mathbb{R}^2$, talk about Gaussian matching and extension of the Gibbs matching with application for the Ginibre ensembles. That also leads to the question when the lower bound on the density is violated.}
%
%%
%\medskip
%%

We conclude this section with open questions that arise in view of our results. Those concern optimality of our convergence rates, extensions to more general costs and different type of correlated point clouds. For the latter, we mention two possible directions concerning the Ginibre ensemble and Coulomb gases that we think are worth investigating.
%
\begin{enumerate}
%
\item \emph{Sharpness of the rates. }The convergence rate in \eqref{ApproximationTransportPlan} and \eqref{ApproximationTransportMap} match with the one obtained for the case of uniformly distributed samples in \cite{ambrosio2019optimal}. However, even in the latter case, it has not been shown whether the rate is optimal and we suspect the opposite. A possible way to track the optimal rate could be to perform a second-order linearization of the Monge-Ampère equation \eqref{eq:MongAmp}. Following the same type of computations leading to \eqref{eq:ApproxPDE} in the case $\rho=1$, a second-order linearization $q^n$ should solve 
%
$$-\Delta q^n=\text{det}(\nabla^2 h^n),$$
%
where we recall that $h^n$ solves \eqref{eq:ApproxPDE}, providing the conjecture 
%
$$\lim_{n\uparrow\infty}\bigg\vert\int_{\mathbb{T}^2}\vert T^n-(\text{Id}+\nabla h^n)\vert^2-\int_{\mathbb{T}^2}\vert\nabla q^n\vert^2\bigg\vert=0.$$
%

%Although the rates of \eqref{eq:approxtranspomapintro} match the ones of the uniform case, no matching lower bounds have been shown yet. We expect the rates to not be sharp \fm{give a simple heuristic argument?}. This might require a higher order linearization of the Monge-Amp\`ere equation or a finer first order estimate. 

%
\medskip
%

\item \emph{Extension to $p$-costs. }A natural question is to investigate if our results hold for different cost functions as $p$-cost functions for $p>1$. The behavior of the cost has been optimally quantified in \cite{AKT84,boutet2002almost,Barthe2013,Fournier2015}. However, to the best of our knowledge, quantitative estimates on the transport plan in the setting of general $p$-costs are not known, even for uniformly distributed samples.
%A possible extension of our result would be to consider other costs than the quadratic one. To the best of our knowledge quantitative estimates for the convergence of the approximating transport map in the setting of $p$-costs are not known for the uniform case. On the other hand, it is well known that similar rates to \eqref{eq:unifmatchingrates} hold in the case of uniformly distributed points. Thus, we would expect that a similar quantitative convergence result for the transport map would hold also in this setting. 
A possible approach would be to revise the linearization ansatz of \cite{caracciolo2014scaling} for general $p$-costs. Indeed, if the transport cost between two points $x,y$ is given by $\frac{1}{p}|x-y|^p$ on the torus, Gangbo-McCann's theorem \cite[Theorem 1.2]{Gangbo1996} ensures that there exists a map $h^n$ such that the optimal transport map $T^n$ takes the form $T^n= \text{Id} +  |\nabla h^n|^{p'-2}\nabla h^n$, where $p'$ denotes the conjugate exponent. Therefore, following the same type of computations leading to \eqref{eq:ApproxPDE}, a first-order linearization should solve the following degenerate $p'$-Laplace equation
%
$$\nabla \cdot \rho|\nabla h^n|^{p'-2}\nabla h^n=\mu^n- \nu^m,$$
%
and we may expect 
%
$$\lim_{n\uparrow\infty}\int_{\mathbb{T}^2}\vert T^n-(\text{Id}+|\nabla h^n|^{p'-2}\nabla h^n)\vert^{p}=0.$$
%
%
See also \cite{Lukas} for a justification of this linearisation ansatz down to mesoscopic scales based on a large-scale regularity theory for the Monge-Amp\`ere equation.


\medskip

%
\item \emph{Ginibre ensemble. }A (complex) Ginibre ensemble is a non-Hermitian random matrix with independent complex Gaussian entries. Given a $n\times n$ Ginibre ensemble $X$, we define its empirical spectral distribution as 
\[
\mu^n = \frac 1n \sum_{i=1}^n \delta_{\lambda_i },
\]
where $\{\lambda_i\}_{i=1}^n$ are the eigenvalues of the matrix $\frac{X}{\sqrt n}$. The so called Circular Law states that, almost-surely, $\mu^n$ weakly converges to the uniform distribution on the complex unit disk $\mu^\infty$ having Lebesgue density $\frac1\pi \mathds{1}_{\bb_1}$, see for instance \cite[Theorem 1.10]{TaoVu}. An interesting question is to quantitatively understand the weak convergence of $\mu^n$ towards $\mu^\infty$ measured Wasserstein distances. A possible approach to this problem would be to employ the linearization argument discussed in Section \ref{sec:lin}. Note that the Ginibre ensemble is not covered by our setting as it is posed on the whole space $\mathbb{R}^2$ and, in general, the eigenvalues $\{\lambda_i\}_{i=1}^n$ possess long-range correlations and therefore do not satisfy our Assumption \ref{Assumptions}. However, in \cite[Th\'eor\`eme 3.1.1]{maxime}, the linearization argument has been shown to be robust enough to derive an upper bound on $\mathbb{E}[W_2(\mu^n, \mu^\infty)]$ (see also \cite[Theorem 1.3]{Jonas} for similar results for $\mathbb{E}[W_1(\mu^n, \mu^\infty)]$). %Though it is far from clear whether the constants in \cite[Th\'eor\`eme 3.1.1]{maxime} and \cite[Theorem 1.3]{Jonas} are sharp, the linearization approach is shown to be robust enough to study some classes of not necessarily weakly correlated point clouds. Indeed, the Ginibre ensemble is not covered by our setting as in general the eigenvalues $(\lambda_i)_{i=1}^n$ do not satisfy the Assumptions \ref{Assumptions}. 
The techniques used in \cite{maxime, Jonas} avoid a quantification of the correlations between the eigenvalues: This is done in \cite{maxime} by a decomposition argument together with concentration estimates of the Wasserstein distance around the Moser's coupling and in \cite{Jonas} by using classical tools from non-Hermitian random matrix theory. A challenging question would be to investigate the exact asymptotics of the transport cost $\mathbb{E}[W_2(\mu^n, \mu^\infty)]$ in the case of the Ginibre ensemble complementing, \cite[Th\'eor\`eme 3.1.1]{maxime} and \cite[Theorem 1.3]{Jonas} with a lower bound and consequently quantify the convergence of the linearized proxies to the optimal transport map. 

%The linearization ansatz initiated in \cite{caracciolo2014scaling} has been shown to be robust also when dealing with certain classes of interacting particle system, see for instance . 
%
%An interesting and challenging question is to investigate the robustness of the linearization initiated in \cite{caracciolo2014scaling} in a more general setting. Motivated by the work \cite{Jonas} on the Ginibre ensemble it would be interesting to study the validity of the two conjectures \eqref{eq:approxcost} in the setting of Coulomb gases. 

\medskip
\item \emph{Planar Coulomb gases. }Planar Coulomb gases are many  particles systems, in which the particles $\{X_i\}_{i=1}^n$ have repulsively Coulomb interactions and are confined by a potential $V: \mathbb{T}^2 \rightarrow \mathbb{R}$. These are modelled by the Hamiltonian
\begin{equation}\label{eq:hamiltonian}
\mathcal{H}_n(\{X_i\}_{i=1}^n) = - \sum_{i \neq j} \log |X_i - X_j| + n \sum_{i=1}^n V(X_i)
\end{equation}
and the Gibbs measure
\begin{equation}\label{eq:gibbsmeasure}
\mathrm{d}\mathbb{P}_{n, \beta} = \frac 1{Z_{n,\beta}} \exp\big(- \beta \mathcal{H}_n  (\{X_i\}_{i=1}^n)\big)\, \mathrm{d}X_1 \dots \mathrm{d}X_n,
\end{equation}
where $Z_{n,\beta}$ is the normalizing constant and $\beta$ denotes the inverse temperature.  In analogy with \eqref{eq:empmeas}, we can define the empirical measure of a Coulomb gas by $\mu^n = \frac1n \sum_{i=1}^n \delta_{X_i}$. In this setting the convergence of the empirical measure exhibits a twofold behavior.  On the one hand, for small temperature $\beta \gg \frac1n$, the empirical measure $\mu^n$ weakly converges to the \emph{equilibrium measure} $\mu^\infty_{\text{low}}$ defined as the minimizer
%
$$\mu^\infty_{\text{low}} := \operatornamewithlimits{argmin}_{\mu \in \mathcal{P}(\mathbb{T}^2)} \bigg\{ - \int_{\mathbb{T}^2\times \mathbb{T}^2} \log (x-y) \,\mathrm{d}\mu \otimes \mu (x,y) + \int_{\mathbb{T}^2} V \,\mathrm{d}\mu\bigg\}. $$
%
On the other hand, for large temperature $\beta\ll \frac{1}{n}$, a correction term is required and the empirical measure $\mu^n$ weakly converges to the \emph{thermal equilibrium measure} $\mu^\infty_{\text{high}}$, that is the minimizer
%
\begin{equation*}
%
\mu^\infty_{\text{high}} :=\operatornamewithlimits{argmin}_{\mu \in \mathcal{P}(\mathbb{T}^2)} \bigg\{- \int_{\mathbb{T}^2\times \mathbb{T}^2} \log (x-y) \,\mathrm{d}\mu \otimes \mu (x,y) + \int_{\mathbb{T}^2} V \,\mathrm{d}\mu + \frac 1{n \beta} \int_{\mathbb{T}^2} \mu \log \mu\bigg\}.
%
\end{equation*}
%
We refer the reader to \cite{serfaty} for a complete exposition on $2$D-Coulomb gases. This setting can be seen as an extension of the previous Ginibre ensemble on the torus. Indeed the law of the spectrum of the Ginibre ensemble is given by the Gibbs measure \eqref{eq:gibbsmeasure} choosing $\beta=2$ and $V(x)=|x|^2$ in \eqref{eq:hamiltonian}. Motivated by this observation and the works \cite{maxime, Jonas} on the Ginibre ensemble, we expect that the linearization approach could be employed also in this setting. First, we could justify the linearization Ansatz in the spirit of Theorem \ref{MainResultMatchingClouds} in both temperature regimes, using available results in the literature. Indeed, concentration inequalities around the equilibrium measure have been derived in \cite{Zelada} (see also \cite{Chafai}), that would replace the Bernstein's type inequality of Proposition \ref{BersteinCorrelated} and the matching cost estimates in Proposition \ref{prop:matchcost}. Second, a more ambitious question would be to use the linearization argument to derive optimal rates of the convergence to the equilibrium measure in both temperature regimes.

\end{enumerate}





\section{Structure of the proof}\label{StrategyOfProofSection}
%
This section is devoted to describe the main ideas and how are organized the proofs of Theorem \ref{MainResultMatchingClouds} and \ref{th1}. We mainly focus on the proof of Theorem \ref{MainResultMatchingClouds} since the proof of Theorem \ref{th1} follows by the same strategy.

%
\medskip
%

\subsubsection*{General strategy.}The proof of Theorem \ref{MainResultMatchingClouds} follows the strategy employed in \cite{ambrosio2019optimal} to deal with independent and uniformly distributed random points. 
The main idea is to use the quantitative stability result for transport maps in \cite[Theorem 3.2]{ambrosio2019optimal}, stating that two transport maps are close in the $\LL^2$-topology if the target measures are close in the $W_2$-topology. We restate below the result for reader's convenience.
%
\begin{theorem}[Stability of transport maps]\label{StabilityResultMap}
%
Let $\nu,\mu_1,\mu_2\in \mathcal{P}(\mathcal{M})$ such that $\nu\ll \m$ and let $T,S: \mathcal{M}\rightarrow\mathcal{M}$ be the optimal transport maps respectively for the pairs of measures $(\nu,\mu_1)$ and $(\nu,\mu_2)$. We assume that $S=\exp(\nabla f)$ for some $f:\mathcal{M}\rightarrow\mathbb{R}$ with $\mathrm{C}^{1,1}$-regularity. 

%
\medskip
%

\noindent There exists a constant $c>0$ depending on $\mathcal{M}$ such that, provided 
%
\begin{equation}\label{RegularityCondition}
%
\|\nabla f\|_{\LL^\infty}+\|\nabla^2 f\|_{\LL^\infty}\leq c,
%
\end{equation}
%
we have 
%
$$\int_{\mathcal{M}}\dd^2(S,T)\,\dd\nu\lesssim W^2_2(\mu_1,\mu_2)+W_2(\mu_1,\mu_2)W_2(\nu,\mu_1).$$
%
\end{theorem}
%
The first step consists of using Theorem \ref{StabilityResultMap} to deduce a stability estimate (in terms of the quadratic Wasserstein distance) of transport plans in the special case where $\mu_1=\nu^m$, $\mu_2=\exp(\nabla h^{n,t})_\#\mu^n$ and $\nu=\mu^n$. In this step, we immediately face the issue of the lack of regularity of $h^{n,t}$ necessary to ensure the condition \eqref{RegularityCondition}: Indeed, recalling that $h^{n,t}$ solves \eqref{eq:f^nmt}, it does not have the $\mathrm{C}^{1,1}$-regularity condition for non-smooth densities $\rho$ that we consider here. We overcome this issue introducing an additional regularization step: We smooth  the operator $-\nabla \cdot \rho \nabla$ and implicitly $\gamma^{n,t}$ in form of 
%
\begin{equation}\label{RegularizationIntro}
%
\gamma^{n,t}_{\delta}:=\big(\text{Id},\exp(\nabla h^{n,t}_\delta)\big)_\#\mu^{n,t}\quad \text{with }\nabla\cdot\rho_\delta\nabla h^{n,t}_\delta=\mu^{n,t}-\nu^{m,t},
%
\end{equation}
%
for a regularization parameter $\delta$ to be optimized and $\rho_\delta:=\text{P}_\delta \rho$. Classical Schauder's theory ensures that $h^{n,t}_\delta$ owns $\cc^\infty$-regularity. 
%
Doing so, we can use Theorem \ref{StabilityResultMap} to deduce, provided that 
%
\begin{equation}\label{SmallnessH}
%
\|\nabla h^{n,t}_\delta\|_{\LL^{\infty}}+\|\nabla^2 h^{n,t}_\delta\|_{\LL^\infty}\ll 1,
%
\end{equation}
%
a stability error estimate which reads
%
\begin{equation}\label{QuantitativeEstiPlanIntro}
%
\begin{aligned}
%
\inf_\pi W^2_2(\pi,\gamma^{n,t}_\delta)\lesssim&\, W^2_2\big(\nu^{m,t},\exp(\nabla h^{n,t}_\delta)_\#\mu^{n,t}\big)+W_2\big(\nu^{m,t},\exp(\nabla h^{n,t}_\delta)_\#\mu^{n,t}\big)W_2(\mu^{n},\nu^{m})\\
&+W^2_2(\nu^{m,t},\nu^{m})+W^2_2(\mu^{n,t},\mu^{n})+\big(W_2(\nu^{m,t},\nu^m)+W_2(\mu^{n,t},\mu^n)\big)W_2(\mu^n,\nu^m),
%
\end{aligned}
\end{equation}
%
where we refer to \eqref{QuantitativeEstiPlan} for further details. Our argument then differs from \cite{ambrosio2019optimal} by splitting the error in two parts
%
\begin{equation}\label{SplittingIntro}
%
\inf_\pi W^2_2(\pi,\gamma^{n,t})\leq \underbrace{W^2_2(\gamma^{n,t},\gamma^{n,t}_\delta)}_{\text{regularization error}}+\underbrace{\inf_{\pi} W^2_2(\pi,\gamma^{n,t}_\delta)}_{\text{stability error}}.
%
\end{equation}
%
Our proof then proceeds in two steps, controlling separately the two terms in \eqref{SplittingIntro}.
%
\subsubsection*{Control of the regularization error}To deal with the regularization error, we look at the difference $e^{n,t}_\delta:=h^{n,t}_\delta-h^{n,t}$ which solves, according to \eqref{eq:f^nmt} and \eqref{RegularizationIntro},
%
\begin{equation}\label{ErrorEquation}
%
-\nabla\cdot\rho_\delta\nabla e^{n,t}_\delta=\nabla\cdot(\rho_\delta-\rho)\nabla h^{n,t}.
%
\end{equation}
%
Using an energy estimate, we get
%
\begin{equation}\label{ErrorIntroRegul}
%
\int_{\mathcal{M}}\vert\nabla e^{n,t}_\delta\vert^2\lesssim \int_{\mathcal{M}}\vert\rho-\rho_\delta\vert^2\vert\nabla h^{n,t}\vert^2.
%
\end{equation}
%
On the one hand, since $\rho\in \LL^\infty$, we have $\rho_\delta\rightarrow\rho$ as $\delta\downarrow 0$ in every $\LL^q$ with $q<\infty$. On the other hand, we learn from Meyers' estimate (recalled in Proposition \ref{Meyers}) that there exists $\bar q>2$ such that $\nabla h^{n,t}\in \LL^{\bar q}$. Consequently, we can treat \eqref{ErrorIntroRegul} using H\"older's inequality which provides
%
\begin{equation}\label{LqControlIntro}
%
\int_{\mathcal{M}}\vert\nabla e^{n,t}_\delta\vert^2\lesssim\|\rho_\delta-\rho\|^2_{\LL^{2(\frac{\bar{q}}{2})'}}\|\nabla h^{n,t}\|^2_{\LL^{\bar{q}}}.
%
\end{equation}
%
The next step is to control the averaged Meyers' norm $\mathbb{E}\big[\|\nabla h^{n,t}\|^2_{\LL^{\bar{q}}}\big]$, that we show in Proposition \ref{prop1} to be of order of 
%
\begin{equation}\label{LqEstimatesIntro}
%
\mathbb{E}\big[\|\nabla h^{n,t}\|^2_{\LL^{\bar{q}}}\big]\lesssim \frac{\vert\log(t)\vert+\log^{\frac{1}{\eta}}(n)}{n},
%
\end{equation} 
%
where we recall that $\eta$ denotes the correlation length, see Assumption \ref{Assumptions}. 

%
\medskip
%

The combination of \eqref{LqControlIntro}, \eqref{LqEstimatesIntro} and local Lipschitzianity of the exponential map yields
%
\begin{equation}\label{regularizedErrorIntro}
%
\mathbb{E}\big[W^2_2(\gamma^{n,t}_\delta,\gamma^{n,t})\big]\lesssim \|\rho_\delta-\rho\|^2_{\LL^{2(\frac{\bar{q}}{2})'}}\frac{\vert \log(t)\vert+\log^{\frac{1}{\eta}}(n)}{n},
%
\end{equation}
%
 where we refer to \eqref{eq:thm1explipW2} for further details. We emphasize here that the stretched exponential decay assumption \eqref{DecayAlphaMixing} of the $\alpha$-mixing coefficient plays a crucial role in the estimate \eqref{LqEstimatesIntro}. Indeed, the additional contribution on the numerator of the r.h.s. of \eqref{LqEstimatesIntro} is due to the correlations and is only of logarithmic type $\log^{\frac{1}{\eta}}(n)$ thanks to the stretched exponential decay \eqref{DecayAlphaMixing}. Finally, the latter appears in the estimate \eqref{regularizedErrorIntro} and can be compensated with the choice of $\delta$ in \eqref{ChoicesIntro} and the regularity assumption on $\rho$ in form of \eqref{ApproximationRhoIntro}.
%
\subsubsection*{Control of the stability error}For the stability error, we first need to ensure \eqref{SmallnessH}. Our strategy follows the idea in \cite{ambrosio2019finer} which consists of showing that \eqref{SmallnessH} is satisfied with very high probability. In our case, a new difficulty comes from our regularization of $\rho$ and the regularization parameter $\delta$ has to be carefully optimized. We show that if $\delta$ is taken as an inverse power of $\log(n)$, \eqref{SmallnessH} becomes very likely as $n\uparrow\infty$. More precisely, we show in Proposition \ref{Fluctuation} that given two exponents $\kappa_1$ and $\upsilon\gg \kappa_1$, there exists $\kappa_2$ such that given the choices
%
\begin{equation}\label{ChoicesIntro}
%
\delta=\frac{1}{\log^{\kappa_1}(n)}\quad \text{and}\quad t=\frac{\log^{\kappa_2}(n)}{n},
%
\end{equation}
%
we have
%
\begin{equation}\label{FluctuationEstimates}
%
\mathbb{P}\Big(\|\nabla h^{n,t}_\delta\|_{\LL^{\infty}}+\|\nabla^2 h^{n,t}_\delta\|_{\LL^\infty}>\tfrac{1}{\log^{\upsilon}(n)}\Big)=o\Big(\frac{1}{n^\ell}\Big)\quad\text{for any $\ell\in \mathbb{N}$}.
%
\end{equation}
%
This result is in the spirit of \cite[Theorem 3.3]{ambrosio2019finer} but, in our setting, the proof relies on Schauder's theory rather than an explicit formula for $h^{n,t}_\delta$ as well as concentration inequalities in form of Proposition \ref{BersteinCorrelated} to treat the correlations. For further details on the strategy, we refer to Section \ref{FluctuationSection}. We emphasize here that, to obtain the super-polynomial behaviour \eqref{FluctuationEstimates}, we crucially use the fact that the concentration inequalities in Proposition \ref{BersteinCorrelated} are of sub-exponential type which is itself ensured by the sub-exponential decay assumption on the $\alpha$-mixing coefficient \eqref{DecayAlphaMixing}. The reason lies in the choices of $\delta$ and $t$ in \eqref{ChoicesIntro}. Indeed, the only room we have when optimizing $t$ is in the logarithmic growth $\log^{\kappa_2}(n)$ (as already mention in Section \ref{sec:lin}, the natural regularization time is $t\sim \frac{1}{n}$ up to logarithmic corrections). Furthermore, the quantity $\|(\nabla h^{n,t}_\delta,\nabla^2 h^{n,t}_\delta)\|_{\LL^{\infty}}$ can be heuristically estimated by an inverse power of $\delta$ and $t^{-1}$ as it involves powers of $\|(\nabla\rho_\delta,\nabla^2\rho_\delta)\|_{\LL^\infty}$ and the norm $\|\mu^{n,t}-1\|_{\LL^{\infty}}$ by Schauder's theory. Hence, in the best case scenario where exponential concentration holds, we expect
%
\begin{align*}
%
\mathbb{P}\Big(\|\nabla h^{n,t}_\delta\|_{\LL^{\infty}}+\|\nabla^2 h^{n,t}_\delta\|_{\LL^\infty}>\lambda\Big)&\lesssim \exp\big(-\lambda n\|(\nabla h^{n,t}_\delta,\nabla^2 h^{n,t}_\delta)\|^{-1}_{\LL^{\infty}}\big)\\
&\leq \exp\big(-\lambda nt\delta^{\tilde{\upsilon}}\big)=\exp\big(-\lambda\log^{\kappa_2}(n)\delta^{\tilde{\upsilon}}\big),
%
\end{align*}
%
for some $\tilde{\upsilon}>0$, which gives a super-polynomial behavior for the choices $\delta=\frac{1}{\log^{\kappa_1}(n)}$ and $\lambda=\frac{1}{\log^{\upsilon}(n)}$ for large $\kappa_2$. Weaker properties, as for instance polynomial concentration, would only lead to a decay given by an inverse power of $\log(n)$ which is not enough for our purpose.
%
\medskip
%

With \eqref{FluctuationEstimates} in hands, we can restrict the analysis to realizations satisfying $\|\nabla h^{n,t}_\delta\|_{\LL^{\infty}}+\|\nabla^2 h^{n,t}_\delta\|_{\LL^\infty}\leq\tfrac{1}{\log^{\upsilon}(n)}$ where, for $n\gg 1$, \eqref{SmallnessH} is satisfied which puts us in the validity of \eqref{QuantitativeEstiPlanIntro}. We then treat each terms appearing in \eqref{QuantitativeEstiPlanIntro} separately:
%
\begin{itemize}
%
\item \textbf{The optimal control of the cost} $W^2_2(\mu^n,\nu^m)$ has been already well studied and is optimally estimated by 
%
\begin{equation}\label{CostIntro}
%
\mathbb{E}\big[W^2_2(\mu^n,\nu^m)\big]\lesssim \frac{\log(n)}{n}.
%
\end{equation}
%
We refer to Appendix \ref{sec:matchcost} for a detailed statement, references and extensions to the cases of Assumption \ref{Assumptions} and Assumption \ref{ClassMarkov}.

%
\medskip
%

\item \textbf{The smoothing errors} $W^2_2(\mu^n,\mu^{n,t})$ and $W^2_2(\nu^m,\nu^{m,t})$. Classical contractivity estimates are known and are usually applied to deal with these errors, see for instance  \cite[Theorem 3]{erbar2015equivalence}, which bound the errors by $t$. However, due to the choice of $t$ in \eqref{ChoicesIntro}, this result is of no use since $t$ is much larger than the magnitude of the cost, namely $t\gg \frac{\log(n)}{n}$. Instead, we follow the approach in \cite{ambrosio2019finer}, where the authors showed that in the particular case of empirical measures in dimension $2$, we can improve the rate and obtain the bound $\frac{\log\log(n)}{n}\ll \frac{\log(n)}{n}$. We extend this result to our setting of non-constant densities and correlated points. In Proposition \ref{Contractivity}, we derive
%
\begin{equation}\label{ContractivityEstimatesIntro}
%
\mathbb{E}\big[W^2_2(\mu^{n},\mu^{n,t})\big]+\mathbb{E}\big[W^2_2(\nu^{m},\nu^{m,t})\big]\lesssim \frac{\log\log(n)}{n}+t\|\rho_{t+\frac{1}{n}}-\rho_{\frac{1}{n}}\|_{\LL^1}.
%
\end{equation}
%
As opposed to \cite{ambrosio2019finer}, our approach uses Fourier analysis together with additional cares to handle the correlations and non-constant densities.

%
\medskip
%

\item \textbf{The error in the Moser coupling} $W^2_2\big(\nu^{m,t},\exp(\nabla h^{n,t}_\delta)_\#\mu^{n,t}\big)$. We follow the strategy in \cite{ambrosio2019finer}. This error can be related to a Moser coupling between $\mu^{n,t}$ and $\nu^{m,t}$ (see for instance \cite[Appendix p. 16]{OldNewVillani}): The equation \eqref{eq:f^nmt} gives a natural coupling between $\mu^{n,t}$ and $\nu^{m,t}$ which can be formulated using Benamou-Brenier's theorem \cite{benamou2000computational},
%
$$\nu^{m,t}=\phi(1,\cdot)_\#\mu^{n,t}\quad\text{with $\phi$ being the flow induced by $s\mapsto \frac{\rho_\delta\nabla h^{n,t}_\delta}{(1-s)\mu^{n,t}+s\nu^{m,t}}$}.$$
%
Then, using the transport plan $\big(\phi(1,\cdot),\exp(\nabla h^{n,t}_\delta)\big)_\#\mu^{n,t}$ as a competitor, we get
%
\begin{align*}
%
W^2_2\big(\nu^{m,t},\exp(\nabla h^{n,t}_\delta)_\#\mu^{n,t}\big)&=W^2_2\big(\phi(1,\cdot)_\#\mu^{n,t},\exp(\nabla h^{n,t}_\delta)_\#\mu^{n,t}\big)\\
&\leq \int_{\mathcal{M}}\dd^2\big(\phi(1,\cdot),\exp(\nabla h^{n,t}_\delta)\big),
%
\end{align*}
%
that we combine with a quantitative stability result for flows of vector fields, proved in \cite[Proposition A.1]{ambrosio2019finer}, leading to
%
$$W^2_2\big(\nu^{m,t},\exp(\nabla h^{n,t}_\delta)_\#\mu^{n,t}\big)\lesssim\big(\|\rho_{\delta}-\rho\|^{2}_{\LL^{2(\frac{\bar q}{2})'}}+\|\rho_{t}-\rho\|^{2}_{\LL^{2(\frac{\bar q}{2})'}}+\tfrac{1}{\log^{\upsilon}(n)}\big)\|\nabla h^{n,t}_\delta\|^2_{\LL^{\bar{q}}},$$
%
%\end{itemize}
%
where we recall that $\bar{q}$ denotes the Meyers' exponent introduced in \eqref{ErrorIntroRegul}. For further details, we refer to \eqref{StabilityFlowProof}. It then remains to appeal to Meyers' estimate, see Proposition \ref{Meyers}, to \eqref{ErrorEquation} together with \eqref{LqEstimatesIntro} in form of 
%
$$\mathbb{E}\big[\|\nabla h^{n,t}_\delta\|^2_{\LL^{\bar{q}}}\big]\lesssim  \mathbb{E}\big[\|\nabla h^{n,t}\|^2_{\LL^{\bar{q}}}\big]\stackrel{\eqref{LqEstimatesIntro}}{\lesssim}\frac{\vert\log(t)\vert+\log^{\frac{1}{\eta}}(n)}{n},$$
%
which finally yields
%
\begin{equation}\label{EstimateMoserCoupling}
%
\mathbb{E}\big[W^2_2\big(\nu^{m,t},\exp(\nabla h^{n,t}_\delta)_\#\mu^{n,t}\big)\big]\lesssim \big(\|\rho_{\delta}-\rho\|^{2}_{\LL^{2(\frac{\bar q}{2})'}}+\|\rho_{t}-\rho\|^{2}_{\LL^{2(\frac{\bar q}{2})'}}+\tfrac{1}{\log^{\upsilon}(n)}\big)\frac{\vert\log(t)\vert+\log^{\frac{1}{\eta}}(n)}{n}.
%
\end{equation}
%
\end{itemize}
%The proof is based on a stability result for transport plans describing how quantitatively (when measuring the error in the $W_2$-Wasserstein distance) optimal plans change when we variate the initial and target measure, see Figure \ref{StabilityFigure}.
%
%\begin{figure}[h]
%
%\includegraphics[scale=2]{Stability.png}
%\caption{Given a optimal transport plan $\pi^n$, the quantitative stability result gives an estimate of $W_2(\pi^n,(\text{Id},\exp(\nabla h^{n,t})_\#\mu^{n,t})$ involving errors between the initial measures $W_2(\mu^n,\mu^{n,t})$ and the target measures $W_2(\nu^n,\exp(\nabla h^{n,t})_\#\mu^{n,t})$.}
%
%\end{figure}
%
%\begin{figure}[h]
%\begin{tikzpicture}
  %\coordinate (A) at (0,0);
  %\coordinate (B) at (6,0);
  %\coordinate (C) at (0,1);
  %\coordinate (D) at (6,1);
  %\draw [->]  (A) node[left]{$\mu^{n,t}$} -- node[above]{$(\mathrm{Id}, \exp(\nabla h^{n,t})_\#\mu^{n,t})$} (B) node[right]{$\exp(\nabla h^{n,t})_\#\mu^{n,t}$};
  %\draw [->]  (C) node[left]{$\mu^n$} -- node[above]{$\pi^n$} (D) node[right]{$\nu^m$};
%\end{tikzpicture}
%\caption{Given a optimal transport plan $\pi^n$, the quantitative stability result gives an estimate of $W_2(\pi^n,(\text{Id},\exp(\nabla h^{n,t})_\#\mu^{n,t})$ involving errors between the initial measures $W_2(\mu^n,\mu^{n,t})$ and the target measures $W_2(\nu^m,\exp(\nabla h^{n,t})_\#\mu^{n,t})$.}\label{StabilityFigure}
%
%\end{figure}
%
%\medskip
%

%Such quantitative stability results have been firstly investigated in \cite{gigli2011holder} for measures with smooth densities and, for instance, more recently in \cite{berman2021convergence} and \cite{gallouet2022strong} for flat geometries. In order to deal with non-flat geometry as Riemannian geometry considered here, we rely on the stability result developed in \cite[Theorem 3.2]{ambrosio2019optimal} that we extend in the case of transport plan. However, applying this result requires Lipschitz regularity of $\nabla h^{n,t}$ and smallness of $\|\nabla h^{n,t}\|_{\LL^{\infty}}+\|\nabla^2 h^{n,t}\|_{\LL^\infty}$. These regularity assumptions are an issue while considering non-smooth densities $\rho$. To overcome this, we regularize $\gamma^{n,t}$ into 
%
%\begin{equation}\label{RegularizationIntro}
%
%\gamma^{n,t}_{\delta}:=\big(\text{Id},\exp(\nabla h^{n,t}_\delta)\big)_\#\mu^{n,t}\quad \text{with }-\nabla\cdot\rho_\delta\nabla h^{n,t}_\delta=\mu^{n,t}-\nu^{m,t},
%
%\end{equation}
%
%for a regularization parameter $\delta$ to optimize and $\rho_\delta:=\text{P}_\delta \rho$. Classical Schauder's theory ensures that $h^{n,t}_\delta$ owns $\cc^\infty$-regularity. We then split the error in two parts: a regularization error $W^2_2(\gamma^{n,t}_\delta,\gamma^{n,t})$ and, provided
%
%\begin{equation}\label{SmallnessH}
%
%\|\nabla h^{n,t}_\delta\|_{\LL^{\infty}}+\|\nabla^2 h^{n,t}_\delta\|_{\LL^\infty}\ll 1,
%
%\end{equation}
%
%a stability error which reads
%
%\begin{equation}\label{QuantitativeEstiPlanIntro}
%
%\begin{aligned}
%
%\inf_\pi W^2_2(\pi,\gamma^{n,t}_\delta)\lesssim&\, W^2_2\big(\nu^{m,t},\exp(\nabla h^{n,t}_\delta)_\#\mu^{n,t}\big)+W_2\big(\nu^{m,t},\exp(\nabla h^{n,t}_\delta)_\#\mu^{n,t}\big)W_2(\mu^{n},\nu^{m})\\
%&+W^2_2(\nu^{m,t},\nu^{m})+W^2_2(\mu^{n,t},\mu^{n})+\big(W_2(\nu^{m,t},\nu^m)+W_2(\mu^{n,t},\mu^n)\big)W_2(\mu^n,\nu^m),
%
%\end{aligned}
%\end{equation}
%
%where the $\inf$ runs over all optimal transport plans $\pi$ between $\mu^n$ and $\nu^m$. For more details on \eqref{QuantitativeEstiPlanIntro} and its proof, see \eqref{QuantitativeEstiPlan}. 

%
%\medskip
%




%
\medskip
%

To conclude, we see that all bounds involve errors in terms of the approximation of $\rho$ using the heat-semigroup, that we need to quantify. This is where the assumption $\rho\in \mathrm{H}^\varepsilon$ plays a role, in form of the quantitative estimate 
%
\begin{equation}\label{ApproximationRhoIntro}
%
\|\rho_s-\rho\|_{\LL^2}\lesssim \|\rho\|_{\mathrm{H}^\varepsilon} \,s^{\varepsilon}\quad\text{for any $s>0$},
%
\end{equation}
%
see \eqref{ApproximationRho} for a proof. Combining \eqref{QuantitativeEstiPlanIntro}, \eqref{SplittingIntro}, \eqref{regularizedErrorIntro}, \eqref{CostIntro}, \eqref{ContractivityEstimatesIntro}, \eqref{EstimateMoserCoupling}, \eqref{ApproximationRhoIntro} with the choices of $\delta$ and $t$ in \eqref{ChoicesIntro}, we obtain Theorem \ref{MainResultMatchingClouds}. The proof of Theorem \ref{th1} is obtained using the same strategy where the first step is simpler, since we apply directly Theorem \ref{StabilityResultMap} with $\mu_1=\rho$, $\nu=\mu^n$ and $\mu_2=\exp(\nabla f^{n,t}_\delta)_\#\mu^n$ where $f^{n,t}$ solves $-\nabla\cdot \rho_\delta\nabla f^{n,t}=\mu^{n,t}-\rho_t$.

%The proof relies on the strategy employed in \cite{ambrosio2019optimal} in the case of uniformly distributed points. This strategy combines both deterministic and probabilistic arguments, that we explain hereafter together with the additional difficulties arising by considering more general densities.

%
%\medskip
%

%\begin{enumerate}
%
%\item \textbf{Deterministic arguments. }The main deterministic ingredient is the quantitative stability result given in \cite[Theorem 3.2]{ambrosio2019optimal}, recalled in \nc{write it somewhere}. This quantitative stability result requires $\cc^2(\mathcal{M})$-regularity of the map $f^{n,t}$, which would require by Schauder's theory $\cc^{1,\alpha}(\mathcal{M})$-regularity of $\rho$, for some $\alpha>0$. To avoid this regularity assumption on $\rho$, we consider its regularization $\rho_\delta:=\text{P}_\delta\rho$, for a parameter $\delta\downarrow 0$ (depending on $n$) which will be optimized later. Doing so, we consider the regularization $f^{n,t}_\delta$ of $f^{n,t}$ defined via
%
%\begin{equation}\label{eq:f^nt_delta}
%-\nabla\cdot \rho_\delta\nabla f^{n,t}_\delta=\mu^{n,t}-\rho_t\quad \text{in $\mathcal{M}$}\quad\text{with}\quad \int_{\mathcal{M}}f^{n,t}_\delta=0.
%\end{equation}
%
%In this process, the regularization error is easily estimated by an energy estimate, since from \eqref{eq:f^nt} and \eqref{eq:f^nt_delta} we deduce
%
%$$-\nabla\cdot \rho_\delta(\nabla f^{n,t}_\delta-f^{n,t})={\color{red}-}\nabla\cdot(\rho_\delta-\rho)\nabla f^{n,t},$$
%
%\fm{I think there is a minus in the line above (the one I added in red)}
%
%which gives, using that $\rho_\delta>c$
%
%\begin{equation}\label{Eq:RegError}
%\int_{\mathcal{M}}\vert\nabla f^{n,t}_\delta-\nabla f^{n,t}\vert^2\lesssim \int_{\mathcal{M}}\vert\rho_\delta-\rho\vert^2\vert\nabla f^{n,t}\vert^2.
%\end{equation}
%
%We then apply the quantitative stability result \nc{cite again} to $\nu:=\rho\,\dd x$, $\mu_1=\hat{\mu}^{n,t}_\delta:= \exp(\nabla f^{n,t}_\delta)\#\rho\,\dd x$ and $\mu_2:=\mu^{n}$,\fm{not $\rho$ but $\rho_t$ and additional approximation term, specify it later}\nc{nop, has to be $\rho$ otherwise you don't have the right transport map} yielding to an estimate involving the two quantities
%
%\begin{equation}\label{Eq:Strategy1}
%W_2(\rho\,\dd x,\mu^{n})\quad\text{and}\quad W_2(\mu^n,\hat{\mu}^{n,t}_\delta),
%\end{equation}
%
%and the regularization error \eqref{Eq:RegError}, holding in the event
%
%\begin{equation}\label{Eq:Strategy2}
%\mathcal{E}_n:=\{\|\nabla^2 f^{n,t}_\delta\|_{\LL^\infty(\mathcal{M})}\leq \log^{-\frac{1}{2}}(n)\}\quad\text{for $n\gg1$.}
%\end{equation}
%

%The control of those quantities are done by probabilistic arguments.

%
%\medskip
%

%\item \textbf{Probabilistic arguments. } To conclude, we now combine the previous quantitative stability result with moment bounds on the quantities define in \eqref{Eq:RegError} and \eqref{Eq:Strategy1} as well as a bound on $\mathbb{P}((\mathcal{E}_n)^c)$.

%
%\medskip
%

%We first discuss about the regularization error \eqref{Eq:RegError}. Meyers' estimate, recalled in Theorem \ref{Meyers}, ensures that there exists $\bar q>2$ such that $\nabla f^{n,t}\in\LL^{\bar q}(\mathcal{M})$. This allow us to split \eqref{Eq:RegError} by
%
%\begin{equation}\label{Eq:RegError2}
%\int_{\mathcal{M}}\vert\nabla f^{n,t}_\delta-\nabla f^{n,t}\vert^2\lesssim \|\rho_\delta-\rho\|^2_{\LL^{\frac{\bar q}{\bar q-2}}(\mathcal{M})}\|\nabla f^{n,t}\|^2_{\LL^{\bar q}(\mathcal{M})},
%\end{equation}
%
%where, since $\rho\in\LL^{\infty}(\mathcal{M})$
%
%$$\|\rho_\delta-\rho\|_{\LL^{q}(\mathcal{M})}\underset{\delta\downarrow 0}{\rightarrow} 0\quad\text{for any $q<\infty$}.$$
%
%Note that if $\rho\in \cc^{0}(\mathcal{M})$, we can also take $q=2$. We then show in Proposition \ref{prop1} that the moments of $\|\nabla f^{n,t}\|^2_{\LL^{\bar q}(\mathcal{M})}$ are optimally controlled yielding, together with \eqref{Eq:RegError2}, the the first term in \eqref{Eq78}.

%
%\medskip
%

%Second, we discuss about the two terms in \eqref{Eq:Strategy1}. For the first quantity it has been well studied in recent work \cite{ambrosio2019pde} and \cite{ambrosio2021quadratic} in the case of \iid random points and in \cite{borda2021empirical} in the case of correlated points. The second term requires more attention. As in \cite{ambrosio2019finer}\nc{write the paragraph}

%
%\medskip
%

%Finally, we discuss about the bound on $\mathbb{P}((\mathcal{E}_n)^c)$, where $\mathcal{E}_n$ is defined in \eqref{Eq:Strategy2}. Compared to \cite{ambrosio2019finer}, we can not rely directly on the heat-kernel representation, due to the regularization process on $\rho$ and the choice of the the heat-kernel for $-\delta$ for the regularization of $\mu^n$. Also, using directly estimates on the heat-kernel for $-\nabla\cdot\rho_\delta\nabla$ would make difficult to track the dependences on $\delta$. Instead, we split the equation \eqref{eq:f^nt_delta} into two parts
%
%$$-\Delta f^{n,t}_\delta=\tfrac{1}{\rho_\delta}(\mu^{n,t}-\rho_t)+\tfrac{1}{\rho_\delta}\nabla\rho_\delta\cdot\nabla f^{n,t}_\delta.$$
%
%We then decompose $f^{n,t}=u^{n,t}_\delta+v^{n,t}_\delta$, where 
%
%\begin{equation}\label{Eq:Strategy3}
%-\Delta u^{n,t}_\delta=\tfrac{1}{\rho_\delta}(\mu^{n,t}-\rho_t)\quad\text{and}\quad -\Delta v^{n,t}_\delta=\tfrac{1}{\rho_\delta}\nabla\rho_\delta\cdot\nabla f^{n,t}_\delta.
%\end{equation}
%
%Using Schauder's theory and the choice $\delta=\frac{1}{\log^\gamma(n)}$, for a well chosen $\gamma>0$, we then show that 
%
%$$\underbrace{\{\|\nabla^2 u^{n,t}_\delta\|_{\LL^{\infty}(\mathcal{M})}\leq \tfrac{1}{\log(n)}\}}_{:=\mathcal{A}_n}\cap\underbrace{\{\|\mu^{n,t}-\rho_t\|_{\LL^{\infty}(\mathcal{M})}\leq \tfrac{1}{\log(n)}\}}_{:=\mathcal{B}_n}\subset \mathcal{E}_n,$$
%
%see \eqref{Eq6} in Proposition \ref{Fluctuation}. We finally show, based on Bernstein's type-inequalities, recalled in Proposition \ref{Berstein} and Proposition \ref{BersteinCorrelated}, that $\mathbb{P}(\mathcal{A}^c_n\cup\mathcal{B}^c_n)$ enjoys exponential type decay estimates, see \eqref{Eq13} and \eqref{Eq76} in Proposition \ref{Fluctuation}.
%
%\end{enumerate}
%\nc{comment more on the correlated case}
\section{Proofs}
\subsection{Notations and preliminary results}\label{sec:preliminary}
%
We provide in this section some notations and preliminary results needed in the proofs of Theorem \ref{MainResultMatchingClouds} and Theorem \ref{th1}. We recall that throughout the paper, we denote by $\mathcal{M}$ a $2$-dimensional compact connected Riemannian manifold (or the square $[0,1]^2$) endowed with the Riemannian distance $\dd$.
%

\medskip
\noindent
\textbf{Wasserstein distance. }Given $\mu, \nu\in\mathcal{P}(\mathcal{M})$, we define the quadratic Wasserstein distance as
%
\begin{equation}\label{Wasserstein}
%
W_2^2(\mu, \nu):= \min_{\pi \in \Pi(\mu, \nu)} \int_{\mathcal{M}\times \mathcal{M}} \dd^2(x,y) \,\dd \pi (x,y),
%
\end{equation}
%
where $\Pi(\mu, \nu)$ is the set of couplings between $\mu$ and $\nu$, that is the set of $\pi\in\mathcal{P}(\mathcal{M}\times \mathcal{M})$ having $\mu$ and $\nu$ as first and second marginal, respectively. We refer the reader to the monographs \cite{Viltop, OldNewVillani} for a detailed overview on the subject of optimal transport. % We recall that the Wasserstein distance admits the following Kantorovich duality formulation, see for instance \cite[Theorem 5.10]{OldNewVillani},
%
%\begin{equation}\label{eq:dualhamjac}
%
%\frac12W_2^2(\mu, \nu) = \sup_{\phi \in \Lip_b}\bigg\{ \int_{\mathcal{M}} Q_1 \phi \, \dd \nu-\int_{\mathcal{M}} \phi \,\dd \mu \bigg\},
%
%\end{equation}
%
%where $\Lip_b$ denotes the set of bounded Lipschitz functions on $\mathcal{M}$ and $Q_t \phi$ is defined by the Hopf-Lax formula
%
%\begin{equation}\label{eq:HopfLax}
%
%Q_t \phi(y) = \inf_{x \in \mathcal{M}} \bigg\{ \phi(x) + \frac1{2t} \dd^2(x,y) \bigg\}\quad\text{for any $t>0$ and $x\in\mathcal{M}$.}
%
%\end{equation}
%
%Using the duality formulation \eqref{eq:dualhamjac}, we prove a Lipschitz property of the Wasserstein metric. 
%
We recall the following simple, but useful Lipschitz contraction property of the Wasserstein distance. 
\begin{lemma}[Lipschitz property of the Wasserstein metric] \label{lem:transportcontraction}
Let $(D, \dd_D)$ be a complete and separable metric space, let $\mu, \nu\in\mathcal{P}(\mathcal{M})$ and let $T: \mathcal{M} \rightarrow D$ be a $L$-Lipschitz map. It holds
%
\begin{equation}\label{Eq:LipEstimateWasser}
W_2^2 (T_{\#}\mu,T_{\#} \nu) \le L^2 W_2^2 (\mu, \nu).
\end{equation}
%
\end{lemma}
\begin{proof}
For any coupling $\pi \in \Pi(\mu, \nu)$, the push-forward $(T,T)_\# \pi$ is a coupling between $T_\# \mu$ and $T_\# \nu$. Moreover, it holds
\[
\begin{split}
\int_{D\times D} \dd^2(x,y) \,\dd ((T,T)_\#\pi)(x,y) & = \int_{\mathcal{M}\times\mathcal{M}} \dd^2(T(x),T(y))\, \dd \pi(x,y)\\
&\le L^2 \int_{\mathcal{M}\times\mathcal{M}} \dd^2(x,y) \, \dd \pi(x,y).
\end{split}
\]
Taking the infimum among all possible couplings $\pi \in \Pi(\mu,\nu)$ leads to \eqref{Eq:LipEstimateWasser}.
\end{proof}
%
%\begin{proof}
%We prove \eqref{Eq:LipEstimateWasser} using the duality formula \eqref{eq:dualhamjac}. We first claim that for any bounded and continuous function $\phi$ on $\mathcal{M}$
%
%\begin{equation}\label{eq:HamJacLip}
%Q_t \phi (T (y)) \le L^2 Q_t  \big(\tfrac{1}{L^2} \phi\circ T\big) (y)\quad \text{for any $t\geq 0$ and $y\in \mathcal{M}$}.
%\end{equation}
%
%where $Q_t$ denotes the Hopf-Lax semigroup
%
%$$Q_t \phi:=\inf_{x\in D}\{\phi(x)+\tfrac{1}{2t}\dd^2(x,\cdot)\}.$$
%
%Indeed, by the definition \eqref{eq:HopfLax} and the Lipschitz property of $T$, we have for any $z\in \mathcal{M}$
%\[
%\begin{split}
%Q_t \phi (T (y))& \le \phi (T (z)) + \tfrac1{2t} \dd^2_D(T (z),T (y))\le L^2 ( \tfrac1{L^2} \phi(T(z)) + \tfrac1{2t} \dd^2(z,y)),
%\end{split}
%\]
%which yields \eqref{eq:HamJacLip} by arbitrariness of $z$. We then deduce from \eqref{eq:HamJacLip} that for any Lipschitz functions $\phi$ on $\mathcal{M}$
%
%\begin{align*}
%
%\int_{\mathcal{M}}Q_1\phi\,\dd T_{\#}\nu-\int_{\mathcal{M}}\phi\, \dd T_{\#}\mu&=\int_{\mathcal{M}}Q_1\phi\circ T\,\dd\nu-\int_{\mathcal{M}}\phi\circ T\,\dd\mu\\
%&\stackrel{\eqref{eq:HamJacLip}}{\leq}L^2\Big(\int_{\mathcal{M}}Q_1(\tfrac{1}{L^2}\phi\circ T)\,\dd \nu-\int_{\mathcal{M}}\tfrac{1}{L^2}\phi\circ T\,\dd \mu\Big)\\
%&\leq \frac{L^2}{2}W^2_2(\mu,\nu),
%
%\end{align*}
%
%which gives \eqref{Eq:LipEstimateWasser} using the duality formula \eqref{eq:dualhamjac}.
%\end{proof}

\medskip
\noindent
\textbf{Heat semigroup and heat kernel. }We recall some basic facts on the heat semigroup and its generator, we refer the reader to \cite[Chapter VI]{chavel1984eigenvalues} for a more detailed overview. For $t>0$, we denote by $p_t$ the fundamental solution of the heat operator $\partial_t-\Delta$ on $\mathcal{M}$, where $\Delta$ denotes the Beltrami-Laplace operator. Classical Schauder's theory ensures that $p_t$ is smooth and it is known, see for instance \cite{stroock1998upper} and \cite[Appendix B]{ambrosio2019finer}, that $p_t$ and its derivatives satisfy for some $C>0$ depending on $\mathcal{M}$
%
\begin{equation}\label{Eq28}
%
\vert\nabla^N p_{t}(x,y)\vert\lesssim t^{-1-\frac{N}{2}}\exp\Big(-\tfrac{1}{C}\tfrac{\dd^2(x,y)}{t}\Big)\quad\text{for any $t>0$, $N\geq 1$ and $x,y\in\mathcal{M}$.}
\end{equation}
%
%It is a well known fact that the generator of $(P_t)_{t\ge0}$ is (half) the Laplace--Beltrami operator $\frac12 \Delta$. The crucial property of the action of $P_t$ is that it transforms singular measures $\mu$ into smooth densities. 
The kernel $p_t$ admits the spectral decomposition
%
\begin{equation}\label{SpectralDecompoHeatKernel}
%
p_{t}(x,y):=\sum_{k\geq 1}e^{-t\lambda_k}\phi_k(x)\phi_k(y)\quad \text{for any $t>0$ and $x,y\in\mathcal{M}$},
%
\end{equation}
%
converging in $\LL^2(\mathcal{M}\times\mathcal{M})$, where we recall that $\{\lambda_k,\phi_k\}$ denotes the eigenvalues and eigenvectors of $-\Delta$ on $\mathcal{M}$. Specifying \eqref{SpectralDecompoHeatKernel} on the diagonal and using $\|\phi_k\|_{\mathrm{L}^2(\mathcal{M})}=1$, we obtain the trace formula
%
\begin{equation}\label{eq:traceformula}
\sum_{k\geq 1}e^{-t\lambda_k}=\int_{\mathcal{M}}p_t(x,x)\dd\m(x)\quad\text{for any $t>0$.}
\end{equation}
%
We recall that $\{\phi_k\}_{k\geq 1}$ can be estimated in terms of the eigenvalues,
%
\begin{equation}\label{eq:eigvaleighfunc}
%
 \|\phi_k\|_{\LL^{\infty}}\lesssim \lambda^{\frac{1}{2}}_k.
%
\end{equation}
%
We briefly recall the argument. Applying the Gagliardo-Nirenberg's interpolation inequality \cite[Theorem 3.70]{aubin1998some}, it holds
%
$$\|\phi_k\|_{\LL^{\infty}}\lesssim \|\nabla^2 \phi_k\|^{\frac{1}{2}}_{\LL^{2}}\|\phi_k\|^{\frac{1}{2}}_{\LL^2}=\|\nabla^2\phi_k\|^{\frac{1}{2}}_{\LL^2}.$$
%
In combination with $-\Delta\phi_k=\lambda_k\phi_k$ and elliptic regularity, in form of
%
$$\|\nabla^2\phi_k\|_{\LL^2}\lesssim \lambda_k\|\phi_k\|_{\LL^2}=\lambda_k,$$
%
we obtain \eqref{eq:eigvaleighfunc}. 

%
\medskip
%

We recall that $(\text{P}_t)_{t>0}$ admits the spectral gap property, that is there exists a constant $C_{\text{sg}}>0$ such that 
%
\begin{equation}\label{SpectralGap}
%
\|\pp_t f\|_{\mathrm{L}^2(\mathcal{M})} \le e^{-C_{\text{sg}} t}\|f\|_{\mathrm{L}^2(\mathcal{M})}\quad\text{for any $f \in \mathrm{L}^2$ with $\int_{\mathcal{M}} f\, \dd \m =0$.}
%
\end{equation}
%
Note that equivalently, \eqref{SpectralGap} can be formulated in terms of the eigenvalues $\{\lambda_k\}_{k\geq 1}$ in form of 
%
\begin{equation}\label{SpectralGapEigenvalues}
%
\inf_{k\geq 1}\lambda_k\geq C_{\text{sg}},
%
\end{equation}
%
simply by specifying \eqref{SpectralGap} on $\{\phi_k\}_{k\geq 1}$. Finally, we recall the Riesz-transform bound
%
\begin{equation}\label{eq:Riesz}
\int_{\mathcal{M}} |\nabla f|^4\,\dd\m  \lesssim \int_{\mathcal{M}}\big|(-\Delta)^\frac12 f\big|^4\,\dd\m\quad\text{for any $f$ such that $\int_{\mathcal{M}}\vert\nabla f\vert^4\,\dd\m<\infty$},
\end{equation}
%
where $(-\Delta)^{\frac{1}{2}}$ can be defined via its spectral representation
%
\begin{equation}\label{Laplacian12}
%
(-\Delta)^{\frac{1}{2}}f=\sum_{k\geq 1} \sqrt{\lambda_k}\,\big(\phi_k,f\big)_{\LL^2}\phi_k\quad\text{for any $f\in \mathrm{H}^{\frac{1}{2}}$,}
%
\end{equation}
%
with $(\cdot,\cdot)_{\LL^2}$ the inner product in $\LL^2$. We refer to the monograph \cite{wangbook} for a discussion of the inequalities \eqref{SpectralGap} and \eqref{eq:Riesz}, see Chapter 1 for the case of a Riemannian manifold without boundary and Chapter 2 for the case of a Riemannian manifold with (convex) boundary. In connection with the Wasserstein metric, the heat semigroup satisfies the following classical contraction property.
\begin{lemma}[Semigroup contraction for absolutely continuous measures]\label{lem:HeatContractionrho}
Let $\rho \in\LL^\infty$ satisfying \eqref{Ellipticity}. Given $\rho_t:=\mathrm{P}_t\rho$, it holds
%
\begin{equation}\label{Eq80}
W_2^2 (\rho_t\,\dd\m, \rho\,\dd\m) \lesssim t \| \rho_t - \rho\|_{\LL^1}\quad\text{for any $t>0$.}
\end{equation}
%
\end{lemma}
\begin{proof}
Using $g$ defined via $-\Delta g=\rho_t-\rho$ together with \eqref{Ellipticity}, \cite[Corollary 4.4]{ambrosio2019finer} yields
%
$$W^2_2(\rho_t\,\dd\m,\rho\,\dd\m)\lesssim \int_{\mathcal{M}}\vert\nabla g\vert^2\,\dd\m=\int_{\mathcal{M}}(\rho_t-\rho)g\,\dd\m.$$
%
Writing $g=\int_{0}^{\infty}\text{P}_\tau(\rho_t-\rho)\,\dd\tau=-\int_{0}^t \rho_\tau\,\dd\tau$ together with $\vert\int_{0}^t \rho_\tau\,\dd\tau\vert\leq \|\rho\|_{\LL^{\infty}} t$ gives \eqref{Eq80}.
%
\end{proof}

\subsection{\texorpdfstring{$\LL^q$}{Lq}-type estimates}
%
As we have seen in \eqref{LqControlIntro}, we need a sharp control of the averaged Meyers' norm $\mathbb{E}\big[\|\nabla h^{n,t}\|^2_{\LL^{\bar q}}\big]$. This will be obtained as an immediate corollary of the following proposition, for more details see \eqref{LqEstihnt}.
%
\begin{proposition}[$\LL^{q}$-estimates]\label{prop1}
%
Let $\{\mu^n\}_n$ be defined in \eqref{eq:empmeas} with point clouds satisfying Assumption \ref{Assumptions}. Let $\bar q$ be the Meyers exponent given in Theorem \ref{Meyers} for the operator $-\nabla\cdot\rho\nabla$. The solution\footnote{with Neumann boundary conditions in case $\mathcal{M}$ has a boundary} $f^{n,t}\in {\rm \dot{H}^1}$ of 
%
\begin{equation}\label{eq:f^nt}
-\nabla\cdot \rho\nabla f^{n,t} =\mu^{n,t}-\rho_t,
\end{equation}
%
satisfies:
%
\begin{equation}\label{Eq22Bis}
\Big(\int_\mathcal{M} \vert\nabla f^{n,t}\vert^{q}\,\dd\m\Big)^\frac{2}{q}\leq \mathcal{C}_{n,t}\frac{\vert\log(t)\vert+\log^{\frac{1}{\eta}}(n)}{n}\quad \text{for any $q\in [2,\min\{\bar q,4\}]$},
\end{equation}
%
and a random variable $\mathcal{C}_{n,t}$ satisfying for some generic constant $C_q$ depending on $q$,
%
\begin{equation}\label{MomentBoundCn}
\sup_{n,t}\mathbb{E}[\tfrac{1}{C_{q}}\mathcal{C}_{n,t}]\leq 1.
\end{equation}
%

%
\medskip
%

Furthermore, if \eqref{DecayAlphaMixing} holds with $\eta\geq 1$ then the assumption \eqref{DecatBetaMixing} can be dropped and the stochastic integrability can be improved up to losing a $\log(n)$ factor, namely
%
\begin{equation}\label{Eq22}
\begin{split}
\lefteqn{\Big(\int_\mathcal{M} \vert\nabla f^{n,t}\vert^{q}\,\dd\m\Big)^\frac{2}{q}}\\ &\leq \mathcal{D}_{n,t}\Big(\frac{\log^{\frac
{1}{\eta}}(n)\vert\log(t)\vert}{n}+\frac{t^{-1}(1+\log^2(n)\mathds{1}_{\eta\neq \infty})}{n^2}\Big)\quad \text{for any $q\in [2,\min\{\bar q,4\}]$},
\end{split}
\end{equation}
%
and a random variable $\mathcal{D}_{n,t}$ satisfying for some generic constant $D_q$ depending on $q$,
%
$$\sup_{n,t}\mathbb{E}[\exp(\tfrac{1}{D_{q}}\mathcal{D}^{\frac{1}{2}}_{n,t})]\leq 2.$$
%
\begin{proof}
%
We proceed in four steps. In the first step, we prove a representation formula for $(-\Delta)^{\frac{1}{2}}$ that we will use as the core tool in the next steps. In the second step, we compare the two operators $-\nabla\cdot \rho\nabla$ and $-\Delta$, with help of Meyers' estimate recalled in Theorem \ref{Meyers}. Doing so, we then have to bound the $\LL^{q}$-norms ($2\leq q\leq\bar q$) of the gradient of the solution\footnote{which belongs to any $\LL^{q}$ for any $q<\infty$ from Calder\'on-Zygmund' theory, see for instance \cite{giaquinta2013introduction}} to the Poisson equation with r.h.s. $\mu^{n,t}-\rho_t$ and Neumann boundary conditions. We control all the norms by the $\LL^{4}$-norm that, in turn, we estimate using the Riesz-transform bound \eqref{eq:Riesz} and following some ideas from \cite[Lemma 3.17]{ambrosio2019pde}. In the third and fourth steps, we control the bound previously obtained in expectation where our main tool is Assumption \ref{Assumptions} and the concentration inequalities in Proposition \ref{BersteinCorrelated}.

%
\medskip
%

{\sc Step 1. A representation formula for $(-\Delta)^{\frac{1}{2}}$. }
%
We show that given $f\in \cc^2$ such that $n_{\mathcal{M}}\cdot \nabla f=0$ on $\partial \mathcal{M}$ we have
%
\begin{equation}\label{Eq45}
(-\Delta)^{\frac{1}{2}} f=\frac{1}{\sqrt{\pi}}\int_{0}^\infty\tau^{-\frac{1}{2}}\Delta \text{P}_{\tau}f\,\dd\tau.
\end{equation}
%
%We first recall that $(-\Delta)^{\frac{1}{2}}f$ can be expressed in terms of the eigenvalues and eigenfunctions $\{\lambda_n,\phi_n\}_n$ of $-\Delta$ in form of
%
%\begin{equation}\label{Laplacian12}
%
%(-\Delta)^{\frac{1}{2}}f=\sum_n \sqrt{\lambda_n}\,\big(\phi_n,f\big)_{\LL^2}\phi_n,
%
%\end{equation}
%
%where $(\cdot,\cdot)_{\LL^2}$ denotes the inner product in $\LL^2$. 
Note that $(-\Delta)^{\frac{1}{2}}f$, defined in \eqref{Laplacian12}, is well defined in $\LL^2$. Indeed, using the fact that $(\phi_k,f)_{\LL^2}=\frac{1}{\lambda_k}(\phi_k,\Delta f)_{\LL^2}$ and \eqref{SpectralGapEigenvalues}, we have for any $N\leq M<\infty$
%
\begin{align*}
%
\bigg\|\sum_{N\leq n\leq M}\sqrt{\lambda_n}\,\big(\phi_n,f)_{\LL^2}\phi_n\bigg\|^2_{\LL^2}=\sum_{N\leq n\leq M}\lambda_n\vert (\phi_n,f)_{\LL^2}\vert^2=&\sum_{N\leq n\leq M}\tfrac{1}{\lambda_n}\vert (\phi_n,\Delta f)_{\LL^2}\vert^2\\
\leq & \tfrac{1}{C_{\text{sg}}}\sum_{n\geq N}\vert (\phi_n,\Delta f)_{\LL^2}\vert^2,
%
\end{align*}
%
which vanishes as $N\uparrow\infty$ uniformly in $M$.

%
\medskip
%

We now justify \eqref{Eq45}. Observe that since $n_\mathcal{M}\cdot \nabla f=0$ on $\partial\mathcal{M}$,
%
\begin{equation}\label{Eq44}
\Delta \text{P}_s f=\text{P}_s\Delta f\quad \text{for any $s\in (0,\infty)$,}
\end{equation}
%
which is a direct consequence of two integration by parts using the heat-kernel representation $\text{P}_s f=\int_{\mathcal{M}}p_s(\cdot,y)f(y)\,\dd \m(y)$. Therefore 
%
\begin{equation}\label{Laplacian12:Eq1}
%
\int_{0}^{\infty}\tau^{-\frac{1}{2}}\Delta {\rm P}_{\tau} f\,\dd\tau=\int_{0}^{\infty}\tau^{-\frac{1}{2}}{\rm P}_{\tau} \Delta f\,\dd\tau,
%
\end{equation}
%
where the last integral is well-defined in $\LL^2$ since from \eqref{SpectralGap} 
%
$$\int_0^\infty t^{-\frac{1}{2}}\|\pp_\tau \Delta f\|_{\LL^2}\,\dd\tau\leq\int_0^\infty t^{-\frac{1}{2}}e^{-C_{sg}\tau}\|\Delta f\|_{\LL^2}\,\dd\tau<\infty.$$
%
We then use the spectral decomposition of the heat semigroup \eqref{SpectralDecompoHeatKernel} to get
%
\begin{equation}\label{Laplacian12:Eq2}
%
\pp_\tau \Delta f=\sum_{n} e^{-\lambda_n\tau}\big(\phi_n,\Delta f\big)_{\LL^2}\phi_n\quad\text{in $\LL^2$.}
%
\end{equation}
%
%We recall the following formula for the trace, see \cite[Chapter 6]{chavel1984eigenvalues},
%
%$$\sum_n e^{-\lambda_n \tau}=\int_{\mathcal{M}}\dd x\, p_{\tau}(x,x).$$
%
The combination of \eqref{Laplacian12:Eq1} and \eqref{Laplacian12:Eq2} yields for any $\eta\in \LL^2$
%
\begin{equation}\label{Laplacian12:Eq3}
%
\begin{aligned}
%
\bigg(\int_{0}^{\infty}\tau^{-\frac{1}{2}}\Delta {\rm P}_{\tau} f,\eta\bigg)_{\LL^{2}}\,\dd\tau=&\int_{0}^\infty \tau^{-\frac{1}{2}}\big(\pp_\tau\Delta f,\eta\big)_{\LL^2}\,\dd\tau\\
\stackrel{\eqref{Laplacian12:Eq2}}{=}&\int_{0}^{\infty}\tau^{-\frac{1}{2}}\sum_{n} e^{-\lambda_n\tau}\big(\phi_n,\Delta f\big)_{\LL^2}\big(\phi_n,\eta\big)_{\LL^{2}}\,\dd\tau.
%
\end{aligned}
%
\end{equation}
%
Using \eqref{SpectralGapEigenvalues}, we have
%
\begin{align*}
%
&\int_{0}^\infty \tau^{-\frac{1}{2}}\sum_n e^{-\lambda_n\tau}\vert\big(\phi_n,\Delta f\big)_{\LL^2}\vert\vert \big(\phi_n,\eta\big)_{\LL^2}\vert\,\dd\tau\\
&\leq \Big(\int_{0}^\infty  \tau^{-\frac{1}{2}}e^{-C_{sg}\tau}\,\dd\tau\Big)\sum_{n}\vert\big(\phi_n,\Delta f\big)_{\LL^2}\vert\vert \big(\phi_n,\eta\big)_{\LL^2}\vert<\infty,
%
\end{align*}
%
so that we can exchange integration and summation in \eqref{Laplacian12:Eq3} to obtain
%
\begin{align*}
%
\bigg(\int_{0}^{\infty}\tau^{-\frac{1}{2}}\Delta {\rm P}_{\tau} f,\eta\bigg)_{\LL^{2}}\,\dd\tau&=\sum_n\Big(\int_{0}^{\infty}\tau^{-\frac{1}{2}}e^{-\lambda_n\tau}\,\dd\tau\Big)\big(\phi_n,\Delta f\big)_{\LL^2}\big(\phi_n,\eta\big)_{\LL^2}\\
&=\sqrt{\pi}\sum_n \tfrac{1}{\sqrt{\lambda_n}}\big(\phi_n,\Delta f\big)_{\LL^2}\big(\phi_n,\eta\big)_{\LL^2}\\
&=\sqrt{\pi}\sum_n \sqrt{\lambda_n}\,\big(\phi_n, f\big)_{\LL^2}\big(\phi_n,\eta\big)_{\LL^2}\stackrel{\eqref{Laplacian12}}{=}\sqrt{\pi}\big((-\Delta f)^{\frac{1}{2}},\eta\big)_{\LL^2},
%
\end{align*}
%
which gives \eqref{Eq45} by arbitrariness of $\eta$. Finally, note that the \rhs integral in \eqref{Eq45} is absolutely convergent thanks to the integration by parts $\Delta \pp_\tau f=\pp_{\tau}\Delta f$ and the bounds on the heat kernel \eqref{Eq28}, so that it defines a function in $\cc^0$.
%
\medskip
%


{\sc Step 2. Comparison with the solution of the Poisson equation. }We claim that for any $q\in [2,\min\{\bar q,4\}]$ and $p<\infty$
%
\begin{equation}\label{Eq51}
\begin{aligned}
%
\mathbb{E}\bigg[\Big(\int_{\mathcal{M}}\vert \nabla f^{n,t}\vert^q\,\dd\m\Big)^{\frac{2p}{q}}\bigg]^{\frac{1}{p}}\lesssim_q \bigg(\int_{\mathcal{M}}\dd\m\bigg(\int_{0}^{\infty}\dd s\,\mathbb{E}\Big[\Big((-s\Delta)^{\frac{1}{2}}\text{P}_{s+t}(\mu^n-\rho)\Big)^{2p}\Big]^{\frac{1}{p}}\bigg)^2\bigg)^{\frac{1}{2}}.
%
\end{aligned}
\end{equation}
%
Let $g^{n,t}\in {\rm \dot{H}^1}$ be the solution of the following Poisson equation 
%
\begin{equation}\label{Eq47}
        -\Delta g^{n,t} =\mu^{n,t}-\rho_t.
\end{equation}
%
Re-expressing the right-hand side of \eqref{eq:f^nt} as $\nabla\cdot \nabla g^{n,t}$, we apply Meyers' estimate recalled in Theorem \ref{Meyers} and Hölder's inequality to obtain: 
%
\begin{equation}\label{Eq41}
\Big(\int_\mathcal{M} \vert\nabla f^{n,t}\vert^{q}\,\dd\m\Big)^\frac{2}{q}\lesssim \Big(\int_\mathcal{M} \vert\nabla g^{n,t}\vert^{q}\,\dd\m\Big)^\frac{2}{q}\lesssim_q\Big(\int_\mathcal{M} \vert\nabla g^{n,t}\vert^{4}\,\dd\m\Big)^\frac{1}{2}\quad\text{for any $q\in [2,\min\{\bar q,4\}]$,}
\end{equation}
%
We now introduce the Paley-Littlewood functional
%
$$\mathcal{L}(g):=\Big(\int_{0}^{\infty}s(\partial_s\text{P}_s g)^2\,\dd s\Big)^{\frac{1}{2}}\quad \text{for any $g\in \LL^{4}$ and $\int_{\mathcal{M}} g=0$.}$$
%
We recall that the inverse of $\mathcal{L}$ is continuous, see \cite{stein2016topics}, namely
%
\begin{equation}\label{Eq40}
\|g\|_{\LL^{4}}\lesssim \|\mathcal{L}(g)\|_{\LL^{4}}.
\end{equation}
%
Combining the Riesz transform bound \eqref{eq:Riesz} with \eqref{Eq40} yields
%
\begin{equation}\label{Eq50}
\begin{aligned}
\int_{\mathcal{M}}\vert\nabla g^{n,t}\vert^4\,\dd\m\lesssim\int_{\mathcal{M}}\Big(\mathcal{L}\Big((-\Delta)^{\frac{1}{2}}g^{n,t}\Big)\Big)^4\dd\m
=\int_{\mathcal{M}}\dd\m\Big(\int_{0}^{\infty}\dd s\, s(\partial_s\text{P}_s(-\Delta)^{\frac{1}{2}} g^{n,t})^2\Big)^{2}.
\end{aligned}
\end{equation}
%
We now claim that
%
\begin{equation}\label{Eq43}
\partial_s\text{P}_s(-\Delta)^{\frac{1}{2}}g^{n,t}=(-\Delta)^{\frac{1}{2}}\text{P}_{s+t}(\mu^n-\rho)\quad\text{for any $s\geq 0$,}
\end{equation}
%
which requires a special care when $\mathcal{M}$ has a boundary. We use the definition of $\text{P}_s$ in form of $\partial_s\text{P}_s=\Delta\text{P}_s$  to get
%
\begin{equation}\label{Eq46}
\partial_s\text{P}_s(-\Delta)^{\frac{1}{2}}g^{n,t}=\Delta\text{P}_s(-\Delta)^{\frac{1}{2}}g^{n,t}.
\end{equation}
% 
Recalling that $n_{\mathcal{M}}\cdot \nabla g^{n,t}=0$, \eqref{Eq44} implies that $\Delta\text{P}_{\tau}g^{n,t}=\text{P}_{\tau}\Delta g^{n,t}$ which, combined with \eqref{Eq47} and \eqref{Eq45}, gives
%
\begin{equation}\label{Eq48}
(-\Delta)^{\frac{1}{2}} g^{n,t}=\frac{1}{\sqrt{\pi}}\int_{0}^\infty \tau^{-\frac{1}{2}}\text{P}_{\tau}(\mu^{n,t}-\rho_t)\,\dd\tau.
\end{equation}
%
In particular, it implies that $n_{\mathcal{M}}\cdot \nabla(-\Delta)^{\frac{1}{2}}g^{n,t}=0$. Therefore, one can use once more \eqref{Eq44} and, together with \eqref{Eq48}, \eqref{Eq46} turns into
%
\begin{equation}\label{Laplace12:Eq5}
%
\partial_s\text{P}_s(-\Delta)^{\frac{1}{2}}g^{n,t}=\frac{1}{\sqrt{\pi}}\int_{0}^\infty\tau^{-\frac{1}{2}}\text{P}_s\Delta\text{P}_{\tau}(\mu^{n,t}-\rho_t)\,\dd\tau.
%
\end{equation}
%
Using a last time \eqref{Eq44} in form of $\text{P}_s\Delta\text{P}_{\tau}(\mu^{n,t}-\rho_t)=\Delta\text{P}_s\text{P}_{\tau}(\mu^{n,t}-\rho_t)$ that we combine with the semigroup property $\text{P}_{t}\text{P}_{t'}=\text{P}_{t+t'}$ yields
%
$$\text{P}_s\Delta\text{P}_{\tau}(\mu^{n,t}-\rho_t)=\Delta\text{P}_{\tau}\text{P}_{s+t}(\mu^{n}-\rho),$$
%
which, together with \eqref{Laplace12:Eq5} and \eqref{Eq45} leads to \eqref{Eq43}.

%
\medskip
%

The combination of \eqref{Eq41}, \eqref{Eq50} and \eqref{Eq43} together with Minkowski's inequality gives
%
\begin{equation*}
\begin{aligned}
%
\mathbb{E}\bigg[\Big(\int_{\mathcal{M}}\vert \nabla f^{n,t}\vert^q\,\dd\m\Big)^{\frac{2p}{q}}\bigg]^{\frac{1}{p}}\stackrel{\eqref{Eq41}}{\lesssim}&\mathbb{E}\bigg[\bigg(\int_{\mathcal{M}}\dd\m\bigg(\int_{0}^{\infty}\dd s\,\Big((-s\Delta)^{\frac{1}{2}}\text{P}_{s+t}(\mu^n-\rho)\Big)^2\bigg)^2\bigg)^{\frac{p}{2}}\bigg]^{\frac{1}{p}}\\
\leq& \bigg(\int_{\mathcal{M}}\dd\m\bigg(\int_{0}^{\infty}\dd s\,\mathbb{E}\Big[\Big((-s\Delta)^{\frac{1}{2}}\text{P}_{s+t}(\mu^n-\rho)\Big)^{2p}\Big]^{\frac{1}{p}}\bigg)^2\bigg)^{\frac{1}{2}},
%
\end{aligned}
\end{equation*}
%
which is \eqref{Eq51}.

%
\medskip
%

{\sc Step 3. Proof of \eqref{Eq22Bis}. }The estimate \eqref{Eq22Bis} is a consequence of \eqref{Eq51} applied with $p=1$ and 
%
\begin{equation}\label{Eq54Bis}
\mathbb{E}\Big[\Big((-s\Delta)^{\frac{1}{2}}\text{P}_{s+t}(\mu^n-\rho)(x)\Big)^2\Big]\lesssim \frac{\zeta(s,t)}{n}(1+\log^{\frac{1}{\eta}}(n)s^{\frac{1}{\log(n)}}\zeta(s,t)^{1-\frac{1}{\log(n)}}) \quad\text{for any  $s\in (0,\infty)$,}
\end{equation}
%
with 
%
\begin{equation}\label{DefZeta}
%
\zeta(s,t):=\min\big\{(s+t)^{-1},(s+t)^{-2}\big\}.
%
\end{equation}
%
Indeed, plugging \eqref{Eq54Bis} in \eqref{Eq51} gives
%
\begin{equation}\label{Eq57Bis}
%
\mathbb{E}\bigg[\Big(\int_{\mathcal{M}}\vert\nabla f\vert^q\,\dd\m\Big)^{\frac{2}{q}}\bigg]\lesssim \frac{1}{n}\int_{0}^{\infty}\zeta(s,t)(1+\log^{\frac{1}{\eta}}(n)s^{\frac{1}{\log(n)}}\zeta(s,t)^{1-\frac{1}{\log(n)}})\,\dd s.
%
\end{equation}
%
Using that 
%
$$\int_{0}^{\infty}\zeta(s,t)\,\dd s\leq t^{-1}\int_{0}^t\dd s+\int_{t}^1s^{-1}\,\dd s+\int_{1}^{\infty}s^{-2}\,\dd s\lesssim \vert\log(t)\vert,$$
%
and analogously
%
$$\int_{0}^{\infty}s^{\frac{1}{\log(n)}}\zeta(s,t)^{1-\frac{1}{\log(n)}}\,\dd s\lesssim 1,$$
%
\eqref{Eq57Bis} implies \eqref{Eq22Bis}.

%
\medskip
%

We now show \eqref{Eq54Bis}. First, using the definition \eqref{eq:empmeas} of $\mu^n$, we have
%
\begin{equation}\label{ExpandSquareLaplace}
%
(-s\Delta)^{\frac{1}{2}}\text{P}_{s+t}(\mu^n-\rho)=\frac{1}{n}\sum_{k=1}^n\omega_{s+t}(\cdot,X_k)\quad\text{with}\quad\omega_{s+t}(\cdot,y):=(-s\Delta)^{\frac{1}{2}}\text{P}_{s+t}(\delta_y-\rho),
%
\end{equation}
%
such that expanding the square gives
%
\begin{align*}
%
\Big((-s\Delta)^{\frac{1}{2}}\text{P}_{s+t}(\mu^n-\rho)(x)\Big)^2=\frac{1}{n^2}\sum_{k=1}^n \big(\omega_{s+t}(x,X_k)\big)^2+\frac{2}{n^2}\sum_{1\leq \ell<\ell'\leq n}\omega_{s+t}(x,X_\ell)\omega_{s+t}(x,X_{\ell'}).
%
\end{align*}
%
The estimate \eqref{Eq54Bis} is then a consequence of
%
\begin{equation}\label{Eq53}
\mathbb{E}[\vert\omega_{s+t}(x,X_1)\vert^2]\lesssim \min\Big\{(s+t)^{-1},(s+t)^{-2}\Big\}\quad\text{and}\quad\sup_{y\in\mathcal{M}}\vert\omega_{s+t}(x,y)\vert\lesssim\sqrt{s} \min\Big\{(s+t)^{-\frac{3}{2}},(s+t)^{-2}\Big\}.
\end{equation}
%
Indeed, the first item of \eqref{Eq53} immediately implies
%
$$\mathbb{E}\bigg[\sum_{k=1}^n \big(\omega_{s+t}(\cdot,X_k)\big)^2\bigg]\lesssim n\zeta(s,t),$$
%
while, using the assumption \eqref{DecayAlphaMixing} and that $t^{-\frac{3}{\log(n)}}\lesssim 1$,
%
\begin{equation}\label{Eq53Bis}
%
\begin{aligned}
%
&\sum_{1\leq \ell<\ell'\leq n}\mathbb{E}[\omega_{s+t}(\cdot,X_\ell)\omega_{s+t}(\cdot,X_{\ell'})]\\
&=\sum_{1\leq \ell<\ell'\leq n}\mathbb{E}[\omega_{s+t}(\cdot,X_\ell)\omega_{s+t}(\cdot,X_{\ell'})]^{\frac{1}{\log(n)}}\mathbb{E}[\omega_{s+t}(\cdot,X_\ell)\omega_{s+t}(\cdot,X_{\ell'})]^{1-\frac{1}{\log(n)}}\\
&\leq \sum_{1\leq \ell<\ell'\leq n}\alpha^{\frac{1}{\log(n)}}_{\ell-\ell'}\big(\sup_{y\in\mathcal{M}}\vert \omega_{s+t}(x,y)\vert\big)^{\frac{2}{\log(n)}}\mathbb{E}[\vert \omega_{s+t}(x,X_1)\vert^2]^{1-\frac{1}{\log(n)}}\\
&\stackrel{\eqref{Eq53},\eqref{DecayAlphaMixing}}\lesssim n\log^{\frac{1}{\eta}}(n)s^{\frac{1}{\log(n)}}t^{-\frac{3}{\log(n)}}\zeta(s,t)^{1-\frac{1}{\log(n)}}\lesssim\log^{\frac{1}{\eta}}(n)s^{\frac{1}{\log(n)}}\zeta(s,t)^{1-\frac{1}{\log(n)}}.
%
\end{aligned}
%
\end{equation}

%
It remains to prove \eqref{Eq53} and we start with the first item. Here, we use the fact that  
%
\begin{equation}\label{UniformContinuityL2}
%
\|(-s\Delta)^{\frac{1}{2}}\text{P}_{\frac{s}{2}}f\|_{\LL^2}\lesssim \|f\|_{\LL^2}\quad\text{for all $f\in\cc^2$ and uniformly in $s$.}
%
\end{equation}
%
This can be seen from \eqref{Laplacian12}: Using that $\pp_{\frac{s}{2}}$ is an auto-adjoint operator in $\LL^2$ and that from \eqref{Laplacian12:Eq2} we learn $\pp_{\frac{s}{2}}\phi_n =e^{-\lambda_n \frac{s}{2}}\phi_n$, we have 
%
\begin{align*}
%
\|(-s\Delta)^{\frac{1}{2}}\text{P}_{\frac{s}{2}}f\|^2_{\LL^2}=\sum_n (s\lambda_n)\vert \big(\phi_n,\pp_{\frac{s}{2}} f\big)_{\LL^2}\vert^2=\sum_n(s\lambda_n)e^{-s\lambda_n }\vert \big(\phi_n,f\big)_{\LL^2}\vert^2\lesssim \|f\|^2_{\LL^2}.
%
\end{align*}
%
Hence, using the semi-group property $\text{P}_{s+t}=\text{P}_{\frac{s}{2}}\text{P}_{\frac{s}{2}+t}$, we deduce
%
\begin{align*}
\mathbb{E}[\vert\omega_{s+t}(x,X_1)\vert^2]\lesssim \|(-s\Delta)^{\frac{1}{2}}\text{P}_{s+t}(\delta_y-\rho)\|^2_{\LL^{2}}&=\|(-s\Delta)^{\frac{1}{2}}\text{P}_{\frac{s}{2}}\text{P}_{\frac{s}{2}+t}(\delta_y-\rho)\|^2_{\LL^{2}}\\
&\stackrel{\eqref{UniformContinuityL2}}{\lesssim} \|\text{P}_{\frac{s}{2}+t}(\delta_y-\rho)\|^2_{\LL^{2}}.
\end{align*}
%
We now bound $\|\text{P}_{\frac{s}{2}+t}(\delta_y-\rho)\|^2_{\LL^{2}}$ in two different ways. First, using the bounds \eqref{Eq28} of the heat-kernel, we have by the triangle inequality
%
$$\|\text{P}_{\frac{s}{2}+t}(\delta_y-\rho)\|^2_{\LL^{2}}\lesssim \|p_{\frac{s}{2}+t}(\cdot,y)\|^2_{\LL^{2}}\stackrel{\eqref{Eq28}}{\lesssim} (s+t)^{-1},$$
%
yielding to the first alternative in the first item of \eqref{Eq53}. Second, applying Poincaré's inequality yields
%
$$\|\text{P}_{\frac{s}{2}+t}(\delta_y-\rho)\|^2_{\LL^{2}}\lesssim \|\nabla p_{\frac{s}{2}+t}(\cdot,y)\|^2_{\LL^{2}}\stackrel{\eqref{Eq28}}{\lesssim} (s+t)^{-2},$$
%
yielding to the second alternative in the first item of \eqref{Eq53}.

%
\medskip
%

We now turn to the second item of \eqref{Eq53}. Here, we make use of the representation formula \eqref{Eq45} applied to $f=\pp_{s+t}(\delta_y-\rho)$. Using the semi-group property $\text{P}_{\tau}\text{P}_{s+t}=\text{P}_{\tau+s+t}$, this takes the form
%
$$\omega_{s+t}(x,y)=\sqrt{\frac{s}{\pi}}\int_{0}^{\infty}\tau^{-\frac{1}{2}}\Delta\text{P}_{\tau+s+t}(\delta_y-\rho)(x)\,\dd\tau.$$
%
A direct application of the heat-kernel bounds \eqref{Eq28} leads to
%
\begin{align*}
\sup_{y\in\mathcal{M}}\vert\omega_{s+t}(x,y)\vert&\lesssim\sqrt{s}\int_{0}^{\infty}\tau^{-\frac{1}{2}}(\tau+s+t)^{-2}\,\dd\tau\lesssim \sqrt{s}\,(s+t)^{-\frac{3}{2}},
\end{align*}
%
which is the first alternative in the second item of \eqref{Eq53}. For the second alternative, we write 
%
$$\Delta \text{P}_{\tau+s+t}(\delta_y-\rho)(x)=\int_{\mathcal{M}}\rho\big(\Delta p_{\tau+s+t}(x,y)-\Delta p_{\tau+s+t}(x,\cdot)\big)\,\dd\m,$$
%
so that, using \eqref{Eq28},
%
\begin{align*}
%
\sup_{y\in\mathcal{M}}\vert\omega_{s+t}(x,y)\vert&\lesssim \sqrt{s}\int_{0}^{\infty}\tau^{-\frac{1}{2}}\sup_{y\in\mathcal{M}}\vert\nabla \Delta p_{\tau+s+t}(x,y)\vert\,\dd\tau\\
&\stackrel{\eqref{Eq28}}{\lesssim}\int_{0}^{\infty}\tau^{-\frac{1}{2}}(\tau+s+t)^{-\frac{5}{2}}\,\dd\tau\lesssim \sqrt{s}\,(s+t)^{-2}.
%
\end{align*}
%
{\sc Step 4. Proof of \eqref{Eq22}. }Following the same computations done in the previous step, the estimate \eqref{Eq22} is a consequence of \eqref{Eq51} and the moment bounds
%
\begin{equation}\label{Eq54}
%
\begin{aligned}
&\mathbb{E}\Big[\Big((-s\Delta)^{\frac{1}{2}}\text{P}_{s+t}(\mu^n-\rho)(x)\Big)^{2p}\Big]^{\frac{1}{p}}\\
&\lesssim p^2\bigg(\frac{\zeta(s,t)}{n}(1+\log^{\frac{1}{\eta}}(n)s^{\frac{1}{2\log(n)}}\zeta(s,t)^{1-\frac{1}{\log(n)}})+\frac{\sup_{y\in\mathcal{M}}\vert \omega_{s+t}(x,y)\vert+\log^2(n)}{n^2}\bigg) \quad\text{for any  $s\in (0,\infty)$,}
\end{aligned}
%
\end{equation}
%
together with the second item of \eqref{Eq53} and Lemma \ref{momentexp}, where we recall that $\zeta$ is defined in \eqref{DefZeta}.

%
\medskip
%

We now show \eqref{Eq54}. Using \eqref{ExpandSquareLaplace} and the assumption $\eta\geq 1$, we apply the concentration inequality in Proposition \ref{BersteinCorrelated} to the effect of: for any $\lambda\geq 0$
%
\begin{equation}\label{BoundProbaTailLqEsti}
%
\begin{aligned}
%
&\mathbb{P}\bigg(\Big\vert(-s\Delta)^{\frac{1}{2}}\text{P}_{s+t}(\mu^n-\rho)(x)\Big\vert\geq \lambda\bigg)\\
&\lesssim \exp\bigg(-\frac{1}{C}\frac{n^2\lambda^2}{nv^2+\big(\sup_{y\in\mathcal{M}}\vert\omega_{s+t}(x,y)\vert\big)^2+n\lambda\big(\sup_{y\in\mathcal{M}}\vert\omega_{s+t}(x,y)\vert\big)\log^2(n)}\bigg),
%
\end{aligned}
%
\end{equation}

%
with 
%
$$v^2:=\mathbb{E}[\vert\omega_{s+t}(x,X_1)\vert^2]+2\sum_{1\leq \ell<\ell'\leq n}\big\vert\mathbb{E}[\omega_{s+t}(x,X_\ell)\omega_{s+t}(x,X_{\ell'})]\big\vert.$$
%
We then obtain \eqref{Eq54} combining \eqref{BoundProbaTailLqEsti} with \eqref{Eq53} and \eqref{Eq53Bis}, together with an application of the Layer-cake formula.
%
%We assume that $(X_k)_{k\in\mathbb{N}^*}$ are i.~i.~d. The estimate \eqref{Eq22} is a consequence of \eqref{Eq51} and
%
%\begin{equation}\label{Eq54}
%\mathbb{E}\Big[\Big((-s\Delta)^{\frac{1}{2}}\text{P}_{s+t}(\mu^n-\rho)\Big)^{2p}\Big]^{\frac{1}{p}}\lesssim p\,\frac{\eta(s,t)}{n}\Big(1+\frac{\eta(s,t)}{n}\Big) \quad\text{for any  $s\in (0,\infty)$,}
%\end{equation}
%
%with
%
%\begin{equation}\label{Eq56}
%\eta(s,t):=\min\Big\{(s+t)^{-1},(s+t)^{-2}\Big\}+\sqrt{s} \min\Big\{(s%+t)^{-\frac{3}{2}},(s+t)^{-2}\Big\}.
%\end{equation}
%
%Indeed, plugging \eqref{Eq54} into \eqref{Eq51} gives
%
%\begin{equation}\label{Eq57}
%\mathbb{E}\bigg[\Big(\int_{\mathcal{M}}\vert \nabla f^{n,t}\vert^q\Big)^{\frac{2p}{q}}\bigg]^{\frac{1}{p}}\lesssim p\frac{1}{n} \int_{0}^{\infty}\dd s\, \eta(s,t)\big(1+\tfrac{\eta(s,t)}{n}\big).
%\end{equation}
%
%Using that 
%
%$$\int_{0}^{\infty}\dd s\, \min\Big\{(s+t)^{-1},(s+t)^{-2}\Big\}\leq t^{-1}\int_{0}^t\dd s+\int_{t}^1\dd s\, s^{-1}+\int_{1}^{\infty}\dd s\, s^{-2}\lesssim \vert\log(t)\vert,$$
%
%and analogously
%
%$$\int_{0}^{\infty}\dd s\, \sqrt{s} \min\Big\{(s+t)^{-\frac{3}{2}},(s+t)^{-2}\Big\}\lesssim \vert\log(t)\vert,$$
%
%we have 
%
%$$\frac{1}{n}\int_{0}^\infty \dd s\, \eta(s,t)\lesssim \frac{\vert \log(t)\vert}{n}.$$
%
%The same way, 
%
%$$\frac{1}{n^2}\int_{0}^{\infty}\dd s\,\eta^2(s,t)\lesssim \frac{t^{-1}}{n^2}.$$
%
%Consequently, \eqref{Eq57} turns into
%
%$$\mathbb{E}\bigg[\Big(\int_{\mathcal{M}}\vert \nabla f^{n,t}\vert^q\Big)^{\frac{2p}{q}}\bigg]^{\frac{1}{p}}\lesssim p\,\Big(\frac{\vert\log(t)\vert}{n}+\frac{t^{-1}}{n^2}\Big),$$
%
%which gives \eqref{Eq22} using Lemma \ref{momentexp}.

%
%\medskip
%

%The estimate \eqref{Eq54} is a consequence of the following bound on the probability tails
%
%\begin{equation}\label{Eq55}
%\mathbb{P}\bigg(\Big\vert(-s\Delta)^{\frac{1}{2}}\text{P}_{s+t}(\mu^n-\rho)(x)\Big\vert\geq \lambda\bigg)\lesssim \exp\Big(-\frac{1}{\eta(s,t)}\frac{n\lambda^2}{1+\lambda}\Big)\quad\text{for any $\lambda\geq 0$ and $x\in\mathcal{M}$,}
%\end{equation}
%
%together with Lemma \ref{pmoments}. 

%
%\medskip
%

%We now prove \eqref{Eq55}. First, using the definition \eqref{DefEmpiricalMeasures} of $\mu^n$, we have
%
%$$(-s\Delta)^{\frac{1}{2}}\text{P}_{s+t}(\mu^n-\rho)=\frac{1}{n}\sum_{k=1}^n\omega_{s+t}(\cdot,X_k)\quad\text{with}\quad\omega_{s+t}(\cdot,y):=(-s\Delta)^{\frac{1}{2}}\text{P}_{s+t}(\delta_y-\rho).$$
%
%Then, applying Bernstein's inequality, recalled in Proposition \ref{Berstein}, we obtain for any $x\in\mathcal{M}$ and $\lambda\geq 0$
%
%\begin{equation}\label{Eq52}
%\begin{aligned}
%&\mathbb{P}\bigg(\Big\vert(-s\Delta)^{\frac{1}{2}}\text{P}_{s+t}(\mu^n-\rho)(x)\Big\vert\geq \lambda\bigg)\leq \exp\bigg(-\frac{n\lambda^2}{2\mathbb{E}[\vert\omega_{s+t}(x,X_1)\vert^2]+\frac{2}{3}\lambda\sup_{y\in\mathcal{M}}\vert\omega_{s+t}(x,y)\vert}\bigg).
%\end{aligned}
%\end{equation}
%
%We now show that 
%
%\begin{equation}\label{Eq53}
%\mathbb{E}[\vert\omega_{s+t}(x,X_1)\vert^2]\lesssim \min\Big\{(s+t)^{-1},(s+t)^{-2}\Big\}\quad\text{and}\quad\sup_{y\in\mathcal{M}}\vert\omega_{s+t}(x,y)\vert\lesssim\sqrt{s} \min\Big\{(s+t)^{-\frac{3}{2}},(s+t)^{-2}\Big\},
%\end{equation}
%
%which, together with \eqref{Eq52}, yields \eqref{Eq55}.

%
%\medskip
%

%We start with the first item in \eqref{Eq53}. Here, we use the fact that  
%
%\begin{equation}\label{UniformContinuityL2}
%
%\|(-s\Delta)^{\frac{1}{2}}\text{P}_{\frac{s}{2}}f\|_{\LL^2}\lesssim \|f\|_{\LL^2}\quad\text{for all $f\in\cc^2$ and uniformly in $s$.}
%
%\end{equation}
%
%This can be seen from \eqref{Laplacian12}: Using that $\pp_{\frac{s}{2}}$ is an auto-adjoint operator in $\LL^2$ and that from \eqref{Laplacian12:Eq1} we learn $\pp_{\frac{s}{2}}\phi_n =e^{-\lambda_n \frac{s}{2}}\phi_n$, we have 
%
%\begin{align*}
%
%\|(-s\Delta)^{\frac{1}{2}}\text{P}_{\frac{s}{2}}f\|^2_{\LL^2}=\sum_n (s\lambda_n)\vert \big(\phi_n,\pp_{\frac{s}{2}} f\big)_{\LL^2}\vert^2=\sum_n(s\lambda_n)e^{-s\lambda_n }\vert \big(\phi_n,f\big)_{\LL^2}\vert^2\lesssim \|f\|^2_{\LL^2}.
%
%\end{align*}
%
%Hence, using the semi-group property $\text{P}_{s+t}=\text{P}_{\frac{s}{2}}\text{P}_{\frac{s}{2}+t}$, we deduce
%
%\begin{align*}
%\mathbb{E}[\vert\omega_{s+t}(x,X_1)\vert^2]\lesssim \|(-s\Delta)^{\frac{1}{2}}\text{P}_{s+t}(\delta_y-\rho)\|^2_{\LL^{2}}&=\|(-s\Delta)^{\frac{1}{2}}\text{P}_{\frac{s}{2}}\text{P}_{\frac{s}{2}+t}(\delta_y-\rho)\|^2_{\LL^{2}}\\
%&\stackrel{\eqref{UniformContinuityL2}}{\lesssim} \|\text{P}_{\frac{s}{2}+t}(\delta_y-\rho)\|^2_{\LL^{2}}.
%\end{align*}
%
%We now bound in two different ways $\|\text{P}_{\frac{s}{2}+t}(\delta_y-\rho)\|^2_{\LL^{2}}$. First, using the bounds \eqref{Eq28} of the heat-kernel we have by the triangle inequality
%
%$$\|\text{P}_{\frac{s}{2}+t}(\delta_y-\rho)\|^2_{\LL^{2}}\lesssim \|p_{\frac{s}{2}+t}(\cdot,y)\|^2_{\LL^{2}}\stackrel{\eqref{Eq28}}{\lesssim} (s+t)^{-1},$$
%
%yielding to the first alternative in the first item of \eqref{Eq53}. Second, applying Poincaré's inequality yields
%
%$$\|\text{P}_{\frac{s}{2}+t}(\delta_y-\rho)\|^2_{\LL^{2}}\lesssim \|\nabla p_{\frac{s}{2}+t}(\cdot,y)\|^2_{\LL^{2}}\stackrel{\eqref{Eq28}}{\lesssim} (s+t)^{-2},$$
%
%yielding to the second alternative in the first item of \eqref{Eq53}.

%
%\medskip
%

%We now turn to the second item of \eqref{Eq53}. Here, we make use of the explicit formula \eqref{Eq45} applied to $f=\pp_{s+t}(\delta_y-\rho)$. Using the semi-group property $\text{P}_{\tau}\text{P}_{s+t}=\text{P}_{\tau+s+t}$, this takes the form
%
%$$\omega_{s+t}(x,y)=\sqrt{\frac{s}{\pi}}\int_{0}^{\infty}\dd \tau\,\tau^{-\frac{1}{2}}\Delta\text{P}_{\tau+s+t}(\delta_y-\rho).$$
%
%A direct application of the heat-kernel bounds \eqref{Eq28} leads to
%
%\begin{align*}
%\sup_{y\in\mathcal{M}}\vert\omega_{s+t}(x,y)\vert&\lesssim\sqrt{s} \min\bigg\{\int_{0}^{\infty}\dd\tau\,\tau^{-\frac{1}{2}}(\tau+s+t)^{-2},\int_{0}^{\infty}\dd\tau\,\tau^{-\frac{1}{2}}(\tau+s+t)^{-\frac{5}{2}}\bigg\}\\
%&\lesssim \sqrt{s}\min\Big\{(s+t)^{-\frac{3}{2}},(s+t)^{-2}\Big\},
%\end{align*}
%
%which is the second item of \eqref{Eq53}.

%
%\medskip
%

%{\sc Step 4. Correlated case. }We assume that $(X_k)_{k\in\mathbb{N}^*}$ are correlated and satisfies \eqref{WeaklyCorrelationBeta}. Arguing as in \eqref{Eq41}, replacing Meyers' estimate by an energy estimate, we have
%
%\begin{equation}\label{Eq73}
%\int_{\mathcal{M}}\vert\nabla f^{n,t}\vert^2\lesssim \int_{\mathcal{M}}\vert \nabla g^{n,t}\vert^2,
%\end{equation}
%
%where we recall that $g^{n,t}$ is defined in \eqref{Eq47}. Then, we express the $\LL^{2}$-norm of $g^{n,t}$ as
%
%\begin{align*}
%\int_{\mathcal{M}}\vert \nabla g^{n,t}\vert^2=\int_{\mathcal{M}}(-\Delta g^{n,t})g^{n,t}=2\int_{t}^{\infty}\int_{\mathcal{M}}\dd x\,\Big(\text{P}_s(\mu^{n}-\rho)(x)\Big)^2,
%\end{align*}
%
%so that, using the definition \eqref{DefEmpiricalMeasures} of $\mu^n$ and \eqref{Eq73}, we obtain
%
%\begin{equation}\label{Eq74}
%\begin{aligned}
%\mathbb{E}\bigg[\int_{\mathcal{M}}\vert\nabla f^{n,t}\vert^2\bigg]\lesssim&\,\frac{1}{n}\int_{t}^{\infty}\int_{\mathcal{M}}\dd x\,\mathbb{E}[(p_s(x,X_1)-\rho_s)^2]\\
%&+\frac{2}{n^2}\sum_{1\leq i<j\leq n}\mathbb{E}\bigg[\int_{t}^{\infty}\dd s\int_{\mathcal{M}}\dd x\,(p_s(x,X_i)-\rho_s(x))(p_s(x,X_j)-\rho_s(x))\bigg].
%\end{aligned}
%\end{equation}
%
%The first right-hand side term of \eqref{Eq74} is directly controlled using the heat-kernel bounds \eqref{Eq28} and give a $\frac{\vert\log(t)\vert}{n}$ contribution. For the second right-hand side term of \eqref{Eq74}, we appeal to the definition of $(\beta_{ij})_{i,j\geq 1}$ given in \eqref{BersteinCorrelated} that we combine with the condition \eqref{WeaklyCorrelationBeta} to get
%
%\begin{equation}\label{Eq75}
%
%\begin{aligned}
%&\sum_{1\leq i<j\leq n}\mathbb{E}\bigg[\int_{t}^{\infty}\int_{\mathcal{M}}(p_s(x,X_i)-\rho_s(x))(p_s(x,X_j)-\rho_s(x))\bigg]\\
%&\lesssim \sum_{1\leq i<j\leq n}\beta_{ij}\sup_{y,z\in\mathcal{M}}\bigg\vert\int_{t}^{\infty}\int_{\mathcal{M}}(p_s(x,y)-\rho_s(x))(p_s(x,z)-\rho_s(x))\bigg\vert\\
%&\stackrel{\eqref{WeaklyCorrelationBeta}}{\lesssim}n\sup_{y,z\in\mathcal{M}}\bigg\vert\int_{t}^{\infty}\int_{\mathcal{M}}\dd x\,(p_s(x,y)-\rho_s(x))(p_s(x,z)-\rho_s(x))\bigg\vert.
%\end{aligned}
%
%\end{equation}
%
%Then, applying the semi-group property $\text{P}_{s} \text{P}_s=\text{P}_{2s}$ and the heat-kernel bounds \eqref{Eq28}, we have for all $y,z\in\mathcal{M}$ and $s\geq 1$
%
%\begin{align*}
%
%&\int_{\mathcal{M}}\dd x\,(p_s(x,y)-\rho_s(x))(p_s(x,z)-\rho_s(x))\\
%&=p_{2s}(y,z)-\int_{\mathcal{M}}\dd x\rho_s(x)p_s(x,z)-\int_{\mathcal{M}}\dd x\rho_s(x)p_s(x,y)+\int_{\mathcal{M}}\dd x\,(\rho_s(x))^2\\
%&=\int_{\mathcal{M}}\dd x\, \rho(x)(p_{2s}(y,z)-p_{2s}(x,z))-\int_{\mathcal{M}}\dd x\, \rho(x) p_{2s}(x,y)+\int_{\mathcal{M}}\dd x\,(\rho_s(x))^2\\
%&\stackrel{\eqref{Eq28}}{\lesssim} \min\{s^{-1},s^{-\frac{3}{2}}\}+ \min\{s^{-1},s^{-2}\},
%
%\end{align*}
%
%so that 
%
%$$\sup_{y,z\in\mathcal{M}}\bigg\vert\int_{t}^{\infty}\int_{\mathcal{M}}\dd x\,(p_s(x,y)-\rho_s(x))(p_s(x,z)-\rho_s(x))\bigg\vert\lesssim \vert\log(t)\vert,$$
%
%and \eqref{Eq75} shows that the second-right hand side term of \eqref{Eq74} give a $\frac{\vert\log(t)\vert}{n}$ contribution.
%{\sc Step 4. Alternative. }It is better to use the spectral gap inequality
%
%\begin{align*}
%\mathbb{E}\Big[\Big((-s\Delta)^{\frac{1}{2}}\text{P}_{s+t}(\mu^n-\rho)(x)\Big)^{2p}\Big]^{\frac{1}{p}}\lesssim& p \frac{1}{n^2} \mathbb{E}\bigg[\bigg(\sum_{k=1}^n\Big((-s\Delta)^{\frac{1}{2}}\text{P}_{s+t}(\delta_{X_k}-\rho)\Big)^2\bigg)^{p}\bigg]^{\frac{1}{p}}\\
%\lesssim& p \frac{1}{n}\sup_{y\mathcal{M}}\vert\mathcal{M} (x,y)\vert.
%\end{align*}
%
%Now note that since the heat-kernel $p_\tau$ and $g^{n,t}$ satisfies the Neumann boundary condition, we have by two integration by parts
%
%$$\Delta \text{P}_{\tau}g^{n,t}= \text{P}_{\tau}\Delta g^{n,t}=\text{P}_{\tau+t}(\mu^{n}-\rho).$$
%
%Then, using that $\partial_s\text{P}_s=\Delta \text{P}_s$, we get 
%
%$$\partial_s\text{P}_s(-\Delta)^{\frac{1}{2}} g^{n,t}=\Delta\text{P}_s(-\Delta)^{\frac{1}{2}}g^{n,t}=\int_{0}^{\infty}\dd\tau\, \tau^{-\frac{1}{2}}\Delta P_{\tau+t+s}(\mu^n-\rho),$$
%and finally
%
%$$\int_{\mathcal{M}}\vert\nabla g^{n,t}\vert^4\lesssim \int_{\mathcal{M}}\bigg(\int_{0}^{\infty}\dd s\,s\Big(\int_{0}^{\infty}\dd\tau\, \tau^{-\frac{1}{2}}\Delta P_{\tau+t+s}(\mu^n-\rho)\Big)^2\bigg)^2.$$
%
%Using then \eqref{Eq41} together with Minkowski's inequality, we arrive at for any $p<\infty$
%
%\begin{equation}\label{Eq42}
%\begin{aligned}
%
%\mathbb{E}\bigg[\Big(\int_{\mathcal{M}}\vert \nabla f^{n,t}\vert^q\Big)^{\frac{2p}{q}}\bigg]^{\frac{1}{p}}\stackrel{\eqref{Eq41}}{\lesssim}&\mathbb{E}\bigg[\bigg(\int_{\mathcal{M}}\bigg(\int_{0}^{\infty}\dd s\,s\Big(\int_{0}^{\infty}\dd\tau\, \tau^{-\frac{1}{2}}\Delta P_{\tau+t+s}(\mu^n-\rho)\Big)^2\bigg)^2\bigg)^{\frac{p}{2}}\bigg]^{\frac{1}{p}}\\
%\leq& \bigg(\int_{\mathcal{M}}\bigg(\int_{0}^{\infty}\dd s\,s\,\mathbb{E}\Big[\Big(\int_{0}^{\infty}\dd\tau\, \tau^{-\frac{1}{2}}\Delta P_{\tau+t+s}(\mu^n-\rho)\Big)^{2p}\Big]^{\frac{1}{p}}\bigg)^2\bigg)^{\frac{1}{2}}.
%
%\end{aligned}
%\end{equation}
%
%We now bound the expectation in \eqref{Eq42} using Bernstein's recalled in Proposition \ref{Berstein}. To do so, we write using the definition of $\mu^n$
%
%$$\int_{0}^{\infty}\dd\tau\, \tau^{-\frac{1}{2}}\Delta P_{\tau+t+s}(\mu^n-\rho)=\frac{1}{n}\sum_{k=1}^n\mathcal{M}(\cdot,X_k)\quad\text{with}\quad\mathcal{M}(\cdot,y):=\int_{0}^{\infty}\dd\tau\, \tau^{-\frac{1}{2}}\Delta (p_{\tau+t+s}(\cdot,y)-\rho),$$
%
%so that for any $\lambda>0$ and $x\in\mathcal{M}$, we apply Proposition \ref{Berstein} to obtain
%
%\begin{align*}
%&\mathbb{P}\bigg(\Big\vert\int_{0}^{\infty}\dd\tau\, \tau^{-\frac{1}{2}}\Delta P_{\tau+t+s}(\mu^n-\rho)\Big\vert\geq \lambda\bigg)\leq \exp\bigg(-\frac{n\lambda^2}{2\mathbb{E}[\vert\mathcal{M}(x,X_1)\vert^2]+\frac{2}{3}\lambda\sup_{y\in\mathcal{M}}\vert\mathcal{M}(x,y)\vert}\bigg).
%\end{align*}
%
%We then bound $\mathbb{E}[\vert\mathcal{M}(x,X_1)\vert^2]$ and $\sup_{y\in\mathcal{M}}\vert \mathcal{M}(x,y)\vert$ separately. First, from Minkowski's inequality
%
%$$\mathbb{E}[\vert\mathcal{M}(x,X_1)\vert^2]\leq \bigg(\int_{0}^{\infty}\dd\tau\,\tau^{-\frac{1}{2}}\mathbb{E}\Big[\vert\Delta(p_{\tau+t+s}(\cdot,X_k)-\rho)\vert^2\Big]^{\frac{1}{2}}\bigg)^2.$$
%
%We then claim that 
%
%$$\mathbb{E}\Big[\vert\Delta(p_{\tau+t+s}(\cdot,X_k)-\rho)\vert^2\Big]^{\frac{1}{2}}\lesssim \min\{$$
%

%using the estimates on the heat-kernel \eqref{Eq28}, we have
%
%\begin{align*}
%\mathbb{E}[\vert\Delta (p_{\tau+t+s}(\cdot,X_k)-\text{P}_{\tau+t}\rho)\vert^2]^{\frac{1}{2}}\leq\, \mathbb{E}[\vert\Delta (p_{\tau+t}(\cdot,X_k)\vert^2]&=\int_{\mathcal{M}}\rho(x)\dd x\,\vert\Delta p_{\tau+t}(\cdot,X_k)\vert^2\\
%&\lesssim t^{-4}
%\end{align*}
%
%
\end{proof}


%
%and
%
%\begin{equation}\label{Eq23}
%\mathbb{E}\bigg[\Big(W^{2}_2(\rho\,\dd x,\mu^{n,t})\Big)^p\bigg]^{\frac{1}{p}}\lesssim p^{\nu}\frac{\vert \log(t)\vert}{n}.
%\end{equation}
%

%
\end{proposition}
%

%
%\begin{proposition}
%Let $X_1, \dots, X_n$ be i.\;i.\;d. $\rho \,\dd x$ distributed points and let $\mu^n$ denote the empirical measure. Then 
%\begin{equation}\label{eq:upper_W2_transportempirical}
%\mathbb{E} \left[ W_2^2(\mu^n, \rho \,\dd x) \right] \lesssim \frac{\log n}n.
%\end{equation}
%\end{proposition}
%\begin{proof}
%\hw{Old: We have to redo the upper-bound since now we assume no regularity on $\rho$. We can follow the argument in \cite[Theorem 4.4]{ambrosio2019pde}.} 
%Fix $\gamma>0$ and let us define \fm{Check again if there is some $\dd x$ missing, i.e. $\mu^{n,t,c}$ is a measure, $u^{n,t,c}$ is a density} 
%
%$$\mu^{n,t,c}:=(1-c)\mu^{n,t}+c\rho_t\,\dd x \quad \text{with}\quad t=c=\gamma n^{-1}\log(n).$$
%
%By the joint convexity of $W_2^2$ \fm{Write it in the preliminaries section and make an inline equation}, it follows 
%$$W^2_2(\mu^{n,t},\mu^{n,t,c})\lesssim c.$$ 
%By the triangle inequality together with Young's product inequality \fm{I think the choice of the parameter $\alpha$ does not matter, so I just use the estimate $(a+b)^2 \le 2 a^2 + b^2$} and the fact that $W_2^2(\nu_t, \nu) \lesssim t$ \fm{Make an inline equation for this and refer to that every time}, we may write
%
%\begin{align*}
%W^2_2(\mu^{n},\rho\,\dd x)\leq& (W_2(\mu^n, \mu^{n,t}) + W_2(\mu^{n,t},\mu^{n,t,c})+W_2(\mu^{n,t,c},\rho_t\,\dd x) +W_2(\rho_t \, \dd x, \rho \, \dd x))^2\\
%\lesssim &W_2^2(\mu^n, \mu^{n,t}) + W_2^2(\mu^{n,t},\mu^{n,t,c})+W_2^2(\mu^{n,t,c},\rho_t\,\dd x) +W_2^2(\rho_t \, \dd x, \rho \, \dd x)\\
%\lesssim & t + c + W_2^2(\mu^{n,t,c},\rho_t\,\dd x).
%\end{align*}
%\leq & (1+\alpha)W^2_2(\rho\,\dd x,\mu^{n,t,c})+(1+\alpha^{-1})W^2_2(\mu^{n,t},\mu^{n,t,c})\\
%\lesssim & (1+\alpha)W^2_2(\rho_t\,\dd x,\mu^{n,t,c}) + (1+\alpha) W_2^2(\rho \,\dd x, \rho_t\, \dd x)+(1+\alpha^{-1})c \\
%\lesssim & (1+\alpha)W^2_2(\rho_t\,\dd x,\mu^{n,t,c}) +(1+\alpha)t + (1+\alpha^{-1})c.
%\end{align*}
%
%\fm{I think there were some errors with the 2s, if I am not wrong the one that I wrote above should be correct, moreover I needed $\rho_t$ in the argument of the Wasserstein distance to apply the upper bound of the proposition the old one was this
%\begin{align*}
%W^2_2(\rho\,\dd x,\mu^{n,t})\leq& (W_2(\mu^{n,t},\mu^{n,t,c})+W_2(\mu^{n,t,c},\rho\,\dd x))^2\\
%\leq & (1+\alpha)W^2_2(\rho\,\dd x,\mu^{n,t,c})+2(1+\alpha^{-1})W^2_2(\mu^{n,t},\mu^{n,t,c})\\
%\lesssim & (1+\alpha)W^2_2(\rho\,\dd x,\mu^{n,t,c})+2(1+\alpha^{-1})c.
%\end{align*}}
%
%Taking the expectation in the expression above, by the choices of $t$ and $c$ 
%\[
%\mathbb{E}[W_2^2 (\mu^n, \rho \, \dd x)] \lesssim \frac{\log n}n + \mathbb{E}[W_2^2(\mu^{n,t,c},\rho_t\,\dd x)].
%\]
%Thus we are left to show
%\begin{equation}\label{eq:expcontributionW2interpolate}
%E[W_2^2(\mu^{n,t,c},\rho_t\,\dd x)] \lesssim \frac{\log n}n
%\end{equation}
%To this end, note that since $f^{n,t}$ solves \eqref{eq:f^nt}, $(1-c) f^{n,t}$ solves
%\[
%-\nabla \cdot \rho \nabla ((1-c) f^{n,t}) = u^{n,t,c} - \rho_t.
%\]
%The latter and \eqref{eq:W_2_upper} combine to 
%
%\begin{equation}\label{eq:W2muntc}
%\begin{split}
%W^2_2(\mu^{n,t,c},\rho_t\,\dd x)&\lesssim (1-c)^2\int_{\mathcal{M}}\frac{\vert \nabla f^{n,t}\vert^2}{M(u^{n,t,c},\rho_t)}\\
%&\leq \int_{\mathcal{M}}\vert \nabla f^{n,t}\vert^2+\int_{\mathcal{M}}\vert \nabla f^{n,t}\vert^2(M(u^{n,t,c},\rho_t)^{-1}-1).
%\end{split}
%\end{equation}
%
%We now estimate the contribution of the first term. Taking the expectation in \eqref{eq:W2muntc}, by Holder's inequality, \eqref{Eq22} and the choice of $t$ \fm{is it ok given the energy estimate?}
%\[
%\mathbb{E}\bigg[ \int_{\mathcal{M}}\vert \nabla f^{n,t}\vert^2 \bigg] \lesssim \mathbb{E} \bigg[ \bigg( \int_\mathcal{M} |\nabla f^{n,t}|^{\bar{q}} \bigg)^{\frac2{\bar{q}}} \bigg] \lesssim \frac{\log t} n \sim \frac{\log n}n.
%\]
%We now estimate the contribution of the second term in \eqref{eq:W2muntc}. By Holder's inequality
%\[
%\mathbb{E}\bigg[\int_{\mathcal{M}}\vert \nabla f^{n,t}\vert^2(M(u^{n,t,c},\rho_t)^{-1}-1)\bigg] \le \mathbb{E}\bigg[\int_\mathcal{M} |\nabla f^{n,t}|^4 \bigg]^\frac12\mathbb{E}[(M(u^{n,t,c},\rho_t)^{-1}-1)^2 ]^\frac12.
%\]
%\fm{I got lost with all those exponent, but here if we get bounds as the one for the L2 moment above we just need to say that the 4 moment is bounded and the other term goes to 0 faster than log n/n to conclude}
%For the second-right hand side term, we use Hölder's inequality together with Meyers' estimate to get
%
%\[
%\begin{split}
%\lefteqn{\mathbb{E}\bigg[\Big(\int_{\mathcal{M}}\vert \nabla f^{n,t}\vert^2(M(u^{n,t,c},\rho_t)^{-1}-1)\Big)^{p}\bigg]^{\frac{1}{p}}}\\
%& \leq \mathbb{E}\bigg[\Big(\int_{\mathcal{M}}\vert \nabla f^{n,t}\vert^{\bar q}\Big)^{\frac{2p}{\bar q}}\bigg]^{\frac{1}{2p}}\mathbb{E}\bigg[\Big(\int_{\mathcal{M}}(M(u^{n,t,c},\rho_t)^{-1}-1)^{\bar q'}\Big)^{\frac{2p}{\bar q'}}\bigg]^{\frac{1}{2p}}.
%\end{split}
%\]
%
%\fm{It is not exactly this one, as one doesn't have the conjugate exponent of $\bar{q}$ but of $\bar{q}/2$}
%The first term on the right-hand side is controlled by \eqref{Eq22}. Moreover by Lemma \ref{lem:integralbound} as $n \rightarrow \infty$
%\[
%\mathbb{E}\bigg[\Big(\int_{\mathcal{M}}(M(u^{n,t,c},\rho_t)^{-1}-1)^{\bar q'}\Big)^{\frac{2p}{\bar q'}}\bigg]^{\frac{1}{2p}} \rightarrow 0,
%\]
%thus \eqref{eq:expcontributionW2interpolate} holds.

%it remains to control the second right-hand side term and to do so we adapt \cite[Lemma 3.13]{ambrosio2019pde}, an inspection of the proof shows that it holds for any moments in probability and space. 
%\end{proof}

%
%\begin{lemma}\label{lem:probabilitybound}
%Let $y \in \mathcal{M}$ and let $u^{n,t}$ \fm{be the right thing}. Then for any $\eta > 0$
%\[
%\mathbb{P}(|u^{n,t}(y) - \rho_t (y)| > \eta) \le 2 \exp \big(- \frac{nt^{\frac d2}}{2 c} F(1, \eta) \big),
%\]
%where
%\[
%F(c, \eta) = \sup_{\lambda} \bigg(\lambda \eta - \frac{\lambda^2}2 \exp(c \lambda)\bigg) > 0.
%\]
%\end{lemma}
%
%\begin{proof}
%\fm{to appear \dots}
%\end{proof}
%
%\begin{lemma}\label{lem:integralbound}
%Let $\mu^{n,t,c} = (1-c) \mu^{n,t} + c \rho_t \, \dd x$ and let $u^{n,t,c} = (1-c) u^{n,t} + c \rho_t \, \dd x$ be its probability density, with $c = c(n) \in (0,1]$. It $t^{\frac d2} = t^{\frac d2} (n) \ge \gamma n^{-1} \log n$ and $n c(n) \rightarrow \infty$ as $n \rightarrow \infty$, then for any $p,q >0$
%\[
%\lim_{n \rightarrow \infty} \mathbb{E} \bigg[ \bigg(\int_\mathcal{M} (M (u^{n,t,c}, \rho_t)^{-1} -1)^q \bigg)^{\frac{2p}q}\bigg]^\frac1{2p}  = 0.
%\]
%\fm{too ambitious?}
%\end{lemma}
%\begin{proof}
%\fm{to appear \dots}
%\end{proof}
%
\subsection{Fluctuation estimates}\label{FluctuationSection}
%
%We define the approximation for a given parameter $\delta>0$
%
%\begin{equation}\label{Eq2}
%\left\{
    %\begin{array}{ll}
        %-\nabla\cdot \rho_\delta\nabla f^{n,t}_\delta =\mu^{n,t}-\rho_t &\quad \text{in $\mathcal{M}$,} \\
        %n_{\mathcal{M}}\cdot \nabla f^{n,t}_\delta = 0  &\quad \text{on $\partial \mathcal{M}$,}
    %\end{array}
%\right.
%\end{equation}
%
%for $\rho_\delta:=\text{P}_\delta\rho$. We first note that the error $f^{n,t}_\delta-f^{n,t}$ satisfies from \eqref{eq:f^nt} and \eqref{Eq2} 
%
%$$- \nabla \cdot\rho_\delta \nabla (f^{n,t}_\delta-f^{n,t})= \nabla \cdot (\rho_\delta - \rho) \nabla f^{n,t}\quad\text{in $\mathcal{M}$},$$
%
%so that from an energy estimate, Hölder's inequality and Proposition \ref{prop1}, we obtain the error bound
%
%\begin{equation}\label{eq:errornablafdelta}
%\mathbb{E}\bigg[\Big(\int_\mathcal{M} | \nabla f^{n,t}_\delta - \nabla f^{n,t}|^2\Big)^p\bigg]^{\frac{1}{p}}\lesssim p \|\rho_\delta-\rho\|^2_{\LL^{2(\frac{\bar q}{2})'}}\frac{\vert\log(t)\vert}{n},
%\end{equation}
%
%where $\bar{q}'$ denotes the conjugate exponent of the Meyers exponent of the operator $-\nabla\cdot \rho\nabla$ given in Theorem \ref{Meyers}. This allow us to focus on $f^{n,t}_\delta$ which owns $\cc^{\infty}(\mathcal{M})$-regularity. The proof of Theorem \ref{th1} relies on a bound on $\|\nabla^2 f^{n,t}_\delta\|_{\LL^{\infty}(\mathcal{M})}$ that we can assume arbitrary small with very high probability when $n\uparrow\infty$. To do so, we decompose the equation into a regular part and a singular part by expanding the equation \eqref{Eq2} thanks to the regularity of $\rho_{\delta}$ as follows
%
%\begin{equation}\label{Eq4}
%-\Delta f^{n,t}_\delta=\tfrac{1}{\rho_\delta}(\mu^{n,t}-\rho_t)+\tfrac{1}{\rho_\delta}\nabla\rho_\delta\cdot \nabla f^{n,t}_\delta.
%\end{equation}
% 
%The second right-hand side term is referred to as the "regular-part" because one can aply Schauder's estimates to get "good" deterministic estimates in $t$. For the first right-hand side term, referred to as the "singular-part", finer estimates via the explicit formulation using the heat semi-group are required. Note that in the case of the cube, rewriting the equation in the form \eqref{Eq4} implies regularity estimate in the form of, for some $\gamma>0$
%
%\begin{equation}\label{RegEstiForCube}
%\|\nabla f^{n,t}_{\delta}\|_{\cc^{0,\alpha}(\mathrm{Q})}\lesssim \|\nabla\rho_\delta\|^\gamma_{\LL^{\infty}(\QQ)}\|\mu^{n,t}-\rho_t\|_{\LL^{\infty}(\QQ)}.
%\end{equation}
%
%\nc{no references available, we can include the short proof for completeness}. The estimate \eqref{RegEstiForCube} is a consequence of
%
%\begin{equation}\label{RegEstiForCube1}
%\|\nabla f^{n,t}_{\delta}\|_{\WW^{1,p}(\QQ)}\lesssim \|\nabla\rho_\delta\|^\gamma_{\LL^{\infty}(\QQ)}\|\mu^{n,t}-\rho_t\|_{\LL^{p}(\QQ)}\quad \text{for $p>d$,}
%\end{equation}
%
%together with Sobolev embedding $\WW^{1,p}(\QQ)\subset \cc^{0,\alpha}(\QQ)$. For \eqref{RegEstiForCube1}, we apply $\LL^p$-regularity estimates to \eqref{Eq4}, which yields
%
%\begin{align*}
%\| f^{n,t}_\delta\|_{\WW^{2,p}(\QQ)}\lesssim& \|\tfrac{1}{\rho_\delta}(\mu^{n,t}-\rho_t)\|_{\LL^{p}(\QQ)}+ \|\tfrac{1}{\rho_\delta}\nabla\rho_{\delta}\cdot\nabla f^{n,t}_\delta\|_{\LL^{p}(\QQ)}\\
%&\lesssim \|\mu^{n,t}-\rho_t\|_{\LL^{p}(\QQ)}+ \|\nabla\rho_{\delta}\|_{\LL^{\infty}(\QQ)}\|\nabla f^{n,t}_\delta\|_{\LL^{p}(\QQ)}.
%\end{align*}
%
%which can be obtained by a classical reflexion argument extending the problem on the torus. We then obtain \eqref{RegEstiForCube1} using the interpolation inequalities, for some $\theta\in (0,1)$
%
%$$\|f^{n,t}_\delta\|_{\WW^{1,p}(\QQ)}\lesssim \|f^{n,t}_\delta\|^{\frac{1}{2}}_{\LL^{p}(\QQ)}\|f^{n,t}_\delta\|^{\frac{1}{2}}_{\WW^{2,p}(\QQ)}\quad\text{and}\quad \|f^{n,t}_{\delta}\|_{\LL^{p}(\QQ)}\leq \|f^{n,t}_\delta\|^{\theta}_{\LL^{p_\star}(\QQ)}\|f^{n,t}_{\delta}\|^{1-\theta}_{\LL^2(\QQ)},$$
%  
%together with Young's inequality and the energy estimate
%
%$$\|f^{n,t}_\delta\|_{\LL^{2}(\QQ)}\lesssim \|\mu^{n,t}-\rho_t\|_{\LL^{2}(\QQ)}.$$

%We show the following.
%
This section is devoted to justify \eqref{FluctuationEstimates} needed to ensure the condition \eqref{SmallnessH} with very high probability. Our result is in the spirit of \cite[Theorem 3.3]{ambrosio2019finer}. However, our strategy differs from \cite[Theorem 3.3]{ambrosio2019finer} and is based on Schauder's theory, with an additional special care on the dependences on $\delta$.
%However, we do not rely on the same proof based on the explicit formula for \eqref{eq:f^nmt} obtained by disintegration using the kernel of the operator $\partial_\tau-\nabla\cdot\rho_\delta\nabla$. On the one hand, because of the way we regularize the \rhs of \eqref{eq:f^nmt}, we would not see direct cancellations since the kernel for the operator $\partial_\tau-\nabla\cdot\rho_\delta\nabla$ and the standard heat-kernel do not match. On the other hand, we need to keep track of the dependences on the parameter $\delta$ that we want to optimize, which are unclear for bounds of type \eqref{Eq28} for the kernel of the operator $\partial_\tau-\nabla\cdot\rho_\delta\nabla$. Instead, we rely on Schauder's theory for which we know that the constant depends at most polynomially on $\|\rho_\delta\|_{\cc^{0,\alpha}}$ and so on $\delta$. 
We briefly sketch the main ingredients of the fluctuation estimates \eqref{FluctuationEstimates}. By linearity, it is enough to show \eqref{FluctuationEstimates} for $f^{n,t}_\delta\in \dot{\mathrm{H}}^1$ being the solution of 
%
\begin{equation}\label{Eq2}
%
-\nabla\cdot\rho_\delta\nabla f^{n,t}_\delta=\mu^{n,t}-\rho_t.
%
\end{equation}
%
We then make use of the chain rule to expand the equation in
%
\begin{equation}\label{Eq3}
%
-\Delta f^{n,t}_\delta=\tfrac{1}{\rho_\delta}\nabla\rho_\delta\cdot\nabla f^{n,t}_\delta+\tfrac{1}{\rho_\delta}(\mu^{n,t}-\rho_t),
%
\end{equation}
%
and we define the auxiliary problem
%
\begin{equation}\label{AuxiliaryProblemFluctu}
%
-\Delta u^{n,t}_\delta=\tfrac{1}{\rho_\delta}(\mu^{n,t}-\rho_t),
%
\end{equation}
%
such that the difference $v^{n,t}_\delta:=f^{n,t}_\delta-u^{n,t}_\delta$ solves $-\Delta v^{n,t}_\delta=\tfrac{1}{\rho_\delta}\nabla\rho_\delta\cdot\nabla f^{n,t}_\delta$. Doing so, on the one hand, we can control $u^{n,t}_\delta$ which can be handled using the explicit formula in terms of the heat-kernel, the explicit bounds on the latter, cf \eqref{Eq28}, and the regularity of $\frac{1}{\rho_\delta}$. On the other hand, we use Schauder's estimates to control $v^{n,t}_\delta$ and $f^{n,t}_\delta$. Using the fact that those estimates depend polynomially on $\|\rho_\delta\|_{\cc^{0,\alpha}}$, we can keep track on the dependences on $\delta$ that we can optimize later on. We finally mention that in the case where $\mathcal{M}$ has a boundary, $u^{n,t}_\delta$ cannot be directly defined by \eqref{AuxiliaryProblemFluctu} since the \rhs does not have zero mean. In order to also include this case, we add a zero order term in the equation, see \eqref{Eq3Bis}.
%
\begin{proposition}[Fluctuation estimates]\label{Fluctuation}
%
Let $\{\mu^n\}_n$ be defined in \eqref{eq:empmeas} with point clouds satisfying Assumption \ref{Assumptions}. For any parameter $\delta\in (0,1)$ and $\upsilon>0$, we define\footnote{with Neumann boundary conditions in case $\mathcal{M}$ has a boundary} $u^{n,t}_{\delta}\in \mathrm{H}^1$ weak solution of
%
\begin{equation}\label{Eq3Bis}
%
u^{n,t}_{\delta} -\Delta u^{n,t}_{\delta}=\tfrac{1}{\rho_\delta}(\mu^{n,t}-\rho_t),
 %
\end{equation}
%
and the two events
%
\begin{equation}\label{Events}
%
\mathcal{A}_n:=\Big\{\|\mu^{n,t}-\rho_t\|_{\LL^{\infty}}\leq \tfrac{1}{\log^{\upsilon}(n)}\Big\}\quad \text{and}\quad\mathcal{B}_{\delta,n}:=\Big\{\big\|(\nabla u^{n,t}_\delta,\nabla^2 u^{n,t}_{\delta})\big\|_{\LL^{\infty}}\leq \tfrac{1}{\log^{\upsilon}(n)}\Big\}.
%
\end{equation}
%
There exists $\kappa>0$ such that for any $\kappa_1>0$ and the choice 
%
\begin{equation}\label{Eq:ChoicesPara}
\delta=\delta_n:=\tfrac{1}{\log^{\kappa_1}(n)},
\end{equation}
%
the solution $f^{n,t}_\delta\in \dot{\mathrm{H}}^1$ of \eqref{Eq2} satisfies
%
\begin{equation}\label{Eq6}
\mathds{1}_{\mathcal{A}_n\cap\mathcal{B}_{\delta,n}}\big\|(\nabla f^{n,t}_\delta,\nabla^2 f^{n,t}_{\delta})\big\|_{\LL^{\infty}}\leq \frac{1}{\log^{\upsilon-(\kappa+2)\kappa_1}(n)}.
\end{equation}
%
Furthermore, there exists $\kappa_2>0$ depending on $\upsilon$ such that for the choice $t=t_n:=\frac{\log^{\kappa_2}(n)}{n}$, we have
%
\begin{equation}\label{Eq13}
\mathbb{P}(\mathcal{A}^{c}_n\cup\mathcal{B}^c_{\delta,n})=o(\tfrac{1}{n^\ell})\quad\text{for any $\ell\in\mathbb{N}$}.
\end{equation}
%


%there exist constants $C,C_1,C_2,C_3,C_4>0$ and an exponent $\alpha$ depending on $\gamma$ such that
%
%\begin{itemize}
%
%\item If $(X_k)_{k\in\{1,\cdots,n\}}$ are i.~i.~d., it holds
%
%\begin{equation}\label{Eq13}
%\mathbb{P}(\mathcal{A}^{c}_n\cup\mathcal{B}^c_{\varepsilon_n,\delta_n,n})\lesssim \frac{n^3}{\log^{\alpha}(n)}\exp(-\tfrac{1}{C}\log^{2}(n))\quad\text{for}\quad t=t_n:=\frac{\log^{4+\frac{4}{\gamma+2}}(n)}{n}.
%\end{equation}
%
%\item If $(X_k)_{k\in\mathbb{N}^*}$ satisfies for $a,b>0$ and $\beta\in (0,1]$
%
%\begin{equation}\label{WeaklyCorrelatedAlpha}
%
%\alpha_\ell\leq a\exp(-b \ell^{\beta}) \quad \text{for any\quad $\ell\geq 1$},
%
%\end{equation}
%where we recall that $(\alpha_\ell)_{\ell\geq 1}$ are defined in \eqref{AlphaMixing}, it holds
%
%\begin{equation}
%\begin{aligned}\label{Eq76}
%\mathbb{P}(\mathcal{A}^{c}_n\cup\mathcal{B}^c_{\varepsilon_n,\delta_n,n})\lesssim&\frac{n^3}{\log^{\alpha}(n)}\bigg(n\exp\Big(-\frac{1}{C_1}\log^{(\nu-1)\beta}(n)\Big)+\exp\Big(-\frac{1}{C_2}\log^2(n)\Big)\\
%&+\exp\bigg(-\frac{1}{C_3}\frac{\log^{2\nu-1}(n)}{n}\exp\Big(\frac{1}{C_4}\tfrac{\log^{(\nu-1)\beta(1-\beta)}(n)}{\log(\log(n))}\Big)\bigg)\bigg)
%\end{aligned}
%\end{equation}
%
%for $t=t_n:=\frac{\log^{\nu}(n)}{n}$ and $\nu:=4+\frac{1}{\beta}+\frac{4}{\gamma+2}$.\nc{add the $\beta=1$ case}
%\end{itemize}
%
\end{proposition}
%
\begin{proof}
%
The proof of Proposition \ref{Fluctuation} is split into two steps. In the first step, we prove \eqref{Eq6} where our main tool is Schauder's theory and an explicit formula for $u^{n,t}_{\delta}$ defined in \eqref{Eq3}. In the second step, we show \eqref{Eq13}, where our main tool is the concentration inequalities in Proposition \ref{BersteinCorrelated} and the explicit formula for $u^{n,t}_{\delta}$ used in the first step.

%
\medskip
%

{\sc Step 1. Proof of \eqref{Eq6}. }We define the difference $v^{n,t}_{\delta}:=f^{n,t}_\delta-\Big(u^{n,t}_{\delta}-\int_{\mathcal{M}} u^{n,t}_\delta\Big)\in \dot{\text{H}}^1$ and note that from \eqref{Eq3} and \eqref{Eq3Bis}, it solves
%
\begin{equation}\label{Eq5}
%
 -\Delta v^{n,t}_{\delta}=\frac{1}{\rho_{\delta}}\nabla\rho_\delta\cdot\nabla f^{n,t}_\delta-\Big(u^{n,t}_{\delta}-\int_{\mathcal{M}}u^{n,t}_\delta\Big).
 %
\end{equation}
%
We claim that there exists $\kappa>0$ such that with the choice of $\delta$ in \eqref{Eq:ChoicesPara}, we have
%
\begin{equation}\label{Eq24}
\mathds{1}_{\mathcal{A}_n\cap\mathcal{B}_{\delta,n}}\big\|(\nabla v^{n,t}_{\delta},\nabla^2 v^{n,t}_{\delta})\big\|_{\LL^{\infty}}\lesssim \frac{1}{\log^{\upsilon-(\kappa+2)\kappa_1}(n)}.
\end{equation}
%
The estimate \eqref{Eq6} then follows from \eqref{Eq24} and the triangle inequality. 

%
\medskip
%

We now prove \eqref{Eq24}. We apply Schauder's estimate (see for instance \cite{gilbarg2015elliptic} and \cite[Chapter 10]{nicolaescu2020lectures}) to \eqref{Eq5} to obtain
%
\begin{equation}\label{Eq11}
\big\|(\nabla v^{n,t}_{\delta},\nabla^2 v^{n,t}_{\delta})\big\|_{\LL^{\infty}}\lesssim \|\tfrac{1}{\rho_{\delta}}\nabla\rho_\delta\cdot\nabla f^{n,t}_\delta\|_{\cc^{0,\alpha}}+\Big\|u^{n,t}_{\delta}-\int_{\mathcal{M}}u^{n,t}_\delta\Big\|_{\cc^{0,\alpha}}.
\end{equation}
%
While the second \rhs term is directly of order of $\frac{1}{\log^{\upsilon}(n)}$ in $\mathcal{B}_{\delta,n}$, the first \rhs term requires additional attention. Using \eqref{Ellipticity} and \eqref{Eq28}, $\rho_\delta$ satisfies
%
\begin{equation}\label{Eq25}
\rho_{\delta}\geq \lambda\quad \text{and}\quad \delta^{-1}\|\nabla\rho_{\delta}\|_{\LL^{\infty}}+\|\nabla^2\rho_{\delta}\|_{\LL^{\infty}}\lesssim \delta^{-2},
\end{equation}
%
and together with the algebraic property of $\|\cdot\|_{\cc^{0,\alpha}}$, we have
%
$$\|\tfrac{1}{\rho_{\delta}}\nabla\rho_\delta\cdot\nabla f^{n,t}_\delta\|_{\cc^{0,\alpha}}\leq \|\tfrac{1}{\rho_{\delta}}\nabla\rho_\delta\|_{\cc^{0,\alpha}}\|\nabla f^{n,t}_\delta\|_{\cc^{0,\alpha}}\lesssim \delta^{-2}\|\nabla f^{n,t}_\delta\|_{\cc^{0,\alpha}}.$$
%
The latter is bounded using Schauder's estimate applied this time on \eqref{Eq2} (knowing that the dependence on $\|\rho_{\delta}\|_{\cc^{0,\alpha}}$ is at most polynomial): there exists $\kappa>0$ such that
%
\begin{equation*}
\begin{aligned}
\|\nabla f^{n,t}_\delta\|_{\cc^{0,\alpha}}&\lesssim \|\rho_{\delta}\|^{\kappa}_{\cc^{0,\alpha}}\|\mu^{n,t}-\rho_t\|_{\LL^{\infty}}\stackrel{\eqref{Eq25}}{\lesssim} \delta^{-\kappa}\|\mu^{n,t}-\rho_t\|_{\LL^{\infty}}\stackrel{\eqref{Eq:ChoicesPara}}{\lesssim} \log^{\kappa\kappa_1}(n)\|\mu^{n,t}-\rho_t\|_{\LL^{\infty}},
\end{aligned}
%
\end{equation*}
%
which yields the following control of the first \rhs term of \eqref{Eq11}
%
\begin{equation}\label{Eq7}
\|\tfrac{1}{\rho_{\delta}}\nabla\rho_\delta\cdot\nabla f^{n,t}_\delta\|_{\cc^{0,\alpha}}\lesssim \log^{(\kappa+2)\kappa_1}(n)\|\mu^{n,t}-\rho_t\|_{\LL^{\infty}},
\end{equation}
%
which is of order of $\frac{1}{\log^{\upsilon-(\kappa+2)\kappa_1}(n)}$ in $\mathcal{A}_n$.

%We now turn to the control of the second right-hand side term of \eqref{Eq11}. 
%Rather than estimating $\|u^{n,t}_{\varepsilon,\delta}\|_{\cc^{0,\alpha}}$, we estimate $\|\nabla u^{n,t}_{\varepsilon,\delta}\|_{\LL^{\infty}}$. To do so, we decompose $u^{n,t}_{\varepsilon,\delta}$ into a regular-part $\mathcal{I}_1$ and a singular-part $\mathcal{I}_2$:
%
%\begin{equation}\label{Eq26}
%u^{n,t}_{\varepsilon,\delta}=\underbrace{\int_{\varepsilon^{-\frac{1}{2}}}^{\infty}\dd s\,e^{-\varepsilon s}\,\text{P}_s\Big(\tfrac{1}{\rho_\delta}(\mu^{n,t}-\rho_t)\Big)}_{=:\mathcal{I}_1}+\underbrace{\int_0^{\varepsilon^{-\frac{1}{2}}}\dd s\,e^{-\varepsilon s}\,\text{P}_s\Big(\tfrac{1}{\rho_\delta}(\mu^{n,t}-\rho_t)\Big)}_{=:\mathcal{I}_2}.
%\end{equation}
%
%For the regular-part $\mathcal{I}_1$, we directly bound the integrand using the heat-kernel bounds \eqref{Eq28}: for any $x\in\mathcal{M}$
%
%\begin{equation*}
%\bigg\vert\nabla\,\text{P}_s\Big(\tfrac{1}{\rho_\delta}(\mu^{n,t}-\rho_t)\Big)(x)\bigg\vert= \bigg\vert\int_\mathcal{M} \dd y\,\nabla\,p_s(x,y)(\tfrac{1}{\rho_\delta}(\mu^{n,t}-\rho_t))(y)\bigg\vert\lesssim s^{-\frac{3}{2}} \|\mu^{n,t}-\rho_t\|_{\LL^{\infty}},
%\end{equation*}
%
%so that
%
%\begin{equation}\label{Eq9}
%\|\nabla\mathcal{I}_1\|_{\LL^{\infty}(\mathcal{M})}\lesssim \|\mu^{n,t}-\rho_t\|_{\LL^{\infty}(\mathcal{M})}\int_{\varepsilon^{-\frac{1}{2}}}^{\infty}\dd s\,s^{-\frac{3}{2}}\lesssim \varepsilon^{\frac{1}{4}}\|\mu^{n,t}-\rho_t\|_{\LL^{\infty}(\mathcal{M})}.
%\end{equation}
%
%For the singular-scale part $\mathcal{I}_2$, we first rewrite it by subtracting and adding back $\frac{1}{\rho_{\delta}(x)}$ in the argument of $\text{P}_s$ in form of, for any $x\in\mathcal{M}$
%
%\begin{align*}
%\mathcal{I}_2(x)=&\int_{0}^{\varepsilon^{-\frac{1}{2}}}\dd s\, e^{-\varepsilon s}\text{P}_s((\tfrac{1}{\rho_{\delta}}-\tfrac{1}{\rho_{\delta}(x)})(\mu^{n,t}-\rho_t))(x)+\frac{1}{\rho_{\delta}(x)}\int_{0}^{\varepsilon^{-\frac{1}{2}}}\dd s\, e^{-\varepsilon s}\text{P}_s(\mu^{n,t}-\rho_t)(x).
%\end{align*}
%
%The second right-hand side term, using the semi-group property $\text{P}_{s}\text{P}_t=\text{P}_{s+t}$, simplifies into
%
%$$\frac{1}{\rho_{\delta}(x)}\int_{0}^{\varepsilon^{-\frac{1}{2}}}\dd s\, e^{-\varepsilon s}\text{P}_s(\mu^{n,t}-\rho_t)(x)=\frac{1}{\rho_{\delta}(x)}\int_{t}^{t+\varepsilon^{-\frac{1}{2}}}\dd s\, e^{-\varepsilon (s-t)}(\mu^{n,s}-\rho_s)(x),$$
%
%so that
%
%\begin{equation}\label{Eq60}
%\mathcal{I}_2(x)=\int_{0}^{\varepsilon^{-\frac{1}{2}}}\dd s\, e^{-\varepsilon s}\text{P}_s((\tfrac{1}{\rho_{\delta}}-\tfrac{1}{\rho_{\delta}(x)})(\mu^{n,t}-\rho_t))(x)+\frac{1}{\rho_{\delta}(x)}\int_{t}^{t+\varepsilon^{-\frac{1}{2}}}\dd s\, e^{-\varepsilon (s-t)}(\mu^{n,s}-\rho_s)(x).
%\end{equation}
%
%Thus, together with \eqref{Eq25}, the heat kernel bounds \eqref{Eq28} and
%
%\begin{equation}\label{Eq29}
%\Big\vert \tfrac{1}{\rho_\delta(y)}-\tfrac{1}{\rho_\delta(x)}\Big\vert\stackrel{\eqref{Eq25}}{\lesssim} \delta^{-1}\vert x-y\vert\quad \text{for any $x, y\in\mathcal{M}$},
%\end{equation}
%
%we obtain
%
%\begin{equation}\label{Eq27}
%\|\nabla \mathcal{I}_2\|_{\LL^{\infty}}\lesssim \delta^{-1}t^{-1}\varepsilon^{-\frac{1}{2}}+\delta^{-1}t^{-\frac{1}{2}}.
%\end{equation}
%
%The combination of \eqref{Eq26}, \eqref{Eq9}, \eqref{Eq27} and \eqref{Eq:ChoicesPara} yields
%
%\begin{equation}\label{Eq12}
%\mathds{1}_{\mathcal{A}_n}\varepsilon\|u^{n,t}_{\varepsilon,\delta}\|_{\cc^{0,\alpha}}\lesssim \mathds{1}_{\mathcal{A}_n}\varepsilon\|\nabla u^{n,t}_{\varepsilon,\delta}\|_{\LL^{\infty}}\lesssim \log^{-\frac{1}{2}}(n).
%\end{equation}
%
%\nc{Continue Here}.
%Finally, the combination of \eqref{Eq11}, \eqref{Eq7} and \eqref{Eq12} implies \eqref{Eq24}.

%
\medskip
%

{\sc Step 2. Proof of \eqref{Eq13}. }We provide the arguments for \eqref{Eq13} in case that \eqref{DecayAlphaMixing} holds for $\eta<1$ and the event $\mathcal{B}^c_{\delta,n}$ that we reduce to $\{\|\nabla^2 u^{n,t}_\delta\|_{\LL^{\infty}}\leq \tfrac{1}{\log^{\upsilon}(n)}\}$, the other cases as well as the case of $\nabla u^{n,t}_\delta$ follow by a straightforward adaptation. The estimate \eqref{Eq13} is a consequence of 
%
\begin{equation}\label{Eq59}
\begin{aligned}
\sup_{x\in\mathcal{M}}\mathbb{P}\Big(\vert\partial^2_{ij} u^{n,t}_{\delta}(x)\vert\geq& \tfrac{1}{2\log^{\upsilon}(n)}\Big)\leq n\exp\bigg(-\frac{1}{C_1}\Big(\frac{nt}{\log^{\upsilon}(n)}\Big)^{\eta}\bigg)+\exp\bigg(-\frac{1}{C_2}\frac{n^2t\log^{-2\upsilon}(n)}{t^{-1}+n\delta^{-4}\log^{\frac{1}{\eta}}(n)}\bigg)\\
&+\exp\bigg(-\frac{1}{C_3}\frac{nt^2}{\log^{\upsilon}(n)}\exp\Big(\frac{1}{C_4}\Big(\frac{nt}{\log^{\upsilon}(n)}\Big)^{\eta(1-\eta)}\log^{-1}(\tfrac{nt}{\log^{\upsilon}(n)})\Big)\bigg).
\end{aligned}
\end{equation}
% 
for some constants $C_1,C_2,C_3,C_4>0$. To see this, let $\eta_n$ be defined by
%
\begin{equation}\label{Eq72}
\eta_n=\frac{1}{2\log^{\upsilon}(n)\|\nabla\partial^2_{ij}u^{n,t}_{\delta}\|_{\LL^{\infty}}}.
\end{equation}
%
By compactness of $\mathcal{M}$, we can find a $\eta_n$-net $\{x_k\}_{1\leq k\leq N}\subset \mathcal{M}$ with $N\lesssim \eta^{-2}_n$. We note that
%
\begin{equation}\label{Eq77}
\Big\{\|\partial^2_{ij}u^{n,t}_{\delta}\|_{\LL^{\infty}}>\tfrac{1}{\log^{\upsilon}(n)}\Big\}\subset\bigcup_{k=1}^{N}\Big\{\vert\partial^2_{ij}u^{n,t}_{\delta}(x_k)\vert>\tfrac{1}{2\log^{\upsilon}(n)}\Big\}.
\end{equation}
%
Indeed, if for any $k\in\{1,\cdots,N\}$, $\vert\partial^2_{ij}u^{n,t}_{\delta}(x_k)\vert\leq\tfrac{1}{2\log^{\upsilon}(n)}$ then for any $x\in \mathcal{M}$ there exists $j\in\{1,\cdots,N\}$ such that
%
$$\vert\partial^2_{ij}u^{n,t}_{\delta}(x)\vert\leq \eta_n\|\nabla\partial^2_{ij}u^{n,t}_{\delta}\|_{\LL^{\infty}}+\vert\partial^2_{ij}u^{n,t}_{\delta}(x_j)\vert\stackrel{\eqref{Eq72}}{\lesssim} \frac{1}{\log^{\upsilon}(n)}.$$
%
Applying $\mathbb{P}$ on \eqref{Eq77} yields 
%
\begin{align*}
\mathbb{P}\Big(\|\partial^2_{ij} u^{n,t}_{\delta}\|_{\LL^{\infty}}\geq \tfrac{1}{\log^{\upsilon}(n)}\Big)&\leq \sum_{k=1}^N\mathbb{P}\Big(\vert\partial^2_{ij} u^{n,t}_{\delta}(x_k)\vert\geq \tfrac{1}{2\log^{\upsilon}(n)}\Big)\\
&\lesssim \eta^{-2}_{n}\sup_{x\in\mathcal{M}}\mathbb{P}\Big(\vert\partial^2_{ij} u^{n,t}_{\delta}(x)\vert\geq \tfrac{1}{2\log^{\upsilon}(n)}\Big).
\end{align*}
%
Using \eqref{Eq59} and 
%
$$\eta^{-2}_n\lesssim \log^{2\upsilon}(n)t^{-3},$$
%
which can be proven following the arguments leading to the second item of \eqref{Eq58} below, this yields $\P(\{\|\nabla^2 u^{n,t}_\delta\|_{\LL^{\infty}}\leq \tfrac{1}{\log^{\upsilon}(n)}\})=o(\tfrac{1}{n^m})$ for any $m\in\mathbb{N}$ for the choice $\kappa_2>4\kappa_1+\frac{1}{\eta}+2\upsilon$.
%
%
\medskip
%

We now prove \eqref{Eq59}. We first exploit the following explicit representation formula\footnote{a simple change of variable gives that the kernel $\tilde{p}_s$ associated to $1-\Delta$ is given by $\tilde{p}_s=e^{-s}p_s$} for $u^{n,t}_{\delta}$, 
%
\begin{equation}\label{eq:stella3}
u^{n,t}_{\delta}=\int_0^{\infty}e^{-s}\,\text{P}_s\Big(\tfrac{1}{\rho_\delta}(\mu^{n,t}-\rho_t)\Big)\,\dd s,
\end{equation}
%
that we expand, using the definition \eqref{eq:empmeas} of $\mu^{n}$, in form of
%
\begin{equation}\label{eq:stella}
u^{n,t}_{\delta}=\frac{1}{n}\sum_{k=1}^n\omega(\cdot,X_k)\quad\text{with } \omega(\cdot,y):=\int_{0}^{\infty}e^{-s}\,\text{P}_s\Big(\tfrac{1}{\rho_{\delta}}(p_t(\cdot,y)-\rho_t)\Big)\,\dd s.
\end{equation}
%
Then applying the concentration inequalities in Proposition \ref{BersteinCorrelated}, we get 
%
\begin{equation}\label{eq:stella2}
\begin{split}
\mathbb{P}\Big(\vert\partial^2_{ij} u^{n,t}_{\delta}(x)\vert &\geq \tfrac{1}{\log^{\upsilon}(n)}\Big)\leq n\exp\bigg(-\frac{1}{C_1}\Big(\frac{n}{M\log^{2}(n)}\Big)^{\beta}\bigg)+\exp\bigg(-\frac{1}{C_2}\frac{n^2\log^{-2\upsilon}(n)}{M^2+n v^2}\bigg)\\
&+\exp\bigg(-\frac{1}{C_3}\frac{n\lambda}{M^2\log^{\upsilon}(n)}\exp\Big(\frac{1}{C_4}\Big(\frac{n}{M\log^{\upsilon}(n)}\Big)^{\beta(1-\beta)}\log^{-1}(\tfrac{n}{M\log^{\upsilon}(n)})\Big)\bigg),
\end{split}
\end{equation}
%
where 
%
\begin{equation}\label{Eq79}
M:=\sup_{y\in\mathcal{M}}\vert\partial^2_{ij}\omega(x,y)\vert\quad\text{and}\quad v^2:=\mathbb{E}[\vert\partial^2_{ij}\omega(x,X_1)\vert^2]+2\sup_{\ell\geq 1}\sum_{k>\ell}\vert\mathbb{E}[\partial^2_{ij}\omega(x,X_\ell)\partial^2_{ij}\omega(x,X_k)]\vert.
\end{equation}
%
The estimate \eqref{Eq59} then follows from the three following estimates
%
\begin{equation}\label{Eq58}
\mathbb{E}[\vert\partial^2_{ij} \omega(x,X_1)\vert^2]\lesssim \delta^{-4}t^{-1},\quad \sup_{y\in\mathcal{M}}\vert\partial^2_{ij} \omega(x,y)\vert\lesssim t^{-1}\quad\text{and}\quad v^2\lesssim \log^{\frac{1}{\eta}}(n)\delta^{-4}t^{-1} ,
\end{equation}
%
that we prove separately in the next three sub-steps.

%
\medskip
%

{\sc Sub-step 2.1. Proof of the first item of \eqref{Eq58}. }Splitting the time integral into $\int_{0}^t+\int_{t}^\infty$ and subtracting and adding back $\tfrac{1}{\rho_\delta(x)}$ in the first integral as well as using the semigroup property of $\{\pp_s\}_{s}$ in form of 
%
$$\int_0^t\dd s\, e^{-s}\,\pp_s(p_t(\cdot,y)-\rho_t)=\int_{t}^{2t}\dd s\, e^{-(s-t)}(p_s(\cdot,y)-\rho_s),$$
%
we decompose $\omega$ into a regular-part $\mathcal{J}_1$ and a singular-part $\mathcal{J}_2$: 
%
\begin{equation}\label{Eq68}
\begin{aligned}
\omega(x,y)=&\underbrace{\frac{1}{\rho_{\delta}(x)}\int_t^{2t} \dd s\, e^{-s}(p_s(\cdot,y)-\rho_{s})(x)
+\int_{t}^{\infty}\dd s\, e^{-s}\text{P}_s\Big(\tfrac{1}{\rho_\delta}(p_t(\cdot,y)-\rho_t)\Big)(x)}_{=:\mathcal{J}_1(x,y)}\\
&+\underbrace{\int_{0}^{t}\dd s\, e^{-s}\text{P}_s\Big((\tfrac{1}{\rho_\delta}-\tfrac{1}{\rho_\delta(x)})(p_t(\cdot,y)-\rho_t)\Big)(x)}_{=:\mathcal{J}_2(x,y)}.
\end{aligned}
\end{equation}
%
For the regular-part $\mathcal{J}_1$, we apply directly the heat-kernel bounds \eqref{Eq28} and the first item of \eqref{Eq25}, to obtain from the triangle inequality and Minkowski's inequality
%
\begin{equation}\label{Eq61}
\begin{aligned}
\mathbb{E}[\vert\partial^2_{ij}\mathcal{J}_1(x,X_1)\vert^2]\lesssim& \|\partial^2_{ij}\mathcal{J}_1(x,\cdot)\|^2_{\LL^{2}}\\
\lesssim &(\|\nabla\rho_{\delta}\|^4_{\LL^{\infty}}+\|\nabla^2\rho_{\delta}\|^2_{\LL^{\infty}})\Big(\int_{t}^{2t}\|p_s(x,\cdot)\|_{\LL^{2}}\,\dd s\Big)^{2}\\
&\underbrace{+\|\nabla\rho_{\delta}\|^2_{\LL^{\infty}}\Big(\int_{t}^{2t}\|\nabla p_s(x,\cdot)\|_{\LL^{2}}\,\dd s\Big)^{2}+\Big(\int_{t}^{2t}\|\nabla^2 p_s(x,\cdot)\|_{\LL^{2}}\,\dd s\Big)^{2}}_{=:\mathcal{R}_1(x)}\\
&+\underbrace{\bigg(\int_{t}^{\infty}\dd s\bigg(\int_{\mathcal{M}}\dd \m(y)\Big(\partial^2_{ij}\text{P}_s\Big(\tfrac{1}{\rho_{\delta}}(p_t(\cdot,y)-\rho_t)\Big)(x)\Big)^2\bigg)^{\frac{1}{2}}\bigg)^2}_{=:\mathcal{R}_2(x)}.
\end{aligned}
\end{equation}
%
The first \rhs term $\mathcal{R}_1(x)$ is dominated using directly the heat-kernel bounds \eqref{Eq28} and \eqref{Eq25}
%
\begin{equation}\label{Eq62}
\mathcal{R}_1(x)\lesssim \delta^{-4}t^{\frac{1}{2}}+\delta^{-2}+t^{-1}\lesssim t^{-1}.
\end{equation}
%
For the second \rhs side term $\mathcal{R}_2(x)$, we first simplify the $y$-integral. Using that 
%
\begin{equation}\label{Eq63}
\int_{\mathcal{M}}\vert p_t(z,w)-\rho_t(z)\vert\,\dd \m(w)\lesssim 1\quad\text{and}\quad\int_{\mathcal{M}}\vert p_t(w,y)-\rho_t(w)\vert\,\dd \m(w)\lesssim 1\quad \text{for any $z,y\in \mathcal{M}$,}
\end{equation}
%
we have by Jensen's inequality and the heat-kernel bounds \eqref{Eq28}
%
\begin{align*}
\lefteqn{\int_{\mathcal{M}}\Big(\partial^2_{ij}\text{P}_s\Big(\tfrac{1}{\rho_{\delta}}(p_t(\cdot,y)-\rho_t)\Big)(x)\Big)^2\,\dd \m(y)}\\ &=\int_{\mathcal{M}}\bigg(\int_{\mathcal{M}}\dd w\,\partial^2_{ij}p_{s}(x,w)\tfrac{1}{\rho_{\delta}(w)}(p_t(w,y)-\rho_t(w))\bigg)^{2}\,\dd \m(y)\\
&\stackrel{\eqref{Eq63}}{\lesssim}\int_{\mathcal{M}}\vert\partial^2_{ij}p_s(x,w)\vert^2\,\dd \m(w)\lesssim s^{-3}.
\end{align*}
%
Thus, 
%
\begin{equation}\label{Eq64}
\mathcal{R}_2(x)\lesssim t^{-1}.
\end{equation}
%
The combination of \eqref{Eq61}, \eqref{Eq62} and \eqref{Eq63} yields
%
\begin{equation}\label{Eq65}
\mathbb{E}[\vert\partial^2_{ij}\mathcal{J}_1(x,X_1)\vert^2]\lesssim t^{-1}.
\end{equation}
%
We now turn to the singular-part $\mathcal{J}_2$. We first apply Minkowski's inequality in form of
%
\begin{equation}\label{Eq66}
\begin{split}
\mathbb{E}[\vert\mathcal{J}_2(x,X_1)\vert^2]&\lesssim \|\partial^2_{ij}\mathcal{J}_2(x,\cdot)\|^2_{\LL^{2}}\\& \lesssim \bigg(\int_{0}^t\dd s\bigg(\int_{\mathcal{M}}\dd \m(y)\Big(\partial^2_{ij}\text{P}_s\Big((\tfrac{1}{\rho_{\delta}}-\tfrac{1}{\rho_{\delta}(x)})(p_t(\cdot,y)-\rho_t)\Big)(x)\Big)^2\bigg)^{\frac{1}{2}}\bigg)^{2}.
\end{split}
\end{equation}
%
We then simplify the $y$-integral. To this aim, we bound the integrand in $\LL^{\infty}$ using the heat-kernel bounds \eqref{Eq28}, \eqref{Eq25} and 
%
\begin{equation}\label{Eq29}
\Big\vert \tfrac{1}{\rho_\delta(y)}-\tfrac{1}{\rho_\delta(x)}\Big\vert\stackrel{\eqref{Eq25}}{\lesssim} \delta^{-1}\dd(x,y)\quad \text{for any $x, y\in\mathcal{M}$},
\end{equation}
%
in form of
%
\begin{equation}\label{Eq69}
\begin{aligned}
& \Big\vert\partial^2_{ij}\text{P}_s\Big((\tfrac{1}{\rho_{\delta}}-\tfrac{1}{\rho_{\delta}(x)})(p_t(\cdot,y)-\rho_t)\Big)(x)\Big\vert\lesssim t^{-1}\int_{\mathcal{M}}\bigg\vert \partial^2_{ij}\Big(p_s(x,w)(\tfrac{1}{\rho_{\delta}(w)}-\tfrac{1}{\rho_{\delta}(x)})\Big)\bigg\vert\,\dd \m(w)\\
\stackrel{\eqref{Eq25},\eqref{Eq29}}{\lesssim}&t^{-1}\delta^{-1}\bigg(\int_{\mathcal{M}}\vert\nabla^2 p_s(x,w)\vert \dd(x,w)\,\dd \m(w)+\int_{\mathcal{M}}\vert\nabla p_s(x,w)\vert\,\dd \m(w)\bigg)\\
& +t^{-1}\delta^{-2}\int_{\mathcal{M}}\vert p_s(x,w)\vert\,\dd \m(w)\\
\stackrel{\eqref{Eq28}}{\lesssim} &t^{-1}\delta^{-1}s^{-\frac{1}{2}}+t^{-1}\delta^{-2}.
\end{aligned}
\end{equation}
%
This yields together with \eqref{Eq66}
%
\begin{equation}\label{Eq67}
\mathbb{E}[\vert\mathcal{J}_2(x,X_1)\vert^2]\lesssim \delta^{-4}t^{-1}.
\end{equation}
%
To conclude, the combination of \eqref{Eq68}, \eqref{Eq65} and \eqref{Eq67} shows the first item of \eqref{Eq58}.

%
\medskip
%

{\sc Sub-step 2.2. Proof of the second item of \eqref{Eq58}. }We use the decomposition \eqref{Eq68}. For the regular-part $\mathcal{J}_1$, we argue as in \eqref{Eq61} for the first term whereas the second-term is estimated using the heat-kernel bounds \eqref{Eq28} in form of
%
\begin{align*}
\bigg\vert\partial^2_{ij}\text{P}_s\Big(\tfrac{1}{\rho_\delta}(p_t(\cdot,y)-\rho_t)\Big)(x)\bigg\vert&=\bigg\vert\int_{\mathcal{M}}\dd \m(w)\,\partial^2_{ij}p_{s}(x,w)\tfrac{1}{\rho_{\delta}(w)}(p_t(w,y)-\rho_t(w))\bigg\vert\\
&\stackrel{\eqref{Eq28},\eqref{Eq25}}{\lesssim} s^{-2}\int_{\mathcal{M}}\dd \m(w)\, \vert p_t(w,y)-\rho_t(w)\vert\\
&\lesssim s^{-2},
\end{align*}
%
so that
%
\begin{align*}
\bigg\vert\int_{t}^{\infty}\dd s\, e^{-s}\partial^2_{ij}\text{P}_s\Big(\tfrac{1}{\rho_\delta}(p_t(\cdot,y)-\rho_t)\Big)(x)\bigg\vert\lesssim t^{-1}.
\end{align*}
%
Hence,
%
\begin{equation}\label{Eq71}
\sup_{y\in\mathcal{M}}\vert\mathcal{J}_1(x,y)\vert\lesssim t^{-1}.
\end{equation}
%
%
\medskip
%

For the singular-part $\mathcal{J}_2$, we use the bound \eqref{Eq69} which directly yields
%
\begin{equation}\label{Eq70}
\sup_{y\in\mathcal{M}}\vert\mathcal{J}_2(x,y)\vert\lesssim \delta^{-1}t^{-\frac{1}{2}}.
\end{equation}
%
The combination of \eqref{Eq68}, \eqref{Eq71} and \eqref{Eq70} gives the second item of \eqref{Eq58}.

%
\medskip
%

{\sc Sub-Step 2.3. Proof of the third item of \eqref{Eq58}. }According to the first item of \eqref{Eq58}, it suffices to give the argument for the second term in the definition \eqref{Eq79} of $v^2$. We use the assumption \eqref{DecayAlphaMixing} together with the two first items of \eqref{Eq58} in form of
%
\begin{align*}
%
\lefteqn{\sum_{k>\ell}\vert\mathbb{E}[\partial^2_{ij}\omega(x,X_i)\partial^2_{ij}\omega(x,X_j)]\vert}\\ &=\sum_{k>\ell}\vert\mathbb{E}[\partial^2_{ij}\omega(x,X_\ell)\partial^2_{ij}\omega(x,X_k)]\vert^{\frac{1}{\log(n)}}\vert\mathbb{E}[\partial^2_{ij}\omega(x,X_\ell)\partial^2_{ij}\omega(x,X_k)]\vert^{1-\frac{1}{\log(n)}}\\
&\lesssim \sum_{k>\ell}\alpha^{\frac{1}{\log(n)}}_{k-\ell}(\sup_{y\in\mathcal{M}}\vert \partial^2_{ij}\omega(x,y)\vert)^{\frac{2}{\log(n)}}(\mathbb{E}[\vert\partial^2_{ij}\omega(x,X_1)\vert^2])^{1-\frac{1}{\log(n)}}\\
&\stackrel{\eqref{Eq58},\eqref{DecayAlphaMixing}}{\lesssim} \delta^{-4(1-\frac{1}{\log(n)})} t^{-1-\frac{1}{\log(n)}}\sum_{k=0}^{\infty}\exp(-b\tfrac{k^{\eta}}{\log(n)})\\
&\lesssim \delta^{\frac{4}{\log(n)}}t^{-\frac{1}{\log(n)}}\log^{\frac{1}{\eta}}(n)\delta^{-4}t^{-1},
%
\end{align*}
%
which concludes since $\limsup_{n\uparrow\infty}\delta^{\frac{4}{\log(n)}}t^{-\frac{1}{\log(n)}}\lesssim 1$. 
%First, using Minkowski's inequality, we have
%
%\begin{equation}\label{Eq14}
%\mathbb{E}[\vert\partial^2_{ij} u_1(x)\vert^2]\leq \left(\int_{0}^{\infty}\dd s\, e^{-\varepsilon s}\,\mathbb{E}\bigg[\Big\vert\partial^2_{ij}\text{P}_s\Big(\tfrac{1}{\rho_{\delta}}(p_t(\cdot,X_1)-\rho_t)\Big)(x)\Big\vert^2\bigg]^{\frac{1}{2}}\right)^2.
%\end{equation}
%

%We then split the time interval into $\int_0^{\infty}=\int_{0}^t+\int_t^{\infty}$. For the regime $s\leq t$, we use that $\partial^2_{ij}$ and $\text{P}_s$ commutes (\nc{As notice below, this two operators do not commute. But we can proceed as below, by subtracting $\tfrac{1}{\rho_{\delta}(x)}$}) as well as once more Minkowski's inequality to obtain
%
%$$\mathbb{E}\bigg[\Big\vert\partial^2_{ij}\text{P}_s\Big(\tfrac{1}{\rho_{\delta}}(p_t(\cdot,X_1)-\rho_t)\Big)(x)\Big\vert^2\bigg]^{\frac{1}{2}}\leq \int_{\mathcal{M}}\dd x\,p_s(x,y)\mathbb{E}\Big[\vert\partial^2_{ij}(\tfrac{1}{\rho_{\delta}}(p_t(\cdot,X_1)-\rho_t))(y)\vert^2\Big]^{\frac{1}{2}}.$$
%
%Then, observe that $\mathbb{E}[\partial^2_{ij}(\tfrac{1}{\rho_\delta}p_t(\cdot, X_1))(y)]=\partial^2_{ij}(\tfrac{1}{\rho_\delta}\rho_t)(y)$, so that
%
%\begin{align*}
%\mathbb{E}\Big[\vert\partial^2_{ij}(\tfrac{1}{\rho_{\delta}}(p_t(\cdot,X_1)-\rho_t)(y))\vert^2\Big]&=\mathbb{E}\Big[\vert\partial^2_{ij}(\tfrac{1}{\rho_{\delta}}(p_t(\cdot,X_1)))(y)\vert^2\Big]-(\partial^2_{ij}(\tfrac{1}{\rho_\delta}\rho_t)(y))^2\\
%&\leq \mathbb{E}\Big[\vert\partial^2_{ij}(\tfrac{1}{\rho_{\delta}}(p_t(\cdot,X_1)))(y)\vert^2\Big].
%\end{align*}
%
%We thus deduce
%
%$$\mathbb{E}\bigg[\Big\vert\partial^2_{ij}\text{P}_s\Big(\tfrac{1}{\rho_{\delta}}(p_t(\cdot,X_1)-\rho_t)\Big)(x)\Big\vert^2\bigg]^{\frac{1}{2}}\leq \int_{\mathcal{M}}\dd x\,p_s(x,y)\mathbb{E}\Big[\vert\partial^2_{ij}(\tfrac{1}{\rho_{\delta}}(p_t(\cdot,X_1)))(y)\vert^2\Big]^{\frac{1}{2}}.$$
%
%Then, using that $X_1\sim \rho\,\dd x$ and 
%
%$$\sup_{y\in\mathcal{M}}\int_\mathcal{M}\,\dd z\,\Big(\vert \nabla^2p_t(z,y)\vert^2+t^{-2}\vert\nabla p_t(z,y)\vert^2\Big)\lesssim t^{-3},$$
%
%we obtain
%\begin{align*}
%\mathbb{E}\Big[\vert\partial^2_{ij}(\tfrac{1}{\rho_{\delta}}(p_t(\cdot,X_1)))(y)\vert^2\Big]=&\int_{\mathcal{M}}\rho(z)\,\dd z\,\vert\partial^2_{ij}(\tfrac{1}{\rho_{\delta}}(p_t(\cdot,z)))(y)\vert^2\\
%\lesssim & \int_{\mathcal{M}}\dd z\vert\partial^2_{ij}p_t(y,z)\vert^2+\delta^{-2}\\
%&+\delta^{-1}\int_{\mathcal{M}}\dd z \,(\vert\partial_i p_t(z,y)\vert^2+\vert\partial_j p_t(z,y)\vert^2)\\
%\lesssim&\, t^{-3}.
%\end{align*}
%
%In sum, 
%
%\begin{equation}\label{Eq15}
%\left(\int_{0}^{t}\dd s\, e^{-\varepsilon s}\,\mathbb{E}\bigg[\Big\vert\partial^2_{ij}\text{P}_s\Big(\tfrac{1}{\rho_{\delta}}(p_t(\cdot,X_1)-\rho_t)\Big)(x)\Big\vert^2\bigg]^{\frac{1}{2}}\right)^2\lesssim t^{-1}.
%\end{equation}
%
%For the regime $s\geq t$, we express by definition
%
%\begin{align*}
%
%&\mathbb{E}\bigg[\Big\vert\partial^2_{ij}\text{P}_s\Big(\tfrac{1}{\rho_{\delta}}(p_t(\cdot,X_1)-\rho_t)\Big)(x)\Big\vert^2\bigg]^{\frac{1}{2}}\\
%&=\bigg(\int_{\mathcal{M}}\rho(z)\,\dd z\,\bigg(\int_{\mathcal{M}}\dd y\,\partial^{2}_{ij}p_s(x,y)(\tfrac{1}{\rho_{\delta}}(p_t(\cdot,z)-\rho_t))(y)\bigg)^2\bigg)^{\frac{1}{2}}.
%
%\end{align*}
%
%Then, using the fact that $\sup_{z\in\mathcal{M}}\Big\vert\int_{\mathcal{M}}\dd y\,(\tfrac{1}{\rho_\delta}(p_t(\cdot,z)-\rho_t))(y)\Big\vert\lesssim 1$ and $\sup_{y\in\mathcal{M}}\Big\vert\int_{\mathcal{M}}\dd z\,(\tfrac{1}{\rho_\delta}(p_t(\cdot,z)-\rho_t))(y)\Big\vert\lesssim 1$, we have by Jensen's inequality
%
%\begin{align*}
%\mathbb{E}\bigg[\Big\vert\partial^2_{ij}\text{P}_s\Big(\tfrac{1}{\rho_{\delta}}(p_t(\cdot,X_1)-\rho_t)\Big)(x)\Big\vert^2\bigg]^{\frac{1}{2}}&\lesssim \bigg(\int_{\mathcal{M}}\dd z\int_{\mathcal{M}}\dd y\,(\partial^{2}_{ij}p_s(x,y))^2(\tfrac{1}{\rho_{\delta}}(p_t(\cdot,z)-\rho_t))(y)\bigg)^{\frac{1}{2}}\\
%&\lesssim \bigg(\int_{\mathcal{M}}\dd y\,(\partial^{2}_{ij}p_s(x,y))^2\bigg)^{\frac{1}{2}}\\
%&\lesssim s^{-\frac{3}{2}}.
%\end{align*}
%
%We thus deduce 
%
%\begin{equation}\label{Eq16}
%\left(\int_{t}^{\infty}\dd s\, e^{-\varepsilon s}\,\mathbb{E}\bigg[\Big\vert\partial^2_{ij}\text{P}_s\Big(\tfrac{1}{\rho_{\delta}}(p_t(\cdot,X_1)-\rho_t)\Big)(x)\Big\vert^2\bigg]^{\frac{1}{2}}\right)^2\lesssim \bigg(\int_{t}^{\infty}\dd s\,s^{-\frac{3}{2}}\bigg)^2\lesssim t^{-1}.
%\end{equation}
%}
%To conclude, the combination of \eqref{Eq14}, \eqref{Eq15} and \eqref{Eq16} shows that 
%

%
%\section{Old proof}
%
%{\sc Step 2. Proof of \eqref{eq:ConvergenceRateMain}. }According to \eqref{eq:thm1explip}, it remains to show that 
%
%\begin{equation}\label{MainResultProof}
%
%\int_{\mathcal{M}}\dd\m\, \dd^2\big(T^n,\exp(\nabla f^{n,t}_\delta)\big)\leq \mathcal{C}_n\frac{\log(n)}{n}\Big(\log^{\gamma_1-1}(n)\|\rho_t-\rho\|_{\LL^1}+\sqrt{\tfrac{\log\log(n)}{\log(n)}}\Big),
%
%\end{equation}
%
%In addition, we only need to show the validity of \eqref{MainResultProof} in the event $\mathcal{A}_n\cap \mathcal{B}_{\delta,n}$ defined in \eqref{Events} thanks to \eqref{Eq13}. Since \eqref{Eq6} holds in $\mathcal{A}_n\cap \mathcal{B}_{\delta,n}$, we can apply for large $n$ the quantitative stability result \cite[Theorem 3.2]{ambrosio2019optimal} to $\mu_1=\mu^n$, $\mu_2=\exp(\nabla f^{n,t}_\delta)_\#\rho$ and $\nu=\rho$ to the effect of 
%
%\begin{equation*}
%
%\int_{\mathcal{M}}\dd\m\,\dd^2 \big(T^{n}, \exp(\nabla f^{n,t}_\delta)\big) \lesssim W_2^2 \big(\mu^n,\exp(\nabla f^{n,t}_\delta)_\#\rho\big)+W_2\big(\mu^n,\exp(\nabla f^{n,t}_\delta)_\#\rho\big) W_2\big(\rho,\mu^n\big).
%
%\end{equation*}
%
%\nc{Continue here.}

%We argue that 
%\begin{equation}\label{eq:thm1estmap}
%\int_{\mathcal{M}}\dd^2 (T^{n}, \exp(\nabla f^{n,t}_\delta))\,\dd m \lesssim W_2^2 (\hat{\mu}_\delta^{n,t}, \mu^n) + W_2^2 (\hat{\mu}_\delta^{n,t}, \mu^n) W_2^2 (\mu_n, \rho\, \dd x),
%\end{equation}
%where $\hat{\mu}_\delta^{n,t} : = \exp(\nabla f_\delta^{n,t})_{\#} \rho \, \dd x$. Indeed, by boundedness of $\rho$ we might smuggle in $\rho \dd m$ and apply \cite[Lemma 2.2]{ambrosio2019optimal} with $\nu = \rho \, \dd x, \mu_1 = \hat{\mu}_\delta^{n,t} $ and $\mu_2 = \mu^n$ \fm{(check if $\mu^n$ has been defined somewhere)} to get 
%\[
%\begin{split}
%\int_{\mathcal{M}}\dd^2 (T^{n}, \exp(\nabla f^{n,t}_\delta))\,\dd m & \lesssim \int_{\mathcal{M}}\dd^2 (T^{n}, \exp(\nabla f^{n,t}_\delta))\rho \,\dd m  \\
%& \lesssim  W_2^2 (, \mu^n) + W_2 (\hat{\mu}_\delta^{n,t}, \mu^n) W_2 (\mu_n, \rho\, \dd x).
%\end{split}
%\]



%Taking the expectation in \eqref{eq:proofThm1-1} we can bound the two terms separately. Let us firstly consider the second term. By Lipschitzianity of the exponential map \nc{on compact set ? Use \cite[Lemma 2.2]{ambrosio2019optimal}}, Minkowski's inequality and \eqref{eq:errornablafdelta} we may write \fm{Check the exponent $p,q$}
%\begin{equation}\label{eq:proofThm1-2}
%\begin{split}
%\mathbb{E} \big[ \int_\mathcal{M} \dd^2(\exp(\nabla f^{n,t}_\delta), \exp(\nabla f^{n,t})) \big] & \le
%\mathbb{E} \big[ \int_\mathcal{M} | \nabla f^{n,t}_\delta - \nabla f^{n,t}|^2 \big]\\
%& \lesssim \mathbb{E} \big[ \big( \int_\mathcal{M} | \nabla f^{n,t}_\delta - \nabla f^{n,t}|^2\big)^p\big]^\frac1p \\
%& \lesssim \| \rho_\delta - \rho\|^2_{L^q} \frac{|\log (t)|}n.
%\end{split}
%\end{equation}
%Let us now turn to the first term of \eqref{eq:proofThm1-1}. We want to apply Theorem \cite[Theorem 3.2]{ambrosio2019optimal} with $\nu = \rho_t \, \dd x, \mu_1 = \hat{\mu}_\delta^{n,t} : = (\text{Id} + \nabla f_\delta^{n,t})_{\#} \rho_t \, \dd x$ and $\mu_2 = \mu^n$ to get 
%\[
%\begin{split}
%\int_{\mathcal{M}}\vert T^{n,t}-(\text{Id}+\nabla f^{n,t}_\delta)\vert^2  & \lesssim \int_{\mathcal{M}}\vert T^{n,t}-(\text{Id}+\nabla f^{n,t}_\delta)\vert^2 \rho_t \, \dd x \\
%& \lesssim  W_2^2 (\hat{\mu}_\delta^{n,t}, \mu^n) + W_2^2 (\hat{\mu}_\delta^{n,t}, \mu^n) W_2^2 (\mu_n, \rho\, \dd x) .
%\end{split}
%\]
%We now need to estimate the right hand side. 
%
%\medskip
%

%\medskip
%{\sc Step 3.} Conclusion. Applying $\mathbb{E}[|\cdot|^p]^\frac1p$ in \eqref{eq:thm1explip} we get
%\begin{equation}\label{eq:thm1avexplip}
%\begin{split}
%\mathbb{E} \bigg[ \bigg(\int_\mathcal{M} \dd^2 (\exp(\nabla f^{n,t}_\delta), \exp(\nabla f^{n,t}))\,\dd m\bigg)^p \bigg]^\frac1p & \lesssim \mathbb{E}[\| \nabla f^{n,t}_\delta - \nabla f^{n,t} \|_{L^2}^{2p}]^\frac1p + C_n^2 \\
%& \stackrel{\eqref{eq:errornablafdelta}, \eqref{eq:defcn}}{\lesssim} p \| \rho_\delta - \rho\|^2_{L^{\bar{q}'}} \frac{|\log (t)|}n + \varepsilon^2 \|\rho_\delta - \rho\|_{L^{\bar{q}'}}^2 \frac{\log n}{n}.
%\end{split}
%\end{equation}
%Applying $\mathbb{E}[|\cdot|^p]^\frac1p$ and Cauchy-Schwarz in \eqref{eq:thm1estmap} we get
%\begin{equation}\label{eq:thm1secondterm}
%\begin{split}
%\lefteqn{\mathbb{E}\bigg[ \bigg(\int_{\mathcal{M}}\dd^2 (T^{n}, \exp(\nabla f^{n,t}_\delta))\,\dd m\bigg)^p\bigg]^\frac1p} \\
%&  \lesssim  \mathbb{E} [W_2^{2p} (\mu^n,\hat{\mu}_\delta^{n,t})]^\frac1p + \mathbb{E} [W_2^{2p} (W_2^{2p} (\mu^n,\hat{\mu}_\delta^{n,t})]^\frac1{2p}]\mathbb{E}[ W_2^{2p} (\mu_n, \rho\, \dd x)]^\frac1{2p}.
%\end{split}
%\end{equation}
%Note that by the triangle inequality we may write
%\begin{equation}\label{eq:thm1triangineq}
%\begin{aligned}
%
 %\mathbb{E} \big[W_2^{2p} (\hat{\mu}_\delta^{n,t}, \mu^n)\big]^\frac1p \lesssim &\,\mathbb{E} [W_2^{2p} (\mu^{n}, \mu^{n,t})]^\frac1p +  \mathbb{E} [W_2^{2p} (\exp(\nabla f^{n,t}_{\delta})\#\rho_t\dd x, \exp(\nabla f^{n,t}_{\delta})\#\rho\dd x)]^\frac1p \\
 %& +\mathbb{E} [W_2^{2p} (\mu^{n,t},\exp(\nabla f^{n,t}_{\delta})\#\rho_t\dd x)]^\frac1p. 
%
%\end{aligned}
%\end{equation}
%
%The first r.h.s term is controlled using \nc{cite the correlated verison}, which reads
%
%$$\mathbb{E} [W_2^{2p} (\mu^{n}, \mu^{n,t})]^\frac1p \lesssim \frac{\log\log n}{ n}.$$
%
%The second r.h.s term is controlled using Lemma \ref{lem:transportcontraction}, together with the fact that $\|\nabla^2 f^{n,t}_\delta\|_{\LL^{\infty}(\mathcal{M})}\leq 1$ in $\mathcal{A}_n \cap \mathcal{B}_{\epsilon, \delta, n}$, combined with Lemma \ref{lem:HeatContractionrho}, leading to
%
%\begin{align*}
%
%\mathbb{E} [W_2^{2p} (\exp(\nabla f^{n,t}_{\delta})\#\rho_t\dd x, \exp(\nabla f^{n,t}_{\delta})\#\rho\dd x)]^\frac1p\lesssim t\|\rho_t-\rho\|_{\LL^{1}(\mathcal{M})}.
%
%\end{align*}
%
%For the third r.h.s term, we note that thanks to Proposition \ref{prop:uppboundBenBre} $\mu^{n,t}$ is the push-forward of $m$ throught the flow at time $1$ of the time dependent vector field
%\[
%Y_s = \frac{\rho_\delta \nabla f_\delta^{n,t}}{(1-s) \rho_t + s \mu^{n,t}}.
%\]
%Since we are on the event $\mathcal{A}_n\cap B_{\epsilon, \delta, n}$ we may assume $n$ to be sufficiently large so that we can apply \cite[Proposition A.1]{ambrosio2019finer} with 
%\[
%X = \nabla f_\delta^{n,t}, \quad Y_s = \frac{\rho_\delta \nabla f_\delta^{n,t}}{(1-s) \rho_t + s \mu^{n,t}}.
%\]
%Note that
%\[
%\begin{split}
%|Y_s - X| & \lesssim (|\rho_\delta - \rho| + |\rho_t - \rho| + |\mu^{n,t} - \rho_t|)|\nabla f^{n,t}_\delta| \\
%& \lesssim \bigg(|\rho_\delta - \rho| + |\rho_t - \rho| + \frac1{\log n}\bigg)|\nabla f^{n,t}_\delta|.
%\end{split}
%\]
%Thus, by the first part of \cite[Proposition A.1]{ambrosio2019finer}\footnote{\fm{it seems to me that it's not necessary the uniform bound $0<\xi<1$}} we may write
%\[
%\begin{split}
%\mathbb{E} [W_2^{2p} (\mu^{n,t},\exp(\nabla f^{n,t}_{\delta})\#\rho_t\dd x)]^\frac1p &\lesssim \mathbb{E} \bigg[\bigg(\int_{\mathcal{M}}\bigg(|\rho_\delta - \rho| + |\rho_t - \rho| + \frac1{\log n}\bigg)^2|\nabla f^{n,t}_\delta|^2\,\dd m\bigg)^p\bigg]^\frac1p \\
%& \lesssim \mathbb{E} \bigg[\bigg(\int_{\mathcal{M}}(|\rho_\delta - \rho|^2 + |\rho_t - \rho|^2 +\frac1{\log n})|\nabla f^{n,t}_\delta|^2\,\dd m\bigg)^p\bigg]^\frac1p.
%\end{split}
%\]
%Thus, by H\"older inequality we get 
%\[
%\begin{split}
%\lefteqn{\mathbb{E} [W_2^{2p} (\mu^{n,t},\exp(\nabla f^{n,t}_{\delta})\#\rho_t\dd x)]^\frac1p}\\ 
%& \lesssim (\|\rho_\delta - \rho \|_{L^{2\big(\frac{\bar{q}}2\big)'}}^2 + \|\rho_t - \rho \|_{L^{2\big(\frac{\bar{q}}2\big)'}}^2 +\frac1{\log n})\mathbb{E}\bigg[\bigg(\int_{\mathcal{M}} |\nabla f^{n,t}_\delta|^{\bar{q}}\,\dd m\bigg)^{\frac2{\bar{q}}p} \bigg]^\frac1p\\
%& \stackrel{\eqref{Eq22}}{\lesssim} (\|\rho_\delta - \rho \|_{L^{2\big(\frac{\bar{q}}2\big)'}}^2 + \|\rho_t - \rho \|_{L^{2\big(\frac{\bar{q}}2\big)'}}^2 +\frac1{\log n}) \frac{\log n}n.
%\end{split}
%\]
%\fm{Case of correlated just square and (ii) of Proposition \ref{prop1}.}
%@@@@@@@@@@@@@@
%to obtain
%\[
%W_2^{2} (\mu^{n,t},\exp(\nabla f^{n,t}_{\delta})\#\rho_t\dd x) \le \frac1{\log n} \int_{\mathcal{M}} |\nabla f_\delta^{n,t}|^2 \,\dd m.
%\]
%\fm{continues}
%\nc{I let you continue !}
%\fm{There I see $t$ and not $\log t$ this should give $\log^\gamma n$ in the numerator.\\@@@@ADD ESTIMATE LAST TERM@@@@}

%\nc{I have the feeling that for the matching of two different set of points, that is the same proof. You define 
%
%$$-\nabla\cdot \rho\nabla f^{n,m,t}=\mu^{n,t}-\mu^{m,t},$$
%
%which is nothing but $f^{n,t}-f^{m,t}$. Then, the quantitative stability result will say
%
%$$\int_{\mathcal{M}}\vert T^{n,m}-(\text{Id}+\nabla f^{n,m,t})\vert^2\lesssim W^2_2(\mu^{m},(\text{Id}+\nabla f^{n,m,t})\#\mu^n)+W_2(\mu^n,\mu^m)W_2(\mu^{m},(\text{Id}+\nabla f^{n,m,t})\#\mu^n),$$
%
%then we argue but adding the quantities in $t$ and previous arguments. Do it by adding also the regularization $\delta$.} 
%\end{proof}
%

%
%\begin{proof}
%Let $u^{n,s} = \frac1n \sum_{k=1}^n p_s(X_i, \,\cdot\,)$ denote the density of $\mu^{n,t}$. By the the continuity of the heat flow with respect to the Wasserstein distance \cite[Proposition 5.1]{ambrosio2019finer} \fm{Maybe better recall the result in some book like EKS15 in the finer estimate paper} we may write
%\[
%\begin{split}
%W_2^2 (\mu_n, \mu^{n,t}) & \lesssim W_2^2 (\mu_n, \mu^{n,\frac1n}) + W_2^2 (\mu_{n,\frac1n}, \mu^{n,t}) \\
%& \lesssim \frac1n + W_2^2 (\mu_{n,\frac1n}, \mu^{n,t}).
%\end{split}
%\]
%Thus we are left to estimate the second term on the right hand side. Let $f$ be the unique null meaan solution to the Poisson equation $- \Delta f = u^{n, \frac1n} - u^{n,t}$, which can be represented by $\int_{\frac1n}^t u^{n,s} - 1 \,\dd s$. Since we are on the event $\mathcal{A}_n \cap \mathcal{B}_{\varepsilon, \delta, n}$ and $\rho \in L^{\infty}(\mathcal{M})$, by \cite[Proposition 4.3]{ambrosio2019finer} we may write
%\[
%W_2^2 (\mu^{n,\frac1n}, \mu^{n,t}) \lesssim \int_{\mathcal{M}} \frac{|\nabla f|^2}{u^{n,t}} \, \dd m \lesssim \int_{\mathcal{M}} | \nabla f|^2 \, \dd m.
%\]
%By the representation formula for $u^{n,t}$ we may write
%\begin{equation}\label{eq:loglognstart}
%\begin{split}
%W_2^2 (\mu^{n,\frac1n}, \mu^{n,t}) & \lesssim \int_{\mathcal{M}} | \nabla f|^2 \, \dd m = \int_{\mathcal{M}} - \Delta f \cdot f \, \dd m \\
%& = \int_{\mathcal{M}} (u^{n, \frac1n} - u^{n,t})\bigg(\int_{\frac1n}^t u^{n,s} - 1\, \dd s\bigg) \, \dd m.
%\end{split}
%\end{equation}

%\medskip
%{\sc Independent case}. Applying $\mathbb{E}[(\ \cdot\ )^p  ]^\frac1p$ to \eqref{eq:loglognstart} by independence we obtain \fm{I am writing it down for $2$}
%\[
%\begin{split}
%\mathbb{E}[W_2^{2} (\mu^{n,\frac1n}, \mu^{n,t})] & \lesssim \frac1n \int_{\frac1n}^t \mathbb{E} \bigg[ \int_{\mathcal{M}} (p_{\frac1n} (X,y) - p_t(X,y))p_s (X, y) \, \dd m (y)\bigg] \, \dd s \\
%& = \frac1n \int_{\frac1n}^t  \int_{\mathcal{M}} (p_{\frac1n + s}(x,x) - p_{t+s}(x,x)) \rho(x) \, \dd m(x) \, \dd s \\
%& \lesssim \frac1n \int_{2 \frac1n}^t \int_{\mathcal{M}} (p_s (x,x) - 1) \, \dd m(x) \, \dd s \lesssim \frac1n \int_{2 \frac1n}^t \frac1s \, \dd s \lesssim \frac{\log \alpha}n.
%\end{split}
%\]


%\medskip
%{\sc Weakly correlated case.} Taking the expectation in \eqref{eq:loglognstart} as before we may write
%\[
%\begin{split}
%\mathbb{E}[W_2^{2} (\mu^{n,\frac1n}, \mu^{n,t})] & \lesssim \frac1{n^2}\sum_{1\le i < j \le n} \int_{\frac1n}^t \mathbb{E} \bigg[ \int_{\mathcal{M}} (p_{\frac1n} (X_i,y) - p_t(X_i,y))p_s (X_j, y) \, \dd m (y)\bigg] \, \dd s \\
%& \lesssim \frac1{n^2}\sum_{1\le i < j \le n} \beta_{ij} \int_{\frac1n}^t  \int_{\mathcal{M}} (p_{\frac1n + s}(x,x) - p_{t+s}(x,x)) \rho(x) \, \dd m(x) \, \dd s \\
%& \stackrel{\eqref{Eq:AssumptionsCorrelatedCase}}{\lesssim} \frac1{n} \int_{\frac1n}^t  \int_{\mathcal{M}} (p_{\frac1n + s}(x,x) - p_{t+s}(x,x)) \rho(x) \, \dd m(x) \, \dd s \lesssim \frac{\log \alpha}n,
%\end{split}
%\]
%where the last inequality follows arguing as in the independent case. 
%\end{proof}
%
%
\end{proof}
%
\subsection{Contractivity estimates}
%
This section is devoted to the control of the smoothing errors $W^2_2(\mu^{n,t},\mu^{n})$ and $W^2_2(\nu^{m,t},\nu^{m})$ for the particular choice of $t$ given in Proposition \ref{Fluctuation}. The first result is in the spirit of \cite[Theorem 5.2]{ambrosio2019finer} that we extend in the case of non-uniformly distributed and correlated points. This extension requires a finer analysis of the error and the proof relies on Berry-Esseen type inequalities in the spirit of \cite[Theorem 5]{borda2021empirical}.
%
\begin{proposition}[Semigroup contraction for empirical measures]\label{Contractivity}
Let $\{\mu^n\}_n$ be defined in \eqref{eq:empmeas} with point clouds satisfying Assumption \ref{Assumptions}. Given $t$ such that Proposition \ref{Fluctuation} holds, we have
%
\begin{equation}\label{ContractivityEstiExpectation}
%
W_2^2(\mu^{n,t}, \mu^{n}) \le \mathcal{C}_n \frac{\log\log(n)}{n}+t\big\|\rho_{t+\frac{1}{n}}-\rho_{\frac{1}{n}}\big\|_{\LL^1},
%
\end{equation}
%
for some random variable $\mathcal{C}_n$ satisfying for $C<\infty$
%
$$\sup_{n \ge 1} \mathbb{E} [\tfrac1{C}\mathcal{C}_n]\leq 1.$$
%
Furthermore, if \eqref{DecayAlphaMixing} holds with $\eta\geq 1$ then the assumption \eqref{DecatBetaMixing} can be dropped and the stochastic integrability can be improved up to losing a $\log(n)$ factor, namely
%
\begin{equation}\label{ContractivityEstiExpMoment}
%
W_2^2(\mu^{n,t}, \mu^{n}) \le \mathcal{D}_n \frac{\log^{\frac{1}{\eta}}(n)\log\log(n)}{n}+t\big\|\rho_{t+\frac{1}{n}}-\rho_{\frac{1}{n}}\big\|_{\LL^1},
%
\end{equation}
%
for some random variable $\mathcal{D}_n$ satisfying for $D<\infty$
%
$$\sup_{n \ge 1} \mathbb{E} \big[\exp(\tfrac1{D}\mathcal{D}^{\frac{1}{2}}_n)\big]\leq 2.$$
%
\end{proposition}
 %
 \begin{proof}
 %
 According to the fluctuation estimates in Proposition \ref{Fluctuation} together with $W^2_2(\mu^{n,t},\mu^n)\leq (\text{diam}(\mathcal{M}))^2$, we can restrict the analysis in $\mathcal{A}_n$ defined in \eqref{Events}. Note that for $n$ large enough, \eqref{Ellipticity} yields
%
\begin{equation}\label{EllipticityAn}
%
\frac{\lambda}{2}\leq \mu^{n,t}\leq \Lambda+1\quad\text{in $\mathcal{A}_n$.}
%
\end{equation}
%

%
\medskip
%

We split the proof into three steps. In the first step, we prove a Berry–Esseen type smoothing inequality for $W^2_2(\mu^{n,t},\mu^n)$ which decomposes the error in a deterministic part involving $\rho$ and a random part involving the Fourier coefficients $\{\widehat{\mu^n}(k)\}_k$ of $\mu^n$. In the second step, we prove \eqref{ContractivityEstiExpectation}. In the third step, we control the fluctuations of $\{\widehat{\mu^n}(k)\}_k$ using the concentration inequalities in Proposition \ref{BersteinCorrelated} and deduce \eqref{ContractivityEstiExpMoment}.

%
\medskip
%

{\sc Step 1. Berry-Esseen type inequality. }Recalling that we denote by $\{\lambda_k,\phi_k\}_k$ the eigenvalues and eigenfunctions of $-\Delta$ respectively, we prove that 
%
\begin{equation}\label{SmoothingIneg}
%
W^2_2(\mu^{n,t},\mu^n)\lesssim \frac{1}{n}+\sum_{k\geq 1} \frac{e^{-\frac{2}{n}\lambda_k}}{\lambda_k}\big(e^{-t\lambda_k}-1\big)^2\vert\widehat{\mu^n}(k)-\widehat{\rho}(k)\vert^2+t\|\rho_{t+\frac{1}{n}}-\rho_{\frac{1}{n}}\|_{\LL^1},
%
\end{equation}
%
where 
%
$$\widehat{\mu^n}(k):=\int_{\mathcal{M}}\phi_k\,\dd\mu^n \quad \text{and}\quad \widehat{\rho}(k):=\int_{\mathcal{M}}\rho\,\phi_k\,\dd\m.$$
%
We first apply the triangle inequality and use the classical contractivity estimate in \cite[Theorem 3]{erbar2015equivalence} to get
%
\begin{equation}\label{EstiPeyre:Eq4}
%
W^2_2(\mu^{n,t},\mu^n)\lesssim W^2_2(\mu^{n,\frac{1}{n}},\mu^{n})+W^2_2(\mu^{n,t+\frac{1}{n}},\mu^{n,t})+W^2_2(\mu^{n,t+\frac{1}{n}},\mu^{n,\frac{1}{n}})\lesssim \frac{1}{n}+W^2_2(\mu^{n,t+\frac{1}{n}},\mu^{n,\frac{1}{n}}).
%
\end{equation}
%
We then apply Peyre's estimate \cite{peyre2018comparison} to the second \rhs, which takes the form
%
\begin{equation}\label{PeyreLemma}
%
W^2_2(\mu^{n,t+\frac{1}{n}},\mu^{n,\frac{1}{n}})\leq 4\sup\bigg\{\Big\vert\int_{\mathcal{M}}(\mu^{n,t+\frac{1}{n}}-\mu^{n,\frac{1}{n}})\,f\,\dd\m\Big\vert^2\quad\text{with}\quad \int_{\mathcal{M}}\mu^{n,t+\frac{1}{n}}\,\vert\nabla f\vert^2\,\dd\m\leq 1\bigg\}.
%
\end{equation}
%
Now, given an arbitrary $f$ such that 
%
\begin{equation}\label{Boundtestfunctionf}
%
\int_{\mathcal{M}}\mu^{n,t+\frac{1}{n}}\,\vert\nabla f\vert^2\,\dd\m\leq 1,
%
\end{equation}
%
we split
%
\begin{equation}\label{EstiPeyre:Eq1}
%
\int_{\mathcal{M}}(\mu^{n,t+\frac{1}{n}}-\mu^{n,\frac{1}{n}})\, f\,\dd\m=\int_{\mathcal{M}}\big(\mu^{n,t+\frac{1}{n}}-\mu^{n,\frac{1}{n}}-(\rho_{t+\frac{1}{n}}-\rho_{\frac{1}{n}})\big)\, f\,\dd\m+\int_{\mathcal{M}}f(\rho_{t+\frac{1}{n}}-\rho_{\frac{1}{n}})\,\dd\m.
%
\end{equation}
%
For the first \rhs term of \eqref{EstiPeyre:Eq1}, we expand the integral using \eqref{SpectralDecompoHeatKernel}. Thus, together with the semigroup property of $\{\mathrm{P}_t\}_{t>0}$ and Cauchy-Schwarz's inequality, we get
%
\begin{align*}
%
&\int_{\mathcal{M}}\big(\mu^{n,t+\frac{1}{n}}-\mu^{n,\frac{1}{n}}-(\rho_{t+\frac{1}{n}}-\rho_{\frac{1}{n}})\big)\, f\,\dd\m\\
&=\int_{\mathcal{M}}\dd\big(\mu^{n,t}-\mu^{n}-(\rho_{t}-\rho)\big)(y)\int_{\mathcal{M}}\dd\m(x)\,f(x)p_{\frac{1}{n}}(x,y)\\
&\stackrel{\eqref{SpectralDecompoHeatKernel}}{=}\sum_{k\geq 1}e^{-\frac{1}{n}\lambda_k}\int_{\mathcal{M}}\dd\big(\mu^{n,t}-\mu^{n}-(\rho_{t}-\rho)\big)(y)\int_{\mathcal{M}}\dd\m(x)\,f(x)\phi_k(x)\phi_k(y)\\
&=\sum_{k\geq 1}e^{-\frac{1}{n}\lambda_k}\widehat{f}(k)\big(\widehat{\mu^{n,t}}(k)-\widehat{\mu^{n}}(k)-(\widehat{\rho_t}(k)-\widehat{\rho}(k))\big)\\
&\leq \Big(\sum_{k\geq 1}\lambda_k\vert\widehat{f}(k)\vert^2\Big)^{\frac{1}{2}}\Big(\sum_{k\geq 1}\frac{e^{-\frac{2}{n}\lambda_k}}{\lambda_k}\vert \widehat{\mu^{n,t}}(k)-\widehat{\mu^{n}}(k)-(\widehat{\rho_t}(k)-\widehat{\rho}(k))\vert^2\Big)^{\frac{1}{2}}
%
\end{align*}
%
Using that from \eqref{EllipticityAn} we have $\mu^{n,t+\frac{1}{n}}\geq \frac{\lambda}{2}$ and recalling \eqref{Boundtestfunctionf}, we get
%
$$ \Big(\sum_{k\geq 1}\lambda_k\vert\widehat{f}(k)\vert^2\Big)^{\frac{1}{2}}\leq \Big(\int_{\mathcal{M}}\vert \nabla f\vert^2\,\dd\m\Big)^{\frac{1}{2}}\leq \frac{2}{\lambda}\Big(\int_{\mathcal{M}}\mu^{n,t+\frac{1}{n}}\, \vert \nabla f\vert^2\,\dd\m\Big)^{\frac{1}{2}}\leq \frac{2}{\lambda}.$$
%
Furthermore, noticing that $\pp_t\phi_k=e^{-t\lambda_k}\phi_k$ which implies that, since $\pp_t$ is self-adjoint
%
$$\widehat{\mu^{n,t}}(k)=e^{-t\lambda_k}\widehat{\mu^n}(k)\quad\text{and}\quad \widehat{\rho_t}(k)=e^{-t\lambda_k}\widehat{\rho}(k),$$
%
we obtain
%
$$\sum_{k\geq 1}\frac{e^{-\frac{2}{n}\lambda_k}}{\lambda_k}\vert \widehat{\mu^{n,t}}(k)-\widehat{\mu^{n}}(k)-(\widehat{\rho_t}(k)-\widehat{\rho}(k))\vert^2=\sum_{k\geq 1}\frac{e^{-\frac{2}{n}\lambda_k}}{\lambda_k}\big(e^{-t\lambda_k}-1\big)^2\vert\widehat{\mu^n}(k)-\widehat{\rho}(k)\vert^2.$$
%
This leads to 
%
\begin{equation}\label{EstiPeyre:Eq2}
%
\int_{\mathcal{M}}\dd\big(\mu^{n,t+\frac{1}{n}}-\mu^{n,\frac{1}{n}}-(\rho_{t+\frac{1}{n}}-\rho_{\frac{1}{n}})\big)\, f\lesssim \Big(\sum_{k\geq 1}\frac{e^{-\frac{2}{n}\lambda_k}}{\lambda_k}\big(e^{-t\lambda_k}-1\big)^2\vert\widehat{\mu^n}(k)-\widehat{\rho}(k)\vert^2\Big)^{\frac{1}{2}}.
%
\end{equation}
%
For the second \rhs of \eqref{EstiPeyre:Eq1}, we introduce $w\in \dot{{\rm H}}^1$ satisfying 
%
$$-\Delta w=\rho_{t+\frac{1}{n}}-\rho_{\frac{1}{n}},$$
%
so that from an integration by parts, Cauchy-Schwarz' inequality and the combination of \eqref{EllipticityAn} and \eqref{Boundtestfunctionf}, we obtain
%
$$\int_{\mathcal{M}}f(\rho_{t+\frac{1}{n}}-\rho_{\frac{1}{n}})\,\dd\m=\int_{\mathcal{M}}\nabla f\cdot\nabla w\,\dd\m\lesssim \Big(\int_{\mathcal{M}}\vert\nabla w\vert^2\,\dd\m\Big)^{\frac{1}{2}}.$$
%
Using then the explicit formula $w=-\int_{\frac{1}{n}}^{t+\frac{1}{n}}\rho_{\tau}\,\dd\tau$ together with $\vert \int_{\frac{1}{n}}^{t+\frac{1}{n}}\rho_{\tau}\,\dd\tau\vert\stackrel{\eqref{Ellipticity}}{\leq} \Lambda t$, we get 
%
$$\int_{\mathcal{M}}\vert\nabla w\vert^2\,\dd\m=\int (\rho_{t+\frac{1}{n}}-\rho_{\frac{1}{n}})w\,\dd\m\lesssim t\|\rho_{t+\frac{1}{n}}-\rho_{\frac{1}{n}}\|_{\LL^1},$$
%
so that 
%
\begin{equation}\label{EstiPeyre:Eq3}
%
\Big\vert\int_{\mathcal{M}}f(\rho_{t+\frac{1}{n}}-\rho_{\frac{1}{n}})\,\dd\m\Big\vert\lesssim \sqrt{t}\|\rho_{t+\frac{1}{n}}-\rho_{\frac{1}{n}}\|^{\frac{1}{2}}_{\LL^1}.
%
\end{equation}
%
The combination of \eqref{EstiPeyre:Eq1}, \eqref{EstiPeyre:Eq2}, \eqref{EstiPeyre:Eq3} and \eqref{EstiPeyre:Eq4} leads to \eqref{SmoothingIneg}.

%
\medskip
%

{\sc Step 2. Proof of \eqref{ContractivityEstiExpectation}. }According to \eqref{SmoothingIneg}, it remains to show that 
%
\begin{equation}\label{ContractExpectation}
%
\mathbb{E}\bigg[\sum_{k\geq 1} \frac{e^{-\frac{2}{n}\lambda_k}}{\lambda_k}\big(e^{-t\lambda_k}-1\big)^2\vert\widehat{\mu^n}(k)-\widehat{\rho}(k)\vert^2\bigg]\lesssim \frac{\log\log(n)}{n}.
%
\end{equation}
%
Writing 
%
\begin{equation}\label{DecomposeFourier}
%
\widehat{\mu^n}(k)-\widehat{\rho}(k)=\frac{1}{n}\sum_{\ell=1}^n (\phi_k(X_\ell)-\mathbb{E}[\phi_k(X_\ell)]),
%
\end{equation}
%
we first expand the square in form of
%
\begin{equation}\label{SquareExpended}
%
\begin{aligned}
%
&\sum_{k\geq 1} \frac{e^{-\frac{2}{n}\lambda_k}}{\lambda_k}\big(e^{-t\lambda_k}-1\big)^2\vert\widehat{\mu^n}(k)-\widehat{\rho}(k)\vert^2\\
=&\frac{1}{n^2}\sum_{\ell=1}^n\sum_{k\geq 1} \frac{e^{-\frac{2}{n}\lambda_k}}{\lambda_k}\big(e^{-t\lambda_k}-1\big)^2\big(\phi_k(X_{\ell})-\mathbb{E}[\phi_k(X_\ell)]\big)^2\\
&+\frac{2}{n^2}\sum_{1\leq \ell<\ell'\leq n}\sum_{k\geq 1} \frac{e^{-\frac{2}{n}\lambda_k}}{\lambda_k}\big(e^{-t\lambda_k}-1\big)^2\big(\phi_k(X_{\ell})-\mathbb{E}[\phi_k(X_\ell)]\big)\big(\phi_k(X_{\ell'})-\mathbb{E}[\phi_k(X_{\ell'})]\big).
%
\end{aligned}
%
\end{equation}

%
For the first \rhs term of \eqref{SquareExpended}, we use the normalisation $\|\phi_k\|_{\LL^2}=1$ together with \eqref{Ellipticity} to the effect of 
%
\begin{equation}\label{BoundsMomentVectors}
%
\mathbb{E}\big[\vert \phi_k(X_1)-\mathbb{E}[\phi_k(X_1)]\vert^2\big]\leq \Lambda,
%
\end{equation}
%
and get
%
\begin{equation}\label{ExpendedSqaure:Est1}
%
\mathbb{E}\bigg[\sum_{\ell=1}^n\sum_{k\geq 1} \frac{e^{-\frac{2}{n}\lambda_k}}{\lambda_k}\big(e^{-t\lambda_k}-1\big)^2(\phi_k(X_{\ell})-\mathbb{E}[\phi_k(X_\ell)])^2\bigg]\leq \Lambda n\sum_{k\geq 1} \frac{e^{-\frac{2}{n}\lambda_k}}{\lambda_k}\big(e^{-t\lambda_k}-1\big)^2.
%
\end{equation}
%
For the second \rhs term of \eqref{SquareExpended}, we use the definition of the $\beta$-mixing coefficient \eqref{BetaMixing} together with the assumption \eqref{DecatBetaMixing} in form of
%
\begin{equation}\label{ExpendedSquare:Est2}
%
\begin{aligned}
%
&\mathbb{E}\bigg[\sum_{1\leq \ell<\ell'\leq n}\sum_{k\geq 1} \frac{e^{-\frac{2}{n}\lambda_k}}{\lambda_k}\big(e^{-t\lambda_k}-1\big)^2\big(\phi_k(X_{\ell})-\mathbb{E}[\phi_k(X_\ell)]\big)\big(\phi_k(X_{\ell'})-\mathbb{E}[\phi_k(X_{\ell'})]\big)\bigg]\\
&\leq \sum_{1\leq \ell<\ell'\leq n}\beta_{\ell\ell'}\sup_{x,y\in\mathcal{M}}\bigg\vert\sum_{k\geq 1} \frac{e^{-\frac{2}{n}\lambda_k}}{\lambda_k}\big(e^{-t\lambda_k}-1\big)^2\phi_k(x)\phi_k(y)\bigg\vert\\
&\stackrel{\eqref{DecatBetaMixing}}{\lesssim}n\sup_{x,y\in\mathcal{M}}\bigg\vert\sum_{k\geq 1} \frac{e^{-\frac{2}{n}\lambda_k}}{\lambda_k}\big(e^{-t\lambda_k}-1\big)^2\phi_k(x)\phi_k(y)\bigg\vert.
%
\end{aligned}
%
\end{equation}
%
The combination of \eqref{SquareExpended}, \eqref{ExpendedSqaure:Est1} and \eqref{ExpendedSquare:Est2} yields
%
\begin{align*}
%
&\mathbb{E}\bigg[\sum_{k\geq 1} \frac{e^{-\frac{2}{n}\lambda_k}}{\lambda_k}\big(e^{-t\lambda_k}-1\big)^2\vert\widehat{\mu^n}(k)-\widehat{\rho}(k)\vert^2\bigg]\\
&\lesssim \frac{1}{n}\bigg(\sum_{k\geq 1} \frac{e^{-\frac{2}{n}\lambda_k}}{\lambda_k}\big(e^{-t\lambda_k}-1\big)^2+\sup_{x,y\in\mathcal{M}}\bigg\vert\sum_{k\geq 1} \frac{e^{-\frac{2}{n}\lambda_k}}{\lambda_k}\big(e^{-t\lambda_k}-1\big)^2\phi_k(x)\phi_k(y)\bigg\vert\bigg).
%
\end{align*}
%
It remains to show that 
%
\begin{equation}\label{ControlExpectationLast}
%
\sum_{k\geq 1} \frac{e^{-\frac{2}{n}\lambda_k}}{\lambda_k}\big(e^{-t\lambda_k}-1\big)^2+\sup_{x,y\in\mathcal{M}}\bigg\vert\sum_{k\geq 1} \frac{e^{-\frac{2}{n}\lambda_k}}{\lambda_k}\big(e^{-t\lambda_k}-1\big)^2\phi_k(x)\phi_k(y)\bigg\vert\lesssim \log\log(n).
%
\end{equation}
%
We only treat the second \lhs term of \eqref{ControlExpectationLast}, the first term is controlled the same way. For any $x,y\in\mathcal{M}$, we expand
%
\begin{align*}
%
\lefteqn{\sum_{k\geq 1}\frac{e^{-\frac{2}{n}\lambda_k}}{\lambda_k}\big(e^{-t\lambda_k}-1\big)^2\phi_k(x)\phi_k(y)}\\ &\quad=\sum_{k\geq 1}\frac{1}{\lambda_k}e^{-2(\frac{1}{n}+t)\lambda_k}\phi_k(x)\phi_k(y)-2\sum_{k\geq 1}\frac{1}{\lambda_k}e^{-(\frac{2}{n}+t)\lambda_k}\phi_k(x)\phi_k(y)+\sum_{k\geq 1}\frac{1}{\lambda_k}e^{-\frac{2}{n}\lambda_k}\phi_k(x)\phi_k(y).
%
\end{align*}
%
We then disintegrate using the spectral decomposition of the heat kernel \eqref{SpectralDecompoHeatKernel} in form of  
%
$$\sum_{k\geq 1}\frac{1}{\lambda_k}e^{-s\lambda_k}\phi_k(x)\phi_k(y)=\int_{s}^1\dd s\, p_s(x,y)-\sum_{k\geq1}\frac{1}{\lambda_k}e^{-\lambda_k}\phi_k(x)\phi_k(y)\quad\text{for any $s>0$, }$$
%
so that we obtain
%
\begin{align*}
%
\sup_{x,y\in\mathcal{M}}\bigg\vert\sum_{k\geq 1}\frac{e^{-\frac{2}{n}\lambda_k}}{\lambda_k}\big(e^{-t\lambda_k}-1\big)^2\phi_k(x)\phi_k(y)\bigg\vert&=\sup_{x,y\in\mathcal{M}}\bigg\vert\int_{2(\frac{1}{n}+t)}^{\frac{2}{n}+t}p_s(x,y)\,\dd s+\int_{\frac{2}{n}}^{\frac{2}{n}+t}p_s(x,y)\,\dd s\bigg\vert\\
&\lesssim \int_{2(\frac{1}{n}+t)}^{\frac{2}{n}+t}s^{-1}\,\dd s+\int_{\frac{2}{n}}^{\frac{2}{n}+t}s^{-1}\,\dd s\\
&\lesssim \log\log(n).
%
\end{align*}

%
\medskip
%

{\sc Step 3. Proof of \eqref{ContractivityEstiExpMoment}. }It is a consequence of the following fluctuation estimates
%
\begin{equation}\label{ContractivityCOM:Eq1}
%
\mathbb{E}\big[\vert\widehat{\mu^n}(k)-\widehat{\rho}(k)\vert^{2p}\big]^{\frac{1}{p}}\lesssim p^2\,\frac{1}{n}\bigg(\log^{\frac{1}{\eta}}(n)\lambda^{\frac{1}{\log(n)}}_k+\frac{\lambda_k(1+\log^2(n))}{n}\bigg)\quad\text{for any $p<\infty$},
%
\end{equation}
%
together with Lemma \ref{momentexp}. Indeed, applying Minkowski's inequality followed by \eqref{ContractivityCOM:Eq1} and $\lambda^{\frac{1}{\log(n)}}_k e^{-\frac{2}{n}\lambda_k}\lesssim e^{-\frac{1}{n}\lambda_k}$ yield
%
\begin{align*}
%
&\mathbb{E}\bigg[\Big(\sum_{k\geq 1}\frac{e^{-\frac{2}{n}\lambda_k}}{\lambda_k}\big(e^{-t\lambda_k}-1\big)^2\vert\widehat{\mu^n}(k)-\widehat{\rho}(k)\vert^2\Big)^p\bigg]^{\frac{1}{p}}\\
&\lesssim p^2\frac{1}{n}\sum_{k\geq 1}\frac{e^{-\frac{1}{n}\lambda_k}}{\lambda_k}\big(e^{-t\lambda_k}-1\big)^2\bigg(\log^{\frac{1}{\eta}}(n)+\frac{\lambda_k(1+\log^2(n))}{n}\bigg),
%
\end{align*}
%
and \eqref{ContractivityEstiExpMoment} follows from
%
\begin{equation}\label{ContractivityConclusion:Eq2}
%
\sum_{k\geq 1}\frac{e^{-\frac{1}{n}\lambda_k}}{\lambda_k}\big(e^{-t\lambda_k}-1\big)^2\lesssim \log\log(n)\quad\text{and}\quad\sum_{k\geq 1}e^{-\frac{1}{n}\lambda_k}\big(e^{-t\lambda_k}-1\big)^2\lesssim t^{-1},
%
\end{equation}
%
which is obtained the same way as in \eqref{ControlExpectationLast} using additionally the trace formula \eqref{eq:traceformula}.


%
\medskip
%

We now prove \eqref{ContractivityCOM:Eq1}. It follows from the estimate on the probability tails
%
\begin{equation}\label{ContractivityCOM:Eq2}
%
\mathbb{P}\big(\vert\widehat{\mu^n}(k)-\widehat{\rho}(k)\vert>\lambda\big)\lesssim \exp\bigg(-\frac{1}{C}\frac{n^2\lambda^2}{n\log^{\frac{1}{\eta}}(n)\lambda^{\frac{1}{\log(n)}}_k+\lambda_k+n\lambda\lambda^{\frac{1}{2}}_k\log^2(n)}\bigg)\quad\text{for any $\lambda>0$},
%
\end{equation}
%
for some $C>0$, together with a simple application of the layer-cake representation. 

%
\medskip
%

To see \eqref{ContractivityCOM:Eq2}, we use \eqref{DecomposeFourier} together with Proposition \ref{BersteinCorrelated} to obtain
%
$$\mathbb{P}\big(\vert\widehat{\mu^n}(k)-\widehat{\rho}(k)\vert>\lambda\big)\lesssim \exp\Big(-\frac{n^2\lambda^2}{nv^2+\|\phi_k\|^2_{\LL^{\infty}}+n\lambda\|\phi_k\|_{\LL^{\infty}}\log^2(n)}\Big),$$
%
with 
%
\begin{equation}\label{vsquaredContra}
%
v^2:=\mathbb{E}[\vert\phi_k(X_1)-\mathbb{E}[\phi_k(X_1)]\vert^2]+2\sum_{i<j}\big\vert\mathbb{E}\big[(\phi_k(X_i)-\mathbb{E}[\phi_k(X_i)])(\phi_k(X_j)-\mathbb{E}[\phi_k(X_j)])\big]\big\vert.
%
\end{equation}
%
The estimate \eqref{ContractivityCOM:Eq2} is then a consequence of 
%
\begin{equation}\label{ContractivityCOM:Eq4}
%
 \|\phi_k\|_{\LL^{\infty}}\lesssim \lambda^{\frac{1}{2}}_k\quad\text{and}\quad v^2\lesssim \log^{\frac{1}{\eta}}(n)\lambda^{\frac{1}{\log(n)}}_k.
%
\end{equation}
%
The first item of \eqref{ContractivityCOM:Eq4} has been treated in \eqref{eq:eigvaleighfunc}. 
%is a direct consequence of Gagliardo-Nirenberg' interpolation inequality \cite[Theorem 3.70]{aubin1998some} which reads
%
%$$\|\phi_k\|_{\LL^{\infty}}\lesssim \|\nabla^2 \phi_k\|^{\frac{1}{2}}_{\LL^{2}}\|\phi_k\|^{\frac{1}{2}}_{\LL^2}=\|\nabla^2\phi_k\|^{\frac{1}{2}}_{\LL^2},$$
%
%combined with $-\Delta\phi_k=\lambda_k\phi_k$ and elliptic regularity
%
%$$\|\nabla^2\phi_k\|_{\LL^2}\lesssim \lambda_k\|\phi_k\|_{\LL^2}=\lambda_k.$$
%
For the second item of \eqref{ContractivityCOM:Eq4}, we use \eqref{BoundsMomentVectors} and, combined with \eqref{DecayAlphaMixing}, we obtain
%
\begin{align*}
%
\sum_{i<j}\big\vert\mathbb{E}\big[(\phi_k(X_i)-\mathbb{E}[\phi_k(X_i)])(\phi_k(X_j)-\mathbb{E}[\phi_k(X_j)])\big]\big\vert\lesssim& \|\phi_k\|^{\frac{1}{\log(n)}}_{\LL^\infty}\sum_{\ell\geq 0}\exp\big(-b\tfrac{\ell^{\eta}}{\log(n)}\big)\\
\lesssim& \log^{\frac{1}{\eta}}(n)\lambda^{\frac{1}{\log(n)}}_k.
%
\end{align*}
%
%, we get from Bernstein's inequality in Proposition \ref{Berstein}
%
%$$\mathbb{P}\big(\vert\widehat{\mu^n}(k)-\widehat{\rho}(k)\vert>\lambda\big)\leq 2\exp\Big(-\frac{n\lambda^2}{2\Lambda+\frac{2}{3}\lambda\|\phi_k\|_{\LL^{\infty}}}\Big).$$
%
%The estimate \eqref{ContractivityCOM:Eq2} then follows from 
%
%\begin{equation}\label{ContractivityCOM:Eq3}
%
%\|\phi_k\|_{\LL^\infty}\lesssim \lambda^\frac{1}{2}_k.
%
%\end{equation}
% 


%{\sc Step 3. Proof of \eqref{Contractivity}. }According to \eqref{SmoothingIneg}, it remains to show that 
%
%\begin{equation}\label{ContractivityConclusion:Eq1}
%
%\mathbb{E}\bigg[\Big(\sum_{k\geq 1} \frac{e^{-\frac{2}{n}\lambda_k}}{\lambda_k}\big(e^{-t\lambda_k}-1\big)^2\vert\widehat{\mu^n}(k)-\widehat{\rho}(k)\vert^2\Big)^p\bigg]^{\frac{1}{p}}\lesssim p^2 \frac{\log\log(n)}{n}.
%
%\end{equation}
%
%Applying \eqref{ContractivityCOM:Eq1} together with Minkowski's inequality yields
%
%$$\mathbb{E}\bigg[\Big(\sum_{k\geq 1}\frac{e^{-\frac{2}{n}\lambda_k}}{\lambda_k}\big(e^{-t\lambda_k}-1\big)^2\vert\widehat{\mu^n}(k)-\widehat{\rho}(k)\vert^2\Big)^p\bigg]^{\frac{1}{p}}\leq p^2\sum_{k\geq 1}\frac{e^{-\frac{2}{n}\lambda_k}}{\lambda_k}\big(e^{-t\lambda_k}-1\big)^2\frac{1}{n}\Big(1+\frac{\lambda_k}{n}\Big).$$
%
%It is enough to prove that
%
%\begin{equation}\label{ContractivityConclusion:Eq2}
%
%\sum_{k\geq 1}\frac{e^{-\frac{2}{n}\lambda_k}}{\lambda_k}\big(e^{-t\lambda_k}-1\big)^2\lesssim \log\log(n)\quad\text{and}\quad\sum_{k\geq 1}e^{-\frac{2}{n}\lambda_k}\big(e^{-t\lambda_k}-1\big)^2\lesssim t^{-1}.
%
%\end{equation}
%
%We only prove the first item of \eqref{ContractivityCOM:Eq2}, the second is treated the same way. We first expand 
%
%$$\sum_{k\geq 1}\frac{e^{-\frac{2}{n}\lambda_k}}{\lambda_k}\big(e^{-t\lambda_k}-1\big)^2=\sum_{k\geq 1}\frac{1}{\lambda_k}e^{-2(\frac{1}{n}+t)\lambda_k}-2\sum_{k\geq 1}\frac{1}{\lambda_k}e^{-(\frac{2}{n}+t)\lambda_k}+\sum_{k\geq 1}\frac{1}{\lambda_k}e^{-\frac{2}{n}\lambda_k}.$$
%
%We then disintegrate in form of  
%
%$$\sum_{k\geq 1}\frac{1}{\lambda_k}e^{-s\lambda_k}=\int_{s}^1\dd s\, \sum_{k\geq 1} e^{-s\lambda_k}-\sum_{k\geq1}\frac{1}{\lambda_k}e^{-\lambda_k}\quad\text{for any $s>0$, }$$
%
%so that, combined with \cite[Chapter VI]{chavel1984eigenvalues}, we obtain
%
%\begin{align*}
%
%\sum_{k\geq 1}\frac{e^{-\frac{2}{n}\lambda_k}}{\lambda_k}\big(e^{-t\lambda_k}-1\big)^2 &=\int_{2(\frac{1}{n}+t)}^{\frac{2}{n}+t}\dd s\int_{\mathcal{M}}\dd\m(x)\, p_s(x,x)+\int_{\frac{2}{n}}^{\frac{2}{n}+t}\dd s\int_{\mathcal{M}}\dd\m(x)\, p_s(x,x)\\
%&\lesssim \int_{2(\frac{1}{n}+t)}^{\frac{2}{n}+t}\dd s\,s^{-1}+\int_{\frac{2}{n}}^{\frac{2}{n}+t}\dd s\,s^{-1}\\
%&\lesssim \log\log(n).
%
%\end{align*}
%
%\subsection{Proof of Lemma \ref{lem:transportcontraction}: Lipschitz property of the Wasserstein metric}
%
%We prove \eqref{Eq:LipEstimateWasser} using the duality formula \cite[Theorem 5.10]{OldNewVillani}. We first claim that for any bounded and continuous function $\phi$ on $D$
%
%\begin{equation}\label{eq:HamJacLip}
%Q_t \phi (T (y)) \le L^2 Q_t  \big(\tfrac{1}{K^2} \phi\circ T\big) (y)\quad \text{for any $t\geq 0$ and $y\in \mathcal{M}$},
%\end{equation}
%
%where $Q_t$ denotes the Hopf-Lax semigroup
%
%$$Q_t \phi:=\inf_{x\in D}\{\phi(x)+\tfrac{1}{2t}\dd^2(x,\cdot)\}.$$
%
%Indeed, by the definition and the Lipschitz property of $T$, we have for any $z\in \mathcal{M}$
%\[
%\begin{split}
%Q_t \phi (T (y))& \le \phi (T (z)) + \tfrac1{2t} \dd^2(T (z),T (y))\le K^2 ( \tfrac1{K^2} \phi(T(z)) + \tfrac1{2t} \dd^2(z,y)),
%\end{split}
%\]
%which yields \eqref{eq:HamJacLip} by arbitrariness of $z$. We then deduce from \eqref{eq:HamJacLip} that for any Lipschitz functions $\phi$ on $D$
%
%\begin{align*}
%
%\int_{\mathcal{M}}Q_1\phi\,\dd T_{\#}\nu-\int_{\mathcal{M}}\phi\, \dd T_{\#}\mu&=\int_{\mathcal{M}}Q_1\phi\circ T\,\dd\nu-\int_{\mathcal{M}}\phi\circ T\,\dd\mu\\
%&\stackrel{\eqref{eq:HamJacLip}}{\leq}L^2\Big(\int_{\mathcal{M}}Q_1(\tfrac{1}{L^2}\phi\circ T)\,\dd \nu-\int_{\mathcal{M}}\tfrac{1}{L^2}\phi\circ T\,\dd \mu\Big)\\
%&\leq \frac{L^2}{2}W^2_2(\mu,\nu),
%
%\end{align*}
%
%which gives \eqref{Eq:LipEstimateWasser} using the duality formula \cite[Theorem 5.10]{OldNewVillani}.
%
%\subsection{Proof of Lemma \ref{lem:HeatContractionrho}: Semigroup contraction for absolutely continuous measures}
%
%Using $g$ defined via $-\Delta g=\rho_t-\rho$ together with \eqref{Ellipticity}, \cite[Corollary 4.4 (4.1)]{ambrosio2019finer} yields
%
%$$W^2_2(\rho_t,\rho)\lesssim \int_{\mathcal{M}}\vert\nabla g\vert^2=\int_{\mathcal{M}}(\rho_t-\rho)g.$$
%
%Writing $g=\int_{0}^{\infty}\dd \tau\,\text{P}_\tau(\rho_t-\rho)=-\int_{0}^t \dd\tau\,\rho_\tau$ together with $\vert\int_{0}^t \dd\tau\,\rho_\tau\vert\leq \|\rho\|_{\LL^{\infty}(\mathcal{M})} t$ gives
%
%$$W^2_2(\rho_t,\rho)\lesssim t \| \rho_t - \rho\|_{\LL^1(\mathcal{M})}.$$
 %
 \end{proof}
%Next, we show a Lipschitz property of the $W_2$-Wasserstein distance.
 %
%\begin{lemma}[Lipschitz property of the Wasserstein metric] \label{lem:transportcontraction}
%Let $(D,\dd)$ be a complete and separable metric space. Consider $\mu, \nu$ be two probability measures on $\mathcal{M}$ and $T: \mathcal{M} \rightarrow D$ be a $L$-Lipschitz map. It holds
%
%\begin{equation}\label{Eq:LipEstimateWasser}
%W_2^2 (T_{\#}\mu,T_{\#} \nu) \le L^2 W_2^2 (\mu, \nu).
%\end{equation}
%
%\end{lemma}
%
%Finally, we show a contractivity estimate between two absolutely continuous measures with respect to bounded densities, needed for the proof of Theorem \ref{th1}.
%
%\begin{lemma}[Semigroup contraction for absolute continuous measures]\label{lem:HeatContractionrho}
%\eqref{Ellipticity}. It holds
%
%\begin{equation}\label{Eq80}
%W_2^2 (\rho_t, \rho) \lesssim t \| \rho_t - \rho\|_{\LL^1},
%\end{equation}
%
%where we recall that $\rho_t=\mathrm{P}_t\rho$.
%\end{lemma}
%



\subsection{Proof of Theorem \ref{MainResultMatchingClouds}: Approximation of the transport plan}\label{ProofTransport}
%
We only give the arguments for \eqref{ApproximationTransportPlan}, \eqref{ApproximationTransportPlanBis} is proved the same way using the corresponding results \eqref{Eq22}, \eqref{ContractivityEstiExpMoment} and \eqref{eq:matchingcostHighMoments} in the case $\eta>2$ (where some additional comments are given if necessary along the proof). We split the proof into four steps. In the first step, we display some preliminary estimates useful all along the proof. In the second step, we deal with the approximation error that occurs in the process of regularizing $\rho$ into $\rho_\delta$. In the third step, we estimate the $W_2$-distance for the regularized quantity using the quantitative stability result in \cite[Theorem 3.2]{ambrosio2019optimal}, splitting the estimates in small pieces that we control in the fourth step. We finally comment on the proof of Remark \ref{FlatGeometry} and Theorem \ref{th1}, which are obtained with similar techniques.

%
\medskip
%

{\sc Step 1. Preliminary estimates. }

%
\medskip
%

\textbf{Heat kernel regularization. }The assumption $\rho\in \mathrm{H}^\varepsilon$ provides 
%
\begin{equation}\label{ApproximationRho}
%
\|\rho_s-\rho\|_{\LL^2}\lesssim \|\rho\|_{\mathrm{H}^\varepsilon} \,s^{\varepsilon}\quad\text{for any $s>0$}.
%
\end{equation}
%
%
Indeed, using the definition of the heat-kernel, Minkowski's inequality and the spectral decomposition \eqref{SpectralDecompoHeatKernel}, we have
%
\begin{align*}
%
\|\rho_{s}-\rho\|^2_{\LL^2}&=\int_{\mathcal{M}}\dd \m(x)\bigg\vert\int_{\mathcal{M}}\rho(y)\dd\m(y)\int_{0}^{s}\dd\tau\,\partial_{\tau}p_{\tau}(x,y)\bigg\vert^2\\
&\leq\bigg(\int_{0}^{s}\dd\tau\,\bigg(\int_{\mathcal{M}}\dd\m(x)\bigg\vert\sum_{k\geq 1}\widehat{\rho}(k)\, \lambda_k e^{-\tau\lambda_k}\phi_k(x)\bigg\vert^2\bigg)^{\frac{1}{2}}\bigg)^2.
%
\end{align*}
%
Noticing that 
%
$$\int_{\mathcal{M}}\dd\m(x)\bigg\vert\sum_{k\geq 1}\widehat{\rho}(k)\, \lambda_k e^{-\tau\lambda_k}\phi_k(x)\bigg\vert^2=\sum_{k\geq 1}\lambda^{2\varepsilon}_k\vert\widehat{\rho}(k)\vert^2\lambda^{2(1-\varepsilon)}_k e^{-2\tau\lambda_k}\lesssim \|\rho\|^2_{{\rm H}^{\varepsilon}}\tau^{2(\varepsilon-1)},$$
%
we get that
%
$$\|\rho_{s}-\rho\|^2_{\LL^2}\lesssim \|\rho\|^2_{{\rm H}^{\varepsilon}}\Big(\int_{0}^{s}\tau^{\varepsilon-1}\,\dd\tau\Big)^2\lesssim_{\varepsilon} \|\rho\|^2_{{\rm H}^{\varepsilon}}s^{2\varepsilon}.$$
%
\textbf{$\LL^\infty$-estimates. }Let $\kappa_1>0$ and $\upsilon>\max\{(\kappa+2)\kappa_1,\frac{1}{\eta}\}$, where $\kappa$ is given in Proposition \ref{Fluctuation}. For the given choices 
%
\begin{equation}\label{Choicestdelta}
%
t:=\frac{\log^{\kappa_2}(n)}{n}\quad\text{and}\quad\delta:=\frac{1}{\log^{\kappa_1}(n)},
%
\end{equation}
%
provided in Proposition \ref{Fluctuation}, we define $h^{n,t}_\delta\in \dot{\mathrm{H}}^1$ the weak solution of 
%
\begin{equation}\label{Equationhdelta}
%
\nabla\cdot\rho_\delta \nabla h^{n,t}_\delta=\mu^{n,t}-\nu^{m,t}.
%
\end{equation}
%
Note that by linearity, one can decompose $h^{n,t}_\delta=h^{(2)n,t}_\delta-h^{(1)n,t}_\delta$ with 
%
\begin{equation}\label{Decompohnt}
%
-\nabla\cdot\rho_\delta \nabla h^{(1)n,t}_\delta=\mu^{n,t}-\rho_t\quad\text{and}\quad-\nabla\cdot\rho_\delta \nabla h^{(2)n,t}_\delta=\nu^{m,t}-\rho_t,
%
\end{equation}
%
so that, considering $u^{n,t}_\delta$ as in \eqref{Eq3Bis} and likewise $v^{n,t}_\delta$ with $\mu^{n,t}$ replaced by $\nu^{m,t}$ and defining 
%
\begin{equation}\label{eq:An}
%
\mathcal{A}_n:=\Big\{\|\mu^{n,t}-\rho_t\|_{\LL^{\infty}}+\|\nu^{m,t}-\rho_t\|_{\LL^{\infty}}\leq \tfrac{1}{\log^{\upsilon}(n)}\Big\}
\end{equation}
as well as
\begin{equation}\label{eq:Bdeltan}
\mathcal{B}_{\delta,n}:=\Big\{\big\|\big(\nabla (u^{n,t}_\delta, v^{n,t}_\delta),\nabla^2 (u^{n,t}_\delta, v^{n,t}_\delta)\big)\big\|_{\LL^\infty}\leq \tfrac{1}{\log^{\upsilon}(n)}\Big\},
%
\end{equation}
%
we deduce from Proposition \ref{Fluctuation} and $m=m(n)\underset{n\uparrow\infty}{\sim} qn$ that
%
\begin{equation}\label{Boundhnt}
%
\big\|(\nabla h^{n,t}_\delta,\nabla^2 h^{n,t}_\delta)\big\|_{\LL^{\infty}}\lesssim \frac{1}{\log^{\upsilon-(\kappa+2)\kappa_1}(n)}\quad\text{in $\mathcal{A}_n\cap \mathcal{B}_{\delta,n}$.}
%
\end{equation}
%
\textbf{$\LL^q$-estimates. }A similar decomposition as in \eqref{Decompohnt} of \eqref{eq:f^nmt} together with \eqref{Eq22Bis} and the choice of $t$ in \eqref{Choicestdelta} yields
%
\begin{equation}\label{LqEstihnt}
%
\Big(\int_{\mathcal{M}}\vert \nabla h^{n,t}\vert^{\bar{q}}\,\dd\m\Big)^{\frac{2}{\bar{q}}}\leq \mathcal{C}_n\frac{\log(n)+\log^{\frac{1}{\eta}}(n)}{n},
%
\end{equation}
%
where $\mathcal{C}_n$ denotes, all along the proof, a random variable which satisfies \eqref{MomentBoundCn} and may change from line to line. 

%
\medskip
%

\textbf{$\LL^2$-regularization error. }Note that from \eqref{Equationhdelta} and \eqref{eq:f^nmt}
%
\begin{equation}\label{DecompositionDeltaProof}
%
-\nabla\cdot\rho_\delta\nabla(h^{n,t}_\delta-h^{n,t})=\nabla\cdot (\rho_\delta-\rho)\nabla h^{n,t},
%
\end{equation}
%
so that from an energy estimate, Hölder's inequality, \eqref{LqEstihnt} and \eqref{Ellipticity}, we obtain 
%
\begin{equation}\label{ErrorApproxL2}
%
\int_{\mathcal{M}}\vert\nabla (h^{n,t}_\delta-h^{n,t})\vert^2\,\dd\m\leq \mathcal{C}_n\|\rho_{\delta}-\rho\|^2_{\LL^{2(\frac{\bar{q}}{2})'}}\frac{\log(n)+\log^{\frac{1}{\eta}}(n)}{n},
%
\end{equation}
%
where $\bar{q}$ denotes the Meyers' exponent of the operator $-\nabla\cdot \rho\nabla$, see Theorem \ref{Meyers}.

%
\medskip
%

Finally, since $\inf_\pi W^2_2(\pi,\gamma^{n,t})\leq (\text{diam}(\mathcal{M}))^2$ and \eqref{Eq13} holds, we can restrict our analysis in $\mathcal{A}_n\cap\mathcal{B}_{\delta,n}$ that we do for the rest of the proof.

%
\medskip
%
{\sc Step 2. Regularization error.}
%
We show that \eqref{ErrorApproxL2} survives when measuring the $W_2$-distance, namely
%
\begin{equation}\label{eq:thm1explipW2}
%
W^2_2(\gamma^{n,t}_\delta,\gamma^{n,t})\leq \mathcal{C}_n \|\rho_\delta-\rho\|^2_{\LL^{2(\frac{\bar{q}}{2})'}}\frac{\log(n)+\log^{\frac{1}{\eta}}(n)}{n}\quad\text{with }\gamma^{n,t}_{\delta}:=\big(\mathrm{Id},\exp(\nabla h^{n,t}_\delta)\big)_\#\mu^{n,t}.
%
\end{equation}
%


%
\medskip
%

Using the coupling $\big((\mathrm{Id},\exp(\nabla h^{n,t}_\delta)),(\mathrm{Id},\exp(\nabla h^{n,t})\big)_\#\mu^{n,t}$ as a competitor in \eqref{Wasserstein} and the fact that $\|\mu^{n,t}\|_{\LL^{\infty}}\lesssim 1$ in $\mathcal{A}_n$, we have
%
\begin{equation}\label{eq:thm1explip}
W^2_2(\gamma^{n,t}_\delta,\gamma^{n,t})\leq \int_{\mathcal{M}}\mu^{n,t}\dd^2 \big(\exp(\nabla h^{n,t}_\delta), \exp(\nabla h^{n,t})\big)\,\dd\m\lesssim \int_{\mathcal{M}}\dd^2 \big(\exp(\nabla h^{n,t}_\delta), \exp(\nabla h^{n,t})\big)\,\dd\m.
\end{equation}
%
We then claim that 
%
\begin{equation}\label{eq:thm1explip2}
\begin{aligned}
\int_\mathcal{M} \dd^2 \big(\exp(\nabla h^{n,t}_\delta), \exp(\nabla h^{n,t})\big)\,\dd\m \lesssim&\,\| \nabla h^{n,t}_\delta - \nabla h^{n,t} \|_{\LL^2}^2\\
&+ \|\rho_\delta-\rho\|^2_{\LL^{(\frac{\bar{q}}{2})'}}\frac{\log(n)}{n}\frac{\| \nabla h^{n,t}_\delta - \nabla h^{n,t}\|_{\LL^2}^2}{\mathbb{E}[ \| \nabla h^{n,t}_\delta - \nabla h^{n,t}\|^{2}_{\LL^2}]}.
\end{aligned}
\end{equation}
%
which, combined with \eqref{eq:thm1explip} and \eqref{ErrorApproxL2} yields \eqref{eq:thm1explipW2}.

%
\medskip
%

We now justify \eqref{eq:thm1explip2}. The difficulty arises from the fact that $\exp$ is not globally Lipschitz. To overcome this, we define
%
\begin{align*}
%
\mathrm{E}_n:=\Big\{\vert\nabla h^{n,t}_\delta-\nabla h^{n,t}\vert &\leq C^{-1}_n\mathbb{E}\big[\|\nabla h^{n,t}_\delta-\nabla h^{n,t}\|^{2}_{\LL^{2}}\big]^{\frac{1}{2}}\Big\}\\
& \text{with }C_n:=\varsigma^{-1}\|\rho_{\delta}-\rho\|_{\LL^{(\frac{\bar{q}}{2})'}}\sqrt{\frac{\log(n)+\log^{\frac{1}{\eta}}(n)}{n}},
%
\end{align*}
%
for a given $\varsigma$ fixed later, and we split
%
\begin{align*}
\lefteqn{\int_\mathcal{M} \dd^2 \big(\exp(\nabla h^{n,t}_\delta), \exp(\nabla h^{n,t})\big)\,\dd \m}\\
&= \int_\mathcal{M} \mathds{1}_{\mathrm{E}_n}\dd^2 \big(\exp(\nabla h^{n,t}_\delta), \exp(\nabla h^{n,t})\big)\,\dd \m+\int_\mathcal{M} \mathds{1}_{\mathrm{E}^c_n}\dd^2 \big(\exp(\nabla h^{n,t}_\delta), \exp(\nabla h^{n,t})\big)\,\dd \m\\
&\leq\int_\mathcal{M}\mathds{1}_{\mathrm{E}_n}\dd^2 \big(\exp(\nabla h^{n,t}_\delta), \exp(\nabla h^{n,t})\big)\,\dd \m+(\text{diam}(\mathcal{M}))^2 \m(\mathrm{E}^c_n).
\end{align*}
%
For the first right-hand side integral, note that from the choice of $C_n$ and \eqref{ErrorApproxL2}, we can choose $\varsigma\ll 1$ uniformly in $n$ such that in $\mathrm{E}_n$ the quantity $\vert\nabla h^{n,t}_{\delta}-\nabla h^{n,t}\vert$ can be made arbitrary small. Since $\exp$ is Lipschitz-continuous in a neighborhood of the null vector, we deduce 
%
$$\int_\mathcal{M}\mathds{1}_{\mathrm{E}_n}\dd^2 \big(\exp(\nabla h^{n,t}_\delta), \exp(\nabla h^{n,t})\big)\,\dd\m\lesssim \|\nabla h^{n,t}_\delta-\nabla h^{n,t}\|^2_{\LL^{2}}.$$
%
For the second right-hand side term, we simply apply Markov's inequality in form of
%
$$\m(\mathrm{E}^c_n)\leq C^2_n\frac{\|\nabla h^{n,t}_\delta-\nabla h^{n,t}\|^2_{\LL^{2}}}{\mathbb{E}\big[\|\nabla h^{n,t}_\delta-\nabla h^{n,t}\|^{2}_{\LL^{2}}\big]}.$$
%
The combination of the two previous estimates gives \eqref{eq:thm1explip2}.

%
\medskip
%

To prove \eqref{ApproximationTransportPlanBis}, we need to control arbitrary $p$-moments, according to Lemma \ref{momentexp}. The argument above can be easily adapted in this case by considering 
%
\begin{align*}
%
\mathrm{E}_n:=\Big\{\vert\nabla h^{n,t}_\delta-\nabla h^{n,t}\vert &\leq C^{-1}_n\mathbb{E}\big[\|\nabla h^{n,t}_\delta-\nabla h^{n,t}\|^{2p}_{\LL^{2}}\big]^{\frac{1}{2p}}\Big\}\\
& \text{with }C_n:=\varsigma\|\rho_{\delta}-\rho\|_{\LL^{(\frac{\bar{q}}{2})'}}\sqrt{\frac{\log^{\frac{1}{\eta}}(n)\log(n)}{n}}.
%
\end{align*}
%
We then follow the same argument, choosing $\varsigma^{-1}=O(\sqrt{p})$.

%\nc{
%For the second term, do like this. Let $\varepsilon\in (0,1)$ and define
%
%$$C_n:=\varepsilon\|\rho_{\delta}-\rho\|_{\LL^{\bar{q}'}}\sqrt{\frac{\log(n)}{n}},$$
%
%and for any $p<\infty$
%
%$$A_n:=\Big\{\vert\nabla f^{n,t}_\delta-\nabla f^{n,t}\vert\leq C^{-1}_n\mathbb{E}\big[\|\nabla f^{n,t}_\delta-\nabla f^{n,t}\|^{2p}_{\LL^{2}}\big]^{\frac{1}{2p}}\Big\}.$$
%
%Now split
%
%\begin{align*}
%\int_\mathcal{M} \dd^2 (\exp(\nabla f^{n,t}_\delta), \exp(\nabla f^{n,t}))\,\dd m=&\int_\mathcal{M} \mathds{1}_{A_n}\dd^2 (\exp(\nabla f^{n,t}_\delta), \exp(\nabla f^{n,t}))\,\dd m\\
%&+\int_\mathcal{M} \mathds{1}_{A^c_n}\dd^2 (\exp(\nabla f^{n,t}_\delta), \exp(\nabla f^{n,t}))\,\dd m\\
%\leq&\int_\mathcal{M} \mathds{1}_{A_n}\dd^2 (\exp(\nabla f^{n,t}_\delta), \exp(\nabla f^{n,t}))\,\dd m+(\text{diam}(\mathcal{M}))^2 m(A^c_n).
%\end{align*}
%
%Now, note that thanks to the choice of $C_n$ and \eqref{eq:errornablafdelta}, we can choose $\varepsilon$ depending on $p$ (explicit) so that in $A_n$, $\vert\nabla f^{n,t}_\delta-\nabla f^{n,t}\vert$ is small so that $\exp$ is Lipschitz. Then, 
%
%$$\int_\mathcal{M} \mathds{1}_{A_n}\dd^2 (\exp(\nabla f^{n,t}_\delta), \exp(\nabla f^{n,t}))\,\dd m\lesssim \|\nabla f^{n,t}_\delta-\nabla f^{n,t}\|^2_{\LL^{2}}.$$
%
%For the second term, we simply apply Markov's inequality in form of
%
%$$m(A^c_n)\leq C^2_n\frac{\|\nabla f^{n,t}_\delta-\nabla f^{n,t}\|^2_{\LL^{2}}}{\mathbb{E}\big[\|\nabla f^{n,t}_\delta-\nabla f^{n,t}\|^{2p}_{\LL^{2}}\big]^{\frac{1}{p}}}.$$
%
%Thus, 
%
%$$\int_\mathcal{M} \dd^2 (\exp(\nabla f^{n,t}_\delta), \exp(\nabla f^{n,t}))\,\dd m\lesssim \|\nabla f^{n,t}_\delta-\nabla f^{n,t}\|^2_{\LL^{2}}+C^2_n\frac{\|\nabla f^{n,t}_\delta-\nabla f^{n,t}\|^2_{\LL^{2}}}{\mathbb{E}\big[\|\nabla f^{n,t}_\delta-\nabla f^{n,t}\|^{2p}_{\LL^{2}}\big]^{\frac{1}{p}}}.$$
%
%}
%
%\fm{
%I think for the Lipschitzianity of $\exp$ you need also that $\nabla f^{n,t}$ is in a neighborhood of the null vector, but in the set $A_n \cap \mathcal{E}_n$ this is true by the simple inequality
%\[
%|\nabla f^{n,t}| \le |\nabla f^{n,t}_\delta - \nabla f^{n,t}| + |\nabla f_\delta^{n,t}|.
%\]
%In \cite[Theorem 1.5]{ambrosio2019optimal} they mention that the exponential map is a global dippheomorphism between a sufficiently small neighborhood of the null vector and its image. I guess we need to be sure that both gradient are small for the Lipschitzianity. Or is there something that you had in mind that only requires the proximity of the two argument of $\exp$?
%}
%

%
\medskip
%

{\sc Step 3. Quantitative stability. }We show that 
%
\begin{equation}\label{QuantitativeEstiPlan}
%
\begin{aligned}
%
\inf_\pi W^2_2(\pi,\gamma^{n,t}_\delta)\lesssim&\, W^2_2\big(\nu^{m,t},\exp(\nabla h^{n,t}_\delta)_\#\mu^{n,t}\big)+W_2\big(\nu^{m,t},\exp(\nabla h^{n,t}_\delta)_\#\mu^{n,t}\big)W_2(\mu^{n},\nu^{m})\\
&+W^2_2(\nu^{m,t},\nu^{m})+W^2_2(\mu^{n,t},\mu^{n})+\big(W_2(\nu^{m,t},\nu^m)+W_2(\mu^{n,t},\mu^n)\big)W_2(\mu^n,\nu^m),
%
\end{aligned}
\end{equation}
%
where we recall that $\gamma^{n,t}_\delta$ is defined in \eqref{eq:thm1explipW2}.

%
\medskip
%

Let $\pi$ be a coupling between $\mu^n$ and $\nu^m$. We introduce a regularization parameter $s<1$ and, smoothing the measure $\mu^n$ into $\mu^{n,s}:=\text{P}_s\mu^n$, the optimal transport plan $\pi^{n,s}$ from $\mu^{n,s}$ to $\nu^m$ is represented by a transport map $T^{n,s}$, according to McCann’s theorem \cite{McCann2001}, that is
%
$$\pi^{n,s}=(\text{Id},T^{n,s})_\#\mu^{n,s}.$$
%
We then apply the triangle inequality in form of
%
\begin{equation}\label{RemarkW2}
\begin{aligned}
W_2(\pi^{n,s},\gamma^{n,t}_\delta)\leq &\, W_2((\mathrm{Id},\exp(\nabla h^{n,t}_\delta)_\#\mu^{n,s},(\mathrm{Id},\exp(\nabla h^{n,t}_\delta)_\#\mu^{n,t})\\
&+W_2(\pi^{n,s},(\mathrm{Id},\exp(\nabla h^{n,t}_\delta)_\#\mu^{n,s}).
\end{aligned}
\end{equation}
%
First, using \eqref{Boundhnt}, $\nabla h^{n,t}_\delta$ is Lipschitz-continuous and $\|\nabla h^{n,t}_\delta\|_{\LL^{\infty}}$ can be made as small as possible for $n$ large. Since $\exp$ is Lipschitz-continuous in a neighborhood of the null vector, we learn from Lemma \ref{lem:transportcontraction} that
%
\begin{equation}\label{EstimateFirstTermRemark}
%
W^{2}_2\big((\mathrm{Id},\exp(\nabla h^{n,t}_\delta))_\#\mu^{n,s},(\mathrm{Id},\exp(\nabla h^{n,t}_\delta))_\#\mu^{n,t}\big)\lesssim W^{2}_2(\mu^{n,t},\mu^{n,s}).
%
\end{equation}
%
Second, we build a competitor for the second right-hand side term of \eqref{RemarkW2}: Defining
%
$$\Gamma:=\big((\text{Id},T^{n,s}),(\text{Id},\exp(\nabla h^{n,t}))\big)_\#\mu^{n,s},$$
%
we have
%
\begin{align*}
W^{2}_2(\pi^{n,s},(\mathrm{Id},\exp(\nabla h^{n,t})_\#\mu^{n,s})&\leq \int_{\mathcal{M}\times\mathcal{M}\times\mathcal{M}\times\mathcal{M}}\delta^2\big((x,z),(y,w)\big)\dd \Gamma\big((x,y),(z,w)\big)\\
&=\int_{\mathcal{M}}\delta^2\big((x,T^{n,s}(x)),(x,\exp(\nabla h^{n,t}(x)))\big)\,\mu^{n,s}(x)\,\dd\m(x)\\
&\stackrel{\eqref{DistanceTransportPlan}}{=}\int_{\mathcal{M}}\dd^2\big(T^{n,s},\exp(\nabla h^{n,t}_{\delta})\big)\mu^{n,s}\,\dd\m.
\end{align*}
%
Using again \eqref{Boundhnt} we can apply, for large $n$, the quantitative stability result of transport maps, Theorem \ref{StabilityResultMap}, to $\mu_1=\nu^m$, $\mu_2=\exp(\nabla h^{n,t}_\delta)_\#\mu^{n,s}$ and $\nu=\mu^{n,s}$ to the effect of 
%
\begin{equation*}%\label{eq:thm1estmap}
%
\begin{aligned}
%
\int_{\mathcal{M}}\dd^2 \big(T^{n,s}, \exp(\nabla h^{n,t}_\delta)\big)\mu^{n,s}\,\dd\m \lesssim&\,W_2^2 \big(\nu^m,\exp(\nabla h^{n,t}_\delta)_\#\mu^{n,s}\big)\\
&+W_2\big(\nu^m,\exp(\nabla h^{n,t}_\delta)_\#\mu^{n,s}\big) W_2\big(\mu^{n,s},\nu^m\big),
%
\end{aligned}
%
\end{equation*}
%
which turns into, using the triangle inequality,
%
\begin{equation}\label{EstimateSecondRemark}
%
\begin{aligned}
%
&\int_{\mathcal{M}}\dd^2 \big(T^{n,s}, \exp(\nabla h^{n,t}_\delta)\big)\mu^{n,s}\,\dd\m\\
&\lesssim\,W_2^2 \big(\nu^{m,t},\exp(\nabla h^{n,t}_\delta)_\#\mu^{n,t}\big)\\
&+W_2\big(\nu^{m,t},\exp(\nabla h^{n,t}_\delta)_\#\mu^{n,t}\big) W_2\big(\mu^{n,s},\nu^m\big)\\
&+W^2_2(\nu^{m,t},\nu^m)+W^2_2\big((\mathrm{Id},\exp(\nabla h^{n,t}_\delta)_\#\mu^{n,s},(\mathrm{Id},\exp(\nabla h^{n,t}_\delta)_\#\mu^{n,t}\big)\\
&+\big(W_2(\nu^{m,t},\nu^m)+W_2((\mathrm{Id},\exp(\nabla h^{n,t}_\delta)_\#\mu^{n,s},(\mathrm{Id},\exp(\nabla h^{n,t}_\delta)_\#\mu^{n,t})\big)W_2(\mu^{n,s},\nu^m).
%
\end{aligned}
%
\end{equation}
%
The combination of \eqref{RemarkW2}, \eqref{EstimateFirstTermRemark} and \eqref{EstimateSecondRemark} yields
%
\begin{equation}\label{EstimateRemarkConclusion}
%
\begin{aligned}
W^2_2(\pi^{n,s},\gamma^{n,t}_\delta)\lesssim&\, W^2_2\big(\nu^{m,t},\exp(\nabla h^{n,t}_\delta)_\#\mu^{n,t}\big)+W_2\big(\nu^{m,t},\exp(\nabla h^{n,t}_\delta)_\#\mu^{n,t}\big)W_2(\mu^{n,s},\nu^{m})\\
&+W^2_2(\nu^{m,t},\nu^{m})+W^2_2(\mu^{n,t},\mu^{n,s})+\big(W_2(\nu^{m,t},\nu^m)+W_2(\mu^{n,t},\mu^{n,s})\big)W_2(\mu^{n,s},\nu^m).
\end{aligned}
%
\end{equation}
%
Since $\mu^{n,s}\underset{s\downarrow 0}{\rightharpoonup}\mu^{n}$, and consequently (up to extracting a subsequence) $\pi^{n,s}\underset{s\downarrow 0}{\rightharpoonup}\pi$, for some optimal transport plan $\pi$, according to the qualitative stability result \cite[Theorem 5.20]{OldNewVillani}, we can pass to the limit as $s\downarrow 0$ in \eqref{EstimateRemarkConclusion} which leads to \eqref{QuantitativeEstiPlan}.

%
\medskip
%

{\sc Step 4. Proof of \eqref{ApproximationTransportPlan}. }
%
We now fix $\kappa_1=\frac{1}{\eta}(\tfrac{\bar q}{2})\frac{1}{2\varepsilon}+1$ such that, applying \eqref{ApproximationRho}, the regularization error \eqref{eq:thm1explipW2} turns into, recalling that $\delta$ is given by \eqref{Choicestdelta}, 
%
\begin{equation}\label{RegularizedErrorGoodScaling}
%
W^2_2(\gamma^{n,t}_\delta,\gamma^{n,t})\leq \mathcal{C}_n\delta^{\frac{2\varepsilon}{(\frac{\bar q}{2})'}}\frac{\log(n)+\log^{\frac{1}{\eta}}(n)}{n}\leq \mathcal{C}_n\frac{1}{n}.
%
\end{equation}
%
It remains to show that 
%
\begin{equation}\label{MainThConclu:Eq1}
%
\begin{aligned}
&\inf_\pi W_2(\pi,\gamma^{n,t}_\delta)\\
&\leq \mathcal{C}_n\frac{\log(n)}{n} \Big(\sqrt{\log^{\kappa_2-1}(n)\|\rho_{t+\frac{1}{n}}-\rho_t\|_{\LL^1}}+\|\rho_\delta-\rho\|^2_{\LL^{2(\frac{\bar{q}}{2})'}}+\|\rho_t-\rho\|^2_{\LL^{2(\frac{\bar{q}}{2})'}}+\tfrac{1}{\log^{\upsilon}(n)}+\sqrt{\tfrac{\log\log n}{\log n}}\Big),
\end{aligned}
%
\end{equation}
%
which together with \eqref{ApproximationRho} and \eqref{RegularizedErrorGoodScaling} leads to \eqref{ApproximationTransportPlan}. To show \eqref{MainThConclu:Eq1}, we control each terms of \eqref{QuantitativeEstiPlan} separately. 

%
\medskip
%

The three last terms are controlled using the contractivity estimate \eqref{ContractivityEstiExpectation} and \eqref{eq:matchingcost} which gives 
%
\begin{align*}
%
&W^2_2(\nu^{m,t},\nu^{m})+W^2_2(\mu^{n,t},\mu^{n})+\big(W_2(\nu^{m,t},\nu^m)+W_2(\mu^{n,t},\mu^n)\big)W_2(\mu^n,\nu^m)\\
&\leq \mathcal{C}_n\frac{\log(n)}{n}\bigg(\sqrt{\frac{\log\log(n)}{\log(n)}}+\sqrt{\log^{\kappa-1}(n)\|\rho_{t+\frac{1}{n}}-\rho_t\|_{\LL^1}}\bigg).
%
\end{align*}
%
For the first two terms, we argue that 
%
\begin{equation}\label{MainResultConclu:Eq2}
%
W^2_2\big(\nu^{m,t},\exp(\nabla h^{n,t}_\delta)_\#\mu^{n,t}\big)\leq \mathcal{C}_n\big(\|\rho_{\delta}-\rho\|^{2}_{\LL^{2(\frac{\bar q}{2})'}}+\|\rho_{t}-\rho\|^{2}_{\LL^{2(\frac{\bar q}{2})'}}+\tfrac{1}{\log^{\upsilon}(n)}\big)\frac{\log(n)}{n},
%
\end{equation}
%
which combined with \eqref{eq:matchingcost} leads to \eqref{MainThConclu:Eq1}.
%

%
\medskip
%

Let us define the curve $\eta: s\in [0,1]\mapsto \eta_s:=(1-s)\mu^{n,t}+s\nu^{m,t}$ and note that from \eqref{Equationhdelta} we have 
%
$$\frac{\dd}{\dd s}\eta_s+\nabla\cdot \Big(\eta_s\frac{\rho_\delta\nabla h^{n,t}_\delta}{\eta_s}\Big)=0.$$
%
Applying Benamou-Brenier' theorem \cite{benamou2000computational}, we learn that 
%
$$\nu^{m,t}=\phi(1,\cdot)_\#\mu^{n,t}\quad\text{with $\phi$ is the flow induced by $s\mapsto \frac{\rho_\delta\nabla h^{n,t}_\delta}{\eta_s}$}.$$
%
Next, using that 
%
\begin{align*}
%
\Big\vert \frac{\rho_\delta\nabla h^{n,t}_\delta}{\eta_s}-\nabla h^{n,t}_\delta\Big\vert&\lesssim \Big(\vert\rho_\delta-\rho\vert+\vert\rho_t-\rho\vert+\vert\mu^{n,t}-\rho_t\vert+\vert\nu^{m,t}-\rho_t\vert\Big)\vert\nabla h^{n,t}_\delta\vert\\
&\lesssim \Big(\vert\rho_\delta-\rho\vert+\vert\rho_t-\rho\vert+\frac{1}{\log^{\upsilon}(n)}\Big)\vert\nabla h^{n,t}_\delta\vert,
%
\end{align*}
%
and applying \cite[Proposition A.1]{ambrosio2019finer} together with Hölder's inequality yields
%
\begin{align}
%
W^2_2\big(\nu^{m,t},\exp(\nabla h^{n,t}_\delta)_\#\mu^{n,t}\big)&=W^2_2\big(\phi(1,\cdot)_\#\mu^{n,t},\exp(\nabla h^{n,t}_\delta)_\#\mu^{n,t}\big)\nonumber\\
&\lesssim \int_{\mathcal{M}}\Big(\vert\rho_\delta-\rho\vert+\vert\rho_t-\rho\vert+\frac{1}{\log^{\upsilon}(n)}\Big)^2\vert\nabla h^{n,t}_\delta\vert^2\nonumber\\
&\leq \big(\|\rho_{\delta}-\rho\|^{2}_{\LL^{2(\frac{\bar q}{2})'}}+\|\rho_{t}-\rho\|^{2}_{\LL^{2(\frac{\bar q}{2})'}}+\tfrac{1}{\log^{\upsilon}(n)}\big)\Big(\int_{\mathcal{M}}\vert\nabla h^{n,t}_\delta\vert^{\bar{q}}\Big)^{\frac{2}{\bar{q}}}.\label{StabilityFlowProof}
%
\end{align}
%
Using Meyers' estimate of Proposition \ref{Meyers} to \eqref{DecompositionDeltaProof} together with \eqref{Ellipticity} and \eqref{LqEstihnt} provides 
%
$$\Big(\int_{\mathcal{M}}\vert\nabla h^{n,t}_\delta\vert^{\bar{q}}\Big)^{\frac{2}{\bar{q}}}\lesssim \Big(\int_{\mathcal{M}}\vert\nabla h^{n,t}\vert^{\bar{q}}\Big)^{\frac{2}{\bar{q}}}\stackrel{\eqref{LqEstihnt}}{\leq} \mathcal{C}_n\frac{\log(n)+\log^{\frac{1}{\eta}}(n)}{n},$$
%
which, combined with \eqref{StabilityFlowProof}, yields \eqref{MainResultConclu:Eq2}.

%
\medskip
%

We finally point out that, in the case $\eta>2$, we use \eqref{ContractivityEstiExpMoment} and \eqref{eq:matchingcostHighMoments} and the same computations lead to \eqref{ApproximationTransportPlanBis}.

%
\medskip
%

{\sc Step 5. Proof of Theorem \ref{th1} and Remark \ref{FlatGeometry}. }The proof of Theorem \ref{th1} follows the same strategy with the main difference that Step 3 is now dropped and Theorem \ref{StabilityResultMap} is directly applied with $\mu_1=\mu^n$, $\nu=\rho$ and $\mu_2=\exp(\nabla f^{n,t}_\delta)_\#\rho\,\dd m$ where $f^{n,t}_\delta$ solves 
%
$$\nabla\cdot \rho_\delta\nabla f^{n,t}=\mu^{n,t}-\rho_t.$$
%
The improvement of Remark \ref{FlatGeometry} follows from the improved contractivity estimate \eqref{ContractivityEstiExpMoment}: Under the assumption \ref{FlatenessAss}, we have (keeping the notations as in Proposition \ref{Contractivity})
%
\begin{equation}\label{InproveContract}
%
W^2_2(\mu^{n,t},\mu^n)\leq \mathcal{D}_n\frac{\log\log(n)}{n},
%
\end{equation}
%
i.e. we do not have the loss $\log^\frac{1}{\eta}(n)$ in \eqref{ContractivityEstiExpMoment}. Inspecting the proof of \eqref{ContractivityEstiExpMoment}, the loss $\log^{\frac{1}{\eta}}(n)$ comes from estimating $v^2$ defined in \eqref{vsquaredContra}. We obtain \eqref{InproveContract} by simply using \eqref{FlatenessAss} and \eqref{AlphaMixing} to upgrade the second item of \eqref{ContractivityCOM:Eq4} into
%
$$v^2\lesssim 1+\sum_{\ell\geq 0}\exp(-b\ell^\eta)\lesssim 1.$$
%
%
\appendix
%
\section{Probabilistic and PDE tools}\label{app:ProbPDE}
%

This section is devoted to recall some probabilistic and analytical tools needed in the proofs. %We first recall Bernstein's inequality to bound averages of i.~i.~d. random variables. For a proof, see \cite[Theorem 3.6 \& 3.7]{chung2006concentration}.
%
%\begin{proposition}\label{Efron-Stein' inequality}
%
%
%\end{proposition}
%
%\begin{proposition}\label{Berstein}
%
%Let $(X_i)_{i\in\mathbb{N}^*}$ be a family of centred random variables such that $\sup_{i\geq 1}\vert X_i\vert\leq M$. Assume that $(X_i)_{i\in\mathbb{N}^*}$ are $i.~i.~d$ random variables with the same distribution $X$. For any $\lambda>0$, it holds
%
%$$\mathbb{P}\bigg(\Big\vert\frac{1}{n}\sum_{i=1}^n X_i\Big\vert>\lambda\bigg)\leq 2\exp\bigg(-\frac{n\lambda^2}{2\mathbb{E}[X^2]+\frac{2}{3}\lambda M}\bigg)\quad\text{for any \quad $n\geq 1$}.$$
%
%\end{proposition}
%
We first recall some concentration inequalities for sequences of random variable satisfying Assumption \ref{Assumptions}. Originally proved for i.i.d. samples, see for instance \cite[Theorem 3.6 \& 3.7]{chung2006concentration}, the proofs in the correlated case can be found in \cite[Theorem 1]{merlevede2011bernstein} and \cite[Theorem 2]{merlevede2009bernstein}.
%
\begin{proposition}\label{BersteinCorrelated}
%
Let $n\in\mathbb{N}$, $M>0$, $\{X_i\}_{i}$ be a family of centred random variables such that $\sup_{i\geq 1}\vert X_i\vert\leq M$ for which \eqref{DecayAlphaMixing} holds.

%
\medskip
%

For any $\lambda>0$, it holds for some constants $(C_i)_{i\in\{1,\cdot,5\}}$ depending on $a$, $b$ :
%
\begin{itemize}
%
\item[(i)]If $\eta<1$,
%
\begin{align*}
\mathbb{P}\bigg(\Big\vert\frac{1}{n}\sum_{i=1}^n X_i\Big\vert>\lambda\bigg)\leq& n\exp\bigg(-\frac{1}{C_1}\Big(\frac{n\lambda}{M}\Big)^{\eta}\bigg)+\exp\bigg(-\frac{1}{C_2}\frac{n^2\lambda^2}{M^2+n v^2}\bigg)\\
&+\exp\bigg(-\frac{1}{C_3}\frac{n\lambda}{M^2}\exp\Big(\frac{1}{C_4}\Big(\frac{n\lambda}{M}\Big)^{\eta(1-\eta)}\log^{-1}(\tfrac{n\lambda}{M})\Big)\bigg),
\end{align*}
%
with
%
$$v^2:=\sup_{i\geq 1}\Big(\mathbb{E}[X^2_i]+2\sum_{j> i}\vert\mathbb{E}[X_iX_j]\vert\Big).$$
%
\item[(ii)]If $\eta=1$,
%
\begin{align*}
\mathbb{P}\bigg(\Big\vert\frac{1}{n}\sum_{i=1}^n X_i\Big\vert>\lambda\bigg)\leq \exp\bigg(-\frac{1}{C_5}\frac{n^2\lambda^2}{nv^2+M^2+n\lambda M(\log(n))^2}\Bigg).
\end{align*}
%
\end{itemize}
%
\end{proposition}
We then recall the link between algebraic moments and exponential moments. The proof is a direct consequence of the Taylor expansion of the exponential function. 
%
\begin{lemma}\label{momentexp}Let $X$ be a non-negative random variable. The following two statements are equivalent: 
\begin{itemize}
\item[(i)]There exists $C_1>0$ such that 
%
$$\mathbb{E}\big[\exp(\tfrac{1}{C_1}X)\big]\leq 2.$$
%
\item[(ii)] There exists $C_2>0$ such that
%
$$\mathbb{E}[X^p]^{\frac{1}{p}}\leq p\,C_2\quad\text{ for any $p<\infty$}.$$
%
\end{itemize}
%
\end{lemma}
%\begin{proof}Let us suppose that there exists $C_2>0$ such that for all $p\geq 1$, $\mathbb{E}[X^p]^{\frac{1}{p}}\leq p C_2$. For all $C_1>0$ we may write
%
%$$\mathbb{E}\bigg[\exp\Big(\frac{1}{C_1}X\Big)\bigg]=\mathbb{E}\bigg[\sum_{n=0}^{+\infty}\frac{X^n}{n!C_1^n}\bigg]\leq \sum_{n=0}^{+\infty}\frac{\left(\frac{C_2}{C_1}n\right)^n}{n!},$$
%
%hen choosing $C_1$ such that $\sum_{n=0}^{+\infty}\frac{\left(\frac{C_2}{C_1}n\right)^n}{n!}\leq 2$ yields the thesis. 

%\medskip
%Let us now suppose that there exists $C_1>0$ such that $\mathbb{E}\bigg[\exp\Big(\frac{1}{C_1}X\Big)\bigg]\leq 2$. This implies that for all $p <\infty$, $\mathbb{E}[ X^p]\leq  C^p_1 p!$. Since, from the Stirling formula, $p!\leq C p^p$ for some $C>0$, we have for all $p\in\mathbb{N}$, $\mathbb{E}[ X^p]^{\frac{1}{p}}\leq C^\frac1p C_1 p$.
%
%\end{proof}
%
%\begin{proposition}[Efron-Stein inequality]\label{EfroStein}
%
%Let $X:=(X_k)_{k\in\{1,\cdots,n\}}$ be $n\geq 1$ independent points in $\mathcal{M}$ with common law $\rho\,\dd x$ and $F: \mathcal{M}^n\rightarrow\mathbb{R}$. Define the vertical derivative 
%
%$$\partial_{X_k} F(X):=F(X)-\mathbb{E}_k[F(X)]\quad\text{for any $k\geq 1$,}$$
%
%where $\mathbb{E}_k[F(X)]$ denotes the conditional expectation
%
%$$\mathbb{E}_k[F(X)]:=\int_{\mathcal{M}}\rho(x)\dd x\, F(X_1,\cdots X_{k-1},x,X_{k+1},\cdots,X_n).$$
%
%We have\nc{check the power of $p$}
%
%$$\mathbb{E}[\vert F(X)-\mathbb{E}[F(X)]\vert^{2p}]^{\frac{1}{p}}\lesssim p\,\mathbb{E}\bigg[\bigg(\sum_{k=1}^n\Big(\partial_{X_k}F(X)\Big)^2\bigg)^{p}\bigg]^{\frac{1}{p}}.$$
%


%\end{proposition}
%
%Furthermore, we state a lemma which gives a relation between probability tails and algebraic moments. The proof is a direct application of the layer cake representation formula.
%
%\begin{lemma}\label{pmoments}
%
%Let $(X_n)_{n\in\mathbb{N}}$ be a sequence of positive random variables satisfying, for some $\eta>0$ 
%
%$$\mathbb{P}(X_n\geq \lambda)\lesssim \exp\Big(-\frac{1}{\eta} \frac{n\lambda^2}{1+\lambda}\Big)\quad \text{for any \quad$\lambda>0$\quad and \quad $n\geq 1$}.$$
%
%Then, 
%
%$$\mathbb{E}[X^p_n]^{\frac{1}{p}}\lesssim p \sqrt{\frac{\eta}{n}}\Big(1+\sqrt{\frac{\eta}{n}}\Big)\quad \text{for any\quad $p<\infty$}.$$
%
%\end{lemma}
%
%\hw{FM: do we want to write this?
%\begin{proof}
%By Cavalieri's formula we may write
%\[
%\mathbb{E} [X^p] = p \int_0^\infty \lambda^{p-1} \mathbb{P}(X_n > \lambda) \dd \lambda \lesssim p \int_0^\infty \lambda^{p-1} \exp\Big(-\frac{An\lambda^2}{1+\lambda}\Big) \dd \lambda.
%\]
%By splitting the latter integral we have
%\[
%\begin{split}
%\int_0^\infty \lambda^{p-1} \exp\Big(-\frac{An\lambda^2}{1+\lambda}\Big) \dd \lambda & = \int_0^{(An)^{-\frac12}} \lambda^{p-1} \exp\Big(-\frac{An\lambda^2}{1+\lambda}\Big) \dd \lambda + \int_{(An)^{-\frac12}}^\infty \lambda^{p-1} \exp\Big(-\frac{An\lambda^2}{1+\lambda}\Big) \dd \lambda \\
%& = I_1 + I_2.
%\end{split}
%\]
%By a brutal estimate we have 
%\begin{equation}\label{eq:cavI1}
%I_1 \le \int_0^{(An)^{-\frac12}} \lambda^{p-1} \dd \lambda = \frac1p (An)^{-\frac p2}.
%\end{equation}
%For $n \gg 1$ we may write
%\[
%I_2 \lessim \int_{(An)^{-\frac12}}^\infty \lambda^{p-1} \exp \Big(-\frac12 An \lambda \Big) \dd \lambda \le (A_n)^{-p} \int_0^\infty t^{p-1} \exp (-t) \dd t \lesssim (A_n)^{-p}.  
%\]
%This and \eqref{eq:cavI1} yields the desired result. (FM: there's a $\frac1p$ on $I_1$, maybe state without a $p$ in front, otherwise also the constant will depend on $p$) 
%\end{proof}
%}
%
%\section{PDE tools}
%
We conclude this section by recalling the standard Meyers' estimate for elliptic equations in divergence form, see for instance the original paper \cite{meyers1963p}.
%
\begin{theorem}[Meyers estimate]\label{Meyers}
%
Let $a : \mathcal{M}\rightarrow \mathbb{R}^{2\times 2}$ be measurable and uniformly elliptic. Consider $u\in \mathrm{H}^1$ the solution of the Neumann boundary problem
%
$$
 \left\{
    \begin{array}{ll}
        -\nabla\cdot a\nabla u=\nabla\cdot g & \text{in $\mathcal{M}$,} \\
        a\nabla u\cdot n_{\mathcal{M}}=0 & \text{on $\partial\mathcal{M}$,}
    \end{array}
\right.
$$
%
for some $g\in \LL^q$ with $q>2$. There exists $2<\bar q<q$ such that 
%
$$\nabla u\in \LL^{\bar q} \quad\text{and}\quad \|\nabla u\|_{\LL^{\bar q}}\lesssim \|g\|_{\LL^{\bar{q}}}.$$
%
\end{theorem}
%
\section{Matching cost for point clouds}\label{sec:matchcost}
%
This section is devoted to recall the upper bounds on the matching cost, results which can be found in \cite[Theorem 2]{borda2021empirical} under mild $\beta$-mixing conditions. The case of Markov chains have been studied in \cite{Riekert, Fournier2015} where sharp upper bounds are obtained. We include a short proof for convenience. %\nc{The proof of \eqref{eq:matchingcost} is already done in \cite{borda2021empirical}, so do not redo it. Just justify why we can can higher moments and why it is true for the Markov chains. In principle, do not rewritte previous estimates, just point out which terms are treated differently and how.}
%

\begin{proposition}[Matching cost]\label{prop:matchcost}
%
Let $\rho$ satisfying \eqref{Ellipticity} and $\{\mu^n\}_n$ be defined in \eqref{eq:empmeas} with point clouds satisfying the Assumption \ref{Assumptions} or in the class of Markov chains satisfying the Assumption \ref{ClassMarkov}. There exists a constant $C>0$ such that 
%
\begin{equation}\label{eq:matchingcost}
%
W_2^2(\mu^n, \rho \,\dd \m) \le \mathcal{C}_n \frac{\log(n)}{n}\quad\text{with $\sup_{n\geq 1}\mathbb{E}\big[\tfrac{1}{C}\mathcal{C}_n\big]\leq 1$.}
%
\end{equation}
%
Furthermore, if \eqref{DecayAlphaMixing} holds with $\eta\geq 1$ then the assumption \eqref{DecatBetaMixing} can be dropped and the stochastic integrability can be improved up to losing a $\log(n)$ factor, namely 
%
\begin{equation}\label{eq:matchingcostHighMoments}
%
W_2^2(\mu^n, \rho \,\dd \m) \le \mathcal{D}_n \frac{\log^{\frac{1}{\eta}}(n)\log(n)}{n}\quad\text{with $\sup_{n\geq 1}\mathbb{E}\big[\exp(\tfrac{1}{C}\mathcal{D}_n)\big]\leq 1$.}
%
\end{equation}
%
\end{proposition}
\begin{proof}
Note that the proof of \eqref{eq:matchingcost} can be found in \cite[Theorem 2]{borda2021empirical} when the point cloud satisfies the Assumption \ref{Assumptions}. We first show how \eqref{eq:matchingcost} can be extended to point clouds which are sampled from a Markov chain satisfying the Assumption \ref{ClassMarkov}. Second, we show how the stochastic integrability can be improved to \eqref{eq:matchingcostHighMoments} when \eqref{DecayAlphaMixing} holds with $\eta \ge 1$.  

\medskip
{\sc Step 1. Markov chains case. }Recall that a Markov chain satisfying the Assumption \ref{ClassMarkov} admits an absolutely continuous invariant measure of the form $\mu_\infty= \rho\, \dd \m$ with $\rho$ satisfying \eqref{Ellipticity}, that is $\lambda\leq \rho\leq \Lambda$. Recalling that we denote by $\{\lambda_n,\phi_n\}_n$ the set of eigenvalues and normalized eigenfunctions of the Laplace-Beltrami operator $-\Delta$ on $\mathcal{M}$, we have by definition \eqref{eq:empmeas} of $\mu^n$, for any $k\geq 1$
%
\begin{equation}\label{ed:decomposefouriermark}
\widehat{\mu^n}(k) - \widehat{\rho}(k) = \frac1n \sum_{\ell=1}^n (\phi_k(X_\ell) - \mathbb{E}[\phi_k(X_\ell)]) + \frac1n \sum_{\ell=1}^n (\mathbb{E}[\phi_k(X_\ell)] - \mu_\infty (\phi_k)),
\end{equation}
%
where we use interchangeably the notation $\widehat{\rho}(k) = \int \phi_k \rho \,\dd \m = \mu_\infty(\phi_k)$. Using the Berry-Esseen smoothing inequality \cite[Theorem 5]{borda2021empirical} together with \eqref{ed:decomposefouriermark}, we get
%
\begin{equation}\label{eq:proofmarkstep1-2}
\begin{split}
\mathbb{E}[W_2^2(\mu^n, \mu_\infty)] \lesssim & \frac1n + \sum_{k\ge1}\frac{e^{-\frac1n\lambda_k}}{\lambda_k} \mathbb{E}\bigg[\bigg(\frac1n \sum_{\ell=1}^n (\phi_k(X_\ell) - \mathbb{E}[\phi_k(X_\ell)])\bigg)^2 \bigg] \\
& + \sum_{k\ge1}\frac{e^{-\frac1n\lambda_k}}{\lambda_k} \mathbb{E} \bigg[ \bigg( \frac1n \sum_{\ell=1}^n (\mathbb{E}[\phi_k(X_\ell)] - \mu_\infty (\phi_k)) \bigg)^2\bigg].
\end{split}
\end{equation}
%
We now estimate the last two terms of \eqref{eq:proofmarkstep1-2} separately and we start with the third one. Using \eqref{eq:lemMark} and \eqref{ContractivityCOM:Eq4}, we have 
\[
|\mathbb{E}[\phi_k(X_\ell)] - \mu_\infty(\phi_k)| \stackrel{\eqref{eq:lemMark}}{\lesssim} \exp(-b\ell^{\eta}) \|\phi_k\|_{\LL^\infty} \stackrel{\eqref{ContractivityCOM:Eq4}}{\lesssim} \exp(-b\ell^{\eta})\lambda_k^\frac12,
\]
Thus, using in addition \eqref{eq:traceformula}, we get
%
\begin{equation}\label{eq:proofmarkstep1-3}
\sum_{k\ge1}\frac{e^{-\frac1n\lambda_k}}{\lambda_k} \mathbb{E}\bigg[\bigg(\frac1n \sum_{\ell=1}^n (\phi_k(X_\ell) - \mathbb{E}[\phi_k(X_\ell)])\bigg)^2 \bigg] \lesssim \frac1{n^2} \sum_{k\ge1} e^{-\frac1n \lambda_k} \stackrel{\eqref{eq:traceformula}}{=} \frac1{n^2}\int_{\mathcal{M}} p_{\frac1n}(x,x)\,\dd\m(x) \lesssim \frac1n.
\end{equation}
%
We now turn to the second term of \eqref{eq:proofmarkstep1-2}. Expanding the square provides
%
\begin{equation}\label{eq:proofmatchcost-2}
\begin{split}
\lefteqn{\sum_{k \ge 1} \frac{e^{-\frac1n \lambda_k}}{\lambda_k}\Big(\frac1n \sum_{\ell=1}^n (\phi_k(X_\ell) - \mathbb{E}[\phi_k(X_\ell)])\Big)^2}\\ 
& = \frac1{n^2}\sum_{k\ge1}\sum_{\ell=1}^n \frac{e^{-\frac1n \lambda_k}}{\lambda_k} |\phi_k(X_\ell) - \mathbb{E}[\phi_k(X_\ell)]|^2 \\
& + \frac2{n^2} \sum_{k\ge1}\sum_{1\le \ell < \ell'\le n} \frac{e^{-\frac1n \lambda_k}}{\lambda_k} (\phi_k(X_\ell) - \mathbb{E}[\phi_k(X_\ell)])(\phi_k(X_{\ell'}) - \mathbb{E}[\phi_k(X_{\ell'})]).
\end{split}
\end{equation}
We now estimate the two terms on the right hand side of \eqref{eq:proofmatchcost-2}. For the first term, an easy induction argument combining \eqref{AbsoluteContinuity} and \eqref{LawOfTheChain} show that for any $n\geq 1$, $\mathbb{P}_{X_n}\ll \m$ with $\lambda\leq\frac{\dd \mathbb{P}_{X_n}}{\dd\m}\leq \Lambda$. Therefore, we have 
%
$$ \mathbb{E}\big[|\phi_k(X_\ell) - \mathbb{E}[\phi_k(X_\ell)]|^2\big]\leq \Lambda,$$
%
and we deduce
%
\begin{equation}\label{eq:proofmatchcost-3}
\frac1{n^2}\sum_{k\ge1}\sum_{\ell=1}^n \frac{e^{-\frac1n \lambda_k}}{\lambda_k} \mathbb{E}\big[|\phi_k(X_\ell) - \mathbb{E}[\phi_k(X_\ell)]|^2\big] \lesssim \frac1n \sum_{k \ge 1}  \frac{e^{-\frac1n \lambda_k}}{\lambda_k}.
\end{equation}
%
For the second term, we use \eqref{eq:markchainbetamixc} to obtain
%
\begin{align*}
&\frac2{n^2} \mathbb{E}\bigg[ \sum_{k\ge1}\sum_{1\le \ell < \ell'\le n} \frac{e^{-\frac1n \lambda_k}}{\lambda_k} (\phi_k(X_\ell) - \mathbb{E}[\phi_k(X_\ell)])(\phi_k(X_{\ell'}) - \mathbb{E}[\phi_k(X_{\ell'})])\bigg]\\
& \lesssim \frac1{n^2}\sum_{1\le \ell < \ell' \le n} \beta_{\ell \ell'} \sup_{x,y \in \mathcal{M}} \bigg\lvert \sum_{k\ge1} \frac{e^{-\frac1n \lambda_k}}{\lambda_k} \phi_k(x)\phi_k(y)\bigg\rvert \\
&\lesssim \frac1n \sup_{x,y \in \mathcal{M}} \bigg\lvert \sum_{k\ge1} \frac{e^{-\frac1n \lambda_k}}{\lambda_k} \phi_k(x)\phi_k(y)\bigg\rvert.
\end{align*}
%
Combining the latter with \eqref{eq:proofmarkstep1-2}, \eqref{eq:proofmarkstep1-3}, \eqref{eq:proofmatchcost-2} and \eqref{eq:proofmatchcost-3} yields 
%
\begin{equation}\label{eq:proofmarkstep1-6}
\mathbb{E}[W_2^2(\mu^n, \mu_\infty)] \lesssim \frac1n + \frac1n \sum_{k\ge1} \frac{e^{-\frac1n\lambda_k}}{\lambda_k} + \frac1n \sup_{x,y \in \mathcal{M}} \bigg\lvert \sum_{k\ge1} \frac{e^{-\frac1n \lambda_k}}{\lambda_k} \phi_k(x)\phi_k(y)\bigg\rvert.
\end{equation}
%
We finally conclude similarly as for \eqref{ControlExpectationLast}.
%Thus we are left to estimate the last two terms on the right hand side of \eqref{eq:proofmarkstep1-6}. To this end we give the argument for the last term, the other being the same. By the spectral decomposition of the heat kernel \eqref{SpectralDecompoHeatKernel} and \eqref{eq:eigvaleighfunc} we get
%\begin{equation}\label{eq:mathccostlastest}
%\begin{split}
%\bigg\lvert\sum_{k\ge1} \frac{e^{-\frac1n \lambda_k}}{\lambda_k} \phi_k(y)\phi_k(z)\bigg\rvert & = \bigg\lvert\int_{\frac1n}^1 \dd s \, p_s(x,y) + \sum_{k\ge 1} \frac{e^{-\lambda_k}}{\lambda_k}\phi_k(x)\phi_k(y)\bigg\rvert  \stackrel{\eqref{eq:eigvaleighfunc}}{\lesssim} \log n + \sum_{k\ge 1} e^{-\lambda_k} \\
%& \stackrel{\eqref{eq:traceformula}}{\lesssim} \log (n) + \int_{\mathcal{M}} \dd \m \, p_1 (x,x) \lesssim \log (n).
%\end{split}
%\end{equation}
%Finally the latter together with \eqref{eq:proofmarkstep1-6} implies \eqref{eq:matchingcost}.
%
\medskip
%

{\sc Step 2. Higher stochastic integrability. }We now prove \eqref{eq:matchingcostHighMoments}. We argue using the moment estimate \eqref{ContractivityCOM:Eq1} which, together with Minkowski's inequality and $\lambda_k^\frac1{\log(n)} e^{-\frac1n \lambda_k} \lessim e^{-\frac1{2n}\lambda_k}$ implies
\[
\mathbb{E}\bigg[\bigg(\sum_{k\ge1} \frac{e^{-\frac1n \lambda_k}}{\lambda_k}|\widehat{\mu^n}(k)-\widehat{\rho}(k)|^2\bigg)^p\bigg]^\frac1p \lesssim p^2 \frac1n \sum_{k\ge 1} \frac{e^{-\frac1{2n} \lambda_k}}{\lambda_k}\bigg( \log^\frac1\eta(n) + \frac{\lambda_k(1+\log^2(n))}{n}  \bigg).
\]
Finally, combining the latter with the Berry-Esseen smoothing inequality \cite[Theorem 5]{borda2021empirical} and arguing similarly as for \eqref{ControlExpectationLast} yields \eqref{eq:matchingcostHighMoments} thanks to Proposition \ref{momentexp}.
\end{proof}
%
%\section{Improvement}
%
%\subsection{Regularity on $\rho$}
%
%{\sc Ideas: }My feeling is that the regularity of $\rho$ is not essential. Let $\rho\in\LL^{\infty}(\mathcal{M})$ such that $\rho\geq c>0$. The only place where I used regularity is for the approximation, but only in $\LL^2(\mathcal{M})$. More precisely
%
%$$\int_{\mathcal{M}}\vert\nabla f^{n,t}_{\delta}-f^{n,t}\vert^2\lesssim \int_{\mathcal{M}}\vert\rho_\delta-\rho\vert^2\vert\nabla f^{n,t}\vert^2.$$
%
%Since $\rho\in \LL^{\infty}(\mathcal{M})$, one has for every $q<\infty$
%
%$$\|\rho_\delta-\rho\|_{\LL^{q}(\mathcal{M})}\underset{\delta\downarrow 0}{\rightarrow} 0.$$
%
%In addition, from Meyers's estimate, we know that there exists $\bar q>1$ such that $\nabla f^{n,t}\in \LL^{2\bar q}(\mathcal{M})$. Therefore, from Hölder's inequality, 
%
%$$\int_{\mathcal{M}}\vert\nabla f^{n,t}_{\delta}-f^{n,t}\vert^2\leq \bigg(\int_{\mathcal{M}}\vert\rho_\delta-\rho\vert^{2\bar q'}\bigg)^{\frac{1}{\bar q'}}\bigg(\int_{\mathcal{M}}\vert\nabla f^{n,t}\vert^{2\bar q}\bigg)^{\frac{1}{\bar q}}.$$
%
%Also, using \eqref{Eq1}, we have from Meyers's estimate
%
%$$\bigg(\int_{\mathcal{M}}\vert\nabla f^{n,t}\vert^{2\bar q}\bigg)^{\frac{1}{\bar q}}\lesssim \bigg(\int_{\mathcal{M}}\vert\nabla u^{n,t}\vert^{2\bar q}\bigg)^{\frac{1}{\bar q}}.$$
%
%Can we still get a bound on 
%
%$$\mathbb{E}\bigg[ \Big(\int_{\mathcal{M}}\vert\nabla u^{n,t}\vert^{2\bar q}\Big)^{\frac{p}{\bar q}}\bigg]^{\frac{1}{p}},$$
%
%for any $p<\infty$ ? Like in Proposition \ref{prop1} ? This should be possible, they did this for $\bar q=4$ in Lemma \cite[LEmma 3.17]{ambrosio2019pde} for instance. Then, if this is possible, is it possible to adapt \cite[Theorem 4.4]{ambrosio2019pde} but only with a bound for the quantity 
%
%$$\mathbb{E}\Big[(W^2_2(\rho\,\dd x,\mu^{n,t}))^{p}\Big]^{\frac{1}{p}},$$
%
%for any $p<\infty$ ? We don't care about the constant $\frac{1}{4\pi}$ here, we just use previous bounds, especially the one on $\bigg(\int_{\mathcal{M}}\vert\nabla f^{n,t}\vert^{2\bar q}\bigg)^{\frac{1}{\bar q}}$.
%

%
%\medskip
%

%{\sc Thoughts: }We can say that $\nabla u^{n,t}$ has actually better integrability and look at 
%
%$$\mathbb{E}\bigg[ \Big(\int_{\mathcal{M}}\vert\nabla u^{n,t}\vert^{4}\Big)^{\frac{p}{2}}\bigg]^{\frac{1}{p}}.$$
%
%Using the arguments of \cite[Lemma 3.17]{ambrosio2019pde}, we have 
%
%$$\int_{\mathcal{M}}\vert\nabla u^{n,t}\vert^4\lesssim\int_{\mathcal{M}} \bigg(\int_{0}^{\infty}\dd s\Big((-s\Delta)^{\frac{1}{2}}\text{P}_{s+t}(\mu^{n}-\rho)\Big)^{2}\bigg)^2.$$
%
%Applying Berstein's inequality would require $\LL^{\infty}(\mathcal{M})\cap \LL^{2}(\mathcal{M})$ bounds on $(-s\Delta)^{\frac{1}{2}}(p_{s+t}(\cdot,X_1)-\rho_{s+t})$. $\LL^{2}(\mathcal{M})$ bounds is already explained in the paper whereas for $\LL^{\infty}(\mathcal{M})$ we use a bit of magic here. We have the following formula for any $f\in\LL^{2}(\mathcal{M})$ such that $\int_{\mathcal{M}} f=0$
%
%$$(-\Delta)^{-\frac{1}{2}}f=\frac{1}{\sqrt{\pi}}\int_{0}^{\infty}\dd\tau\, \tau^{-\frac{1}{2}}\text{P}_\tau f.$$
% 
%Applying this with $f=p_{s+t}(\cdot,X_1)-\rho_{s+t}$ combined with the semi-group property yields
%
%$$(-\Delta)^{-\frac{1}{2}}(p_{s+t}(\cdot,X_1)-\rho_{s+t})=\frac{1}{\sqrt{\pi}}\int_{s+t}^{\infty}\dd\tau\, \tau^{-\frac{1}{2}}(p_{\tau}(\cdot,X_1)-\rho_{\tau}).$$
%
%Therefore, for any $x\in\mathcal{M}$
%
%\begin{align*}
%\vert(-\Delta)^{\frac{1}{2}}(p_{s+t}(\cdot,X_1)-\rho_{s+t})(x)\vert=&\vert(-\Delta)(-\Delta)^{-\frac{1}{2}}(p_{s+t}(\cdot,X_1)-\rho_{s+t})(x)\vert\\
%\lesssim&\int_{s+t}^{\infty}\dd \tau\, \tau^{-\frac{1}{2}}\vert\Delta (p_{\tau}(x,X_1)-\rho_\tau(x))\vert.
%\end{align*}
%
%Then, we use that
%
%$$\vert\Delta (p_{\tau}(x,X_1)-\rho_\tau)\vert\lesssim \min\{\tau^{-2},\tau^{-\frac{5}{2}}\},$$
%
%so that 
%
%$$\|(-s\Delta)^{\frac{1}{2}}(p_{s+t}(\cdot,X_1)-\rho_{s+t})\|_{\LL^{\infty}(\mathcal{M})}\lesssim s^{\frac{1}{2}}\min\{(s+t)^{-\frac{3}{2}},(s+t)^{-2}\}.$$
%
%We will have to deal with integrals like
%
%$$\int_{0}^{\infty}\dd s\,s^{\frac{1}{2}}\min\{(s+t)^{-\frac{3}{2}},(s+t)^{-2}\},$$
%
%which is of the right order since
%
%\begin{align*}
%\int_{0}^{\infty}\dd s\,s^{\frac{1}{2}}\min\{(s+t)^{-\frac{3}{2}},(s+t)^{-2}\}&=\Big(\int_{0}^t+\int_{t}^1+\int_{1}^{\infty}\Big)\dd s\,s^{\frac{1}{2}}\min\{(s+t)^{-\frac{3}{2}},(s+t)^{-2}\}\\
%&\lesssim t^{-\frac{3}{2}}\int_{0}^t\dd s\, s^{\frac{1}{2}}+\int_{t}^1\dd s\, s^{-1}+\int_{1}^{\infty}\dd s\, s^{-\frac{3}{2}}\\
%&\lesssim \vert\log(t)\vert.
%\end{align*}
%






\section{Proof for the class of Markov chains}\label{ClassMarkovAppendix}
%
We provide in this Section the arguments for extending Theorem \ref{MainResultMatchingClouds} and Theorem \ref{th1} to the class of Markov chains introduced in Section \ref{Examples}. The proof follows the lines of the proof of Theorem \ref{MainResultMatchingClouds}, where the main difference is that we drop the assumption that the point clouds is identically distributed. That affects the proofs of the main ingredients (we recall that the scaling of the cost has already be proven in Proposition \ref{prop:matchcost}), namely the $\LL^q$ estimates in Proposition \ref{prop1}, the fluctuation estimates in Proposition \ref{Fluctuation} and the contractivity estimates in Proposition \ref{Contractivity}. We show in the following how to adapt the proofs for a given Markov chain $\{X_n\}_{n\geq 1}$ satisfying Assumption \ref{ClassMarkov}. In the following, we recall that $\mu_\infty=\rho\,\dd\m$ denotes the unique invariant measure of the chain. We split the proof into three steps.

%
\medskip
%

{\sc Step 1. $\LL^q$ estimates. }We have to understand the extra error term coming from the deviation of $\mathbb{E}[\mu^n]$ from $\mu_\infty$. In view of \eqref{Eq51}, it is 
%
$$\bigg(\int_{\mathcal{M}}\dd\m\Big(\int_{0}^\infty \dd s\,\big((-s\Delta)^{\frac{1}{2}}\text{P}_{s+t}(\mathbb{E}[\mu^n]-\mu_\infty)\big)^2\Big)^2\bigg)^{\frac{1}{2}}.$$
%
Using the definition \eqref{eq:empmeas} of $\mu^n$, the convergence to equilibrium \eqref{eq:lemMark} applied to $f=(-s\Delta)^{\frac{1}{2}}p_{s+t}(x,\cdot)$ and the heat-kernel estimates \eqref{Eq28}, we have for any $s\geq 0$ and $x\in\mathcal{M}$
\begin{equation}\label{eq:estLqmark}
%
\begin{aligned}
(-s\Delta)^\frac12 \text{P}_{s+t}\big(\mathbb{E}[\mu^n] - \rho\big)(x)
% & = \frac1n \sum_{\ell=1}^n (-s\Delta)^\frac12 (\mathbb{E}[p_{s+t}(x,X_\ell)] - \mu_\infty(p_{s+t}(x,\cdot))\\
& = \frac1n \sum_{\ell=1}^n  (\mathbb{E}[(-s\Delta)^\frac12 p_{s+t}(x,X_\ell)] - \mu_\infty((-s\Delta)^\frac12 p_{s+t}(x,\cdot))\\ & \stackrel{\eqref{eq:lemMark}}{\lesssim}  \frac{\|(-s \Delta)^\frac12 p_{s+t}(x,\cdot)\|_{\LL^\infty}}n\stackrel{\eqref{Eq28}}{\lesssim}\frac{s^\frac12(s+t)^{-\frac32}}n,
\end{aligned}
%
\end{equation}
%
so that, recalling $t=\frac{\log^{\kappa}(n)}{n}$,
%
$$\bigg(\int_{\mathcal{M}}\dd\m\Big(\int_{0}^\infty \dd s\,\big((-s\Delta)^{\frac{1}{2}}\text{P}_{s+t}(\mathbb{E}[\mu^n]-\mu_\infty)\big)^2\Big)^2\bigg)^{\frac{1}{2}}\lesssim \frac{1}{n^2}\int_0^\infty s(s+t)^{-3}\lesssim \frac{1}{n\log^{\kappa}(n)}\ll\frac{\log(n)}{n}.$$
%
{\sc Step 2. Fluctuation estimates. }Here, the distribution of the Markov chain affects the concentration estimate \eqref{Eq59}. We show that, defining 
%
\begin{equation}\label{ExtensionMarkovFluctu3}
%
\bar{u}^{n,t}_\delta:=\int_{0}^\infty e^{-s}\text{P}_s\Big(\tfrac{1}{\rho_\delta}(\mu^{n,t}-\mathbb{E}[\mu^{n,t}])\Big)\,\dd s,
%
\end{equation}
%
we have
%
\begin{equation}\label{ExtensionMarkovFluctu}
%
\mathbb{P}\Big(\vert\partial^2_{ij} u^{n,t}_\delta(x)\vert\geq \tfrac{1}{2\log^{\nu}(n)}\Big)\leq \mathbb{P}\Big(\vert\partial^2_{ij} \bar{u}^{n,t}_\delta(x)\vert\geq \tfrac{1}{4\log^{\nu}(n)}\Big)\quad\text{for any $x\in\mathcal{M}$,}
%
\end{equation}
%
where the r.h.s can be estimated following the lines of the proof of \eqref{Eq59}. As before, we investigate the extra term coming from the deviation of $\mathbb{E}[\mu^n]$ from $\mu_\infty$. The estimate \eqref{ExtensionMarkovFluctu} follows from 
%
\begin{equation}\label{ExtensionMarkovFluctu2}
%
\|\partial^2_{ij}(u^{n,t}_\delta-\bar{u}^{n,t}_\delta)\|_{\LL^{\infty}}\ll \frac{1}{\log^{\nu}(n)}.
%
\end{equation}
%
We argue as in \eqref{Eq68}, decomposing $u^{n,t}_\delta-\bar{u}^{n,t}_\delta$ into a regular-part and a singular part: for any $x\in\mathcal{M}$
%
\begin{equation}\label{MarkovExtensionFluctu4}
%
\begin{aligned}
%
(u^{n,t}_\delta-\bar{u}^{n,t}_\delta)(x)=&\frac1{\rho_\delta(x)}\int_0^\infty e^{-s}\,\big(\mathbb{E}[\mu^{n,t+s}] - \rho_{t+s}\big)(x)\,\dd s\\
&+ \int_0^\infty e^{-s}\,\text{P}_s \Big(\big(\tfrac1{\rho_\delta}-\tfrac1{\rho_\delta(x)}\big)\big(\mathbb{E}[\mu^{n,t}] - \rho_t\big)\Big)(x)\,\dd s.
%
\end{aligned}
%
\end{equation}
%
To estimate the second r.h.s integral of \eqref{MarkovExtensionFluctu4}, we use \eqref{Eq69}. For the first r.h.s integral, that we denote by $\mathcal{J}$, we use the definition \eqref{eq:empmeas} of $\mu^n$, the convergence to equilibrium \eqref{eq:lemMark} and the heat-kernel bounds \eqref{Eq28} to obtain 
%
\begin{align*}
%
\vert\partial^2_{ij}\mathcal{J}(x)\vert=&\bigg\vert\frac{1}{n}\sum_{k=1}^n\partial^2_{ij}\bigg(\frac1{\rho_\delta(\cdot)}\int_0^\infty e^{-s}\,\big(\mathbb{E}[p_{t+s}(\cdot,X_k)] - \mu_\infty(p_{t+s}(x,\cdot))\big)\,\dd s\bigg)(x)\bigg\vert\\
\lesssim& \frac{1}{n}\sum_{k=1}^n\bigg(\|\nabla^2\tfrac{1}{\rho_\delta}\|_{\LL^\infty}\int_{0}^\infty \vert\mathbb{E}[p_{s+t}(\cdot,X_k)]-\mu_\infty(p_{t+s}(x,\cdot))\vert\\
&+\int_{0}^\infty \vert\mathbb{E}[\nabla^2 p_{s+t}(\cdot,X_k)]-\mu_\infty(\nabla^2 p_{t+s}(x,\cdot))\vert\\
&+\|\nabla\tfrac{1}{\rho_\delta}\|_{\LL^\infty}\int_{0}^\infty \vert\mathbb{E}[\nabla p_{s+t}(\cdot,X_k)]-\mu_\infty(\nabla p_{t+s}(x,\cdot))\vert\Big)\\
\stackrel{\eqref{eq:lemMark},\eqref{Eq28}\eqref{Eq25}}{\lesssim}&\frac{1}{n}\bigg(\delta^{-2}\int_{0}^\infty \min\{(s+t)^{-1},(s+t)^{-\frac{3}{2}}\}\,\dd s+\int_{0}^\infty (s+t)^{-2}\,\dd s+\delta^{-1}\int_{0}^{\infty} (s+t)^{-\frac{3}{2}}\, \dd s\bigg)\\
\lesssim & \frac{1}{n}(\delta^{-2}t^{-\frac{1}{2}}+t^{-1}+\delta^{-1}t^{-\frac{1}{2}})\ll \frac{1}{\log^{\nu}(n)}.
%
\end{align*}
%
{\sc Step 3. Contractivity estimate. }Here, the law of the Markov chain affects the estimate \eqref{ContractExpectation}. The extra error term coming from the deviation of $\mathbb{E}[\mu^n]$ from $\mu_\infty$ reads
%
$$\sum_{k\geq 1} \frac{e^{-\frac{2}{n}\lambda_k}}{\lambda_k}\big(e^{-t\lambda_k}-1\big)^2\vert\mathbb{E}[\widehat{\mu^n}(k)]-\mu_\infty(\phi_k)\vert^2.$$
%
Using the definition \eqref{eq:empmeas} of $\mu^n$ and the convergence to equilibrium \eqref{eq:lemMark} applied with $f=\phi_k$ and the bound on the eigenfunctions \eqref{eq:eigvaleighfunc}, we have for any $k\leq n$
%
\begin{align*}
%
\vert\mathbb{E}[\widehat{\mu^n}(k)]-\mu_\infty(\phi_k)\vert=&\frac{1}{n}\bigg\vert\sum_{\ell=1}^n(\mathbb{E}[\phi_k(X_{\ell})]-\mu_\infty(\phi_k))\bigg\vert^2\stackrel{\eqref{eq:lemMark}}{\lesssim} \frac{1}{n}\|\phi_k\|_{\LL^{\infty}}\stackrel{\eqref{eq:eigvaleighfunc}}{\lesssim}\frac{1}{n}\lambda^{\frac{1}{2}}_k,
%
\end{align*}
%
so that, using the trace formula \eqref{eq:traceformula} and the heat-kernel estimates \eqref{Eq28}, we deduce 
%
$$\sum_{k\geq 1} \frac{e^{-\frac{2}{n}\lambda_k}}{\lambda_k}\big(e^{-t\lambda_k}-1\big)^2\vert\mathbb{E}[\widehat{\mu^n}(k)]-\mu_\infty(\phi_k)\vert^2\lesssim \frac{1}{n^2}\sum_{k\geq 1}e^{-\frac{2}{n}\lambda_k}(e^{-t\lambda_k}-1)^2\stackrel{\eqref{eq:traceformula},\eqref{Eq28}}{\lesssim} \frac{1}{n} \ll \frac{\log\log(n)}{n}.$$
%
%\nc{A bit like same remarks as previously, do not rewritte previous estimates, just point out the terms you treat differently and how.}
%\nc{ATTENTION ! There are arguments missing... The justification of the $\LL^q$-estimates require a priori some comments. Indeed, see for instance in \eqref{ExpandSquareLaplace}, $\omega_{s+t}$ is not of mean zero, you have to substract the mean and argue that you don't add an higher order term. Please check carefully that all the arguments are there.}
%In this section we show that \eqref{ApproximationTransportPlan} and \eqref{ApproximationTransportMap} extend to the class of Markov chains satisfying the Assumption \ref{ClassMarkov}. The proof follows the same lines as the ones of Theorems \ref{MainResultMatchingClouds} and \ref{th1}, thus we just highlight the main differences in the Markovian case of Assumption \ref{ClassMarkov}. As in the setting of Theorem \ref{MainResultMatchingClouds}, the key ingredient for the proof is the quantitative stability result \cite[Theorem 3.2]{ambrosio2019optimal}. We will make use of the same notation of the previous sections with the understanding that $\rho$ is now the density of the stationary distribution $\pi$ of the Markov chain as in \eqref{AbsoluteContinuity} and $\mu^n, \nu^m$ will be referring to the empirical measures associated to the Markov chains $(X_n)_{n \ge 0}, (Y_n)_{n \ge 0}$ satisfying the Assumption \ref{ClassMarkov} with the same initial distribution $\lambda$ and the same transition kernel $K$. In particular, the two Markov chains will have the same invariant distribution $\pi$. Moreover, recall that $h^{n,t}$ will denote the solution to \eqref{eq:f^nt} and $h^{n,t}_\delta$ will denote the solution to \eqref{Equationhdelta}.

%\medskip
%We first notice that the $\beta$-mixing condition of the Assumption \ref{ClassMarkov} is stronger than the $\alpha$-mixing condition of the Assumption \ref{Assumptions}. Indeed, in general the $\beta$-mixing coefficient (recall \eqref{BetaMixing} for a definition) provides an upper bound for the $\alpha$-mixing coefficient (recall $\eqref{AlphaMixing}$ for a definition), thus \eqref{eq:markchainbetamixc} implies the decay \eqref{DecayAlphaMixing}. As in Theorem \ref{MainResultMatchingClouds} we can restrict our analysis to the set $\mathcal{A}_n \cap B_{\delta,n}$ (recall \eqref{eq:An} and \eqref{eq:Bdeltan} for a definition). Indeed, the arguments of Proposition \ref{Fluctuation} can be adapted to the Markovian case. The main issue is proving an analogous of \eqref{Eq59}. Indeed, since the states of a Markov chain are not identically distributed we cannot directly apply the bound \eqref{eq:markchainbetamixc} to obtain the decay \eqref{eq:stella2} as the quantity $\omega$ defined in \eqref{eq:stella} do not have zero mean. However, we can overcome this difficulty by adding and subtracting the expectation of the empirical measure on the integral representation \eqref{eq:stella3} to get
%\[
%\begin{split}
%u^{n,t}_\delta & = \int_0^\infty \dd s \, e^{-s}P_s \bigg(\frac1{\rho_\delta}(\mu^{n,t} - \rho_t)\bigg) \\
%& = \underbrace{\int_0^\infty \dd s \, e^{-s}P_s \bigg(\frac1{\rho_\delta}(\mu^{n,t} - \mathbb{E}[\mu^{n,t}])\bigg)}_{\tilde{u}^{n,t}_\delta} + \underbrace{\int_0^\infty \dd s \, e^{-s}P_s \bigg(\frac1{\rho_\delta}(\mathbb{E}[\mu^{n,t}] - \rho_t)\bigg)}_{\bar{u}^{n,t}_\delta}.
%\end{split}
%\]
%Along the same lines of the proof of \eqref{Eq59} we can show that the second derivative $\tilde{u}^{n,t}_\delta$ satisfy an analogous decay in probability. We now show that the the spatial second derivative of the deterministic error term $\bar{u}_\delta^{n,t}$ are bounded by an higher order term. We first decompose $\bar{u}_\delta^{n,t}$ into a ``regular'' and an ``singular'' part. Adding and subtracting $\frac1{\rho_\delta(x)}$ in the integral we get
%\[
%\begin{split}
%\lefteqn{\int_0^\infty \dd s \, e^{-s}P_s \bigg(\frac1{\rho_\delta}(\mathbb{E}[\mu^{n,t}] - \rho_t)\bigg)} \\
%& = \underbrace{\int_0^\infty \dd s \, e^{-s}P_s \bigg(\frac1{\rho_\delta(x)}(\mathbb{E}[\mu^{n,t}] - \rho_t)\bigg)}_{:=\tilde{\mathcal{J}}_1} + \underbrace{\int_0^\infty \dd s \, e^{-s}P_s \bigg(\bigg(\frac1{\rho_\delta}-\frac1{\rho_\delta(x)}\bigg)(\mathbb{E}[\mu^{n,t}] - \rho_t)\bigg)}_{:=\tilde{\mathcal{J}}_2}
%\end{split}
%\]
%By the same argument as in the proof of Proposition \ref{prop1}, the ``singular'' term $\tilde{\mathcal{J}}_2$ satisfies the estimate \eqref{Eq69}. We now focus on the ``regular'' term. By the definition of the empirical measure \eqref{eq:empmeas}, by Lemma \ref{lem:Markblock} and by \eqref{Eq28} we get
%\[
%\begin{split}
%|\partial_{ij}^2 \tilde{\mathcal{J}}_1(x)| & = \frac1n \partial_{ij}^2 \bigg( \frac 1{\rho_\delta}\bigg) \sum_{k=1}^n \int_0^\infty \dd s \, e^{-s}(\mathbb{E}[p_{s+t}(x, X_k)] - \pi(p_{s+t}(x,\cdot))\\
%& + \frac1n  \frac 1{\rho_\delta} \sum_{k=1}^n \int_0^\infty \dd s \, e^{-s}(\mathbb{E}[\partial_{ij}^2p_{s+t}(x, X_k)] - \pi(\partial_{ij}^2p_{s+t}(x,\cdot))\\
%& +\frac1n \partial_i \bigg( \frac 1{\rho_\delta}\bigg) \sum_{k=1}^n \int_0^\infty \dd s \, e^{-s}(\mathbb{E}[\partial_j p_{s+t}(x, X_k)] - \pi(\partial_j p_{s+t}(x,\cdot))\\
%& +\frac1n \partial_j \bigg( \frac 1{\rho_\delta}\bigg) \sum_{k=1}^n \int_0^\infty \dd s \, e^{-s}(\mathbb{E}[\partial_i p_{s+t}(x, X_k)] - \pi(\partial_i p_{s+t}(x,\cdot))\\
%& \stackrel{\eqref{Eq28},\eqref{Eq25},\eqref{eq:lemMark}}{\lesssim} \frac{\delta^{-2}}n \int_0^\infty \dd s\, (s+t)^{-1} e^{-s} + \frac1n \int_0^\infty \dd s\, (s+t)^{-2} + \delta^{-1}\int_0^\infty \dd s\,(s+t)^{-\frac32} \\
%& \lesssim \frac{\delta^{-2}}n|\log(t)|+ \frac{t^{-1}}n+\frac{\delta^{-1}}nt^{-\frac12} \ll \log^{-\nu}(n). 
%\end{split}
%\] 
%Therefore, an analogous of Proposition \ref{prop1} holds and we can apply \cite[Theorem 3.2]{ambrosio2019optimal}. By the same argument as in {\sc Step 3} of the proof of Theorem \ref{MainResultMatchingClouds} we get the quantitative stability estimate  
%\begin{equation}\label{eq:MarkQuantStab}
%%
%\begin{aligned}
%%
%\inf_q W^2_2(q,\gamma^{n,t}_\delta)\lesssim&\, W^2_2\big(\nu^{n,t},\exp(\nabla h^{n,t}_\delta)_\#\mu^{n,t}\big)+W_2\big(\nu^{n,t},\exp(\nabla h^{n,t}_\delta)_\#\mu^{n,t}\big)W_2(\mu^{n},\nu^{n})\\
%&+W^2_2(\nu^{n,t},\nu^{n})+W^2_2(\mu^{n,t},\mu^{n})+\big(W_2(\nu^{n,t},\nu^n)+W_2(\mu^{n,t},\mu^n)\big)W_2(\mu^n,\nu^n),
%%
%\end{aligned}
%\end{equation}
%where the $\inf$ runs over all optimal transport plans $q$ between $\mu^n$ and $\nu^n$. Therefore, it remains to estimate the terms on the right hand side of \eqref{eq:MarkQuantStab}. 
%
%\medskip
%The matching cost between the two point clouds is readily estimated by \eqref{eq:matchingcost} which, together with the triangle inequality for the $W_2$-distance, implies
%\begin{equation}\label{eq:estwasmarkcloud}
%W_2^2 (\mu^n, \nu^n) \lesssim \mathcal{C}_n \frac{\log (n)}n \quad\text{with $\sup_{n\ge 1} \mathbb{E}[\frac1C \mathcal{C}_n \le1]$.}
%\end{equation}
%
%\medskip
%Moreover the refined heat contraction estimate
%\begin{equation}\label{eq:refinedcontractionmark}
%W_2^2(\mu^{n,t}, \mu^n) \le \mathcal{C}_n \frac{\log \log (n)}{n}\quad\text{with $\sup_{n\ge 1} \mathbb{E}[\frac1C \mathcal{C}_n \le1]$.}
%\end{equation}
%holds in the Markovian case as well. Indeed, by Lemma \ref{lem:Markblock} the proof of Proposition \ref{Contractivity} can be adapted to the Markovian case. As in the setting of identically distributed point clouds, we can restrict our analysis to the set $\mathcal{A}_n$ where the elliptic condition \eqref{EllipticityAn} holds. Arguing as in {\sc Step 1} of Proposition \ref{Contractivity} we get the smoothing inequality \eqref{SmoothingIneg}. As opposed to the i.i.d. case the decomposition formula \eqref{DecomposeFourier} does not hold since the $X_n$ are not identically distributed. However, we can argue as in Proposition \ref{prop:matchcost} and treat the term as in \eqref{ed:decomposefouriermark}. Adding and subtracting the expectation of the empirical measure $\mu^n$ in \eqref{SmoothingIneg} and using the triangle inequality yields the upper bound
%\begin{equation}\label{eq:proofmark11}
%\begin{split}
%W^2_2 (\mu^{n,t},\mu^n) & \lesssim \frac1n +\sum_{k\geq 1} \frac{e^{-\frac{2}{n}\lambda_k}}{\lambda_k}\big(e^{-t\lambda_k}-1\big)^2\vert\widehat{\mu^n}(k)-\mathbb{E}[\widehat{\mu^n}(k)]\vert^2\\
%& +\sum_{k\geq 1} \frac{e^{-\frac{2}{n}\lambda_k}}{\lambda_k}\big(e^{-t\lambda_k}-1\big)^2\vert\mathbb{E}[\widehat{\mu^n}(k)]-\widehat{\rho}(k)\vert^2 +t\|\rho_{t+\frac{1}{n}}-\rho_{\frac{1}{n}}\|_{\LL^1}.
%\end{split}
%\end{equation}
%We note that the second term on the right hand side of \eqref{eq:proofmark11} can be estimated as in the proof of Proposition \ref{Contractivity}, thus the Proposition \ref{Contractivity} can be adapted to the Markovian case once we show that the additional error given by the third term is of higher order. By Lemma \ref{lem:Markblock} and by \eqref{ContractivityCOM:Eq4} we obtain
%\begin{equation}\label{eq:markproof12}
%\begin{split}
%&\lefteqn{\sum_{k\geq 1} \frac{e^{-\frac{2}{n}\lambda_k}}{\lambda_k}\big(e^{-t\lambda_k}-1\big)^2\vert\mathbb{E}[\widehat{\mu^n}(k)]-\widehat{\rho}(k)\vert^2}
%\\
%& = \sum_{k\geq 1} \frac{e^{-\frac{2}{n}\lambda_k}}{\lambda_k}\big(e^{-t\lambda_k}-1\big)^2\bigg(\frac1n\sum_{\ell=1}^n\vert\mathbb{E}[\phi_k(X_\ell)]-\pi(\phi_k(X_\ell)\vert\bigg)^2  \\
%& \stackrel{\eqref{eq:lemMark}}{\lesssim} \frac1{n^2}\sum_{k \ge 1}\frac{e^{-\frac{2}{n}\lambda_k}}{\lambda_k}\big(e^{-t\lambda_k}-1\big)^2 \|\phi_k\|_{\LL^\infty} \\ 
%&\stackrel{\eqref{ContractivityCOM:Eq4}}{\lesssim} \frac1{n^2} \sum_{k \ge 1} e^{-\frac2n \lambda_k}(e^{-t \lambda_k} - 1)^2 \lesssim \sum_{k \ge 1} e^{-\frac2n \lambda_k} \stackrel{\eqref{eq:traceformula}}{=} \int_{\mathcal{M}} \dd \mathrm{m}\, p_{\frac2n}(x,x) \lesssim \frac1n. 
%\end{split}
%\end{equation}
%Finally, \eqref{eq:proofmark11} together with \eqref{eq:markproof12} implies \eqref{eq:refinedcontractionmark}. 
%
%%which leads to the estimate\nc{this looks complicated, I would explain that the difference is that we need to substract the expectation. Thanks to the assumption \eqref{ConvergenceRate}, the remaing term is of higher order (show the proof) and we then proceed the same way}
%%\[
%%\begin{split}
%%\lefteqn{\mathbb{E}\bigg[\sum_{k\geq 1} \frac{e^{-\frac{2}{n}\lambda_k}}{\lambda_k}\big(e^{-t\lambda_k}-1\big)^2\vert\widehat{\mu^n}(k)-\widehat{\rho}(k)\vert^2\bigg]}\\
%%& \lesssim \frac{1}{n}\bigg(\sum_{k\geq 1} \frac{e^{-\frac{2}{n}\lambda_k}}{\lambda_k}\big(e^{-t\lambda_k}-1\big)^2+\sup_{x,y\in\mathcal{M}}\bigg\vert\sum_{k\geq 1} \frac{e^{-\frac{2}{n}\lambda_k}}{\lambda_k}\big(e^{-t\lambda_k}-1\big)^2\phi_k(x)\phi_k(y)\bigg\vert\bigg).
%%\end{split}
%%\]
%%Finally, \eqref{eq:refinedcontractionmark} follows from \eqref{ControlExpectationLast}.
%
%
%
%\medskip
%Moreover, we have the estimate
%\begin{equation}\label{eq:firsttermmark}
%%
%W^2_2\big(\nu^{n,t},\exp(\nabla h^{n,t}_\delta)_\#\mu^{n,t}\big)\leq \mathcal{C}_n\frac{\log(n)}{n}\quad\text{with $\sup_{n\ge 1} \mathbb{E}[\frac1C \mathcal{C}_n \le2]$.}
%%
%\end{equation}
%The proof \eqref{eq:firsttermmark} will be identical to the one of \eqref{MainResultConclu:Eq2} provided an analogous of Proposition \ref{prop1} holds when the samples are drawn from not identically distributed point clouds. The main issue introduced by the Markovian case is that we cannot directly consider the representation given by \eqref{ExpandSquareLaplace} to get a $\LL^q$ type estimate, but we have to add and subtract the expectation of the empirical measure $\mu^n$ in order to apply the $\beta$-mixing bound \eqref{eq:markchainbetamixc}. In particular this provides us the upper bound for \eqref{Eq54Bis} given by 
%\begin{equation}\label{eq:proofmark14}
%\begin{split}
%\lefteqn{\mathbb{E} [ ((-s\Delta)^\frac12 P_{s+t}(\mu^n - \rho)(x))^2 ]}\\
%&  \lesssim \mathbb{E} [ ((-s\Delta)^\frac12 P_{s+t}(\mu^n - \mathbb{E}[\mu^n])^2 ] + ((-s\Delta)^\frac12 P_{s+t}(\mathbb{E}[\mu^n] - \rho)(x))^2.
%\end{split}
%\end{equation}
%The first term on the right hand side of \eqref{eq:proofmark14} can be estimated in the same way as in the proof of Proposition \ref{prop1}, while the second deterministic term provides a higher order term. Indeed, by the definition of the empirical measure $\mu^n$ \eqref{eq:empmeas}, by the linearity of the operator $(-s\Delta)^\frac12$, by Lemma \ref{lem:Markblock} and by \eqref{Eq28} we have
%\begin{equation}\label{eq:estLqmark}
%\begin{split}
%((-s\Delta)^\frac12 P_{s+t}(\mathbb{E}[\mu^n] - \rho)(x))  
%% & = \frac1n \sum_{\ell=1}^n (-s\Delta)^\frac12 (\mathbb{E}[p_{s+t}(x,X_\ell)] - \pi(p_{s+t}(x,\cdot))\\
%& = \frac1n \sum_{\ell=1}^n  (\mathbb{E}[(-s\Delta)^\frac12 p_{s+t}(x,X_\ell)] - \pi((-s\Delta)^\frac12 p_{s+t}(x,\cdot))\\ & \stackrel{\eqref{eq:lemMark}}{\lesssim}  \frac{\|(-s \Delta)^\frac12 p_{s+t}(x,\cdot)\|_{L^\infty(\mathcal{M})}}n\stackrel{\eqref{Eq28}}{\lesssim}\frac{s^\frac12(s+t)^{-\frac32}}n,
%\end{split}
%\end{equation}
%which is of higher order than the right hand side of \eqref{Eq54Bis}. 
%
%\medskip
%Finally, combining \eqref{eq:MarkQuantStab}, \eqref{eq:estwasmarkcloud}, \eqref{eq:refinedcontractionmark}, \eqref{eq:firsttermmark} and recalling the assumption $\rho \in \mathrm{H}^\varepsilon$ yields the desired estimate:
%\[
%\inf_{q}W^2_2(q,\gamma^{n,t})\leq \mathcal{C}_n\frac{\log n}{n}\sqrt{\frac{\log\log(n)}{\log(n)}}\quad\text{with \,$\sup_{n\geq 1}\mathbb{E}[\tfrac{1}{C}\mathcal{C}_n]\leq 1$,}
%\]
%%
%where the $\inf$ runs over all optimal transport plans $q$ between $\mu^n$ and $\nu^n$.


\section*{Acknowledgments} The authors warmly thank Lorenzo Dello Schiavo, Antonio Agresti and Martin Huesmann for useful discussions and fruitful comments.

\bibliographystyle{plain}
\bibliography{references}
\end{document}
